% Generated by build_chapter.py from the GY1966 files. Do not edit by hand; edit the files and rebuild.
\chapter{The academic year 1966--67}\label{ch:1966}
This chapter covers all 38 MCCME records numbered 92--129 and all 68 supplied scans, P000140--P000207. The catalogue order is preserved even when dates are not monotone. Its overlapping page ranges identify records; each occurrence instead points to the pages containing that actual statement. The year heading, prospective plan, and Parshin heading remain present without invented mathematical content. The outcome is a critical occurrence index of 277 items, 21 indexed problems (one continuing an earlier problem identifier), 108 correction or qualification entries, and 78 checked bibliographic entries. Status assessments are dated 9 November 2026.

\begin{center}
\begin{tabular}{lrrrrrr}
\toprule
Class & A & B & C & D & E & F\\
\midrule
Items & 22 & 45 & 3 & 0 & 185 & 22\\
\bottomrule
\end{tabular}
\end{center}
Here A denotes an authorial publication soon after the talk, B subsequent development by others, C a precisely delimited partly developed question, D an undeveloped identifiable item, E prior mathematics or exposition, and F failure of exact identification. These are historical classifications, not correctness scores. An item may be A and false as written; a reconstruction may be correct without proving what the speaker intended. The separate assessment fields preserve these distinctions.

Three lines dominate. Bernstein’s equivariant K-theory and elliptic-family constructions connect representation rings, completion, and index bundles. The recorded claim that every representation ring embeds in its completion fails for a disconnected compact group; a concrete cyclic-group counterexample is supplied. Minlos’s lectures move from canonical models and factorial correlations to thermodynamic quantities and limiting distributions. The later published lectures explicitly trace their 1965 and 1966 origins \citep[Introduction]{minlos1968lectures}. The year also follows the expansion of singularity theory: jets, Boardman symbols, intrinsic derivatives, preparation, and real local algebras. Arnold’s two November lectures are singularity-theory lectures despite the catalogue’s hydrodynamics label \citep{arnold1968singularities,boardman1967singularities}.

The strongest publication matches include Kazhdan’s property-\(T\) announcement, Shilov’s regular Hilbert-space analysis, Seeley’s complex-power kernel statements, and Vishik--Eskin’s family-index criterion for completion of elliptic boundary systems \citep{kazhdan1967dual,shilov1967hilbert1,seeley1968powers,vishikeskin1967normally}. In the last case the manuscript received in April 1967 resolves the March question within its stated setting; it is not counted as an orphan problem. Dates of talks, receipt, preprint circulation and printing are kept separate whenever the evidence distinguishes them.

The three C items concern the algebraic topological classification of Anosov diffeomorphisms, the arbitrary-smooth-metric remainder of Helgason’s exact geodesic-sphere support question, and the smooth type of equal-prime-period geodesic manifolds. Each is matched with its own hypotheses; a geodesic-ray theorem, an analyticity theorem, or an entire-Grauert-tube theorem is not substituted for the recorded problem \citep[Conjecture 2.3]{spatzier2016rigidity}\citep[Problem 4.1]{quinto2008helgason}\citep[Main Theorem]{lisong2026zoll}. No D item is justified, and none of these C items yielded a plausible bounded research advance in this run.

The full mathematical annotation therefore audits Minlos’s variance argument. It proves an explicit conditional variance estimate, an \(L^1\) truncation conclusion, and counterexamples to finite variance under a bare \(L^1\) hypothesis. It also identifies the exact canonical volume correction and the legitimate dependence of the effective volume on correlation order. These are complete, independently checked editorial derivations, with code and outputs, and carry no claim of a new Gibbs-state theorem. The larger thermodynamic programme’s undefined exceptional parameters are retained as a formulation limitation (F).

Continuity with earlier chapters includes lattice rigidity and cohomology (1964-8 and 1964-24), family indices (1964-15.4), the Maslov thread (1964-21.3), and the pseudodifferential calculus of 1965. The perturbed complex square reuses GP-1965-03 and its earlier occurrence 1965-72.5. The corrections distinguish changed readings, translations, unresolved ink, provenance, and faithful mathematical errors: the Boardman sign and positive-sequence convention, the elliptic local-algebra monomial, entropy normalization, and the heat-kernel factor four are consequential examples. Every source was opened, but a verified link does not mean that every proof was read; the item assessments and bibliography annotations state the actual verification level.

\GStalk{92}{1966-09-05}{I. M. Gelfand}{P000140-P000142}{Programme; Brieskorn links and exotic spheres}
The actual opening heading and programme occupy P000140: statistical mechanics, pseudodifferential operators, algebraic groups, and deformations. A report attributed to Hirzebruch then describes exotic spheres as links of isolated hypersurface singularities. The catalogue range overlaps subsequent talks; items below follow the handwritten headings. The printed exponent requires an emendation, not merely removal of an uncertainty marker.

\begin{GSitem}{1966-92.1}{E}{A seven-dimensional hypersurface link}
\GSfield{Kind; source}{construction; \GSpage{P000140}}
\GSfield{The notes}{The notebook intersects \(z_1^3+z_2^{3k-1}+z_3^2+z_4^2+z_5^2=0\) with \(\sum|z_i|^2=\epsilon\), and counts \(10-3=7\) real dimensions.}
\GSfield{Statement}{For integers \(k\geq1\), the corrected Brieskorn link \(\Sigma(3,6k-1,2,2,2)\subset\mathbb C^5\), cut out on a sufficiently small positive sphere, is a smooth closed real seven-manifold and a homotopy sphere. The written exponent \(3k-1\) does not give homology spheres for every \(k\).}
\GSfield{Reconstruction}{\(6k-1\) is our mathematical emendation. In the exponent graph for \((3,3k-1,2,2,2)\), odd \(k\) leaves an even component of four vertices, which fails the homology-sphere criterion. Even \(k=2\ell\) gives the standard family with parameter \(\ell\).}
\GSfield{Fate}{Brieskorn published the construction in 1966; Hirzebruch explained its sphere criterion in his November 1966 expos\'e, printed in 1968 \citep[\S\S2--3]{brieskorn1966beispiele,hirzebruch1968singularities}.}
\GSfield{Present status}{As of 9 November 2026: established. The notebook is an early report of an existing construction, with a substantive exponent error.}
\GSfield{Assessment}{literal statement: false as written; publication match: exact; verification: read the result; development: annotated; confidence: high}
\end{GSitem}
\begin{GSitem}{1966-92.2}{E}{Twenty-eight oriented seven-sphere classes}
\GSfield{Kind; source}{theorem; \GSpage{P000140}}
\GSfield{The notes}{As \(k\) varies, all 28 “Milnor spheres” are said to occur.}
\GSfield{Statement}{The group of oriented smooth homotopy seven-spheres is \(\Theta_7=bP_8\cong\mathbb Z/28\). With the corrected exponent, 28 consecutive parameters represent its 28 oriented diffeomorphism classes, including the standard sphere.}
\GSfield{Reconstruction}{The orientation convention and inclusion of the standard sphere are supplied explicitly.}
\GSfield{Fate}{The signature computation and the enumeration are given in \citet[\S3]{hirzebruch1968singularities}, reporting the 1966 construction.}
\GSfield{Present status}{As of 9 November 2026: established; 28 counts all classes, not 28 exotic classes.}
\GSfield{Assessment}{literal statement: correct with missing hypotheses; publication match: exact; verification: read the result; development: none; confidence: high}
\end{GSitem}
\begin{GSitem}{1966-92.3}{E}{The analogous eleven-dimensional enumeration}
\GSfield{Kind; source}{theorem; \GSpage{P000140}}
\GSfield{The notes}{The report adds that the same method produces 992 spheres in dimension 11.}
\GSfield{Statement}{Adding two quadratic variables gives \(\Sigma(3,6k-1,2,2,2,2,2)\), a homotopy eleven-sphere. Its classes exhaust \(\Theta_{11}=bP_{12}\cong\mathbb Z/992\), with orientations retained.}
\GSfield{Reconstruction}{The explicit seven-variable polynomial is reconstructed from the phrase “the same method.”}
\GSfield{Fate}{\citet[\S3]{hirzebruch1968singularities} gives the corresponding higher-dimensional signature calculation and the order 992.}
\GSfield{Present status}{As of 9 November 2026: established. This numerical claim survives correction of the preceding family.}
\GSfield{Assessment}{literal statement: correct; publication match: exact; verification: read the result; development: none; confidence: high}
\end{GSitem}

\GStalk{93}{1966-09-05}{D. A. Kazhdan}{P000140-P000143}{Bass rigidity and an integral cohomology reduction}
P000141--P000142 report Bass’s work on rigidity of \(\mathrm{SL}_n(\mathbb Z)\), relating deformations to first cohomology and proposing reduction modulo a prime. This resumes the cohomological rigidity and congruence-subgroup discussions indexed at 1964-8.1--2 and 1964-24.1--2. The established conclusion and the incompletely recorded proof reduction must be distinguished.

\begin{GSitem}{1966-93.1}{B}{Local rigidity of the standard higher-rank lattice}
\GSfield{Kind; source}{theorem; \GSpage{P000141}}
\GSfield{The notes}{For \(n>2\), a continuous family of homomorphisms \(\varphi_t:\mathrm{SL}_n(\mathbb Z)\to\mathrm{SL}_n(\mathbb R)\) starting with inclusion has the form \(g(t)\gamma g(t)^{-1}\).}
\GSfield{Statement}{For \(n\geq3\), the standard inclusion is locally rigid: all homomorphisms sufficiently close on a finite generating set are conjugate to it. A continuous parameter path starting at inclusion is locally conjugate; a path of conjugators can be chosen on an interval by lifting through the conjugation orbit.}
\GSfield{Reconstruction}{Closeness on generators and the local-to-path interpretation are made explicit.}
\GSfield{Fate}{\citet[\S16, Corollaries 16.4--16.5 and the following remark]{bassmilnorserre1967congruence} give the cohomological rigidity consequences; publication was in 1967 by Bass, Milnor and Serre, whose work Kazhdan is reporting.}
\GSfield{Present status}{As of 9 November 2026: established. This is local deformation rigidity, not the assertion that arbitrary homomorphisms belong to one component.}
\GSfield{Assessment}{literal statement: correct with missing hypotheses; publication match: exact; verification: read the result; development: none; confidence: high}
\end{GSitem}
\begin{GSitem}{1966-93.2}{E}{The first-cohomology criterion}
\GSfield{Kind; source}{method; \GSpage{P000141}}
\GSfield{The notes}{The derivative of a deformation is a 1-cocycle with values in the adjoint module; conjugation produces a coboundary. Vanishing \(H^1\) is used to prove rigidity.}
\GSfield{Statement}{For a finitely generated group \(\Gamma\), Lie group \(G\), and homomorphism \(\rho\), a differentiable path has infinitesimal class in \(H^1(\Gamma,\mathfrak g_{\operatorname{Ad}\rho})\). If that group vanishes, Weil’s local rigidity criterion applies. Nonzero cohomology alone need not integrate to a nontrivial deformation.}
\GSfield{Reconstruction}{The warning about nonintegrable cocycles is an editorial clarification.}
\GSfield{Fate}{The criterion predates this talk and is explicitly used in \citet[\S16, remark on rigidity]{bassmilnorserre1967congruence}; it recurs from 1964-8.2.}
\GSfield{Present status}{As of 9 November 2026: established; the obstruction to integrating arbitrary cocycles must not be omitted.}
\GSfield{Assessment}{literal statement: correct with missing hypotheses; publication match: partial; verification: read the result; development: none; confidence: high}
\end{GSitem}
\begin{GSitem}{1966-93.3}{E}{Passing from an integral adjoint lattice to real coefficients}
\GSfield{Kind; source}{method; \GSpage{P000141}}
\GSfield{The notes}{The notes choose \(L=\mathfrak{sl}_n(\mathbb Z)\) and identify real first cohomology with its integral counterpart tensored with \(\mathbb R\).}
\GSfield{Statement}{For a finitely presented \(\Gamma\) and a finite-rank free integral \(\Gamma\)-module \(L\), \(H^1(\Gamma,L)\otimes\mathbb R\cong H^1(\Gamma,L\otimes\mathbb R)\). Thus finiteness of the integral group proves real vanishing.}
\GSfield{Reconstruction}{Finite presentation makes the cocycle relations a finite matrix calculation and justifies tensoring their kernels with the flat module \(\mathbb R\).}
\GSfield{Fate}{The integral conclusion appears in \citet[Corollary 16.5]{bassmilnorserre1967congruence}; the change-of-scalars step is elementary homological algebra.}
\GSfield{Present status}{As of 9 November 2026: correct with finiteness hypotheses. No assertion about all degrees or arbitrary coefficient modules is intended.}
\GSfield{Assessment}{literal statement: correct with missing hypotheses; publication match: partial; verification: read the result; development: none; confidence: high}
\end{GSitem}
\begin{GSitem}{1966-93.4}{B}{Reduction modulo a prime and the missing restriction term}
\GSfield{Kind; source}{method; \GSpage{P000141-P000142}}
\GSfield{The notes}{The proof attempts to make multiplication by \(p\) surjective on \(H^1(\Gamma,L)\), then reduces to a finite quotient and to \(H^1(\Gamma/\Gamma_p,L/pL)=0\).}
\GSfield{Statement}{From \(0\to L\xrightarrow{p}L\to L/pL\to0\), vanishing \(H^1(\Gamma,L/pL)\) implies \(p\)-surjectivity. For finitely generated abelian \(H^1(\Gamma,L)\), this excludes a free summand. But inflation--restriction also has \(H^1(\Gamma_p,L/pL)^{\Gamma/\Gamma_p}\); vanishing on the quotient does not alone remove that term.}
\GSfield{Reconstruction}{Our reconstruction supplies the exact-sequence obstruction and deliberately does not infer the missing vanishing.}
\GSfield{Fate}{The final finiteness theorem matches \citet[Corollary 16.5]{bassmilnorserre1967congruence}, by the authors whose work is being reported; the exact prime-reduction proof in these notes was not matched to a complete published argument.}
\GSfield{Present status}{As of 9 November 2026: the conclusion is established; the recorded proof remains incomplete, rather than an independent open theorem.}
\GSfield{Assessment}{literal statement: correct with missing hypotheses; publication match: partial; verification: read the result; development: annotated; confidence: medium}
\end{GSitem}

\GStalk{94}{1966-09-05}{V. I. Arnold}{P000141-P000147}{Anosov systems, structural stability, and Smale questions}
The handwritten Arnold heading is on P000142; his discussion ends on P000145. Uniform hyperbolicity, an arithmetic nilmanifold, Morse--Smale systems, Axiom A, and spectral decomposition are interspersed with explicit conjectures. Positive and negative later answers belong to different questions; in particular hyperbolicity of the nonwandering set alone does not force dense periodic points.

\begin{GSitem}{1966-94.1}{E}{A U-system is uniformly hyperbolic on the whole manifold}
\GSfield{Kind; source}{definition; \GSpage{P000142}}
\GSfield{The notes}{The tangent bundle splits into \(X\oplus Y\), with one part contracted and the other expanded; such an automorphism is a \(U\)-system.}
\GSfield{Statement}{A \(C^1\) diffeomorphism \(f\) of a closed smooth manifold is Anosov if \(TM=E^s\oplus E^u\) is a continuous invariant splitting and \(\|Df^nv^s\|,\|Df^{-n}v^u\|\leq C\lambda^n\|v^{s,u}\|\) for \(n\geq0\), \(C<\infty\), \(0<\lambda<1\).}
\GSfield{Reconstruction}{The modern constants and bundle regularity explicate the shorthand.}
\GSfield{Fate}{This is the established Anosov notion reviewed in \citet[\S\S1--2]{smale1967dynamical}; the survey followed the 1966 Moscow congress.}
\GSfield{Present status}{As of 9 November 2026: standard. Continuity, invariance and uniform constants are essential.}
\GSfield{Assessment}{literal statement: correct with missing hypotheses; publication match: exact; verification: read the result; development: none; confidence: high}
\end{GSitem}
\begin{GSitem}{1966-94.2}{E}{Hyperbolic toral automorphisms}
\GSfield{Kind; source}{construction; \GSpage{P000142}}
\GSfield{The notes}{A linear torus automorphism with no unit-modulus eigenvalues is the first example.}
\GSfield{Statement}{If \(A\in\mathrm{GL}(d,\mathbb Z)\) has no eigenvalue of modulus one, \(x+\mathbb Z^d\mapsto Ax+\mathbb Z^d\) is Anosov. Its stable and unstable bundles descend from the generalized eigenspaces.}
\GSfield{Fate}{The example is included in \citet[\S1]{smale1967dynamical}, and was already discussed in the earlier index at 1964-33.4.}
\GSfield{Present status}{As of 9 November 2026: established and explicit.}
\GSfield{Assessment}{literal statement: correct; publication match: exact; verification: read the result; development: none; confidence: high}
\end{GSitem}
\begin{GSitem}{1966-94.3}{E}{An arithmetic six-dimensional nilmanifold example}
\GSfield{Kind; source}{construction; \GSpage{P000142-P000143}}
\GSfield{The notes}{Two real upper-unitriangular three-by-three groups are combined using the two embeddings of \(\mathbb Z[\sqrt3]\); the quotient is compact. No actual automorphism formula is retained.}
\GSfield{Statement}{Let \(H\) have coordinates \((a,b,c)\) and law \((a,b,c)(a_2,b_2,c_2)=(a+a_2,b+b_2+ac_2,c+c_2)\). Embed \(H(\mathbb Z[\sqrt3])\) diagonally in \(H(\mathbb R)^2\) by conjugate embeddings. This is a uniform lattice. Multiplication of \(a,c\) by \(u=2+\sqrt3\) and of \(b\) by \(u^2\), with the conjugate action in the second factor, descends to an Anosov automorphism.}
\GSfield{Reconstruction}{\ours\ The chosen map respects the cross-term because \((ua)(uc_2)=u^2ac_2\), preserves the integral lattice since \(u\) is a unit, and has eigenvalues \(u,u,u^2,u^{-1},u^{-1},u^{-2}\). Compactness follows by reduction in the two-dimensional Minkowski lattice successively in \(a,c,b\).}
\GSfield{Fate}{The nilmanifold family is part of the algebraic examples reviewed by \citet[\S1]{smale1967dynamical}; the formula supplied here is a compatible reconstruction, not a claimed verbatim match.}
\GSfield{Present status}{As of 9 November 2026: established construction; the scan determines the lattice more clearly than the missing map.}
\GSfield{Assessment}{literal statement: correct with missing hypotheses; publication match: partial; verification: read the result; development: annotated; confidence: high}
\end{GSitem}
\begin{GSitem}{1966-94.4}{C}{Are all Anosov examples algebraic?}
\GSfield{Kind; source}{conjecture; \GSpage{P000143}}
\GSfield{The notes}{After the nilmanifold construction the notes suggest that all such systems may be of this kind.}
\GSfield{Statement}{A precise modern version asks whether every Anosov diffeomorphism of a closed connected manifold is topologically conjugate to an affine Anosov automorphism of an infranilmanifold. The weaker manifold question asks only whether the manifold is homeomorphic to an infranilmanifold.}
\GSfield{Reconstruction}{Topological conjugacy and infranilmanifolds are modern precisifications, not words legible in the notebook.}
\GSfield{Fate}{The strong topological-conjugacy formulation is explicit in \citet[Conjecture 2.3]{spatzier2016rigidity}; \citet[Introduction]{deremainkar2023anosov} retain the weaker manifold question. The latter circulated in 2020, appeared online in 2022, and in print in 2023.}
\GSfield{Present status}{As of 9 November 2026: no general resolution was found in the dated primary-source search. Smooth conjugacy and diffeomorphism to a standard smooth infranilmanifold are stronger, generally invalid replacements.}
\GSfield{Assessment}{literal statement: correct with missing hypotheses; publication match: partial; verification: read the result; development: none; confidence: medium; problem \hyperref[prob:GP-1966-01]{GP-1966-01}}
\GSfield{Worth developing}{Further work is worthwhile in a specified algebraic or dimension-restricted class; the unrestricted classification has no plausible bounded proof from these notes.}
\end{GSitem}
\begin{GSitem}{1966-94.5}{E}{Structural stability and the choice of equivalence}
\GSfield{Kind; source}{definition; \GSpage{P000143-P000144}}
\GSfield{The notes}{Nearby systems should be equivalent, initially described as diffeomorphic and subsequently as related by an \(\epsilon\)-homeomorphism.}
\GSfield{Statement}{A diffeomorphism is \(C^r\) structurally stable if all sufficiently \(C^r\)-close diffeomorphisms are topologically conjugate. For vector fields use an orbit homeomorphism preserving time orientation, usually allowing reparameterization. A near-identity conjugacy is available in the hyperbolic stability setting.}
\GSfield{Reconstruction}{The precise topology and the flow time-change convention are supplied.}
\GSfield{Fate}{\citet[Introduction and \S2]{smale1967dynamical} uses topological equivalence. The same distinction is relevant to 1965-74.1.}
\GSfield{Present status}{As of 9 November 2026: standard; smooth conjugacy would wrongly preserve eigenvalue data under perturbation.}
\GSfield{Assessment}{literal statement: false as written; publication match: exact; verification: read the result; development: annotated; confidence: high}
\end{GSitem}
\begin{GSitem}{1966-94.6}{E}{The two-dimensional structural-stability model}
\GSfield{Kind; source}{theorem; \GSpage{P000143}}
\GSfield{The notes}{Andronov--Pontryagin theory is cited as the two-dimensional precedent.}
\GSfield{Statement}{For flows on a closed orientable surface, the classical criterion requires finitely many hyperbolic equilibria and periodic orbits, no saddle connections, and no other nonwandering recurrence. Such fields are structurally stable; the relevant density theorem is a \(C^1\) surface theorem.}
\GSfield{Reconstruction}{The compact-surface formulation combines the precedent with the standard complete criterion; it is not attributed entirely to the original planar theorem.}
\GSfield{Fate}{The historical surface theory and Peixoto’s extension are reviewed by \citet[\S2]{smale1967dynamical}.}
\GSfield{Present status}{As of 9 November 2026: established. The notes do not justify transferring the density conclusion to higher dimensions.}
\GSfield{Assessment}{literal statement: correct with missing hypotheses; publication match: partial; verification: read the result; development: none; confidence: high}
\end{GSitem}
\begin{GSitem}{1966-94.7}{E}{Morse--Smale systems}
\GSfield{Kind; source}{definition; \GSpage{P000143-P000144}}
\GSfield{The notes}{Finitely many hyperbolic stationary and periodic trajectories, convergence toward them, and transverse stable/unstable intersections are listed.}
\GSfield{Statement}{For a smooth flow on a closed manifold, the nonwandering set must be a finite union of hyperbolic equilibria and periodic orbits, and their stable and unstable manifolds must meet transversely. For diffeomorphisms the nonwandering set consists of finitely many hyperbolic periodic points with the same transversality condition.}
\GSfield{Reconstruction}{The nonwandering-set condition makes the informal convergence clause precise.}
\GSfield{Fate}{The finite-hyperbolic/transverse formulation is the setting of \citet[Theorem 1.2 and \S5]{palissmale1970stability}, circulated in 1968.}
\GSfield{Present status}{As of 9 November 2026: standard; this definition is much narrower than Axiom A.}
\GSfield{Assessment}{literal statement: correct with missing hypotheses; publication match: exact; verification: read the result; development: none; confidence: high}
\end{GSitem}
\begin{GSitem}{1966-94.8}{B}{Does Morse--Smale imply structural stability?}
\GSfield{Kind; source}{question; \GSpage{P000143}}
\GSfield{The notes}{The heading says four assertions. Only the first, that Morse--Smale conditions imply stability, is explicitly marked as a question.}
\GSfield{Statement}{Is every Morse--Smale diffeomorphism or flow on a closed manifold structurally stable in the \(C^1\) topology?}
\GSfield{Fate}{Yes: \citet[Theorem 1.2; flow version in \S5]{palissmale1970stability}. The paper circulated in 1968, appeared in Russian in 1969, and in the cited proceedings in 1970.}
\GSfield{Present status}{As of 9 November 2026: solved affirmatively.}
\GSfield{Assessment}{literal statement: correct; publication match: exact; verification: read the result; development: none; confidence: high; problem \hyperref[prob:GP-1966-02]{GP-1966-02}}
\end{GSitem}
\begin{GSitem}{1966-94.9}{E}{Must a structurally stable system be Morse--Smale?}
\GSfield{Kind; source}{question; \GSpage{P000143}}
\GSfield{The notes}{The second assertion reverses the implication. The notebook groups assertions 2--4 under \ru{Опровергнуты Смейлом}: refuted by Smale.}
\GSfield{Statement}{Does structural stability force only finitely many hyperbolic recurrent critical elements?}
\GSfield{Fate}{The Anosov examples already refute this: see \citet[\S1]{smale1967dynamical}. A hyperbolic toral automorphism is structurally stable and has infinitely many periodic orbits; its suspension gives the flow analogue.}
\GSfield{Present status}{As of 9 November 2026: false in general, with examples available in the theory discussed at the talk.}
\GSfield{Assessment}{literal statement: correct; publication match: exact; verification: read the result; development: none; confidence: high; problem \hyperref[prob:GP-1966-03]{GP-1966-03}}
\end{GSitem}
\begin{GSitem}{1966-94.10}{E}{Are Morse--Smale systems dense?}
\GSfield{Kind; source}{question; \GSpage{P000143}}
\GSfield{The notes}{The third assertion proposes density of Morse--Smale systems; the notebook explicitly includes it among the three assertions refuted by Smale.}
\GSfield{Statement}{Are Morse--Smale systems dense among all smooth systems on each closed manifold in a specified \(C^r\) topology?}
\GSfield{Fate}{An open Anosov neighbourhood with infinitely many periodic orbits prevents this general density statement \citep[\S1]{smale1967dynamical}.}
\GSfield{Present status}{As of 9 November 2026: false in general; manifold- and dimension-specific density results do not repair the universal assertion.}
\GSfield{Assessment}{literal statement: correct; publication match: exact; verification: read the result; development: none; confidence: high; problem \hyperref[prob:GP-1966-04]{GP-1966-04}}
\end{GSitem}
\begin{GSitem}{1966-94.11}{E}{Are structurally stable systems dense?}
\GSfield{Kind; source}{question; \GSpage{P000143}}
\GSfield{The notes}{The fourth assertion proposes density of structurally stable systems; the notebook already reports Smale’s refutation.}
\GSfield{Statement}{Is the set of structurally stable \(C^r\) vector fields dense on every compact smooth manifold?}
\GSfield{Fate}{\citet[Main theorem]{smale1966nondensity}, received in June 1965 and published in 1966, constructs a compact four-manifold with an open set containing no structurally stable fields.}
\GSfield{Present status}{As of 9 November 2026: false in general. The negative result was already published and reported at this seminar.}
\GSfield{Assessment}{literal statement: correct; publication match: exact; verification: read the result; development: none; confidence: high; problem \hyperref[prob:GP-1966-05]{GP-1966-05}}
\end{GSitem}
\begin{GSitem}{1966-94.12}{E}{A hyperbolic invariant set}
\GSfield{Kind; source}{definition; \GSpage{P000144}}
\GSfield{The notes}{Uniform hyperbolicity is imposed on a closed invariant set, with a point, the whole space, and a Cantor-type set as examples.}
\GSfield{Statement}{For a diffeomorphism, a compact invariant set \(\Lambda\) is hyperbolic when \(T_\Lambda M\) admits a continuous invariant stable/unstable splitting with uniform exponential estimates. For nonsingular flows add the one-dimensional flow direction. A horseshoe illustrates a proper, totally disconnected hyperbolic set.}
\GSfield{Reconstruction}{The distinction between maps and flows is supplied.}
\GSfield{Fate}{These examples and the invariant-set formulation appear in \citet[\S\S1--2]{smale1967dynamical}.}
\GSfield{Present status}{As of 9 November 2026: standard; the flow-direction summand cannot be omitted from the flow definition.}
\GSfield{Assessment}{literal statement: correct with missing hypotheses; publication match: exact; verification: read the result; development: none; confidence: high}
\end{GSitem}
\begin{GSitem}{1966-94.13}{E}{Nonwandering points and Axiom A}
\GSfield{Kind; source}{definition; \GSpage{P000144-P000145}}
\GSfield{The notes}{The nonwandering set \(\Omega\) is closed; “Smale systems” satisfy hyperbolicity of \(\Omega\) and density in it of periodic points.}
\GSfield{Statement}{For a homeomorphism of a compact space, \(x\in\Omega(f)\) if every neighbourhood \(U\) meets \(f^nU\) for some \(n>0\). A \(C^1\) diffeomorphism satisfies Axiom A when \(\Omega(f)\) is hyperbolic and \(\overline{\operatorname{Per}(f)}=\Omega(f)\).}
\GSfield{Reconstruction}{The modern name Axiom A is used retrospectively.}
\GSfield{Fate}{The two independent clauses are those used by \citet[\S\S2--3]{smale1967dynamical}.}
\GSfield{Present status}{As of 9 November 2026: standard. On a flow, equilibria must also be included among the recurrent critical elements.}
\GSfield{Assessment}{literal statement: correct with missing hypotheses; publication match: exact; verification: read the result; development: none; confidence: high}
\end{GSitem}
\begin{GSitem}{1966-94.14}{E}{Nonwandering does not imply periodic}
\GSfield{Kind; source}{remark; \GSpage{P000145}}
\GSfield{The notes}{The discussion mentions a system with every point nonwandering but no periodic trajectories.}
\GSfield{Statement}{An irrational rotation \(x\mapsto x+\alpha\) of \(\mathbb R/\mathbb Z\), \(\alpha\notin\mathbb Q\), has \(\Omega=S^1\) and no periodic point. This example is not hyperbolic.}
\GSfield{Reconstruction}{The irrational rotation is our explicit model for the brief example.}
\GSfield{Fate}{This elementary distinction is consistent with the separate hypotheses in \citet[\S2]{smale1967dynamical}.}
\GSfield{Present status}{As of 9 November 2026: established counterexample to a claim without hyperbolicity.}
\GSfield{Assessment}{literal statement: correct; publication match: partial; verification: read the result; development: none; confidence: high}
\end{GSitem}
\begin{GSitem}{1966-94.15}{B}{Generic density of periodic points in the nonwandering set}
\GSfield{Kind; source}{theorem; \GSpage{P000145}}
\GSfield{The notes}{A generic periodic-point density theorem is attributed to a difficult-to-read name with a 1964 date.}
\GSfield{Statement}{For a residual subset of \(\operatorname{Diff}^1(M)\) of a closed smooth manifold, \(\Omega(f)=\overline{\operatorname{Per}(f)}\). The analogous generic flow statement includes equilibria. This is a \(C^1\), not an unrestricted \(C^r\), conclusion.}
\GSfield{Fate}{\citet[Theorem 6.5 in the Russian version]{pugh1967density}, published in 1967, states the generic flow result; the diffeomorphism formulation above is its standard counterpart. The paper refers to earlier announcements, but the handwritten name/date remains insecure.}
\GSfield{Present status}{As of 9 November 2026: established in \(C^1\). The ambiguous marginal attribution is retained as ambiguous.}
\GSfield{Assessment}{literal statement: correct with missing hypotheses; publication match: partial; verification: read the result; development: none; confidence: medium}
\end{GSitem}
\begin{GSitem}{1966-94.16}{B}{Density of Axiom A systems}
\GSfield{Kind; source}{conjecture; \GSpage{P000145}}
\GSfield{The notes}{A conjecture asks whether systems satisfying the two Smale conditions are dense.}
\GSfield{Statement}{Are Axiom A diffeomorphisms \(C^1\)-dense on every closed smooth manifold?}
\GSfield{Fate}{The general conjecture failed through the robust examples of \citet{abrahamsmale1970nongenericity}; this is a different assertion from Smale’s earlier nondensity of structural stability.}
\GSfield{Present status}{As of 9 November 2026: false in general. Density in restricted classes is a separate question.}
\GSfield{Assessment}{literal statement: correct; publication match: exact; verification: publisher metadata checked; development: none; confidence: medium; problem \hyperref[prob:GP-1966-06]{GP-1966-06}}
\end{GSitem}
\begin{GSitem}{1966-94.17}{B}{Does hyperbolicity of the nonwandering set imply dense periodic points?}
\GSfield{Kind; source}{conjecture; \GSpage{P000145}}
\GSfield{The notes}{The possibility that condition 1 implies condition 2 is assigned “70\%” confidence.}
\GSfield{Statement}{If the nonwandering set of a smooth diffeomorphism of a compact manifold is hyperbolic, must periodic points be dense in that set?}
\GSfield{Fate}{No: \citet[main counterexample]{kurata1979hyperbolic} constructs a hyperbolic nonwandering set without dense periodic points.}
\GSfield{Present status}{As of 9 November 2026: refuted in general. Local maximality and shadowing cannot be silently supplied to the arbitrary nonwandering set.}
\GSfield{Assessment}{literal statement: correct; publication match: exact; verification: read the result; development: none; confidence: high; problem \hyperref[prob:GP-1966-07]{GP-1966-07}}
\end{GSitem}
\begin{GSitem}{1966-94.18}{B}{Spectral decomposition into basic sets}
\GSfield{Kind; source}{theorem; \GSpage{P000145}}
\GSfield{The notes}{The nonwandering set decomposes into finitely many invariant transitive pieces.}
\GSfield{Statement}{For an Axiom A diffeomorphism on a closed manifold, \(\Omega(f)\) is a disjoint union of finitely many compact invariant topologically transitive basic sets. Periodic density belongs among the hypotheses, even if a preceding speculative implication would have removed it.}
\GSfield{Reconstruction}{Axiom A is stated explicitly rather than inferred from condition 1.}
\GSfield{Fate}{The spectral decomposition theorem is given in \citet[\S3]{smale1967dynamical}, published in 1967. The notebook attests its circulation here in September 1966; an independently checked earlier public-circulation date for this exact statement is not asserted.}
\GSfield{Present status}{As of 9 November 2026: established under Axiom A. The two-clause assumption matters because the preceding conjecture is false.}
\GSfield{Assessment}{literal statement: correct with missing hypotheses; publication match: exact; verification: read the result; development: none; confidence: high}
\end{GSitem}

\GStalk{95}{1966-09-12}{G. E. Shilov}{P000145-P000148}{Regular Hilbert-space functions and a first-order mean operator}
P000146--P000147 introduce a uniform algebra between weak and norm continuity, limiting finite-dimensional spherical means, and a mean-defined differential operator. Linear boundary problems and a fixed-point approach to a semilinear equation follow. Shilov’s 1967 papers closely match this announcement, but the operator is not the usual Hilbert-space trace Laplacian.

\begin{GSitem}{1966-95.1}{A}{The regular uniform algebra on a Hilbert ball}
\GSfield{Kind; source}{definition; \GSpage{P000146}}
\GSfield{The notes}{For the unit ball \(S\subset H\), \(L(S)\subset\overline R\subset C(S)\), where \(R\) consists of continuous functions of finitely many coordinates and distances \(\|x-a_j\|\).}
\GSfield{Statement}{Let \(H\) be an infinite-dimensional separable real Hilbert space. In the supremum norm on its closed unit ball, let \(A\) be the uniform closure of the unital algebra generated by the linear coordinates and \(\|x\|^2\). This equals the stated coordinate-and-distance algebra. One has \(C(S_{\rm weak})\subset A\subset C_b(S_{\rm norm})\).}
\GSfield{Reconstruction}{The identity \(\|x-a\|^2=\|x\|^2-2\langle x,a\rangle+\|a\|^2\) explains the equivalent generators. The subscript \(b\) is an explicit clarification.}
\GSfield{Fate}{The regular algebra is defined in \citet[\S1]{shilov1967hilbert1}, received 4 January 1967.}
\GSfield{Present status}{As of 9 November 2026: established. The ambient norm-continuous algebra must be bounded if a uniform Banach-algebra norm is intended.}
\GSfield{Assessment}{literal statement: correct with missing hypotheses; publication match: exact; verification: read the result; development: none; confidence: high}
\end{GSitem}
\begin{GSitem}{1966-95.2}{A}{A compactification and its nonevaluation characters}
\GSfield{Kind; source}{construction; \GSpage{P000146}}
\GSfield{The notes}{The maximal-ideal space of \(L(S)\) is the Hilbert ball; that of \(\overline R\) is a compact extension.}
\GSfield{Statement}{The character space of \(C(S_{\rm weak})\) is the weakly compact ball. For \(A\), an explicit model is \(K=\{(z,s):\|z\|^2\leq s\leq1\}\subset S_{\rm weak}\times[0,1]\), with evaluation points \((x,\|x\|^2)\). Points with \(s>\|z\|^2\) are additional characters.}
\GSfield{Reconstruction}{\ours\ The set \(K\) is compact because norm squared is weakly lower semicontinuous. The generators extend continuously and separate its points. Evaluation points are dense: use \(x_j=z+\sqrt{s-\|z\|^2}\,e_j\), with orthonormal \(e_j\perp z\). Stone--Weierstrass then identifies \(A\) with \(C(K)\).}
\GSfield{Fate}{\citet[\S1 and the reference to Part II]{shilov1967hilbert1} connects the regular algebra to its maximal ideals. The concrete model here is supplied as a reconstruction.}
\GSfield{Present status}{As of 9 November 2026: the compactification is well defined; it is not the norm ball itself.}
\GSfield{Assessment}{literal statement: correct with missing hypotheses; publication match: partial; verification: read the result; development: annotated; confidence: high}
\end{GSitem}
\begin{GSitem}{1966-95.3}{A}{Limiting spherical and ball means coincide}
\GSfield{Kind; source}{theorem; \GSpage{P000146}}
\GSfield{The notes}{Finite-dimensional means on spheres and balls have the same infinite-dimensional limit on \(\overline R\).}
\GSfield{Statement}{For admissible centres \(z\) and radii \(r\), normalized means over increasing finite-dimensional balls or spheres converge for every \(f\in A\), independently of the orthonormal exhaustion. On a generator \(F(\langle e_1,x\rangle,\ldots,\langle e_m,x\rangle,\|x\|^2)\), the limit is \(F(\langle e_1,z\rangle,\ldots,\langle e_m,z\rangle,\|z\|^2+r^2)\).}
\GSfield{Reconstruction}{The displayed formula spells out the finite-coordinate concentration calculation; admissibility means the averaging sets lie in the domain.}
\GSfield{Fate}{The mean construction is in \citet[\S2]{shilov1967hilbert1}.}
\GSfield{Present status}{As of 9 November 2026: established on this algebra; it is not a countably additive uniform measure on an infinite-dimensional sphere.}
\GSfield{Assessment}{literal statement: correct with missing hypotheses; publication match: exact; verification: read the result; development: none; confidence: high}
\end{GSitem}
\begin{GSitem}{1966-95.4}{A}{The mean is multiplicative}
\GSfield{Kind; source}{theorem; \GSpage{P000146}}
\GSfield{The notes}{\(M_r\) is a multiplicative functional, usually evaluation only in the compact extension.}
\GSfield{Statement}{For \(f,g\in A\), \(M_r(fg,z)=M_r(f,z)M_r(g,z)\). For \(r>0\) this is evaluation at \((z,\|z\|^2+r^2)\in K\), not at a Hilbert-space point.}
\GSfield{Fate}{\citet[\S2]{shilov1967hilbert1} establishes multiplicativity.}
\GSfield{Present status}{As of 9 November 2026: established; finite-dimensional averaging is not multiplicative, and the limiting assertion relies on the restricted regular algebra.}
\GSfield{Assessment}{literal statement: correct; publication match: exact; verification: read the result; development: none; confidence: high}
\end{GSitem}
\begin{GSitem}{1966-95.5}{A}{The mean-defined first-order derivation}
\GSfield{Kind; source}{definition; \GSpage{P000146}}
\GSfield{The notes}{\(\Delta f(z)=\lim_{r\to0}[M_r(f,z)-f(z)]/r^2\), with the emphatic comment that it is first order and obeys Leibniz’s rule.}
\GSfield{Statement}{On the domain where this limit exists in the stipulated topology, the mean generator \(L\) is a derivation: \(L(fg)=fLg+gLf\). On smooth regular functions expressed in \((z,s)\), it differentiates in the \(s\)-direction at \(s=\|z\|^2\). In particular \(L\|x\|^2=1\), while it annihilates a smooth finite-coordinate function.}
\GSfield{Reconstruction}{The notation \(L\) avoids confusion with a trace Laplacian. The examples follow directly from the notebook’s normalization.}
\GSfield{Fate}{The operator and its derivation property are in \citet[\S\S2,5]{shilov1967hilbert1}.}
\GSfield{Present status}{As of 9 November 2026: established for its specified function class and closure. It should not be identified with the ordinary sum of second derivatives.}
\GSfield{Assessment}{literal statement: correct with missing hypotheses; publication match: exact; verification: read the result; development: annotated; confidence: high}
\end{GSitem}
\begin{GSitem}{1966-95.6}{A}{A regular Dirichlet problem}
\GSfield{Kind; source}{theorem; \GSpage{P000146-P000147}}
\GSfield{The notes}{For a star-shaped surface \(\sum_{k>m}\xi_k^2=\varphi(\xi_1,\ldots,\xi_m)\) with bounded curvature, the Dirichlet problem has a unique, continuously dependent solution given by an angular mean.}
\GSfield{Statement}{In Shilov’s regular boundary class, the mean-defined operator has a unique regular harmonic extension of regular continuous boundary data, with continuous dependence in the uniform norm. An unqualified assertion for every weakly closed star-shaped surface is not supported. Even the norm unit sphere is not weakly closed in infinite dimension.}
\GSfield{Reconstruction}{The scope is restricted to the published regular-domain construction instead of guessing a general geometric theorem from “bounded curvature.”}
\GSfield{Fate}{The linear boundary construction and uniqueness are in \citet[\S\S3,6]{shilov1967hilbert1}.}
\GSfield{Present status}{As of 9 November 2026: established in the published regular-domain framework; the notebook does not retain all the domain hypotheses.}
\GSfield{Assessment}{literal statement: correct with missing hypotheses; publication match: partial; verification: read the result; development: none; confidence: medium}
\end{GSitem}
\begin{GSitem}{1966-95.7}{A}{A Poisson inverse for the regular operator}
\GSfield{Kind; source}{construction; \GSpage{P000147}}
\GSfield{The notes}{The inhomogeneous problem \(\Delta u=g\in\overline R\) also has an explicit formula.}
\GSfield{Statement}{For the same regular boundary problem, separate a harmonic extension of the boundary data and a zero-boundary solution operator \(P\) satisfying \(LPg=g\) on the allowed regular class. Boundedness of \(P\), with the relevant domain norm, is needed for continuous dependence.}
\GSfield{Reconstruction}{The operator notation \(P\) is introduced for the next fixed-point step.}
\GSfield{Fate}{\citet[\S\S4--5]{shilov1967hilbert1} gives the inhomogeneous construction.}
\GSfield{Present status}{As of 9 November 2026: established within that framework.}
\GSfield{Assessment}{literal statement: correct with missing hypotheses; publication match: exact; verification: read the result; development: none; confidence: high}
\end{GSitem}
\begin{GSitem}{1966-95.8}{A}{Semilinear equations by a fixed point}
\GSfield{Kind; source}{method; \GSpage{P000147}}
\GSfield{The notes}{\(\Delta u=\varphi(x,u(x))\) is said to be solved by a fixed-point method.}
\GSfield{Statement}{A sufficient abstract version is: on a closed ball of the regular Banach function space, \(T(u)=h+P\varphi(\cdot,u)\) preserves the ball and \(\varphi\) is uniformly Lipschitz in \(u\) with \(\|P\|\operatorname{Lip}_u\varphi<1\). Then \(T\) has a unique fixed point, solving the specified boundary problem.}
\GSfield{Reconstruction}{The ball-invariance and contraction assumptions are our transparent sufficient version, not a claim that they reproduce the strongest published theorem.}
\GSfield{Fate}{The nonlinear programme was published in \citet[Theorems 1--2]{shilov1967hilbert3}, received 16 February 1967.}
\GSfield{Present status}{As of 9 November 2026: established under appropriate existence hypotheses. The bare phrase does not prove global solvability for arbitrary nonlinearities.}
\GSfield{Assessment}{literal statement: correct with missing hypotheses; publication match: partial; verification: read the result; development: annotated; confidence: high}
\end{GSitem}

\GStalk{96}{1966-09-12}{I. M. Gelfand}{P000146-P000149}{Uniqueness, admissibility, and viscosity for hyperbolic systems}
The middle and lower parts of P000147 ask when nonlinear first-order systems have unique solutions. Scalar convexity, characteristic eigenvector curves, conservation form, viscosity, and travelling waves are juxtaposed. The equations and solution classes in the hydrodynamic remarks are too incomplete for an exact publication match.

\begin{GSitem}{1966-96.1}{E}{Hyperbolicity requires more than a real spectrum}
\GSfield{Kind; source}{definition; \GSpage{P000147}}
\GSfield{The notes}{\(u_t=A(u)u_x\) is called hyperbolic when all eigenvalues of \(A(u)\) are real.}
\GSfield{Statement}{For a standard one-dimensional strictly hyperbolic quasilinear system, \(A\) is smooth on an open state set and has \(n\) distinct real eigenvalues there. For non-strict notions, diagonalizability and uniform control of eigenvectors must be specified; a real spectrum alone allows defective matrices.}
\GSfield{Reconstruction}{Strict hyperbolicity is a sufficient precise scope, not an assertion that every physical example has distinct speeds.}
\GSfield{Fate}{The modern strict-hyperbolicity setting of \citet[Theorem 1]{bianchinibressan2005viscosity} makes the assumptions needed for the later uniqueness question explicit.}
\GSfield{Present status}{As of 9 November 2026: the notebook’s definition is incomplete for standard well-posedness theorems.}
\GSfield{Assessment}{literal statement: correct with missing hypotheses; publication match: partial; verification: read the result; development: none; confidence: high}
\end{GSitem}
\begin{GSitem}{1966-96.2}{E}{Scalar convexity does not select a unique weak solution}
\GSfield{Kind; source}{claim; \GSpage{P000147}}
\GSfield{The notes}{For \(n=1\), uniqueness is said always to hold when \(f\) is convex.}
\GSfield{Statement}{For \(u_t+F(u)_x=0\) with smooth scalar flux and bounded initial data, entropy admissibility yields the standard unique solution. Convexity alone does not make arbitrary distributional solutions unique. For Burgers flux \(F(u)=u^2/2\) and data 0 left, 1 right, both an expansion shock of speed \(1/2\) and the rarefaction fan are weak solutions.}
\GSfield{Reconstruction}{\ours\ The shock satisfies Rankine--Hugoniot since \((1/2-0)/(1-0)=1/2\). The fan is \(u=0,x\leq0\), \(u=x/t,0<x<t\), \(u=1,x\geq t\); it solves the equation in each region and is continuous across its interfaces. Their common initial trace proves nonuniqueness without admissibility.}
\GSfield{Fate}{Entropy selection was present in \citet{oleinik1957discontinuous}; the general scalar uniqueness framework is \citet{kruzhkov1970first}.}
\GSfield{Present status}{As of 9 November 2026: the literal weak-solution claim is false; the entropy formulation is established.}
\GSfield{Assessment}{literal statement: false as written; publication match: partial; verification: read the result; development: annotated; confidence: high}
\end{GSitem}
\begin{GSitem}{1966-96.3}{F}{Hydrodynamics, magnetohydrodynamics, and an unspecified Godunov class}
\GSfield{Kind; source}{remark; \GSpage{P000147}}
\GSfield{The notes}{Two-component hydrodynamics and apparently magnetohydrodynamics are said to have uniqueness; a narrow class treated by Godunov is mentioned.}
\GSfield{Statement}{No single uniqueness theorem can be recovered without the exact equations, state restrictions, domain, data, and admissibility criterion. These fragments do not establish uniqueness for every Euler or magnetohydrodynamic system.}
\GSfield{Fate}{The checked scalar and strictly hyperbolic literature does not identify these particular fragments: \citet{oleinik1957discontinuous,bianchinibressan2005viscosity} are comparisons only.}
\GSfield{Present status}{As of 9 November 2026: unidentified at theorem level; not evidence of an undeveloped problem.}
\GSfield{Assessment}{literal statement: not determinable; publication match: thematic only; verification: unresolved lead; development: none; confidence: low}
\end{GSitem}
\begin{GSitem}{1966-96.4}{E}{Convexity along an eigenvector field}
\GSfield{Kind; source}{definition; \GSpage{P000147}}
\GSfield{The notes}{Draw eigenvectors in state space; convexity may mean growth of the corresponding eigenvalue along their integral curves.}
\GSfield{Statement}{If \(A(u)r_i(u)=\lambda_i(u)r_i(u)\) is a smooth simple eigenpair, the genuine-nonlinearity condition is \(D\lambda_i(u)\cdot r_i(u)\ne0\). Linear degeneracy means that derivative vanishes identically. These are properties of characteristic fields in state space, distinct from the space-time characteristics of one solution.}
\GSfield{Reconstruction}{The derivative condition and distinction of the two kinds of curves are supplied.}
\GSfield{Fate}{This is the classical characteristic-field distinction from Lax’s 1957 theory, discussed in the historical setup of \citet{bianchinibressan2005viscosity}.}
\GSfield{Present status}{As of 9 November 2026: established; the notebook gives the correct geometric idea without specifying nonvanishing or normalization.}
\GSfield{Assessment}{literal statement: correct with missing hypotheses; publication match: partial; verification: read the result; development: none; confidence: high}
\end{GSitem}
\begin{GSitem}{1966-96.5}{E}{Conservation form and the meaning of a shock}
\GSfield{Kind; source}{method; \GSpage{P000147}}
\GSfield{The notes}{Writing \(u_t=\partial_x f(u)\) is said to be better than general quasilinear form.}
\GSfield{Statement}{For a conservative system \(u_t+F(u)_x=0\), weak solutions and jump conditions are defined distributionally. For a general nonconservative matrix \(A(u)\), \(A(u)u_x\) is not canonically defined at jumps; a regularization or prescribed nonconservative product is additional data.}
\GSfield{Reconstruction}{The distributional explanation is supplied.}
\GSfield{Fate}{The conservative/nonconservative distinction is explicit in \citet[Introduction and Theorem 1]{bianchinibressan2005viscosity}.}
\GSfield{Present status}{As of 9 November 2026: standard; the notebook’s plus-sign convention is converted by \(F=-f\).}
\GSfield{Assessment}{literal statement: correct with missing hypotheses; publication match: partial; verification: read the result; development: none; confidence: high}
\end{GSitem}
\begin{GSitem}{1966-96.6}{B}{Uniqueness selected by identity viscosity}
\GSfield{Kind; source}{question; \GSpage{P000147}}
\GSfield{The notes}{When does adding \(\epsilon u_{xx}\) give uniqueness?}
\GSfield{Statement}{A definite interpretation asks whether, for smooth strictly hyperbolic \(A\) and common initial data, solutions of \(u_t+A(u)u_x=\epsilon u_{xx}\) converge as \(\epsilon\downarrow0\) to one limit independent of subsequence. Small total variation, values in a fixed hyperbolicity neighbourhood, and the identity viscosity matrix specify an established scope.}
\GSfield{Reconstruction}{The vanishing-viscosity interpretation is distinguished from the easier uniqueness of a fixed strictly parabolic equation.}
\GSfield{Fate}{\citet[Theorem 1]{bianchinibressan2005viscosity} prove this small-BV statement, with uniform BV and Lipschitz dependence; the preprint circulated in November 2001 and publication followed in 2005.}
\GSfield{Present status}{As of 9 November 2026: solved in the stated small-BV setting. The unrestricted wording does not specify a single larger well-posedness problem and is not claimed solved.}
\GSfield{Assessment}{literal statement: correct with missing hypotheses; publication match: partial; verification: read the result; development: none; confidence: high; problem \hyperref[prob:GP-1966-08]{GP-1966-08}}
\end{GSitem}
\begin{GSitem}{1966-96.7}{E}{Travelling-wave reduction is local information}
\GSfield{Kind; source}{method; \GSpage{P000147}}
\GSfield{The notes}{The viscosity question is said to reduce to ordinary differential equations.}
\GSfield{Statement}{For a conservative viscous system, a travelling profile \(U((x-st)/\epsilon)\) satisfies \(U^{\prime}=F(U)-F(u_-)-s(U-u_-)\) after integration, with end states at equilibria. For the notebook sign convention in a nonconservative system, the second-order equation is \(U^{\prime\prime}=-(A(U)+sI)U^{\prime}\).}
\GSfield{Reconstruction}{\ours\ Both ODEs are derived here by substitution. They reconstruct a plausible local meaning of the last sentence, not the cancelled word.}
\GSfield{Fate}{Travelling profiles are part of the analysis in \citet{bianchinibressan2005viscosity}; a full Cauchy-problem uniqueness proof needs substantially more.}
\GSfield{Present status}{As of 9 November 2026: the profile reduction is valid; equivalence to global uniqueness of arbitrary solutions is false.}
\GSfield{Assessment}{literal statement: correct with missing hypotheses; publication match: partial; verification: read the result; development: annotated; confidence: high}
\end{GSitem}

\GStalk{97}{1966-09-12}{I. N. Bernstein}{P000147-P000150}{The representation-ring completion theorem and the circle}
P000148 explicitly reports part of what Atiyah had explained. The augmentation-ideal completion theorem is illustrated by \(S^1\) and \(\mathbb{CP}^{\infty}\). Two literal slips deserve separate treatment: calling the uncompleted map an embedding for every compact group, and giving the wrong fibre and total space for the tautological bundle.

\begin{GSitem}{1966-97.1}{E}{The representation ring and associated universal bundles}
\GSfield{Kind; source}{construction; \GSpage{P000148}}
\GSfield{The notes}{\(R_G\) is the commutative ring of finite-dimensional representations and their formal differences; a representation gives a bundle over \(B_G\).}
\GSfield{Statement}{For compact Lie \(G\), use finite-dimensional complex continuous representations. Direct sum and tensor product define \(R(G)\). The associated-bundle construction sends \([V]\) to \([EG\times_GV]\in K^0(BG)\) and is a ring homomorphism.}
\GSfield{Fate}{This construction underlies \citet[Introduction]{atiyahsegal1969completion}, whose introduction traces the finite-group and connected-group cases to 1961.}
\GSfield{Present status}{As of 9 November 2026: standard; complex coefficients and virtual bundles are made explicit.}
\GSfield{Assessment}{literal statement: correct with missing hypotheses; publication match: exact; verification: publisher metadata checked; development: none; confidence: high}
\end{GSitem}
\begin{GSitem}{1966-97.2}{E}{The uncompleted map need not be an embedding}
\GSfield{Kind; source}{claim; \GSpage{P000148}}
\GSfield{The notes}{The arrow \(R_G\to K(B_G)\) is labelled \ru{Вложение}.}
\GSfield{Statement}{The map need not be injective for disconnected compact groups. In \(R(C_6)=\mathbb Z[x]/(x^6-1)\), let \(a=(x-1)(x+1)(x^2+x+1)\ne0\). Then \((x^2-x+1)a=0\) and \((x-1)(-x)a=a\), so \(a\in\bigcap_{n\geq1}(x-1)^nR(C_6)\); its completion image vanishes.}
\GSfield{Reconstruction}{\ours\ The displayed polynomial has degree four and is nonzero modulo the monic degree-six relation. Since \(1+(x-1)x=x^2-x+1\), multiplication by \(-x\) is an inverse to \(x-1\) on the submodule generated by \(a\). Iteration proves membership in every augmentation power. The completion theorem then identifies its image in \(K^0(BC_6)\) with zero.}
\GSfield{Fate}{The correct general statement is completion, as in \citet{atiyahsegal1969completion}; the elementary cyclic-group calculation here checks the stronger notebook wording.}
\GSfield{Present status}{As of 9 November 2026: the literal embedding assertion is false in this generality.}
\GSfield{Assessment}{literal statement: false as written; publication match: partial; verification: publisher metadata checked; development: annotated; confidence: high}
\end{GSitem}
\begin{GSitem}{1966-97.3}{E}{Augmentation completion gives K of the classifying space}
\GSfield{Kind; source}{theorem; \GSpage{P000148}}
\GSfield{The notes}{Completing \(R_G\) at virtual representations of dimension zero gives \(K(B_G)\).}
\GSfield{Statement}{For compact Lie \(G\), \(I_G=\ker(\dim:R(G)\to\mathbb Z)\) and \(\widehat{R(G)}_{I_G}=\varprojlim R(G)/I_G^n\). The associated-bundle map induces \(\widehat{R(G)}_{I_G}\cong K^0(BG)\); also \(K^1(BG)=0\).}
\GSfield{Fate}{\citet[Introduction]{atiyahsegal1969completion} records the 1961 cases and supplies a unified 1969 proof; the seminar itself documents prior circulation of the general statement.}
\GSfield{Present status}{As of 9 November 2026: established. Class E describes Bernstein’s explicit exposition of Atiyah, rather than a claim of a 1966 Bernstein discovery.}
\GSfield{Assessment}{literal statement: correct; publication match: exact; verification: publisher metadata checked; development: none; confidence: high}
\end{GSitem}
\begin{GSitem}{1966-97.4}{E}{The circle completion is a formal power-series ring}
\GSfield{Kind; source}{construction; \GSpage{P000148}}
\GSfield{The notes}{\(R_{S^1}=\sum n_ix^i\), only finitely many nonzero coefficients; \(y=x-1\) generates the augmentation ideal and \(\widehat R_{S^1}=\mathbb Z[[y]]\).}
\GSfield{Statement}{\(R(S^1)=\mathbb Z[x,x^{-1}]\), where \(x\) is the weight-one character. Modulo \(y^n\), \(x=1+y\) is invertible; inverse limits therefore give \(\mathbb Z[[y]]\). The notebook’s geometric series \(1/(1-y)\) is an example of a formal series, not the expansion of \(x^{-1}\).}
\GSfield{Reconstruction}{The inverse-limit justification and distinction of the two inverses are supplied.}
\GSfield{Fate}{This is the standard circle example of the completion theorem \citep{atiyahsegal1969completion}; Atiyah’s course circulated in 1964 and appeared as \citet{atiyah1967ktheory}.}
\GSfield{Present status}{As of 9 November 2026: established elementary calculation.}
\GSfield{Assessment}{literal statement: correct; publication match: exact; verification: publisher metadata checked; development: none; confidence: high}
\end{GSitem}
\begin{GSitem}{1966-97.5}{E}{Projective space and its tautological line bundle}
\GSfield{Kind; source}{construction; \GSpage{P000148}}
\GSfield{The notes}{\(B_{S^1}=\mathbb{CP}^{\infty}\); on \(\mathbb{CP}^n\) the canonical bundle is described as having fibre \(\mathbb{CP}^1\) and total space \(\mathbb{CP}^{n+1}\).}
\GSfield{Statement}{The tautological complex line bundle \(L\to\mathbb{CP}^n\) has fibre the line \(\ell\cong\mathbb C\) above \([\ell]\). Its unit-circle bundle is \(S^{2n+1}\to\mathbb{CP}^n\), with fibre \(S^1\). Neither the line bundle nor this circle bundle has total space \(\mathbb{CP}^{n+1}\).}
\GSfield{Reconstruction}{The two actual bundles are distinguished to avoid guessing an unstated projectivization.}
\GSfield{Fate}{The standard construction is part of Atiyah’s 1964 course, published as \citet{atiyah1967ktheory}.}
\GSfield{Present status}{As of 9 November 2026: the model for \(BS^1\) is correct; the finite-dimensional bundle description is false as written.}
\GSfield{Assessment}{literal statement: false as written; publication match: partial; verification: publisher metadata checked; development: annotated; confidence: high}
\end{GSitem}
\begin{GSitem}{1966-97.6}{E}{K of finite and infinite complex projective space}
\GSfield{Kind; source}{theorem; \GSpage{P000148}}
\GSfield{The notes}{Every class on \(\mathbb{CP}^n\) is a polynomial in \(x-1\), and \((x-1)^{n+1}=0\); passage to the limit gives the infinite case.}
\GSfield{Statement}{For \(x=[L]\), \(K^0(\mathbb{CP}^n)=\mathbb Z[x]/((x-1)^{n+1})\), \(K^1(\mathbb{CP}^n)=0\), and the compatible inverse limit is \(K^0(\mathbb{CP}^{\infty})=\mathbb Z[[x-1]]\).}
\GSfield{Reconstruction}{The odd K group is supplied to justify the uncomplicated inverse limit.}
\GSfield{Fate}{This is the projective-space calculation in \citet{atiyah1967ktheory}, based on the 1964 lectures.}
\GSfield{Present status}{As of 9 November 2026: established. The exponent depends on the finite stage \(n\).}
\GSfield{Assessment}{literal statement: correct; publication match: exact; verification: publisher metadata checked; development: none; confidence: high}
\end{GSitem}

\GStalk{98}{1966-09-19}{I. N. Bernstein}{P000148-P000152}{Equivariant completion, induction, and a noncompact programme}
P000149--P000150 continue Atiyah’s programme: skeletal limits, equivariant vector bundles, the Borel construction, and reductions to a unitary group and maximal torus. Gelfand asks for a noncompact analogue and suggests a Plancherel replacement for dimension, but does not specify a representation ring or ideal. The family-index explanation is continued at record 104.

\begin{GSitem}{1966-98.1}{E}{Skeletal inverse limits}
\GSfield{Kind; source}{construction; \GSpage{P000149}}
\GSfield{The notes}{Take \(K((BG)_1)\leftarrow K((BG)_2)\leftarrow\cdots\) and its projective limit.}
\GSfield{Statement}{For this completion problem, skeletal K-groups determine the completed theory. For a general countable CW filtration \(X=\bigcup X_n\), however, the comparison is governed by \(0\to\varprojlim{}^1K^{-1}(X_n)\to K^0(X)\to\varprojlim K^0(X_n)\to0\). The derived-limit term must be shown to vanish or the pro-system formulation retained.}
\GSfield{Reconstruction}{The derived-limit caveat is supplied to prevent an overgeneralization.}
\GSfield{Fate}{The filtration/completion distinction is explicit in \citet{atiyahsegal1969completion}; the simple projective-space tower has vanishing odd groups.}
\GSfield{Present status}{As of 9 November 2026: valid in the appropriate completed interpretation, not an unconditional definition of topological K-theory for every infinite complex.}
\GSfield{Assessment}{literal statement: correct with missing hypotheses; publication match: partial; verification: publisher metadata checked; development: none; confidence: high}
\end{GSitem}
\begin{GSitem}{1966-98.2}{F}{Gelfand’s noncompact completion programme}
\GSfield{Kind; source}{question; \GSpage{P000149}; I. M. Gelfand}
\GSfield{The notes}{Generalize the K-functor and classifying space to locally compact Lie groups, using infinite-dimensional representations, so that the theorem remains valid.}
\GSfield{Statement}{This specifies a programme but not a single theorem: it leaves the representation category, Grothendieck construction, completion ideal, and space unspecified. One precise later analogue treats finite proper \(G\)-CW complexes for discrete \(G\); this does not settle all locally compact Lie groups in one canonical interpretation.}
\GSfield{Fate}{\citet[Theorems 3.2 and 4.3; \S5]{luckoliver2001completion} provide a proper-action analogue and explain difficulties for positive-dimensional noncompact groups.}
\GSfield{Present status}{As of 9 November 2026: partly realized through several inequivalent theories, with no exact notebook conjecture to classify as solved or still open.}
\GSfield{Assessment}{literal statement: not determinable; publication match: thematic only; verification: read the result; development: none; confidence: medium}
\end{GSitem}
\begin{GSitem}{1966-98.3}{F}{A Plancherel factor in place of dimension}
\GSfield{Kind; source}{remark; \GSpage{P000149}; I. M. Gelfand}
\GSfield{The notes}{Dimension is compared with an uncertain “factor” in the Plancherel formula.}
\GSfield{Statement}{For a noncompact group, an irreducible unitary representation may be infinite-dimensional; formal degree or Plancherel density can play an analogue of dimension in some formulas. The fragment does not specify a multiplication law, an augmentation homomorphism, or even a class of representations.}
\GSfield{Fate}{The noncompact K-theory comparison in \citet[\S5]{luckoliver2001completion} does not identify this particular suggestion. No exact theorem match was found.}
\GSfield{Present status}{As of 9 November 2026: unidentifiable beyond the heuristic analogy.}
\GSfield{Assessment}{literal statement: not determinable; publication match: none found; verification: unresolved lead; development: none; confidence: low}
\end{GSitem}
\begin{GSitem}{1966-98.4}{E}{Equivariant K-theory of vector bundles}
\GSfield{Kind; source}{definition; \GSpage{P000149}}
\GSfield{The notes}{A \(G\)-bundle is a vector bundle with compatible linear action on fibres; take sums, tensor products, and formal differences. At a point the answer is \(R_G\).}
\GSfield{Statement}{For compact Lie \(G\) acting on compact \(X\), \(K_G^0(X)\) is the Grothendieck group of finite-rank complex \(G\)-equivariant vector bundles. Fibre maps are linear and cover the given base action. Pullback from a point makes it an \(R(G)\)-module, and \(K_G^0(*)=R(G)\).}
\GSfield{Fate}{The construction was already being taught by Atiyah; \citet{atiyahsegal1969completion} uses the theory systematically.}
\GSfield{Present status}{As of 9 November 2026: standard for compact groups. Noncompact groups require additional qualifications.}
\GSfield{Assessment}{literal statement: correct; publication match: exact; verification: publisher metadata checked; development: none; confidence: high}
\end{GSitem}
\begin{GSitem}{1966-98.5}{E}{The Borel construction}
\GSfield{Kind; source}{construction; \GSpage{P000149-P000150}}
\GSfield{The notes}{\(X_G=E\times_GX\), with \(E\) a free right \(G\)-space over \(B_G\); for a point one recovers \(B_G\).}
\GSfield{Statement}{Choose a contractible free right \(G\)-space \(EG\). For a left \(G\)-space \(X\), set \(X_G=(EG\times X)/((eg,x)\sim(e,gx))\). An equivariant bundle induces a bundle over \(X_G\), producing \(K_G^*(X)\to K^*(X_G)\).}
\GSfield{Reconstruction}{Contractibility and the explicit quotient convention are supplied.}
\GSfield{Fate}{This is the comparison map of \citet{atiyahsegal1969completion}.}
\GSfield{Present status}{As of 9 November 2026: standard; contractibility is required of the universal free space.}
\GSfield{Assessment}{literal statement: correct with missing hypotheses; publication match: exact; verification: publisher metadata checked; development: none; confidence: high}
\end{GSitem}
\begin{GSitem}{1966-98.6}{B}{The equivariant completion theorem}
\GSfield{Kind; source}{theorem; \GSpage{P000150}}
\GSfield{The notes}{\(\widehat K_G(X)\cong K(X_G)\), with completion in the dimension-zero ideal of \(R_G\).}
\GSfield{Statement}{For compact Lie \(G\) and a finite \(G\)-CW complex \(X\), completion at \(I_G\) gives \(\widehat{K_G^*(X)}_{I_G}\cong K^*(EG\times_GX)\). A pro-system version handles the filtration explicitly; a statement for an arbitrary \(G\)-space without finiteness assumptions is too broad.}
\GSfield{Reconstruction}{The finite \(G\)-CW hypothesis is a safe standard formulation of the announced result.}
\GSfield{Fate}{\citet{atiyahsegal1969completion} is the exact general theorem publication. The notebook proves its circulation through Atiyah by September 1966.}
\GSfield{Present status}{As of 9 November 2026: established in the stated finite equivariant-CW setting.}
\GSfield{Assessment}{literal statement: correct with missing hypotheses; publication match: exact; verification: publisher metadata checked; development: none; confidence: high}
\end{GSitem}
\begin{GSitem}{1966-98.7}{E}{Completion by tensoring with the completed coefficient ring}
\GSfield{Kind; source}{claim; \GSpage{P000150}}
\GSfield{The notes}{\(\widehat K_G(X)=K_G(X)\otimes_{R_G}\widehat R_G\) is offered as an exact definition.}
\GSfield{Statement}{For a Noetherian ring \(R\), ideal \(I\), and finitely generated \(R\)-module \(M\), \(\widehat M_I\cong M\otimes_R\widehat R_I\). Here compact Lie \(G\) and a finite \(G\)-CW complex ensure the relevant finiteness. This tensor formula is not a universal definition for arbitrary modules.}
\GSfield{Reconstruction}{The commutative-algebra hypotheses are supplied.}
\GSfield{Fate}{This algebra is used in \citet{atiyahsegal1969completion}.}
\GSfield{Present status}{As of 9 November 2026: correct with finite-generation assumptions.}
\GSfield{Assessment}{literal statement: correct with missing hypotheses; publication match: partial; verification: publisher metadata checked; development: none; confidence: high}
\end{GSitem}
\begin{GSitem}{1966-98.8}{E}{Induction to a unitary group}
\GSfield{Kind; source}{construction; \GSpage{P000150}}
\GSfield{The notes}{For \(G\subset U\), \(K_G(X)=K_U(U\times_GX)\), and \(X_G=(U\times_GX)_U\).}
\GSfield{Statement}{For a closed inclusion of compact Lie groups \(G\subset U(N)\), induction gives an equivalence of equivariant vector-bundle categories over \(X\) and \(U(N)\times_GX\), hence the K-theory isomorphism. The corresponding Borel constructions are naturally equivalent using a common universal free \(U(N)\)-space.}
\GSfield{Reconstruction}{Closed inclusion and common universal models make the displayed identities precise.}
\GSfield{Fate}{This induction mechanism belongs to the completion proof in \citet{atiyahsegal1969completion}.}
\GSfield{Present status}{As of 9 November 2026: established; the equality of classifying constructions should be interpreted with the chosen models.}
\GSfield{Assessment}{literal statement: correct with missing hypotheses; publication match: exact; verification: publisher metadata checked; development: none; confidence: high}
\end{GSitem}
\begin{GSitem}{1966-98.9}{E}{Maximal-torus splitting by a family index}
\GSfield{Kind; source}{method; \GSpage{P000150}}
\GSfield{The notes}{For a maximal torus \(T\subset U\), restriction embeds a direct summand, with a dashed reverse map to be explained analytically.}
\GSfield{Statement}{For \(U=U(N)\), the flag manifold \(U/T\) is complex. The fibrewise Dolbeault index along \(U\times_TX\to X\) defines a retraction of restriction \(K_U(X)\to K_T(X)\): the projection formula and \(\operatorname{ind}_{U}(\bar\partial_{U/T})=1\) give \(p_!p^*=1\). The same mechanism compares the completed/Borel diagram.}
\GSfield{Reconstruction}{The Dolbeault operator is our precise realization of the dashed arrow; the notebook postpones its definition to the family-index discussion.}
\GSfield{Fate}{The direct-image strategy is discussed by Atiyah in \citet{atiyah1988commentary}; the theorem is \citet{atiyahsegal1969completion}.}
\GSfield{Present status}{As of 9 November 2026: established. An arbitrary elliptic operator, or the de Rham Euler index, need not give 1.}
\GSfield{Assessment}{literal statement: correct with missing hypotheses; publication match: partial; verification: publisher metadata checked; development: none; confidence: high}
\end{GSitem}

\GStalk{99}{1966-09-19}{R. A. Minlos}{P000150-P000156}{Canonical ensembles and factorial correlation functions}
The first lecture starts at P000151 and ends near the top of P000155. It defines continuum and lattice models, the thermodynamic limit, canonical Gibbs weights, and correlation densities. The 1968 lectures explicitly identify the autumn 1966 Gelfand seminar as a source, while tracing the earlier course to Khumsan in September 1965. Much of this first lecture is foundational exposition, not a new result. The ordered-pair convention is retained consistently.

\begin{GSitem}{1966-99.1}{E}{Classical phase space and a Hamiltonian}
\GSfield{Kind; source}{definition; \GSpage{P000151}}
\GSfield{The notes}{A physical system is a state space and a real Hamiltonian; the classical states are \((p_1,\ldots,p_N,x_1,\ldots,x_N)\), with \(p_i\in\mathbb R^\nu\), \(x_i\in\Omega\).}
\GSfield{Statement}{For bounded measurable \(\Omega\subset\mathbb R^d\), the labelled phase space is \(\mathbb R^{dN}\times\Omega^N\), possibly restricted by \(|x_i-x_j|>\delta\). With \(m>0\), \(H=\sum_i|p_i|^2/(2m)+W_N(x_1,\ldots,x_N)\). Forbidden configurations may equivalently have \(H=+\infty\).}
\GSfield{Reconstruction}{Measurability and labelled-particle convention are explicit.}
\GSfield{Fate}{Expository model setup: \citet[Lecture I]{minlos1968lectures}.}
\GSfield{Present status}{As of 9 November 2026: standard; allowing \(+\infty\) reconciles hard cores with the earlier phrase real-valued.}
\GSfield{Assessment}{literal statement: correct with missing hypotheses; publication match: exact; verification: read the result; development: none; confidence: high}
\end{GSitem}
\begin{GSitem}{1966-99.2}{E}{Ordered pair interactions and identical systems}
\GSfield{Kind; source}{definition; \GSpage{P000151-P000152}}
\GSfield{The notes}{\(W_N=\sum_{i\ne j}U(|x_i-x_j|)\); systems with fixed potential and mass but different \(N,\Omega\) are identical.}
\GSfield{Statement}{The notebook sums over ordered distinct pairs. Thus \(\sum_{i\ne j}U=\sum_{i<j}\phi\) with the usual unordered pair potential \(\phi=2U\). A thermodynamic sequence fixes \(U,m,\beta\) while changing \(N\) and the vessel.}
\GSfield{Reconstruction}{The conversion \(\phi=2U\) prevents silently importing a different normalization.}
\GSfield{Fate}{The conventions are compared with \citet[Lecture I]{minlos1968lectures}.}
\GSfield{Present status}{As of 9 November 2026: consistent if the factor two is retained throughout; it is not a transcription error.}
\GSfield{Assessment}{literal statement: correct; publication match: exact; verification: read the result; development: none; confidence: high}
\end{GSitem}
\begin{GSitem}{1966-99.3}{E}{Integrating out momenta}
\GSfield{Kind; source}{construction; \GSpage{P000152-P000154}}
\GSfield{The notes}{The configurational system omits kinetic energy; the full Gibbs density factors into momentum and position distributions.}
\GSfield{Statement}{For \(\beta,m>0\), Gaussian integration contributes \((2\pi m/\beta)^{dN/2}\). The position marginal is \(Z_{N,\Omega}^{-1}e^{-\beta W_N}\), independently of the momentum variables. A fixed \(N!\) factor cancels from normalized labelled probabilities but matters when defining thermodynamic partition functions.}
\GSfield{Reconstruction}{The Gaussian factor and normalization qualification are supplied.}
\GSfield{Fate}{The reduction is in \citet[Lecture I]{minlos1968lectures}.}
\GSfield{Present status}{As of 9 November 2026: established exact factorization.}
\GSfield{Assessment}{literal statement: correct; publication match: exact; verification: read the result; development: none; confidence: high}
\end{GSitem}
\begin{GSitem}{1966-99.4}{E}{A hard-core lattice gas}
\GSfield{Kind; source}{definition; \GSpage{P000152}}
\GSfield{The notes}{States are distinct sites of a cubic lattice; \(U(x)\) is even, without requiring spherical symmetry.}
\GSfield{Statement}{On a finite \(\Omega\subset\mathbb Z^d\), configurations are \(N\)-element subsets, or labelled distinct tuples with the corresponding symmetry factor. An even pair interaction \(U(z)=U(-z)\) defines \(W_N=\sum_{i\ne j}U(x_i-x_j)\). Lattice symmetry does not imply rotational invariance.}
\GSfield{Fate}{This recurs from 1964-3.1; see \citet[Lecture I]{minlos1968lectures}.}
\GSfield{Present status}{As of 9 November 2026: standard.}
\GSfield{Assessment}{literal statement: correct; publication match: exact; verification: read the result; development: none; confidence: high}
\end{GSitem}
\begin{GSitem}{1966-99.5}{E}{Repulsion, decay, and the Lennard--Jones example}
\GSfield{Kind; source}{remark; \GSpage{P000152}}
\GSfield{The notes}{\(U(r)=a r^{-12}-b r^{-6}\), \(a,b>0\), illustrates short-range repulsion and decay; finite-range and hard-core versions are also drawn.}
\GSfield{Statement}{For rigorous limits one needs quantitative assumptions, for example stability \(W_N\geq-BN\) and integrable tails or \(\int|e^{-\beta\phi(x)}-1|\,dx<\infty\). In physical dimension three the Lennard--Jones form has the appropriate repulsive core and integrable far tail. Mere growth at zero and decay at infinity, without rates or stability, are insufficient.}
\GSfield{Reconstruction}{The precise stability and regularity assumptions are editorial additions.}
\GSfield{Fate}{The stability issue was treated before this lecture by \citet{fisherruelle1966stability}; later density-expansion hypotheses are explicit in \citet[\S2]{kunatsagkarogiannis2018density}.}
\GSfield{Present status}{As of 9 November 2026: established framework. The aside that physicists use only this potential is not a mathematical theorem.}
\GSfield{Assessment}{literal statement: correct with missing hypotheses; publication match: partial; verification: read the result; development: none; confidence: high}
\end{GSitem}
\begin{GSitem}{1966-99.6}{E}{Thermodynamic scaling and vessel regularity}
\GSfield{Kind; source}{definition; \GSpage{P000153}}
\GSfield{The notes}{\(\Omega_N\) fills space, \(|\Omega_N|/N\to v\), and the vessel shape is usually immaterial.}
\GSfield{Statement}{Take \(N\to\infty\), \(|\Omega_N|\to\infty\), \(N/|\Omega_N|\to\rho=1/v\in(0,\infty)\), with a van Hove exhaustion: for each fixed \(a>0\), the volume of the \(a\)-boundary strip is \(o(|\Omega_N|)\). Shape-independent limits require hypotheses on the interaction and on the quantity being limited.}
\GSfield{Reconstruction}{The van Hove condition replaces an unrestricted exhaustion.}
\GSfield{Fate}{This continues 1964-3.2; the boundary-strip condition is explicit in \citet[Lecture III, \S4, p.~164]{minlos1968lectures}.}
\GSfield{Present status}{As of 9 November 2026: standard scaling, not a theorem that every correlation limit is unique at every density.}
\GSfield{Assessment}{literal statement: correct with missing hypotheses; publication match: exact; verification: read the result; development: none; confidence: high}
\end{GSitem}
\begin{GSitem}{1966-99.7}{E}{The finite canonical Gibbs law}
\GSfield{Kind; source}{definition; \GSpage{P000153}}
\GSfield{The notes}{\(p(a)=Q^{-1}e^{-\beta H(a)}\), \(Q=\int e^{-\beta H}\), is postulated in equilibrium; for a lattice replace density by probabilities.}
\GSfield{Statement}{Given a reference measure and \(0<Q<\infty\), the canonical law at fixed \(N,\Omega,\beta>0\) has that density. Stationarity alone does not characterize it: many functions of the energy yield invariant laws for Hamiltonian dynamics.}
\GSfield{Reconstruction}{Normalization and the logical status of the postulate are supplied.}
\GSfield{Fate}{The canonical ensemble is introduced in \citet[Lecture I]{minlos1968lectures}.}
\GSfield{Present status}{As of 9 November 2026: established equilibrium model; the Gibbs form is an assumption here, not deduced merely from time-independence.}
\GSfield{Assessment}{literal statement: correct with missing hypotheses; publication match: exact; verification: read the result; development: none; confidence: high}
\end{GSitem}
\begin{GSitem}{1966-99.8}{E}{Factorial correlation densities}
\GSfield{Kind; source}{definition; \GSpage{P000154}}
\GSfield{The notes}{\(\tau_k=(N!/(N-k)!)p(x_1,\ldots,x_k)\), where \(p\) is a marginal density.}
\GSfield{Statement}{For a symmetric labelled \(N\)-particle density, \(r_k=(N)_kp_k\), \((N)_k=N(N-1)\cdots(N-k+1)\). It represents ordered distinct particle tuples: \[\mathbb E\sum_{i_1,\ldots,i_k\ {\rm distinct}}h(x_{i_1},\ldots,x_{i_k})=\int h r_k,\] whenever the integral is meaningful.}
\GSfield{Reconstruction}{We use \(r_k\) as a typographically unambiguous synonym for the notebook’s \(\tau_k\).}
\GSfield{Fate}{The factorial convention is explicit in \citet[Lecture III, \S1]{minlos1968lectures}.}
\GSfield{Present status}{As of 9 November 2026: standard. These densities are not normalized probability densities.}
\GSfield{Assessment}{literal statement: correct; publication match: exact; verification: read the result; development: none; confidence: high}
\end{GSitem}
\begin{GSitem}{1966-99.9}{E}{Expectations of summatory observables}
\GSfield{Kind; source}{theorem; \GSpage{P000154}}
\GSfield{The notes}{First- and second-order summatory functions have means \(\int f\tau_1\) and \(\iint f\tau_2\).}
\GSfield{Statement}{For an integrable test function, \(\mathbb E\sum_if(x_i)=\int_\Omega f(x)r_1(x)\,dx\). For an integrable symmetric pair kernel, \(\mathbb E\sum_{i\ne j}h(x_i,x_j)=\int_{\Omega^2}h(x,y)r_2(x,y)\,dx\,dy\); no additional factor \(1/2\) occurs.}
\GSfield{Fate}{These identities are used in \citet[Lecture III, \S1 and \S4]{minlos1968lectures}.}
\GSfield{Present status}{As of 9 November 2026: exact consequences of the definition.}
\GSfield{Assessment}{literal statement: correct; publication match: exact; verification: read the result; development: none; confidence: high}
\end{GSitem}
\begin{GSitem}{1966-99.10}{E}{Local mean particle density}
\GSfield{Kind; source}{theorem; \GSpage{P000154-P000155}}
\GSfield{The notes}{For the particle count \(n_C\), \(\mathbb E n_C=\int_C\tau_1\); small regions identify \(\tau_1\) as local mean density.}
\GSfield{Statement}{For measurable \(C\subset\Omega\), \(\mathbb E n_C=\int_Cr_1\). If \(r_1\) is locally integrable, averages over shrinking balls converge to \(r_1(x)\) at its Lebesgue points; if it is continuous, this holds at every interior point.}
\GSfield{Reconstruction}{Lebesgue points specify the sense of the limiting density.}
\GSfield{Fate}{The interpretation is given in \citet[Lecture III, \S1]{minlos1968lectures}.}
\GSfield{Present status}{As of 9 November 2026: established, with the usual differentiation qualifications.}
\GSfield{Assessment}{literal statement: correct with missing hypotheses; publication match: exact; verification: read the result; development: none; confidence: high}
\end{GSitem}
\begin{GSitem}{1966-99.11}{E}{Probability of exactly one point in a small region}
\GSfield{Kind; source}{claim; \GSpage{P000155}}
\GSfield{The notes}{\(P(n_{\Delta C}=1)=\tau_1(x_0)\Delta C+o(\Delta C)\).}
\GSfield{Statement}{For shrinking neighbourhoods \(C\) of \(x_0\), continuity of \(r_1\) there and a locally bounded \(r_2\) imply \(P(n_C=1)=r_1(x_0)|C|+o(|C|)\). Indeed \(0\leq\mathbb E n_C-P(n_C=1)\leq\mathbb E[n_C(n_C-1)]=\int_{C^2}r_2=O(|C|^2)\).}
\GSfield{Reconstruction}{The explicit factorial-moment estimate is supplied.}
\GSfield{Fate}{The local interpretation occurs in \citet[Lecture III, \S1]{minlos1968lectures}.}
\GSfield{Present status}{As of 9 November 2026: correct with control of multiple occupancy; hard cores make that control immediate in sufficiently small cells.}
\GSfield{Assessment}{literal statement: correct with missing hypotheses; publication match: exact; verification: read the result; development: annotated; confidence: high}
\end{GSitem}

\GStalk{100}{1966-09-26}{R. A. Minlos}{P000154-P000159}{Bulk correlations, clustering, and representative means}
The second lecture begins midway down P000155 and ends on P000157. It deliberately labels several thermodynamic assertions as hypotheses, allows an effective specific volume depending on the correlation order, and uses clustering to argue that pair observables concentrate. The 1968 version restores an integrable decay requirement; nevertheless the displayed variance proof needs more than the stated \(L^1\) condition on the pair kernel. Our bounded development isolates the exact probabilistic implication and its limits.

\begin{GSitem}{1966-100.1}{E}{Two-cell probabilities and the second correlation density}
\GSfield{Kind; source}{claim; \GSpage{P000155}}
\GSfield{The notes}{\(\tau(x_1,x_2)\Delta x_1\Delta x_2\) is the probability of one particle in each of two small cells.}
\GSfield{Statement}{For disjoint shrinking cells \(C_1,C_2\) around distinct points, continuous \(r_2\) and locally bounded \(r_3,r_4\) yield \(P(n_{C_1}=n_{C_2}=1)=r_2(x_1,x_2)|C_1||C_2|+o(|C_1||C_2|)\). The exact identity without small-cell qualification is \(\mathbb E(n_{C_1}n_{C_2})=\int_{C_1\times C_2}r_2\).}
\GSfield{Reconstruction}{Disjointness and multiple-occupancy bounds are supplied.}
\GSfield{Fate}{Factorial interpretation: \citet[Lecture III, \S1]{minlos1968lectures}.}
\GSfield{Present status}{As of 9 November 2026: correct as a leading-order interpretation, not an exact finite-cell probability.}
\GSfield{Assessment}{literal statement: correct with missing hypotheses; publication match: exact; verification: read the result; development: none; confidence: high}
\end{GSitem}
\begin{GSitem}{1966-100.2}{F}{Bulk-correlation hypotheses with an unspecified exceptional set}
\GSfield{Kind; source}{conjecture; \GSpage{P000155-P000156}}
\GSfield{The notes}{Hypotheses 1--5 assert limiting \(\bar\tau_k\), a boundary-distance error, translation invariance, continuity away from exceptional parameters, and boundedness. The effective \(v_N\) may depend on \(k\).}
\GSfield{Statement}{For a specified stable finite-range continuum pair potential and fixed \((\rho,\beta)\), ask for a canonical bulk limit \(r_{k,N}\to r_k\), independent of a regular vessel exhaustion, and quantitative boundary/mixing estimates strong enough to integrate products of local observables. The dilute regime supplies such results in explicit ranges; the notebook proposes an extension through thermodynamic regimes, without specifying the exceptional set.}
\GSfield{Reconstruction}{The conditional moment implication is audited separately under item 1966-100.11. Its completed proof is not presented as an open-problem advance. The broader programme cannot be converted into a precise C question by silently defining its exceptional set.}
\GSfield{Fate}{The 1968 account retains hypotheses and asks for precise inequalities \citep[Lecture III, \S3 and Problem 3, p.~164]{minlos1968lectures}. Dilute canonical results appear in \citet[Theorem 2.5 and Corollary 2.6]{pulvirentitsagkarogiannis2015finite} and \citet[Theorems 2.1,2.7]{kunatsagkarogiannis2018density}.}
\GSfield{Present status}{As of 9 November 2026: the programme has rigorous dilute-regime realizations, but no precise all-density residual conjecture is certified by this source. The exceptional set and admissible phase selection are unspecified. F records this failure of exact formulation, not a failure to recognize the subject.}
\GSfield{Assessment}{literal statement: not determinable; publication match: partial; verification: read the result; development: none; confidence: medium}
\end{GSitem}
\begin{GSitem}{1966-100.3}{E}{An order-dependent effective specific volume}
\GSfield{Kind; source}{claim; \GSpage{P000155}}
\GSfield{The notes}{For every \(\tau_k\) there is \(v_N\), depending on the vessel and \(k\), that improves the boundary approximation; \(v_N\to v\).}
\GSfield{Statement}{In the ideal canonical gas of volume \(V\), \(r_{k,N}=(N)_k/V^k\). The infinite-volume ideal density is \(\rho^k\). For each fixed \(k\leq N\), choosing \(\rho_{N,k}=(N)_k^{1/k}/V\), equivalently \(v_{N,k}=V/(N)_k^{1/k}\), matches that order exactly and \(v_{N,k}\to v\). A single \(\rho_N\) cannot match both orders 1 and 2.}
\GSfield{Reconstruction}{\ours\ The explicit \(k\)-dependent formula is calculated here and checked by the supplied code. It verifies only the ideal gas, not interacting boundary rates.}
\GSfield{Fate}{The canonical adjustment is discussed in \citet[Lecture III, \S3, p.~161]{minlos1968lectures}; the exact ideal-gas calculation is elementary.}
\GSfield{Present status}{As of 9 November 2026: the notebook’s \(k\)-dependent quantifier passes the ideal-gas test. The stronger common-parameter version would fail.}
\GSfield{Assessment}{literal statement: correct with missing hypotheses; publication match: partial; verification: read the result; development: annotated; confidence: high}
\end{GSitem}
\begin{GSitem}{1966-100.4}{F}{The exponential wall-rate claim and its missing regime}
\GSfield{Kind; source}{claim; \GSpage{P000155-P000156}}
\GSfield{The notes}{For compactly supported potentials the decay function is indicated as \(e^{-\lambda r}\). Continuity may fail at exceptional \((v,\beta)\).}
\GSfield{Statement}{The displayed \(\gamma\) controls error by distance to the vessel wall, after the order-dependent parameter adjustment. This is distinct from spatial decay between two particle clusters. A precise theorem would specify the potential, thermodynamic regime, boundary conditions, and adjustment before asserting an exponential wall estimate. The notes explicitly allow exceptional parameters and failure of the rate at some volumes per particle.}
\GSfield{Reconstruction}{The canonical correction in the development concerns unadjusted pair covariance. The ideal gas actually supports the separate k-dependent wall-parameter adjustment, so that calculation does not refute this source claim.}
\GSfield{Fate}{\citet[Theorem 2.5, Corollary 2.6]{pulvirentitsagkarogiannis2015finite} gives a related dilute two-point estimate with a finite-volume term; it is not an exact match to this adjusted wall-error assertion.}
\GSfield{Present status}{As of 9 November 2026: no exact theorem with the notebook’s undeclared regime was identified. Its qualified claim is not refuted by imposing a universal quantifier it never states.}
\GSfield{Assessment}{literal statement: not determinable; publication match: thematic only; verification: unresolved lead; development: none; confidence: medium}
\end{GSitem}
\begin{GSitem}{1966-100.5}{E}{Translation invariance, continuity, and the density parameter}
\GSfield{Kind; source}{claim; \GSpage{P000155-P000156}}
\GSfield{The notes}{Translation invariance gives \(\bar\tau_1=\rho\). The notes say the hypotheses are proved for small \(v\) and \(\beta\), crediting Ruelle.}
\GSfield{Statement}{A translation-invariant state has a constant one-point density; identifying it with the imposed canonical density additionally uses bulk convergence without loss of mass. The perturbative gas regime is small density/activity, hence large specific volume \(v=1/\rho\), not small \(v\). Within a convergent expansion, correlations depend analytically on the appropriate parameter.}
\GSfield{Reconstruction}{The relation \(v=1/\rho\) is already explicit on P000155 and determines the correction.}
\GSfield{Fate}{Ruelle’s gas-phase results appeared in 1963--1964 \citep{ruelle1963correlation,ruelle1964cluster}; canonical density expansions are made precise by \citet[Theorem 2.7]{kunatsagkarogiannis2018density}.}
\GSfield{Present status}{As of 9 November 2026: the qualitative gas-phase conclusion is established; the source’s volume/density wording is inverted.}
\GSfield{Assessment}{literal statement: false as written; publication match: partial; verification: read the result; development: none; confidence: high}
\end{GSitem}
\begin{GSitem}{1966-100.6}{F}{Can arbitrary initial distributions lead to Gibbs laws?}
\GSfield{Kind; source}{question; \GSpage{P000156}; I. M. Gelfand; intervention by V. I. Arnold}
\GSfield{The notes}{Gelfand suggests that any distribution might approach Gibbs as particle number increases; Arnold objects that an ergodic hypothesis is missing.}
\GSfield{Statement}{Without a rule associating initial laws and dynamics for different \(N\), this is not one convergence theorem. Invariance or ergodicity on an energy surface concerns a microcanonical measure, whereas a canonical law mixes energy levels. Deterministic Hamiltonian evolution preserves energy and cannot send every initial energy distribution to one prescribed canonical distribution.}
\GSfield{Fate}{The ensemble distinction is treated in \citet[Lecture II]{minlos1968lectures}; no exact publication is identified for this exchange.}
\GSfield{Present status}{As of 9 November 2026: the unrestricted universality suggestion fails; no definite residual conjecture is recoverable.}
\GSfield{Assessment}{literal statement: not determinable; publication match: thematic only; verification: unresolved lead; development: none; confidence: medium}
\end{GSitem}
\begin{GSitem}{1966-100.7}{E}{Representativity is concentration of a normalized sequence}
\GSfield{Kind; source}{definition; \GSpage{P000156}}
\GSfield{The notes}{\(P(|F-\langle F\rangle|>\epsilon)\to0\); the notes immediately call this imprecise because there is no one function \(F\).}
\GSfield{Statement}{For an extensive observable \(F_N\), the intended property is \((F_N-\mathbb EF_N)/N\to0\) in probability. If \(\mathbb EF_N/N\to m\), this gives \(F_N/N\to m\). Unnormalized fluctuations need not converge to zero.}
\GSfield{Reconstruction}{The triangular sequence of observables and probability laws is made explicit.}
\GSfield{Fate}{The normalization is explicit in \citet[Lecture III, \S4]{minlos1968lectures}.}
\GSfield{Present status}{As of 9 November 2026: established definition; the notebook itself flags the shorthand.}
\GSfield{Assessment}{literal statement: correct with missing hypotheses; publication match: exact; verification: read the result; development: none; confidence: high}
\end{GSitem}
\begin{GSitem}{1966-100.8}{E}{An integrable pair observable}
\GSfield{Kind; source}{construction; \GSpage{P000156}}
\GSfield{The notes}{\(F_N=\sum_{i\ne j}f(x_i-x_j)\), with one fixed \(f\) and \(\int|f|<\infty\).}
\GSfield{Statement}{Use a real even measurable \(f\in L^1(\mathbb R^d)\) and ordered pairs. A uniform bound \(r_{2,N}\leq C_2\) implies \(\mathbb E|F_N|\leq C_2|\Omega_N|\|f\|_1\). It does not imply a finite second moment. The symmetry condition is inherited from the preceding lecture.}
\GSfield{Reconstruction}{Explicit evenness and the moment distinction are supplied.}
\GSfield{Fate}{The same integrability assumption is (3.24) in \citet[Lecture III, \S4]{minlos1968lectures}.}
\GSfield{Present status}{As of 9 November 2026: the expectation is controlled; variance needs additional short-distance integrability.}
\GSfield{Assessment}{literal statement: correct with missing hypotheses; publication match: exact; verification: read the result; development: annotated; confidence: high}
\end{GSitem}
\begin{GSitem}{1966-100.9}{A}{The mean per particle has a correlation-integral limit}
\GSfield{Kind; source}{claim; \GSpage{P000156}}
\GSfield{The notes}{\(\langle F_N\rangle=|\Omega_N|\int\bar\tau_2(\xi)f(\xi)\,d\xi+o(|\Omega_N|)\), hence the mean per particle tends to \(v\) times that integral.}
\GSfield{Statement}{A sufficient precise hypothesis is \(r_{2,N}\leq C_2\) and almost-everywhere convergence of \(A_N(z)=|\Omega_N|^{-1}\int\mathbf1_{\Omega_N}(x)\mathbf1_{\Omega_N}(x-z)r_{2,N}(x,x-z)\,dx\) to a bulk function \(R(z)\). Then dominated convergence gives \(\mathbb EF_N/|\Omega_N|\to\int fR\), and \(\mathbb EF_N/N\to v\int fR\) if \(|\Omega_N|/N\to v\). Pointwise convergence near a fixed origin alone is insufficient.}
\GSfield{Reconstruction}{Uniform integrability and the vessel condition are supplied; the final subscript must remain 2.}
\GSfield{Fate}{The calculation is published in \citet[Lecture III, \S4, pp.~162--163]{minlos1968lectures}, from the 1965/1966 lectures.}
\GSfield{Present status}{As of 9 November 2026: valid under the stated uniform bulk hypotheses.}
\GSfield{Assessment}{literal statement: correct with missing hypotheses; publication match: partial; verification: read the result; development: none; confidence: high}
\end{GSitem}
\begin{GSitem}{1966-100.10}{E}{The multiplicative cluster property}
\GSfield{Kind; source}{conjecture; \GSpage{P000157}}
\GSfield{The notes}{\(|\bar\tau_{k+s}-\bar\tau_k\bar\tau_s|\leq C_{k,s}\gamma(\min|x_i-y_j|)\) is said to follow, with difficulty, from preceding hypotheses.}
\GSfield{Statement}{This is a mixing property of the limiting state and requires more than existence, boundedness, translation invariance and continuity of correlations. For example a stationary mixture of Poisson processes of two distinct intensities has bounded factorial densities of each order but a nonzero, distance-independent two-point connected correlation.}
\GSfield{Reconstruction}{\ours\ If the random intensity is \(Z\in\{a,b\}\), \(r_k=\mathbb EZ^k\), and \(r_2-r_1^2=\operatorname{Var}Z>0\). This example refutes the logical implication between qualitative correlation properties; it is not claimed to be a canonical limit for the notebook’s fixed potential.}
\GSfield{Fate}{The gas-phase cluster theorem predates the talk \citep{ruelle1964cluster}. The general inference is not supplied by \citet[Lecture III, \S3]{minlos1968lectures}.}
\GSfield{Present status}{As of 9 November 2026: established in suitable pure/uniqueness regimes; false as an automatic consequence of the qualitative properties listed.}
\GSfield{Assessment}{literal statement: false as written; publication match: partial; verification: read the result; development: annotated; confidence: high}
\end{GSitem}
\begin{GSitem}{1966-100.11}{A}{The claimed linear variance and its integrability gap}
\GSfield{Kind; source}{claim; \GSpage{P000157}}
\GSfield{The notes}{\(DF_N\sim dN\), so \(D(F_N/N)\sim d/N\) and Chebyshev gives representativity.}
\GSfield{Statement}{With \(f\in L^1\cap L^2\), bounded second and third factorial densities, and an integrable finite-volume two-pair connected bound allowing an \(O(V^{-1})\) canonical term, one obtains \(\operatorname{Var}F_N=O(V)=O(N)\). These assumptions alone do not give a positive asymptotic coefficient. With only \(L^1\), variance can be infinite; concentration can nevertheless follow by truncation under the same correlation bounds.}
\GSfield{Reconstruction}{\ours\ The repair, sharp distinction of moments, and truncation implication are proved in the development. They clarify the bounded implication within GP-1966-09, rather than solve its general thermodynamic part.}
\GSfield{Fate}{The 1968 proof adds integrable decay and a van Hove condition in its p.~164 remark \citep[Lecture III, \S4]{minlos1968lectures}, but still displays \(f^2\) after assuming only (3.24).}
\GSfield{Present status}{As of 9 November 2026: the notebook’s stated assumptions do not prove the variance asymptotic. The chapter supplies a complete conditional repair and a counterexample, without claiming novelty for the probabilistic method.}
\GSfield{Assessment}{literal statement: false as written; publication match: partial; verification: read the result; development: developed in this chapter; confidence: high; problem \hyperref[prob:GP-1966-09]{GP-1966-09}}
\end{GSitem}
\begin{GSitem}{1966-100.12}{E}{The exact ordered-pair square expansion}
\GSfield{Kind; source}{theorem; \GSpage{P000157}}
\GSfield{The notes}{\(F_N^2=2\sum_{i\ne j}f_{ij}^2+4\sum_{i,j,k\ {\rm distinct}}f_{ij}f_{kj}+\sum_{i,j,k,s\ {\rm distinct}}f_{ij}f_{ks}\).}
\GSfield{Statement}{For \(f_{ij}=f_{ji}\), this identity partitions two ordered pairs by whether their union has two, three, or four labels. Integration gives coefficients 2, 4, and 1 against \(r_2,r_3,r_4\), respectively. All three sums are over ordered distinct tuples as displayed.}
\GSfield{Reconstruction}{\oursComputed\ The evenness requirement is explicit; the exact combinatorial verification is supplied with its output.}
\GSfield{Fate}{The same combinatorial calculation occurs in \citet[Lecture III, \S4, p.~163]{minlos1968lectures}.}
\GSfield{Present status}{As of 9 November 2026: exact. The supplied code checks the identity on exhaustive small labelled configurations independently of the proof.}
\GSfield{Assessment}{literal statement: correct with missing hypotheses; publication match: exact; verification: read the result; development: annotated; confidence: high}
\end{GSitem}

\GStalk{101}{1966-09-26}{D. A. Kazhdan}{P000157-P000159}{Isolation of the trivial representation and lattice finiteness}
Two finiteness conclusions for lattices are announced on P000158, followed by the representation-theoretic mechanism now called property \(T\). Kazhdan’s short paper, received two weeks after the seminar, is an unusually close match. A separate sentence about groups over \(\mathbb Q\) retains no conclusion and cannot be identified with a theorem.

\begin{GSitem}{1966-101.1}{A}{Finite abelianization of the lattice}
\GSfield{Kind; source}{theorem; \GSpage{P000158}}
\GSfield{The notes}{If the simple real or \(p\)-adic group \(G\) contains \(\mathrm{SL}(3,P)\), and \(\Gamma\) is discrete with finite-volume quotient, then \(\Gamma/[\Gamma,\Gamma]\) is finite.}
\GSfield{Statement}{Let \(K\) be a nondiscrete locally compact field. Every lattice in \(\mathrm{SL}_3(K)\) has finite abelianization. The same conclusion holds for a locally compact group \(G\) with property \(T\), including the simple algebraic groups covered by Kazhdan’s Theorem 6. In that extension the algebraic subgroup and compact-generation hypotheses of that theorem must be retained.}
\GSfield{Reconstruction}{The explicit \(\mathrm{SL}_3(K)\) case is exact. The larger-group wording is restricted to the published theorem rather than extrapolated from the notebook’s abbreviation.}
\GSfield{Fate}{\citet[Theorems 1, 3, 5--6]{kazhdan1967dual}; received 10 October 1966, published 1967.}
\GSfield{Present status}{As of 9 November 2026: established. Containing an arbitrary property-\(T\) subgroup, without the ambient hypotheses, does not transfer property \(T\) to every group.}
\GSfield{Assessment}{literal statement: correct with missing hypotheses; publication match: exact; verification: read the result; development: none; confidence: high}
\end{GSitem}
\begin{GSitem}{1966-101.2}{A}{Finite generation of the lattice}
\GSfield{Kind; source}{theorem; \GSpage{P000158}}
\GSfield{The notes}{\(\Gamma\) has finitely many generators, even extractable from any generating set.}
\GSfield{Statement}{A countable discrete group with property \(T\) is finitely generated; hence the lattices of the preceding item are finitely generated. Any generating set of a finitely generated group contains a finite generating subset, since finitely many fixed generators use only finitely many words in that set.}
\GSfield{Reconstruction}{The parenthetical strengthening is the elementary finite-word argument, not a stronger finiteness property.}
\GSfield{Fate}{\citet[Theorems 2--3]{kazhdan1967dual}, announced here in September 1966 and printed in 1967.}
\GSfield{Present status}{As of 9 November 2026: established. Finite presentation is a separate assertion and is not stated in these notes.}
\GSfield{Assessment}{literal statement: correct; publication match: exact; verification: read the result; development: none; confidence: high}
\end{GSitem}
\begin{GSitem}{1966-101.3}{A}{Isolation of the trivial unitary representation}
\GSfield{Kind; source}{method; \GSpage{P000158}}
\GSfield{The notes}{The proof uses \ru{изолированность единичного представления}, isolation of the trivial representation.}
\GSfield{Statement}{For a locally compact group, property \(T\) means that every unitary representation with almost invariant unit vectors has a nonzero invariant vector. Equivalently, the trivial representation is isolated in the appropriate unitary-dual topology. Kazhdan proves this for \(\mathrm{SL}_3(K)\) and transfers it to finite-covolume closed subgroups.}
\GSfield{Reconstruction}{Almost invariant vectors give an explicit modern formulation of the isolated-representation shorthand.}
\GSfield{Fate}{\citet[\S1 and Theorems 3, 5]{kazhdan1967dual}. The terminology and mechanism occur in the original 1967 paper.}
\GSfield{Present status}{As of 9 November 2026: established; the notes record the central innovation, not merely its later name.}
\GSfield{Assessment}{literal statement: correct; publication match: exact; verification: read the result; development: none; confidence: high}
\end{GSitem}
\begin{GSitem}{1966-101.4}{F}{A rational-group representation fragment}
\GSfield{Kind; source}{remark; \GSpage{P000158}}
\GSfield{The notes}{An algebraic group over \(\mathbb Q\) containing \(\mathrm{SL}(3,\mathbb Q)\) is followed only by \ru{что-то о представлениях}, something about representations.}
\GSfield{Statement}{No asserted property of representations, topology on \(G(\mathbb Q)\), or class of representations survives. Neither a local-field property-\(T\) theorem nor an automorphic theorem can be inferred from this sentence.}
\GSfield{Fate}{The nearby results match \citet{kazhdan1967dual}; the rational-group fragment has no exact statement to match.}
\GSfield{Present status}{As of 9 November 2026: unidentified as a mathematical assertion, even though its subject is clear.}
\GSfield{Assessment}{literal statement: not determinable; publication match: thematic only; verification: unresolved lead; development: none; confidence: low}
\end{GSitem}

\GStalk{102}{1966-09-26}{Lee Lorch}{P000157-P000160}{Higher monotonicity of zeros and wave integrals}
The lecturer’s handwritten name on P000158 is uncertain; the MCCME entry and the joint work with Peter Szego identify Lee Lorch. The page treats Sturm comparison, higher finite differences of zeros, Bessel and Airy examples, and integrals between consecutive zeros. The strict inequality cannot follow from the weak derivative hypotheses exactly as written. Gelfand’s endpoint remark is mathematically suggestive but not completely recorded.

\begin{GSitem}{1966-102.1}{E}{Monotonicity of zero spacing and its alleged converse}
\GSfield{Kind; source}{claim; \GSpage{P000158}}
\GSfield{The notes}{Under the label Sturm (1836), increasing \(f\) in \(y^{\prime\prime}+fy=0\) makes consecutive zeros closer, “and conversely.”}
\GSfield{Statement}{For a continuous increasing coefficient on an oscillatory interval, Sturm comparison gives nonincreasing consecutive zero gaps, with strictness under an appropriate strict comparison. Decreasing gaps of one solution do not conversely force its coefficient to be increasing.}
\GSfield{Reconstruction}{\ours\ For \(x>1\), put \(y=x^{-1/2}e^{\epsilon\cos(2x^2)}\sin(x^2)\), \(\epsilon=1/100\). Its zeros \(\sqrt{k\pi}\) have strictly decreasing gaps. Direct differentiation gives \(f=x^2[4+16\epsilon+32\epsilon\cos(2x^2)-16\epsilon^2\sin^2(2x^2)]-3/(4x^2)\). This is smooth and positive on \((1,\infty)\), but \(f^{\prime}<0\) at all sufficiently large points with \(2x^2=\pi/2+2k\pi\). Thus even positive oscillatory coefficients do not satisfy the proposed converse. The identity is checked symbolically in the supplied code.}
\GSfield{Fate}{The comparison-theory setting precedes the talk and is the starting point of \citet[Introduction]{lorchszego1963higher}.}
\GSfield{Present status}{As of 9 November 2026: the forward comparison is standard; the unrestricted converse is false.}
\GSfield{Assessment}{literal statement: false as written; publication match: partial; verification: read the result; development: annotated; confidence: high}
\end{GSitem}
\begin{GSitem}{1966-102.2}{E}{Higher monotonicity needs a nondegeneracy condition}
\GSfield{Kind; source}{theorem; \GSpage{P000158}}
\GSfield{The notes}{The assumptions are \((-1)^jf^{(j+1)}\geq0\), \(j=1,\ldots,N\), and \(f(\infty)>0\); the conclusion is \((-1)^n\Delta^{n+1}x_k>0\).}
\GSfield{Statement}{The written statement fails for \(f\equiv1\), whose zero gaps are constant. A checked sufficient theorem is formulated through independent real solutions \(y_1,y_2\) and \(p=y_1^2+y_2^2\): if \(p>0\), \(-p^{\prime}>0\), and \((-1)^jp^{(j)}\geq0\) for \(2\leq j\leq N\), with the derivatives continuous up to order \(N-1\) and the last derivative existing on the interval, then the alternating finite differences of zero gaps are strictly positive through the stated orders wherever the zeros exist.}
\GSfield{Reconstruction}{\ours\ The constant-coefficient counterexample also survives adding the missing \(j=0\) weak inequality. Accordingly no silent change of an index repairs strictness.}
\GSfield{Fate}{\citet[Theorem 2.1]{lorchszego1963higher} has the strict phase-amplitude conditions; \citet[Theorem, pp.~447--448]{muldoon1968extension} weakens the higher ones while retaining \(-p^{\prime}>0\).}
\GSfield{Present status}{As of 9 November 2026: the theorem family is established, but the exact recorded coefficient-only formulation is false. The later phase-amplitude theorem is a partial, explicitly different match.}
\GSfield{Assessment}{literal statement: false as written; publication match: partial; verification: read the result; development: annotated; confidence: high}
\end{GSitem}
\begin{GSitem}{1966-102.3}{E}{The Bessel and Airy examples}
\GSfield{Kind; source}{remark; \GSpage{P000158}}
\GSfield{The notes}{Bessel functions with \(\nu>1/2\), Airy functions, and others are listed as examples.}
\GSfield{Statement}{Writing \(u(x)=\sqrt{x}\,J_\nu(x)\) puts the Bessel equation in normal form \(u^{\prime\prime}+[1-(\nu^2-1/4)x^{-2}]u=0\). The oscillatory Airy equation is \(u^{\prime\prime}+xu=0\), for example \(u=\operatorname{Ai}(-x)\). Normal-form dependent variables and the oscillatory branch must be specified before importing zero or wave-integral assertions.}
\GSfield{Reconstruction}{The normalizing factor and the sign of the Airy argument are supplied.}
\GSfield{Fate}{These special-function applications belong to \citet[applications]{lorchszego1963higher}; the transformed equations identify the notebook’s examples.}
\GSfield{Present status}{As of 9 November 2026: established examples. A statement about wave areas of \(J_\nu\) itself is not identical to one about \(\sqrt{x}J_\nu\).}
\GSfield{Assessment}{literal statement: correct with missing hypotheses; publication match: partial; verification: read the result; development: none; confidence: high}
\end{GSitem}
\begin{GSitem}{1966-102.4}{E}{Higher monotonicity of wave integrals}
\GSfield{Kind; source}{theorem; \GSpage{P000158}}
\GSfield{The notes}{For consecutive zeros, \(M_k=\int_{x_k}^{x_{k+1}}|y|^\lambda\,dx\), \(\lambda>-1\), has the same property; \(\lambda=0\) gives gaps and \(\lambda=1\) areas.}
\GSfield{Statement}{Under the hypotheses of item 1966-102.2, \((-1)^n\Delta^nM_k>0\) for \(0\leq n\leq N\), whenever all required consecutive zeros lie in the interval. Nontrivial solutions have simple zeros, and \(\lambda>-1\) is the exact local integrability condition there.}
\GSfield{Reconstruction}{The range of differences is explicitly limited by the available zeros.}
\GSfield{Fate}{\citet[Theorem 2.1]{lorchszego1963higher}; the weakened hypotheses and unchanged conclusion are in \citet[Theorem]{muldoon1968extension}.}
\GSfield{Present status}{As of 9 November 2026: established with the stated hypotheses; the finite-difference order for gaps is one higher because \(M_k=x_{k+1}-x_k\) when \(\lambda=0\).}
\GSfield{Assessment}{literal statement: correct with missing hypotheses; publication match: exact; verification: read the result; development: none; confidence: high}
\end{GSitem}
\begin{GSitem}{1966-102.5}{F}{The divergent endpoint and reciprocal derivatives}
\GSfield{Kind; source}{remark; \GSpage{P000158}; I. M. Gelfand}
\GSfield{The notes}{Gelfand’s \(\lambda=-1\) remark contains \(1/y^{\prime}(x_k)\) and an unresolved abbreviation; it is not a complete limiting formula.}
\GSfield{Statement}{The ordinary integral at \(\lambda=-1\) diverges at every simple endpoint zero. A precise nearby identity, not asserted to be the intended one, is \(\lim_{\lambda\downarrow-1}(\lambda+1)\int_a^b|y(x)|^\lambda dx=|y^{\prime}(a)|^{-1}+|y^{\prime}(b)|^{-1}\) for consecutive simple zeros of a \(C^1\) function with no zero inside.}
\GSfield{Reconstruction}{\ours\ Near \(a\), for \(0<\delta<|y^{\prime}(a)|\), \((|y^{\prime}(a)|-\delta)(x-a)\leq|y(x)|\leq(|y^{\prime}(a)|+\delta)(x-a)\). Integrating these bounds with \(-1<\lambda<0\), multiplying by \(\lambda+1\), taking the limit and then \(\delta\downarrow0\) gives the first reciprocal. Repeat at \(b\); on the remaining compact interval the scaled integral tends to zero. The sine example is independently checked with the beta-integral formula.}
\GSfield{Fate}{The integral theorem in \citet{lorchszego1963higher} assumes \(\lambda>-1\); no checked source identifies the exact endpoint remark.}
\GSfield{Present status}{As of 9 November 2026: the literal ordinary integral is divergent; the intended regularization or signed difference remains unidentified.}
\GSfield{Assessment}{literal statement: not determinable; publication match: none found; verification: unresolved lead; development: annotated; confidence: low}
\end{GSitem}

\GStalk{103}{1966-09-26}{I. M. Gelfand}{P000158-P000160}{A recapitulation of equivariant K-theory and completion}
On P000159 Gelfand explicitly repeats Bernstein’s and Atiyah’s material: equivariant bundles, the Borel construction, the representation-ring action, and the approximate identification with ordinary K-theory. These are new notebook occurrences of the same constructions, not additional discovery claims.

\begin{GSitem}{1966-103.1}{E}{Equivariant K-theory as a contravariant ring functor}
\GSfield{Kind; source}{definition; \GSpage{P000159}}
\GSfield{The notes}{For compact \(G\), equivariant bundles and formal differences define \(K_G(X)\); the functor is contravariant and receives \(R(G)\) from \(X\to*\).}
\GSfield{Statement}{For a compact Lie group and compact \(G\)-space in the finite \(G\)-CW setting, complex equivariant vector bundles form the Grothendieck ring \(K_G^0(X)\). Pullback gives contravariance, and \(K_G^0(*)=R(G)\) acts by tensoring with representation bundles.}
\GSfield{Reconstruction}{Finite \(G\)-CW hypotheses align the next completion statement with its checked classical scope.}
\GSfield{Fate}{The definitions precede this reprise; see \citet[Introduction]{atiyahsegal1969completion} and the earlier occurrences 1966-97.1 and 1966-98.4.}
\GSfield{Present status}{As of 9 November 2026: standard. Being an \(R(G)\)-algebra does not require the structural homomorphism to be injective.}
\GSfield{Assessment}{literal statement: correct with missing hypotheses; publication match: exact; verification: publisher metadata checked; development: none; confidence: high}
\end{GSitem}
\begin{GSitem}{1966-103.2}{E}{The Borel construction is not the ordinary orbit space}
\GSfield{Kind; source}{construction; \GSpage{P000159}}
\GSfield{The notes}{\(X_G=E\times_GX\) is called a generalization of the classifying space; for a point it is \(BG\). The accompanying phrase says the orbits are contracted to points.}
\GSfield{Statement}{For a contractible free \(G\)-space \(EG\), the diagonal quotient \(EG\times_GX\) is the homotopy quotient and is a bundle over \(BG\) with fibre \(X\). Only for free suitable actions is it homotopy equivalent to \(X/G\). For trivial actions it is \(BG\times X\).}
\GSfield{Reconstruction}{The distinction between homotopy and ordinary quotient is made explicit.}
\GSfield{Fate}{This is the construction used by \citet{atiyahsegal1969completion}; it repeats item 1966-98.5.}
\GSfield{Present status}{As of 9 November 2026: standard. The orbit phrase must refer to the diagonal action, not the loss of all isotropy information from \(X\).}
\GSfield{Assessment}{literal statement: correct with missing hypotheses; publication match: exact; verification: publisher metadata checked; development: none; confidence: high}
\end{GSitem}
\begin{GSitem}{1966-103.3}{B}{The approximate identification means I-adic completion}
\GSfield{Kind; source}{theorem; \GSpage{P000159}}
\GSfield{The notes}{\(K_G(X)\) approximately coincides with \(K(X_G)\).}
\GSfield{Statement}{For a finite \(G\)-CW complex and compact Lie \(G\), the comparison identifies the augmentation-ideal completion of \(K_G^*(X)\) with \(K^*(EG\times_GX)\), interpreted through the skeletal inverse system. It is generally not an isomorphism from the uncompleted ring.}
\GSfield{Reconstruction}{The ideal, completion, grading, and finiteness assumptions are supplied.}
\GSfield{Fate}{\citet{atiyahsegal1969completion}, printed 1969, supplies the general completion theorem. Atiyah’s earlier special cases and Bernstein’s exposition account for the result circulating before this occurrence.}
\GSfield{Present status}{As of 9 November 2026: established in the stated finite setting; this is the same theorem family as item 1966-98.6.}
\GSfield{Assessment}{literal statement: correct with missing hypotheses; publication match: exact; verification: publisher metadata checked; development: none; confidence: high}
\end{GSitem}

\GStalk{104}{1966-09-26}{I. N. Bernstein}{P000158-P000162}{Families of elliptic operators and the index bundle}
P000159--P000160 supply the analytic construction behind the K-theoretic direct image, then a finite-dimensional modification of a Fredholm family. The essential issue is that kernel dimensions jump. The index is a virtual bundle after stabilization, not generally the difference of two pre-existing kernel and cokernel bundles. This continues 1964-15.4.

\begin{GSitem}{1966-104.1}{E}{Twisting a vertical elliptic family}
\GSfield{Kind; source}{construction; \GSpage{P000159}}
\GSfield{The notes}{For a fibration \(Y\to X\) with compact manifold fibre \(M_x\), a family \(D_x\) acts between restrictions of two bundles; twisting by a bundle \(V\) on \(Y\) gives an elliptic family.}
\GSfield{Statement}{Let \(p:Y\to X\) be a smooth fibre bundle with closed fibres and compact base, and let \(D\) be a continuous family of vertical elliptic differential operators between complex vector bundles. A chosen connection defines a twist \(D\otimes V\) with symbol \(\sigma(D)\otimes1_V\); its analytic index defines an additive homomorphism \(p_!^D:K^0(Y)\to K^0(X)\).}
\GSfield{Reconstruction}{Closed fibres, compact base, operator continuity, and a connection are supplied. The symbol determines the index.}
\GSfield{Fate}{The family-index construction was discussed in Atiyah’s 1966 congress lecture, printed in 1968, and fully developed in \citet[\S2, Proposition 2.2 and Definition 2.3]{atiyahsinger1971families}; see \citet{atiyah1968global}.}
\GSfield{Present status}{As of 9 November 2026: established. Independence of auxiliary choices is an index assertion, not the claim that twisting a differential operator is connection-free.}
\GSfield{Assessment}{literal statement: correct with missing hypotheses; publication match: exact; verification: read the result; development: none; confidence: high}
\end{GSitem}
\begin{GSitem}{1966-104.2}{E}{The direct image obeys the projection formula}
\GSfield{Kind; source}{theorem; \GSpage{P000159}}
\GSfield{The notes}{The reverse map is required to be a map of \(K(X)\)-modules; the de Rham complex is mentioned as an example.}
\GSfield{Statement}{For the family index, \(p_!^D(p^*a\cdot b)=a\cdot p_!^D(b)\). A left inverse to \(p^*\) follows only if \(p_!^D(1)=1\). A de Rham family has fibrewise numerical index \(\chi(M)\), which need not equal one; the Dolbeault operator on the flag fibre used in the completion argument has index one.}
\GSfield{Reconstruction}{The condition on the index of 1 is the missing algebraic step that connects this construction to the preceding completion argument.}
\GSfield{Fate}{The family-index framework is in \citet[\S2, Proposition 2.2 and Definition 2.3]{atiyahsinger1971families}; Atiyah discusses the flag-manifold direct image in \citet{atiyah1988commentary}.}
\GSfield{Present status}{As of 9 November 2026: established. An arbitrary elliptic family is not automatically a retraction of pullback.}
\GSfield{Assessment}{literal statement: correct with missing hypotheses; publication match: exact; verification: read the result; development: none; confidence: high}
\end{GSitem}
\begin{GSitem}{1966-104.3}{E}{Kernels and cokernels need not form vector bundles}
\GSfield{Kind; source}{claim; \GSpage{P000159-P000160}}
\GSfield{The notes}{Finite-dimensional kernels and cokernels are subtracted to obtain an element of K-theory; after a projected deformation the kernel dimension is said to be locally constant.}
\GSfield{Statement}{A norm-continuous family of Fredholm operators has locally constant numerical index but need not have locally constant kernel or cokernel dimension. For \(A_t:\mathbb C\to\mathbb C\), \(A_tz=tz\), both dimensions jump at \(t=0\). Thus the virtual index requires stabilization or local patching.}
\GSfield{Reconstruction}{\ours\ The one-dimensional family is a direct counterexample to the unstabilized bundle claim: its kernel is zero for \(t\ne0\) and \(\mathbb C\) at zero.}
\GSfield{Fate}{This distinction underlies the family construction of \citet[\S2, Proposition 2.2 and Definition 2.3]{atiyahsinger1971families}; it also explains the Fredholm discussion in \citet{atiyah1968global}.}
\GSfield{Present status}{As of 9 November 2026: established. The abbreviated difference is valid only after the required index-bundle construction.}
\GSfield{Assessment}{literal statement: false as written; publication match: exact; verification: read the result; development: annotated; confidence: high}
\end{GSitem}
\begin{GSitem}{1966-104.4}{E}{A rigorous finite-dimensional stabilization}
\GSfield{Kind; source}{method; \GSpage{P000160}}
\GSfield{The notes}{A map \(X\to F(H,H)\) is modified through \(A_x(1-t)+tA_xP_V\), where the projection’s precise definition is abbreviated.}
\GSfield{Statement}{For a compact Hausdorff base \(X\), complex Hilbert spaces \(H_0,H_1\), and norm-continuous Fredholm family \(A_x:H_0\to H_1\), choose a finite-dimensional \(V\subset H_1\) with \(\operatorname{im}A_x+V=H_1\) for every \(x\). The augmented operators \(B_x:H_0\oplus V\to H_1\), \(B_x(u,v)=A_xu+v\), are surjective and their kernels form a finite-rank vector bundle. Define \(\operatorname{Ind}(A)=[\ker B]-[X\times V]\).}
\GSfield{Reconstruction}{\ours\ For each \(x\), a complement to the range gives a surjective augmentation on a neighbourhood; a finite cover supplies a common \(V\). Locally a continuous right inverse exists by the Neumann-series argument, so the kernels form a bundle. Enlarging \(V\) to \(W\) yields \(0\to\ker B_V\to\ker B_W\to X\times(W/V)\to0\), preserving the virtual class. A common enlargement compares any two choices. Apply the same argument over compact \(X\times[0,1]\) for homotopy invariance.}
\GSfield{Fate}{This is the standard family-index construction, contemporary with the 1966 lecture and recorded in \citet[\S2, Proposition 2.2 and Definition 2.3]{atiyahsinger1971families}. The specific projected formula in the notebook is not asserted to be an exact transcription of that argument.}
\GSfield{Present status}{As of 9 November 2026: established. The resulting class is independent of the auxiliary finite-dimensional subspace and homotopy invariant.}
\GSfield{Assessment}{literal statement: correct with missing hypotheses; publication match: partial; verification: read the result; development: annotated; confidence: high}
\end{GSitem}

\GStalk{105}{1966-10-03}{Louis Nirenberg}{P000160-P000164}{Complex powers and an analytic index formula}
The actual notes occupy P000161--P000162, after the preceding talk on P000160. Nirenberg reports Seeley's and H\"ormander's analytic approach to the index: positive-spectrum cancellation, complex powers, continuation of kernels and a local expression at exponent zero. Gelfand comments on analytic descriptions of topological quantities. The catalogue's wider interval also contains Minlos's next lecture. Both formula-critical pages were inspected against their scans and Seeley's contemporary preprint in its published Russian translation \citep[pp.~96--98]{seeley1968powers}.

\begin{GSitem}{1966-105.1}{E}{Analytic expressions for topological index data}
\GSfield{Kind; source}{remark; \GSpage{P000161}; I. M. Gelfand}
\GSfield{The notes}{Nirenberg's report is analytic without explicit topology. The scan says the connection to topology and the disappearance of lower-order terms are unclear; the supplied transcription instead says “no connections” and “leading terms”. Gelfand replies that Atiyah--Singer and Pontryagin classes also have analytic expressions.}
\GSfield{Statement}{The Fredholm index of an elliptic operator on a closed manifold depends on the topological class of its principal symbol. Characteristic classes have differential-form representatives constructed from connections, allowing analytic integral formulae for topological invariants.}
\GSfield{Reconstruction}{The corrected readings are \ru{неясна связь} and \ru{младшие члены}. No particular degree map or characteristic polynomial is supplied. Independence of the index under lower-order perturbations must be distinguished from dependence of a chosen local density representative.}
\GSfield{Fate}{The index theorem predates this seminar; the analytic construction actually developed next is Seeley's \citep[pp.~96--98]{seeley1968powers}.}
\GSfield{Present status}{As of 9 November 2026: established background; no separate conjecture is specified.}
\GSfield{Assessment}{literal statement: correct with missing hypotheses; publication match: partial; verification: read the result; development: annotated; confidence: high}
\end{GSitem}
\begin{GSitem}{1966-105.2}{E}{Positive spectra of the two Laplacians}
\GSfield{Kind; source}{theorem; \GSpage{P000161}}
\GSfield{The notes}{For \(\lambda>0\), \(E_\lambda(L^*L)\) and \(E_\lambda(LL^*)\) have common dimension \(d_\lambda\); the identification explicitly excludes zero.}
\GSfield{Statement}{Let \(L:C^\infty(M;E)\to C^\infty(M;F)\) be elliptic of positive order on a closed smooth manifold with Hermitian bundles and a smooth positive density. For \(\lambda>0\), \(L:E_\lambda(L^*L)\to E_\lambda(LL^*)\) is an isomorphism with inverse \(\lambda^{-1}L^*\). Moreover \(\ker L^*L=\ker L\) and \(\ker LL^*=\ker L^*\).}
\GSfield{Reconstruction}{The inverse and analytic setting are supplied. Direct composition proves the eigenspace assertion; \(\langle L^*Lu,u\rangle=\|Lu\|^2\) proves the kernel assertion.}
\GSfield{Fate}{This elementary spectral argument is the input to the complex-power index calculation \citep[p.~96]{seeley1968powers}.}
\GSfield{Present status}{As of 9 November 2026: standard and correct under the stated domain hypotheses.}
\GSfield{Assessment}{literal statement: correct; publication match: exact; verification: read the result; development: annotated; confidence: high}
\end{GSitem}
\begin{GSitem}{1966-105.3}{E}{Spectral complex powers and traces}
\GSfield{Kind; source}{construction; \GSpage{P000161}}
\GSfield{The notes}{For \(\operatorname{Re}s\ll0\), \((I+L^*L)^s\) has an eigenfunction kernel whose integrated diagonal gives \(\sum d_\lambda(1+\lambda)^s\).}
\GSfield{Statement}{If \(L\) has order \(r>0\), \(A=I+L^*L\) has order \(2r\). For \(\operatorname{Re}s<-n/(2r)\), \(A^s\) is trace class and \(\operatorname{Tr}A^s=\sum_j(1+\lambda_j)^s\), counting multiplicity. The kernel is distributionally \(\sum_j(1+\lambda_j)^s\Psi_j(x)\Psi_j(y)^*\); sufficiently negative real part gives continuity and the diagonal trace integral.}
\GSfield{Reconstruction}{One may sum over distinct eigenvalues with multiplicities, or over a full eigenbasis without an additional multiplicity factor. The convergence threshold and bundle adjoint are supplied.}
\GSfield{Fate}{Spectral calculus is the known starting point for the continuation in \citet[pp.~96--97]{seeley1968powers}.}
\GSfield{Present status}{As of 9 November 2026: standard with convergence qualifications.}
\GSfield{Assessment}{literal statement: correct with missing hypotheses; publication match: exact; verification: read the result; development: annotated; confidence: high}
\end{GSitem}
\begin{GSitem}{1966-105.4}{E}{Regularized trace difference and index}
\GSfield{Kind; source}{theorem; \GSpage{P000161}}
\GSfield{The notes}{All positive-eigenvalue terms cancel in \(\operatorname{Tr}(I+L^*L)^s-\operatorname{Tr}(I+LL^*)^s\), leaving \(\dim\ker L-\dim\ker L^*\), independent of \(s\).}
\GSfield{Statement}{Initially for \(\operatorname{Re}s<-n/(2r)\), the difference of the two traces equals \(\operatorname{ind}L\). Their meromorphic continuations therefore have the same constant difference, in particular at a common regular point such as zero.}
\GSfield{Reconstruction}{At zero the separately unregularized traces are not finite. The continuation identity follows from the identity theorem applied to the trace-class half-plane.}
\GSfield{Fate}{The cancellation is elementary and underlies the local formula \citep[pp.~96--98]{seeley1968powers}.}
\GSfield{Present status}{As of 9 November 2026: established; regularization must remain explicit.}
\GSfield{Assessment}{literal statement: correct; publication match: exact; verification: read the result; development: annotated; confidence: high}
\end{GSitem}
\begin{GSitem}{1966-105.5}{B}{Contour powers with a spectral cut}
\GSfield{Kind; source}{construction; \GSpage{P000162}}
\GSfield{The notes}{An angular gap about the negative ray is assumed for the principal symbol; powers use a resolvent contour and multiplication by integer powers. A parenthesis proposes reducing to order one by multiplication by an invertible operator.}
\GSfield{Statement}{Let \(A\) be an invertible classical elliptic operator of positive order \(m\) on a closed manifold, with a ray avoided by its principal-symbol spectra and operator spectrum. Fix a branch and oriented contour. Then \(A^s=(2\pi i)^{-1}\int_\Gamma\lambda^s(\lambda-A)^{-1}d\lambda\), initially for negative real part, extends by integer powers to a holomorphic family satisfying \(A^sA^t=A^{s+t}\). The parameter resolvent estimate applies along the contour.}
\GSfield{Reconstruction}{The symbol condition concerns eigenvalues in the bundle case; invertibility and an operator spectral cut are added. Order reduction can simplify an index problem but does not preserve individual powers or kernels under arbitrary multiplication.}
\GSfield{Fate}{Seeley's circulated preprint contains this construction; its Russian translation appeared in 1968 \citep[pp.~96--97]{seeley1968powers}. Nirenberg is reporting another author's work.}
\GSfield{Present status}{As of 9 November 2026: standard spectral-cut calculus.}
\GSfield{Assessment}{literal statement: correct with missing hypotheses; publication match: exact; verification: read the result; development: annotated; confidence: high}
\end{GSitem}
\begin{GSitem}{1966-105.6}{B}{Holomorphic kernels in a left half-plane}
\GSfield{Kind; source}{theorem; \GSpage{P000162}}
\GSfield{The notes}{Assertion (i) says \(s\mapsto K_s\) is holomorphic into \(C^k\) when \(\operatorname{Re}(ms)<-n-k\).}
\GSfield{Statement}{Under the spectral-cut hypotheses of the preceding item, the distribution kernel of \(A^s\) is a holomorphic family; for \(\operatorname{Re}(ms)<-n-k\) its values are in \(C^k(M\times M)\), holomorphically in that topology.}
\GSfield{Reconstruction}{The operator acts on a smooth bundle over a closed \(n\)-manifold. The stated threshold ensures enough negative order for the kernel derivatives to be continuous.}
\GSfield{Fate}{This is the first kernel assertion of \citet[p.~97]{seeley1968powers}.}
\GSfield{Present status}{As of 9 November 2026: established for classical elliptic operators with a spectral cut.}
\GSfield{Assessment}{literal statement: correct; publication match: exact; verification: read the result; development: annotated; confidence: high}
\end{GSitem}
\begin{GSitem}{1966-105.7}{B}{Entire continuation off the diagonal}
\GSfield{Kind; source}{theorem; \GSpage{P000162}}
\GSfield{The notes}{Assertion (ii) takes a compact set in \(M\times M\) disjoint from the diagonal and continues the restricted kernels holomorphically over the whole \(s\)-plane.}
\GSfield{Statement}{For every compact subset \(C\subset M\times M\setminus\operatorname{diag}M\), the kernel of \(A^s\) depends entirely on \(s\) as a smooth function on a neighbourhood of \(C\), with every \(C^\infty\) seminorm controlled locally in \(s\).}
\GSfield{Reconstruction}{Smoothness in \(x,y\) is not real analyticity; the notebook itself warns that the coefficients are only infinitely differentiable.}
\GSfield{Fate}{This is the off-diagonal kernel assertion of \citet[p.~97]{seeley1968powers}.}
\GSfield{Present status}{As of 9 November 2026: established in the classical elliptic calculus.}
\GSfield{Assessment}{literal statement: correct; publication match: exact; verification: read the result; development: annotated; confidence: high}
\end{GSitem}
\begin{GSitem}{1966-105.8}{B}{Meromorphic continuation of the diagonal density}
\GSfield{Kind; source}{theorem; \GSpage{P000162}}
\GSfield{The notes}{Assertion (iii) gives possible simple poles at \((j-n)/m\), omitting zero because it is regular.}
\GSfield{Statement}{The initially defined diagonal \(K_s(x,x)\) has meromorphic continuation, with at most simple poles among \((j-n)/m\), \(j\in\mathbb N_0\). The candidate pole at zero is removable; other possible poles may cancel. This is a continued density-valued endomorphism.}
\GSfield{Reconstruction}{Continuation is performed after taking the diagonal in a half-plane where it is defined. It is not a general distributional restriction operation.}
\GSfield{Fate}{The candidate poles and regularity at zero are in \citet[p.~97]{seeley1968powers}.}
\GSfield{Present status}{As of 9 November 2026: established; the distinction from the full zero-power kernel is treated next.}
\GSfield{Assessment}{literal statement: correct; publication match: exact; verification: read the result; development: annotated; confidence: high}
\end{GSitem}
\begin{GSitem}{1966-105.9}{E}{Exponent zero versus diagonal continuation}
\GSfield{Kind; source}{claim; \GSpage{P000162}}
\GSfield{The notes}{The manuscript says \ru{при \(s=0\) это не единичный оператор}, with the explanation that the kernel is not a delta function.}
\GSfield{Statement}{As an operator \(A^0=I\), whose full distribution kernel is the identity delta distribution. The finite value at zero of the meromorphically continued diagonal \(K_s(x,x)\) results from restricting first in a left half-plane and then continuing. It is not the diagonal restriction of the full kernel of \(A^0\).}
\GSfield{Reconstruction}{This corrects a conflation in the handwritten source, not a mistranscription. Acting on each eigenvector verifies \(A^0=I\) directly.}
\GSfield{Fate}{The functional calculus and continued diagonal in \citet[pp.~96--97]{seeley1968powers} distinguish these operations.}
\GSfield{Present status}{As of 9 November 2026: the literal operator assertion is false; the intended regularized-diagonal distinction is valid.}
\GSfield{Assessment}{literal statement: false as written; publication match: partial; verification: read the result; development: annotated; confidence: high}
\end{GSitem}
\begin{GSitem}{1966-105.10}{B}{Local resolvent coefficient and index density}
\GSfield{Kind; source}{method; \GSpage{P000162}}
\GSfield{The notes}{The notes write \((2\pi)^nK_0(x,x)=m^{-1}\int_{|\xi|=1}\int_{-\infty}^0b_{-m-n}(x,\xi;\lambda)\,d\lambda\,dS\), then take the difference of such expressions for the index.}
\GSfield{Statement}{Continue the diagonal densities of \((I+L^*L)^s\) and \((I+LL^*)^s\) to zero. The integral of the difference of their fibre traces is \(\operatorname{ind}L\). Each local continued density is computable from the parameter-resolvent symbol term of degree \(-m-n\), with its resolvent-sign and contour conventions fixed consistently.}
\GSfield{Reconstruction}{The coefficient formula is recorded in Seeley's convention. Replacing \((A-\lambda)^{-1}\) by \((\lambda-A)^{-1}\) changes the coefficient sign; the contour must change consistently. The invariant assertion is the integrated difference.}
\GSfield{Fate}{The coefficient expression is \citet[p.~97, formula (II)]{seeley1968powers}; it is the close publication match.}
\GSfield{Present status}{As of 9 November 2026: established analytic index machinery.}
\GSfield{Assessment}{literal statement: correct with missing hypotheses; publication match: exact; verification: read the result; development: annotated; confidence: high}
\end{GSitem}
\begin{GSitem}{1966-105.11}{B}{Complex powers remain pseudodifferential}
\GSfield{Kind; source}{theorem; \GSpage{P000162}}
\GSfield{The notes}{The final paragraph says \(A^s\) is pseudodifferential, with a symbol obtained by inserting the symbol in the contour formula.}
\GSfield{Statement}{For an invertible classical elliptic operator of order \(m>0\) with a spectral cut, \(A^s\) is classical of complex order \(ms\). Its homogeneous symbol coefficients are obtained by contour integration of the parameter-resolvent symbol expansion.}
\GSfield{Reconstruction}{The substitution is in the full symbolic calculus: lower-order derivative corrections are present. It is generally false that the exact symbol is just the pointwise scalar power \(a(x,\xi)^s\).}
\GSfield{Fate}{This is the symbol construction of \citet[pp.~97--99]{seeley1968powers}.}
\GSfield{Present status}{As of 9 November 2026: established; its symbol calculation supplies the preceding local density formula.}
\GSfield{Assessment}{literal statement: correct with missing hypotheses; publication match: exact; verification: read the result; development: annotated; confidence: high}
\end{GSitem}

\GStalk{106}{1966-10-03}{R. A. Minlos}{P000162-P000165}{Infinite configurations and correlation axioms}
The actual lecture is P000163--P000164. It constructs the measurable space of locally finite configurations, discusses limiting Gibbs measures and mixing, invokes ergodicity for mean density and introduces thermodynamic functions. Gelfand asks for intrinsic axioms on correlations. Minlos's expanded lectures trace their first delivery to Khumsan in September 1965 and their amplification to the autumn 1966 Gelfand seminar \citep[pp.~137--138]{minlos1968lectures}. That publication distinguishes results and hypotheses more carefully than this notebook.

\begin{GSitem}{1966-106.1}{E}{Unlabelled finite configurations}
\GSfield{Kind; source}{construction; \GSpage{P000163}}
\GSfield{The notes}{Collision-free labelled configurations map to \(N\)-point subsets \(C_{\Omega,N}\). Symmetry of \(H\) allows the Hamiltonian and Gibbs measure to descend; the discarded collision set has measure zero.}
\GSfield{Statement}{For bounded Borel \(\Omega\subset\mathbb R^\nu\), quotient the collision-free subset of \(\Omega^N\) by permutations. A symmetric measurable Hamiltonian with positive finite partition function defines a Gibbs probability, whose pushforward is a law on simple \(N\)-point configurations. Lebesgue reference measure gives all collision diagonals measure zero.}
\GSfield{Reconstruction}{The \(1/N!\) factor belongs consistently to the quotient reference measure or to both normalized integrals. The null-set claim presupposes a non-atomic reference measure.}
\GSfield{Fate}{The construction is explicit in \citet[\S1, pp.~60--62]{minlos1967limiting} and \citet[Lecture II]{minlos1968lectures}.}
\GSfield{Present status}{As of 9 November 2026: standard configuration-space probability.}
\GSfield{Assessment}{literal statement: correct; publication match: exact; verification: read the result; development: annotated; confidence: high}
\end{GSitem}
\begin{GSitem}{1966-106.2}{A}{Counting-cylinder sigma algebra}
\GSfield{Kind; source}{definition; \GSpage{P000163}}
\GSfield{The notes}{\(C_\infty\) consists of finite or countable subsets without accumulation points; the sigma algebra is generated by \(n(d,U_i)=k_i\) for finite collections of disjoint sets.}
\GSfield{Statement}{Let \(\Gamma\) be the locally finite simple configurations in \(\mathbb R^\nu\), with the sigma algebra generated by \(N_U(\gamma)=|\gamma\cap U|\) for bounded Borel \(U\). Finite intersections of counting events generate it; disjoint observation sets suffice by a finite Boolean refinement.}
\GSfield{Reconstruction}{Bounded Borel sets and ambient local finiteness specify the intended measurable space.}
\GSfield{Fate}{This is the space used in Minlos's forthcoming construction \citep[\S1]{minlos1967limiting}; A marks the exact authorial publication, not a claim of novelty for counting cylinders.}
\GSfield{Present status}{As of 9 November 2026: the standard measurable configuration space.}
\GSfield{Assessment}{literal statement: correct; publication match: exact; verification: read the result; development: annotated; confidence: high}
\end{GSitem}
\begin{GSitem}{1966-106.3}{A}{Dilute limiting Gibbs distribution}
\GSfield{Kind; source}{theorem; \GSpage{P000163}}
\GSfield{The notes}{The notes ask for a limiting space with the finite-volume asymptotic properties and say existence is proved at low density.}
\GSfield{Statement}{A verified precise version is grand canonical: for stable pair interactions with integrable Mayer function, and sufficiently small activity at fixed \(\beta>0\), the low-activity expansion constructs a translation-invariant probability law on locally finite configurations with the limiting local distributions. A concrete subclass consists of hard-core finite-range potentials bounded below outside the core.}
\GSfield{Reconstruction}{Stability is \(H_N\ge-BN\); Mayer integrability is \(\int|e^{-\beta\Phi(x)}-1|dx<\infty\). The ensemble and small parameter are supplied. This does not alone establish every canonical fixed-density limit or equivalence of ensembles.}
\GSfield{Fate}{\citet[\S\S1--3]{minlos1967limiting}, received 20 December 1966, is the close construction. The canonical limiting-measure hypothesis is separately identified by \citet[Lecture II, pp.~150--152]{minlos1968lectures}.}
\GSfield{Present status}{As of 9 November 2026: proved in dilute regimes; unrestricted density, uniqueness and complete phase diagrams are different questions.}
\GSfield{Assessment}{literal statement: correct with missing hypotheses; publication match: partial; verification: read the result; development: annotated; confidence: high}
\end{GSitem}
\begin{GSitem}{1966-106.4}{A}{Uniform mixing and its quantifiers}
\GSfield{Kind; source}{claim; \GSpage{P000163}}
\GSfield{The notes}{Separated cylinder events satisfy \(\mu(A\cap B)-\mu(A)\mu(B)\to0\), with the gloss “uniformly over all pairs \(U\)”; the limiting measure is said to have this property.}
\GSfield{Statement}{The later lecture formulation fixes bounded \(U\) and \(A\in\mathcal F_U\), then asks for \(\sup_{V:\operatorname{dist}(U,V)\ge r}\sup_{B\in\mathcal F_V}|\mu(A\cap B)-\mu(A)\mu(B)|\to0\). Uniformity over unrestricted sizes of both windows is stronger. Pointwise correlation convergence does not by itself supply either uniformity or a quantitative variance bound.}
\GSfield{Reconstruction}{The fixed-event quantifier is taken from the expanded lectures, not silently substituted into the diplomatic transcription. The 1967 existence paper explicitly defers the mixing/ergodicity discussion.}
\GSfield{Fate}{\citet[Lecture II, formula (2.5), pp.~151--152]{minlos1968lectures} states the careful version and its general-regime hypothesis status; low-density cluster theory already appears in \citet{ruelle1964cluster}.}
\GSfield{Present status}{As of 9 November 2026: qualified dilute clustering is established, but the notebook's unrestricted gloss is unsupported.}
\GSfield{Assessment}{literal statement: correct with missing hypotheses; publication match: partial; verification: read the result; development: annotated; confidence: medium}
\end{GSitem}
\begin{GSitem}{1966-106.5}{E}{Ergodic mean density}
\GSfield{Kind; source}{theorem; \GSpage{P000164}}
\GSfield{The notes}{Almost every configuration has a mean density, by Birkhoff's theorem and ergodicity.}
\GSfield{Statement}{For a translation-invariant ergodic point-process law with finite intensity \(\rho=\mathbb E N_{[0,1)^\nu}\), growing cubes satisfy \(N_{Q_R}/|Q_R|\to\rho\) almost surely. Without ergodicity the limit is an invariant conditional expectation and can be random.}
\GSfield{Reconstruction}{Finite intensity and a specified averaging sequence are added. The assertion is a law of large numbers, not a rate or central limit theorem.}
\GSfield{Fate}{This application of the spatial ergodic theorem is restated in \citet[Lecture II, p.~152]{minlos1968lectures}.}
\GSfield{Present status}{As of 9 November 2026: established; mixing is stronger than necessary.}
\GSfield{Assessment}{literal statement: correct with missing hypotheses; publication match: exact; verification: read the result; development: annotated; confidence: high}
\end{GSitem}
\begin{GSitem}{1966-106.6}{B}{Characterizing correlation hierarchies}
\GSfield{Kind; source}{question; \GSpage{P000164}; I. M. Gelfand}
\GSfield{The notes}{Gelfand proposes starting with axioms on correlation functions and asks how to single them out.}
\GSfield{Statement}{Given a full hierarchy of symmetric locally finite factorial correlation measures with \(\rho_0=1\), characterize when it is realized by a probability law on locally finite configurations. The key positivity condition tests finite configuration polynomials \(KG(\gamma)=\sum_{\eta\subset\gamma,\ |\eta|<\infty}G(\eta)\): if \(KG\ge0\) pointwise, then the corresponding integral against the correlation measure must be nonnegative. Symmetry, this positivity and normalization characterize realizability in Lenard's correlation-measure framework; uniqueness is separate.}
\GSfield{Reconstruction}{This is a precise reconstruction of the one-line proposal. Entrywise nonnegativity of the correlations is insufficient. Full-hierarchy realizability differs from a truncated pair-correlation problem.}
\GSfield{Fate}{\citet[abstract and main characterization]{lenard1975characterization} gives necessary and sufficient conditions and explicitly permits nonunique realizations; \citet{lenard1973correlation} treats determinacy.}
\GSfield{Present status}{As of 9 November 2026: established at the level of an infinite family of positivity conditions.}
\GSfield{Assessment}{literal statement: correct with missing hypotheses; publication match: exact; verification: publisher metadata checked; development: annotated; confidence: high; problem \hyperref[prob:GP-1966-10]{GP-1966-10}}
\end{GSitem}
\begin{GSitem}{1966-106.7}{E}{Partition function and thermodynamic limit}
\GSfield{Kind; source}{theorem; \GSpage{P000164}}
\GSfield{The notes}{The notes define \(Z=N!^{-1}\int e^{-\beta H}da\), say \(E,P,S,F\) are expressed through \(\log Z\), and promise \(\lim N^{-1}\log Z\).}
\GSfield{Statement}{For stable tempered classical pair potentials, regular expanding vessels and fixed admissible density \(N/|\Omega|\to\rho>0\), the thermodynamic limit defines configurational free energy per particle through \(-(\beta N)^{-1}\log Z_{\Omega,N}\). Stability excludes collapse, decay controls long-range tails, and admissibility excludes impossible hard-core densities. Existence of this limit does not automatically imply convergence of every derivative.}
\GSfield{Reconstruction}{The potential assumptions and regular vessel sequence are absent from this page. Full phase-space normalization adds a known momentum contribution. No all-shapes assertion is inferred.}
\GSfield{Fate}{The expanded account is \citet[Lecture IV, \S\S1--4, pp.~164--176]{minlos1968lectures}; the thermodynamic-limit programme predates this talk.}
\GSfield{Present status}{As of 9 November 2026: established under stability and decay assumptions, compatible with phase-boundary nondifferentiability.}
\GSfield{Assessment}{literal statement: correct with missing hypotheses; publication match: partial; verification: read the result; development: annotated; confidence: high}
\end{GSitem}

\GStalk{107}{1966-10-10}{R. A. Minlos}{P000163-P000166}{Energy, wall pressure and entropy}
The fourth lecture occupies the upper part of P000165. It records the energy derivative, a separate moving-wall pressure conjecture, and entropy per particle. The manuscript itself has an entropy sign error and an erroneous reciprocal-temperature factor. Arnold starts below the separator. The expanded treatment is \citet[Lecture IV]{minlos1968lectures}.

\begin{GSitem}{1966-107.1}{E}{Mean energy and temperature differentiation}
\GSfield{Kind; source}{theorem; \GSpage{P000165}}
\GSfield{The notes}{\(E=\langle H\rangle=-\partial_\beta\log Z\), with \(\langle H\rangle\approx Ne(v,\beta)\) for pair interactions; representativeness is said to have been proved.}
\GSfield{Statement}{For a \(\beta\)-independent Hamiltonian and fixed reference measure, finiteness of \(Z\) and of the first energy moment in a neighbourhood permits \(-\partial_\beta\log Z=\mathbb EH\). If convex functions \(f_N=N^{-1}\log Z_N\) converge on an open interval to \(f\), then \(f_N'(\beta)\to f'(\beta)\) at differentiability points of \(f\), and \(N^{-1}\mathbb EH_N\to-f'(\beta)\).}
\GSfield{Reconstruction}{Numerator and denominator must use the same \(1/N!\) convention. Mean convergence is weaker than concentration; the latter requires fluctuation control.}
\GSfield{Fate}{The identities and thermodynamic use are in \citet[Lecture IV, \S\S1--3]{minlos1968lectures}.}
\GSfield{Present status}{As of 9 November 2026: standard, with derivative limits qualified at phase transitions.}
\GSfield{Assessment}{literal statement: correct with missing hypotheses; publication match: partial; verification: read the result; development: annotated; confidence: high}
\end{GSitem}
\begin{GSitem}{1966-107.2}{F}{Limiting wall pressure}
\GSfield{Kind; source}{question; \GSpage{P000165}}
\GSfield{The notes}{\(p(x)=\tau_1(x)\), \(P=S^{-1}\int p\,d\sigma\); a hypothesis says wall pressure tends to a constant. The notes correctly warn that fixed-window convergence of \(\tau_1\) does not imply this.}
\GSfield{Statement}{A definite theorem would require hard-wall or other boundary conditions, a family of growing vessels, a moving boundary observation point, normalization of pressure, and the exact uniformity sought. These are unspecified, so a unique modern proposition cannot be extracted.}
\GSfield{Reconstruction}{For a classical hard wall the contact relation is \(\beta p_{\rm wall}=\rho_{\rm contact}\), with the appropriate mechanical stress. It is not a universal identity of pressure with a bulk correlation. Moving boundary points are outside the earlier fixed-window limit.}
\GSfield{Fate}{The thermodynamic framework is present in \citet[Lecture IV]{minlos1968lectures}, but no publication establishing this unrestricted moving-wall assertion was verified.}
\GSfield{Present status}{As of 9 November 2026: insufficiently specified to label solved or open; the distinction of limiting procedures is valid.}
\GSfield{Assessment}{literal statement: not determinable; publication match: thematic only; verification: unresolved lead; development: annotated; confidence: medium}
\end{GSitem}
\begin{GSitem}{1966-107.3}{E}{Pressure normalization}
\GSfield{Kind; source}{claim; \GSpage{P000165}}
\GSfield{The notes}{The manuscript writes \(p=\partial_V\log Z\).}
\GSfield{Statement}{Physical pressure with \(k_B=1\) is \(p=\beta^{-1}(\partial_V\log Z)_{\beta,N}\), whenever a differentiable family of vessel deformations defines the derivative. The written equation instead holds for reduced pressure \(\beta p\).}
\GSfield{Reconstruction}{The units are unspecified; the neighbouring contact-density formula suggests reduced pressure. This is a normalization ambiguity, not a certain transcription error.}
\GSfield{Fate}{These conventions belong to the identities reviewed by \citet[Lecture IV]{minlos1968lectures}.}
\GSfield{Present status}{As of 9 November 2026: established after fixing units and the volume variation.}
\GSfield{Assessment}{literal statement: correct with missing hypotheses; publication match: partial; verification: read the result; development: annotated; confidence: high}
\end{GSitem}
\begin{GSitem}{1966-107.4}{E}{Entropy sign and temperature factor}
\GSfield{Kind; source}{claim; \GSpage{P000165}}
\GSfield{The notes}{The scan has \(S=\int\pi\log\pi\) and \(S/N=N^{-1}\log Z-\beta^{-1}\partial_\beta(N^{-1}\log Z)\).}
\GSfield{Statement}{For \(\pi=Z^{-1}e^{-\beta H}\), Gibbs entropy is \(S=-\int\pi\log\pi=\log Z+\beta\mathbb EH=\log Z-\beta\partial_\beta\log Z\). Thus \(S/N=f_N-\beta f_N'\). An opposite entropy sign convention changes both terms consistently.}
\GSfield{Reconstruction}{Substitution of \(\log\pi=-\beta H-\log Z\) verifies the correction without a thermodynamic-limit theorem. Both discrepancies are present in the handwriting.}
\GSfield{Fate}{The corrected identity is standard in \citet[Lecture IV, \S1]{minlos1968lectures}; the two notebook formulae cannot share a consistent convention.}
\GSfield{Present status}{As of 9 November 2026: false as written; the correction is completely determined.}
\GSfield{Assessment}{literal statement: false as written; publication match: partial; verification: read the result; development: annotated; confidence: high}
\end{GSitem}
\begin{GSitem}{1966-107.5}{E}{Entropy per particle}
\GSfield{Kind; source}{claim; \GSpage{P000165}}
\GSfield{The notes}{The limit of entropy per particle is asserted.}
\GSfield{Statement}{If \(f_N=N^{-1}\log Z_N\) converges on an open interval to differentiable \(f\), convex convergence gives \(S_N/N\to f(\beta)-\beta f'(\beta)\). At a nondifferentiability point this argument provides no unique limit without a state or one-sided prescription.}
\GSfield{Reconstruction}{This sufficient formulation covers a regular single-phase regime without inferring differentiability merely from existence of free energy.}
\GSfield{Fate}{The relevant expanded account is \citet[Lecture IV, \S\S1--3]{minlos1968lectures}.}
\GSfield{Present status}{As of 9 November 2026: standard with the phase qualification.}
\GSfield{Assessment}{literal statement: correct with missing hypotheses; publication match: partial; verification: read the result; development: annotated; confidence: high}
\end{GSitem}

\GStalk{108}{1966-10-10}{V. I. Arnold}{P000164-P000168}{Lagrangian planes and the Maslov class}
Arnold begins in lower P000165 and continues through P000166. He identifies \(\Lambda(n)=U(n)/O(n)\), constructs determinant squared and pulls the circle class back by the tangent map of a Lagrangian manifold. The exact authorial match is \citet{arnold1967characteristic}, submitted 28 September 1966, before this talk, and published January 1967. It clarifies the geometric class behind earlier occurrence 1964-21.3.

\begin{GSitem}{1966-108.1}{E}{Three compatible linear structures}
\GSfield{Kind; source}{definition; \GSpage{P000165}}
\GSfield{The notes}{The Euclidean metric, multiplication by \(i\), and \([x,y]=(Ix,y)\) are introduced. The pairwise intersections of their preserving groups all equal \(U(n)\).}
\GSfield{Statement}{If \(g\) is positive definite and \(J^2=-I\), \(g(Jv,Jw)=g(v,w)\), then \(\omega(v,w)=g(Jv,w)\) is nondegenerate skew symmetric. Preserving any two of \(g,J,\omega\) preserves the third. Thus the pairwise intersections of \(O(2n)\), \(GL(n,\mathbb C)\), and \(Sp(2n,\mathbb R)\) equal \(U(n)\).}
\GSfield{Reconstruction}{Compatibility of \(g\) and \(J\) is explicit; arbitrary metric and complex structures need not satisfy the assertion.}
\GSfield{Fate}{This linear algebra is Lemma 1.1 of \citet[\S1.1]{arnold1967characteristic}.}
\GSfield{Present status}{As of 9 November 2026: standard.}
\GSfield{Assessment}{literal statement: correct; publication match: exact; verification: read the result; development: annotated; confidence: high}
\end{GSitem}
\begin{GSitem}{1966-108.2}{E}{Isotropic and Lagrangian spaces}
\GSfield{Kind; source}{definition; \GSpage{P000165-P000166}}
\GSfield{The notes}{An \(n\)-plane with vanishing symplectic form is Lagrangian. Real and imaginary planes are examples; planes are said to have dimension at most \(n\), and every curve for \(n=1\) is Lagrangian.}
\GSfield{Statement}{An isotropic subspace \(W\subset V^{2n}\) satisfies \(\omega|_W=0\), hence \(\dim W\le n\). It is Lagrangian when \(\dim W=n\), equivalently \(W=W^\omega\). An immersed \(n\)-manifold is Lagrangian when all tangent spaces are. Every immersed curve in a symplectic surface satisfies this.}
\GSfield{Reconstruction}{The inequality concerns isotropic spaces; “Lagrangian” means maximal dimension. A singular arbitrary curve need not be a submanifold.}
\GSfield{Fate}{These are the definitions used in \citet[\S1.2]{arnold1967characteristic}.}
\GSfield{Present status}{As of 9 November 2026: standard with the terminology clarified.}
\GSfield{Assessment}{literal statement: correct with missing hypotheses; publication match: exact; verification: read the result; development: annotated; confidence: high}
\end{GSitem}
\begin{GSitem}{1966-108.3}{E}{Generating functions for Lagrangian graphs}
\GSfield{Kind; source}{theorem; \GSpage{P000166}}
\GSfield{The notes}{The graph \(p=p(q)\) is Lagrangian if and only if \(p=\partial S/\partial q\).}
\GSfield{Statement}{For a smooth one-form \(\alpha=\sum p_i(q)dq_i\) on \(U\subset\mathbb R^n\), its cotangent graph is Lagrangian exactly when \(d\alpha=0\). Locally this is equivalent to \(\alpha=dS\). Globally it is equivalent to exactness, and holds in particular for simply connected \(U\).}
\GSfield{Reconstruction}{The local qualification is missing. A closed angular one-form on an annulus has a Lagrangian graph but no globally single-valued primitive.}
\GSfield{Fate}{This is the standard generating-function example in the framework of \citet[\S1]{arnold1967characteristic}.}
\GSfield{Present status}{As of 9 November 2026: correct locally; exactness is required globally.}
\GSfield{Assessment}{literal statement: correct with missing hypotheses; publication match: partial; verification: read the result; development: annotated; confidence: high}
\end{GSitem}
\begin{GSitem}{1966-108.4}{A}{Lagrangian Grassmannian as a homogeneous space}
\GSfield{Kind; source}{theorem; \GSpage{P000166}}
\GSfield{The notes}{\(U(n)\) acts transitively and \(O(n)\) stabilizes the real plane, so \(\Lambda(n)=U(n)/O(n)\).}
\GSfield{Statement}{A real orthonormal basis of any Lagrangian plane extends complex linearly to a unitary map. The stabilizer of the standard real plane is exactly \(O(n)\), yielding a smooth identification \(\Lambda(n)\cong U(n)/O(n)\).}
\GSfield{Reconstruction}{A compatible metric presents the homogeneous space; Lagrangian subspaces themselves are defined by the symplectic form.}
\GSfield{Fate}{This is Lemma 1.2 in \citet[\S1.2]{arnold1967characteristic}, already submitted when the lecture was given.}
\GSfield{Present status}{As of 9 November 2026: established; A records the matched authorial exposition rather than priority for elementary linear algebra.}
\GSfield{Assessment}{literal statement: correct; publication match: exact; verification: read the result; development: annotated; confidence: high}
\end{GSitem}
\begin{GSitem}{1966-108.5}{A}{Determinant-square fibration}
\GSfield{Kind; source}{construction; \GSpage{P000166}}
\GSfield{The notes}{Determinant squared descends to \(\Lambda(n)\) because orthogonal determinants are \(\pm1\); a diagram relates \(SO,O,S^0\), \(SU,U,S^1\), and \(S\Lambda,\Lambda,S^1\).}
\GSfield{Statement}{The map \(d([U])=(\det U)^2\) from \(U(n)/O(n)\) to \(S^1\) is well defined. Its fibre over one is \(SU(n)/SO(n)\), since a determinant-minus-one representative can be adjusted by a real reflection. Quotient maps and squaring on \(S^1\) give the commutative diagram of fibrations.}
\GSfield{Reconstruction}{The reflection argument identifies the handwritten \(S\Lambda\). The vertical maps are not all inclusions.}
\GSfield{Fate}{This is the construction of \citet[\S1.3]{arnold1967characteristic}.}
\GSfield{Present status}{As of 9 November 2026: the standard universal Maslov construction.}
\GSfield{Assessment}{literal statement: correct; publication match: exact; verification: read the result; development: annotated; confidence: high}
\end{GSitem}
\begin{GSitem}{1966-108.6}{A}{Maslov cohomology class}
\GSfield{Kind; source}{construction; \GSpage{P000166}}
\GSfield{The notes}{The tangent map \(M\to\Lambda(n)\to S^1\) assigns a loop the degree of its final circle map, yielding a class in \(H^1(M,\mathbb Z)\).}
\GSfield{Statement}{For a Lagrangian immersion \(i:M\to\mathbb R^{2n}\), use the ambient trivialization and tangent Gauss map \(\tau(x)=di_x(T_xM)\). Then \(\mu_M=(d\circ\tau)^*\eta\in H^1(M;\mathbb Z)\), for the chosen generator \(\eta\) on the circle. Evaluation on a loop is its winding number under \(d\circ\tau\).}
\GSfield{Reconstruction}{The ambient trivialization is available in this phase space. A general symplectic ambient manifold requires the appropriate relative formulation. The doubtful “katsikl” means cocycle/class in this mathematical context.}
\GSfield{Fate}{\citet[\S\S1--2]{arnold1967characteristic} develops the class already implicit in occurrence 1964-21.3; submission preceded the talk and publication followed it.}
\GSfield{Present status}{As of 9 November 2026: established, with an orientation-dependent sign convention.}
\GSfield{Assessment}{literal statement: correct; publication match: exact; verification: read the result; development: annotated; confidence: high}
\end{GSitem}
\begin{GSitem}{1966-108.7}{A}{Morse index and quantization}
\GSfield{Kind; source}{remark; \GSpage{P000166}}
\GSfield{The notes}{The Maslov index occurs in asymptotic formulae; the Morse index is obtained by counting conjugate points of a geodesic. Another geometric definition by Maslov is mentioned.}
\GSfield{Statement}{For a nondegenerate fixed-endpoint geodesic segment, the index of the second variation equals the sum of interior conjugate-point multiplicities. The relative Maslov index of the path of Jacobi-data Lagrangian subspaces gives the same integer under consistent sign and endpoint conventions. In semiclassical quantization it measures caustic phase changes.}
\GSfield{Reconstruction}{A relative path index is required; an arbitrary closed-loop winding number is not automatically a Morse index. The alternative definition is not actually recorded.}
\GSfield{Fate}{These applications belong to \citet[\S\S2--3]{arnold1967characteristic}; the Morse index theorem itself is older.}
\GSfield{Present status}{As of 9 November 2026: established with multiplicities and endpoint conventions.}
\GSfield{Assessment}{literal statement: correct with missing hypotheses; publication match: partial; verification: read the result; development: annotated; confidence: high}
\end{GSitem}

\GStalk{109}{1966-10-17}{I. M. Gelfand}{P000166-P000168}{Recovering a compact group from tensor intertwiners}
The actual report is the upper-middle portion of P000167. Gelfand explains reconstruction from irreducible representation spaces and their tensor intertwiners. Multiplicities alone and arbitrary matrix solutions require corrections. The top of the same page separately says “Minlos, Lecture 5: Phase transitions”, with no mathematical content following that heading; no additional MCCME record or phase-transition theorem is invented. Nirenberg's heat-kernel report below belongs to record 110.

\begin{GSitem}{1966-109.1}{E}{Tensor decomposition and intertwiner data}
\GSfield{Kind; source}{construction; \GSpage{P000167}}
\GSfield{The notes}{For a finite or compact \(G\), irreducible finite-dimensional spaces \(A_i\) are given, together with \(A_i\otimes A_k\cong\sum A_j\) and component maps \(d_{ik}^j\).}
\GSfield{Statement}{For a compact topological group, finite-dimensional continuous unitary representations are completely reducible. Every tensor product \(V_i\otimes V_k\) decomposes as a finite direct sum of irreducibles with multiplicity spaces, and a choice of unitary decomposition gives intertwining maps. A general compact group can have infinitely many irreducible isomorphism classes, even though each individual tensor product is finite dimensional.}
\GSfield{Reconstruction}{The notebook's finite list \(A_1,\ldots,A_n\) is appropriate to a finite group or a selected finite list, not to all irreducibles of an arbitrary compact group. The phrase “each 1500 times” illustrates a sufficiently large ambient direct sum rather than a universal multiplicity bound.}
\GSfield{Fate}{This is classical compact-group representation theory, the concrete category recalled by \citet[summary]{doplicherroberts1989duality}.}
\GSfield{Present status}{As of 9 November 2026: standard; multiplicity spaces and actual intertwining maps are different levels of information.}
\GSfield{Assessment}{literal statement: correct with missing hypotheses; publication match: exact; verification: publisher metadata checked; development: annotated; confidence: high}
\end{GSitem}
\begin{GSitem}{1966-109.2}{E}{Fusion multiplicities do not determine even a finite group}
\GSfield{Kind; source}{claim; \GSpage{P000167}}
\GSfield{The notes}{The crossed-out and inserted line is transcribed as saying that properties of the tensor multiplicities recover a finite group, although not a compact group.}
\GSfield{Statement}{If only the fusion coefficients and dimensions of irreducibles are meant, the assertion is false. The dihedral group of order eight and the quaternion group each have four one-dimensional complex irreducibles forming \(C_2\times C_2\), and one two-dimensional irreducible \(W\), with \(\chi\otimes W=W\) and \(W\otimes W=\bigoplus_\chi\chi\). They are not isomorphic: their numbers of involutions differ.}
\GSfield{Reconstruction}{This elementary countercheck \ours\ is supplied here. It tests precisely the reading “multiplicities alone”; extra tensor coherence or intertwiner data could change the claim. The overwritten sentence remains a textual ambiguity.}
\GSfield{Fate}{Classical Tannaka reconstruction uses more than a representation ring. The required categorical data are explicit in the comparison drawn by \citet[summary]{doplicherroberts1989duality}.}
\GSfield{Present status}{As of 9 November 2026: false on the multiplicity-only reading. No unsolved finite-group reconstruction problem is inferred.}
\GSfield{Assessment}{literal statement: false as written; publication match: partial; verification: publisher metadata checked; development: annotated; confidence: high}
\end{GSitem}
\begin{GSitem}{1966-109.3}{E}{Tannaka reconstruction requires unit and reality conditions}
\GSfield{Kind; source}{theorem; \GSpage{P000167}}
\GSfield{The notes}{Replace all representation matrices \(T_i(g)\) by arbitrary linear maps \(C_i\) satisfying the displayed tensor intertwining equations; the notes claim every solution comes from a group element.}
\GSfield{Statement}{Let \(\mathcal R(G)\) be the rigid symmetric category of finite-dimensional continuous unitary representations of compact \(G\), and \(F\) its forgetful functor to Hilbert spaces. The group of unitary tensor-natural automorphisms of \(F\), with \(C_{\mathbf1}=1\), is naturally \(G\), with its compact topology. The naturality conditions include all intertwiners and dual compatibility.}
\GSfield{Reconstruction}{Unqualified equations admit the zero family unless unit normalization is imposed. Merely complex linear tensor automorphisms need not be unitary; even for \(S^1\) the character family \(C_k=z^k\), \(z\in\mathbb C^\times\), exhibits the missing compact reality condition.}
\GSfield{Fate}{The classical concrete duality predates the seminar. \citet[summary and references 18, 27]{doplicherroberts1989duality} explicitly contrasts it with their later abstract symmetric \(C^*\)-category reconstruction.}
\GSfield{Present status}{As of 9 November 2026: established with unitary tensor-natural data; the literal unrestricted matrix formulation is false.}
\GSfield{Assessment}{literal statement: false as written; publication match: partial; verification: publisher metadata checked; development: annotated; confidence: high}
\end{GSitem}
\begin{GSitem}{1966-109.4}{F}{An unspecified extension to SL(2,R)}
\GSfield{Kind; source}{remark; \GSpage{P000167}}
\GSfield{The notes}{The report ends “There is a generalization to \(SL(2,\mathbb R)\)”.}
\GSfield{Statement}{No class of representations, topology, dual object, admissible tensor operations or reconstruction statement is given. Consequently the particular noncompact duality theorem intended cannot be identified from this sentence.}
\GSfield{Reconstruction}{Finite-dimensional representations alone, all unitary representations, and operator-algebraic duality lead to materially different statements. A general later noncompact duality theorem cannot be treated as an exact match.}
\GSfield{Fate}{The later literature cited in \citet[references 28--29]{doplicherroberts1989duality} supplies possible noncompact leads, not an identification of this report.}
\GSfield{Present status}{As of 9 November 2026: unresolved identification rather than a demonstrated open problem.}
\GSfield{Assessment}{literal statement: not determinable; publication match: thematic only; verification: unresolved lead; development: annotated; confidence: low}
\end{GSitem}

\GStalk{110}{1966-10-17}{Louis Nirenberg}{P000166-P000168}{Varadhan's short-time heat-kernel asymptotic}
The lower portion of P000167 records Nirenberg reporting Varadhan's result: the leading exponential of the heat kernel measures squared geodesic distance. The handwritten general formula omits the factor four that is present in its own Euclidean example. A final allusion to elliptic eigenvalues contains no identifiable theorem. The close publication is \citet{varadhan1967fundamental}, published in May 1967.

\begin{GSitem}{1966-110.1}{B}{Short-time heat kernel and intrinsic distance}
\GSfield{Kind; source}{theorem; \GSpage{P000167}}
\GSfield{The notes}{For \(\partial_tu=\sum a^{ij}(x)u_{x_ix_j}\) with point-source initial data, the notes give \(u=\exp[-g(x,y)/t+o(t^{-1})]\) and identify \(g=d^2\) for the metric generated by \(a^{ij}\).}
\GSfield{Statement}{For a smooth symmetric uniformly positive definite coefficient matrix \(a(x)\), with standard boundedness assumptions ensuring the positive fundamental solution on \(\mathbb R^n\), put \(g_x(v,v)=v^\top a(x)^{-1}v\). Then \(\lim_{t\downarrow0}-4t\log p(t,x,y)=d_g(x,y)^2\). Equivalently \(p(t,x,y)=\exp[-d_g(x,y)^2/(4t)+o(t^{-1})]\) for fixed \(x,y\).}
\GSfield{Reconstruction}{The inverse matrix defines the metric. Uniform ellipticity and a standard global coefficient setting are supplied as sufficient hypotheses. The normalization is for the displayed generator with no factor \(1/2\).}
\GSfield{Fate}{\citet[pp.~431--455]{varadhan1967fundamental} is the exact subsequent paper by Varadhan, reported here by Nirenberg. Publisher metadata was checked; the normalization is independently fixed by the notebook's Gaussian.}
\GSfield{Present status}{As of 9 November 2026: established. The literal identification of \(g\) in the exponent is off by a factor four.}
\GSfield{Assessment}{literal statement: false as written; publication match: exact; verification: publisher metadata checked; development: annotated; confidence: high}
\end{GSitem}
\begin{GSitem}{1966-110.2}{E}{Gaussian normalization and subleading powers of time}
\GSfield{Kind; source}{theorem; \GSpage{P000167}}
\GSfield{The notes}{The Euclidean example is \(Ct^{-n/2}\exp[-|x-y|^2/(4t)]\); the polynomial factor is said to enter the \(o(t^{-1})\) term.}
\GSfield{Statement}{For \(\partial_tu=\Delta u\), \(p(t,x,y)=(4\pi t)^{-n/2}\exp[-|x-y|^2/(4t)]\). Since \(t\log t\to0\), \(\log(t^{-n/2})=o(t^{-1})\), so this prefactor does not affect \(-4t\log p\).}
\GSfield{Reconstruction}{The constant is supplied from mass-one normalization. Direct differentiation and Gaussian integration independently check the formula and its factor four.}
\GSfield{Fate}{This classical example is the normalization check for the general result \citep{varadhan1967fundamental}.}
\GSfield{Present status}{As of 9 November 2026: standard and correct.}
\GSfield{Assessment}{literal statement: correct; publication match: exact; verification: publisher metadata checked; development: annotated; confidence: high}
\end{GSitem}
\begin{GSitem}{1966-110.3}{F}{Unspecified elliptic-eigenvalue analogy}
\GSfield{Kind; source}{remark; \GSpage{P000167}}
\GSfield{The notes}{The last sentence says that elliptic eigenvalues of the Laplacian in a domain can be treated “in the same spirit”.}
\GSfield{Statement}{The sentence supplies no domain assumptions, asymptotic parameter, eigenvalue regime or proposed formula. It cannot uniquely identify Weyl asymptotics, a small-noise principal-eigenvalue estimate or another spectral result.}
\GSfield{Reconstruction}{No precise theorem has been inserted from the speaker's presumed intentions.}
\GSfield{Fate}{The heat-kernel paper \citep{varadhan1967fundamental} is relevant context only; it does not identify this unspecific extra sentence.}
\GSfield{Present status}{As of 9 November 2026: unidentified allusion, not a definite open problem.}
\GSfield{Assessment}{literal statement: not determinable; publication match: thematic only; verification: unresolved lead; development: annotated; confidence: low}
\end{GSitem}

\GStalk{111}{1966-10-24}{M. I. Vishik}{P000166-P000168}{Pseudodifferential operators: two lecture headings}
At the top of P000168, Vishik's first lecture has only “I. Kohn--Nirenberg definition” and “II. Boundedness theorem”. No hypotheses, formulas or proof survive. A separate unindexed heading attributes a Sturm--Liouville topic to a professor whose surname is uncertain (possibly Lorch); it has no mathematical statement either. The entries below distinguish contextual reconstructions from actual recorded assertions.

\begin{GSitem}{1966-111.1}{F}{Kohn--Nirenberg definition: heading only}
\GSfield{Kind; source}{definition; \GSpage{P000168}}
\GSfield{The notes}{The first outline entry names a Kohn--Nirenberg definition, without a formula or symbol class.}
\GSfield{Statement}{A contextual modern model is \(a(x,D)u(x)=(2\pi)^{-n}\int e^{ix\cdot\xi}a(x,\xi)\widehat u(\xi)\,d\xi\), initially for Schwartz \(u\), with \(a\) satisfying specified symbol estimates. The heading alone does not identify which symbol estimates were taught.}
\GSfield{Reconstruction}{The formula is explanatory reconstruction, not a transcription of a missing lecture. In particular no \(S^m_{\rho,\delta}\) class is attributed to the speaker.}
\GSfield{Fate}{\citet{kohnnirenberg1965algebra} is the natural earlier source, but only its publisher metadata and the contextual match were verified.}
\GSfield{Present status}{As of 9 November 2026: the named framework is identifiable, while the precise definition used in this lecture is not.}
\GSfield{Assessment}{literal statement: not determinable; publication match: thematic only; verification: publisher metadata checked; development: annotated; confidence: medium}
\end{GSitem}
\begin{GSitem}{1966-111.2}{F}{Boundedness theorem: heading only}
\GSfield{Kind; source}{theorem; \GSpage{P000168}}
\GSfield{The notes}{The second outline entry is merely “Theorem on boundedness”.}
\GSfield{Statement}{One standard theorem is \(a(x,D):H^s\to H^{s-m}\) for a properly supported symbol \(a\in S^m_{1,0}\), with the usual uniform derivative estimates and localization on a manifold. The heading does not specify whether this, an order-zero \(L^2\) theorem, or a different boundedness criterion was intended.}
\GSfield{Reconstruction}{This example supplies context only and is not counted as a verified formulation of Vishik's theorem. Earlier occurrences 1965-81.8--9 discuss Sobolev order and angular regularity, but do not fill this gap.}
\GSfield{Fate}{The relevant pre-lecture sources include \citet{kohnnirenberg1965algebra,hormander1965pseudo}; their existence does not determine the missing hypotheses.}
\GSfield{Present status}{As of 9 November 2026: no unique theorem can be assigned to the heading.}
\GSfield{Assessment}{literal statement: not determinable; publication match: thematic only; verification: publisher metadata checked; development: annotated; confidence: medium}
\end{GSitem}

\GStalk{112}{1966-10-31}{M. I. Vishik}{P000166-P000170}{Composition and formal symbols}
Vishik's second lecture begins on P000168 and ends in the upper half of P000169. Unlike the preceding outline, it supplies the composition expansion, its Fourier/Taylor proof, formal asymptotic symbols, uniqueness modulo smoothing and realization of a formal symbol. These are expository parts of the calculus already published in 1965 by Kohn--Nirenberg and H\"ormander \citep{kohnnirenberg1965algebra,hormander1965pseudo}. Oleinik begins below and belongs to record 113.

\begin{GSitem}{1966-112.1}{E}{Composition expansion with a controlled remainder}
\GSfield{Kind; source}{theorem; \GSpage{P000168}}
\GSfield{The notes}{\(AB=C_0+\cdots+C_{s-1}+T_s\), with \(\sigma(C_k)=\sum_{|\alpha|=k}(\alpha!)^{-1}\partial_\xi^\alpha a\,D_x^\alpha b\); \(T_s\) has order \(\sigma+\lambda-s\).}
\GSfield{Statement}{For properly supported classical or \(S^{m}_{1,0}\) operators \(A,B\) of orders \(m_1,m_2\), the Kohn--Nirenberg product symbol satisfies \(a\# b-\sum_{|\alpha|<N}(\alpha!)^{-1}\partial_\xi^\alpha a\,D_x^\alpha b\in S^{m_1+m_2-N}_{1,0}\), with \(D_x=(1/i)\partial_x\). In particular the principal symbol is \(a_{m_1}b_{m_2}\).}
\GSfield{Reconstruction}{The symbol class, support/localization conditions and Fourier convention are supplied. Symbol estimates are essential; a formal Taylor polynomial alone does not bound the remainder.}
\GSfield{Fate}{This is the already published calculus of \citet{kohnnirenberg1965algebra,hormander1965pseudo}. The notebook's formula was checked directly against P000168; external verification here is bibliographic rather than a full rereading of those papers.}
\GSfield{Present status}{As of 9 November 2026: standard pseudodifferential composition.}
\GSfield{Assessment}{literal statement: correct with missing hypotheses; publication match: exact; verification: publisher metadata checked; development: annotated; confidence: high}
\end{GSitem}
\begin{GSitem}{1966-112.2}{E}{Fourier-space proof of the composition formula}
\GSfield{Kind; source}{method; \GSpage{P000168-P000169}}
\GSfield{The notes}{The localized Fourier kernels \(\widetilde h(\xi-\eta,\eta)\), \(\widetilde g(\eta-\tau,\tau)\) are composed; Taylor expansion of the second variable of \(\widetilde h\) at \(\eta=\tau\) gives derivatives and convolution products.}
\GSfield{Statement}{With compact spatial localization and the preceding symbol estimates, write the composition as an oscillatory double integral. Taylor-expand \(a(x,\xi+\eta)\) in \(\eta\); transfer each \(\eta^\alpha\) to \(D_x^\alpha b\) by Fourier inversion. The integral Taylor remainder and integration by parts yield the order \(m_1+m_2-N\) estimate.}
\GSfield{Reconstruction}{The notebook gives the algebraic part of this proof but omits remainder estimates and convergence regularization. They are hypotheses and standard analytic steps required for the theorem, not additional recorded calculations.}
\GSfield{Fate}{The method is part of the original symbol calculus \citep{kohnnirenberg1965algebra,hormander1965pseudo}.}
\GSfield{Present status}{As of 9 November 2026: established; the recorded proof is a sketch of the composition theorem.}
\GSfield{Assessment}{literal statement: correct with missing hypotheses; publication match: partial; verification: publisher metadata checked; development: annotated; confidence: high}
\end{GSitem}
\begin{GSitem}{1966-112.3}{E}{Formal asymptotic symbol}
\GSfield{Kind; source}{definition; \GSpage{P000169}}
\GSfield{The notes}{For homogeneous \(A_k\) of orders \(\tau_k\searrow-\infty\), \(A\sim\sum A_k\) means the remainder after \(N\) terms has order below \(\tau_N\); the formal sum of \(a_k\) is called the symbol.}
\GSfield{Statement}{In the classical calculus a full formal symbol is a sequence \(a_{m-j}(x,\xi)\), homogeneous of degree \(m-j\) for \(|\xi|\ge1\), such that \(a-\sum_{j<N}\chi(\xi)a_{m-j}\in S^{m-N}_{1,0}\) for every \(N\), with a cutoff near zero. More generally one may specify a decreasing order sequence tending to \(-\infty\) with corresponding remainder bounds.}
\GSfield{Reconstruction}{This distinguishes a formal homogeneous expansion from an exact total symbol. The notebook's weaker “order below \(\tau_N\)” does not itself specify the sharper conventional next-order bound.}
\GSfield{Fate}{Asymptotic symbols were part of the 1965 calculus \citep{kohnnirenberg1965algebra,hormander1965pseudo}.}
\GSfield{Present status}{As of 9 November 2026: standard after a remainder convention is fixed.}
\GSfield{Assessment}{literal statement: correct with missing hypotheses; publication match: exact; verification: publisher metadata checked; development: annotated; confidence: high}
\end{GSitem}
\begin{GSitem}{1966-112.4}{E}{Zero formal symbol and smoothing operators}
\GSfield{Kind; source}{theorem; \GSpage{P000169}}
\GSfield{The notes}{“Symbol zero” means \(A\) has order \(-\infty\).}
\GSfield{Statement}{For a properly supported classical operator, vanishing of every homogeneous formal-symbol component is equivalent to membership in \(\Psi^{-\infty}\). Such an operator has a smooth kernel and maps compactly supported distributions to smooth functions; on a closed manifold it improves every Sobolev order.}
\GSfield{Reconstruction}{A zero principal symbol merely lowers order; it is the entire formal expansion that must vanish.}
\GSfield{Fate}{This is the smoothing ideal of the established calculus \citep{kohnnirenberg1965algebra,hormander1965pseudo}.}
\GSfield{Present status}{As of 9 November 2026: standard and correct with the meaning of “symbol” fixed.}
\GSfield{Assessment}{literal statement: correct; publication match: exact; verification: publisher metadata checked; development: annotated; confidence: high}
\end{GSitem}
\begin{GSitem}{1966-112.5}{E}{Uniqueness of the formal homogeneous expansion}
\GSfield{Kind; source}{theorem; \GSpage{P000169}}
\GSfield{The notes}{The theorem says an operator's symbol is uniquely determined.}
\GSfield{Statement}{For a fixed coordinate system and quantization, the homogeneous components of a classical symbol are uniquely determined by the operator modulo smoothing. Two distinct homogeneous leading terms cannot differ by a lower-order symbol. An exact total symbol depends on localization and quantization, and transforms nontrivially under coordinate changes.}
\GSfield{Reconstruction}{The uniqueness is of the formal expansion in a fixed calculus, not of a coordinate-invariant scalar full symbol.}
\GSfield{Fate}{This is the uniqueness convention in the calculus of \citet{kohnnirenberg1965algebra,hormander1965pseudo}.}
\GSfield{Present status}{As of 9 November 2026: standard after conventions are fixed.}
\GSfield{Assessment}{literal statement: correct with missing hypotheses; publication match: exact; verification: publisher metadata checked; development: annotated; confidence: high}
\end{GSitem}
\begin{GSitem}{1966-112.6}{E}{Realization of a formal symbol}
\GSfield{Kind; source}{theorem; \GSpage{P000169}}
\GSfield{The notes}{The notes assert that for every symbol there exists an operator.}
\GSfield{Statement}{Given smooth homogeneous components \(a_{m-j}\) satisfying the symbol bounds on each compact coordinate set, an asymptotic summation using cutoffs at successively larger frequencies constructs \(a\in S^m_{1,0}\) with that formal expansion. Quantization and localization give an operator, unique modulo smoothing for the prescribed expansion.}
\GSfield{Reconstruction}{“Every symbol” means an admissible formal classical symbol, not an arbitrary collection without smoothness or local estimates. The cutoff construction is a Borel-type asymptotic summation.}
\GSfield{Fate}{This realization theorem is already part of the 1965 calculus \citep{kohnnirenberg1965algebra,hormander1965pseudo}.}
\GSfield{Present status}{As of 9 November 2026: established.}
\GSfield{Assessment}{literal statement: correct with missing hypotheses; publication match: exact; verification: publisher metadata checked; development: annotated; confidence: high}
\end{GSitem}
\begin{GSitem}{1966-112.7}{E}{Filtered algebra of pseudodifferential operators}
\GSfield{Kind; source}{theorem; \GSpage{P000169}}
\GSfield{The notes}{Pseudodifferential operators form an algebra and their formal product is \(\sum_\alpha(\alpha!)^{-1}\partial_\xi^\alpha\sigma(A)D_x^\alpha\sigma(B)\).}
\GSfield{Statement}{The properly supported classical operators satisfy \(\Psi^{m_1}\Psi^{m_2}\subset\Psi^{m_1+m_2}\), with the written asymptotic symbol product. The smoothing operators form a two-sided ideal, so this product descends to complete formal symbols modulo smoothing.}
\GSfield{Reconstruction}{The conclusion is the algebraic organization of the preceding composition and realization theorems; it requires the same support and symbol-class hypotheses.}
\GSfield{Fate}{This is the central already published result named in the title of \citet{kohnnirenberg1965algebra}.}
\GSfield{Present status}{As of 9 November 2026: standard calculus; no undeveloped research proposal is present.}
\GSfield{Assessment}{literal statement: correct with missing hypotheses; publication match: exact; verification: publisher metadata checked; development: annotated; confidence: high}
\end{GSitem}

\GStalk{113}{1966-10-31}{O. A. Oleinik}{P000168-P000172}{The unsteady Prandtl boundary layer}
Oleinik begins in lower P000169 and finishes P000170. She contrasts viscous and ideal-fluid boundary conditions, writes the Prandtl system and data, and announces uniqueness, short-time existence, existence near the inflow for all time, and stability. The monotonicity and compatibility assumptions needed for her theorems are absent from the notebook. The announcement had already appeared in 1966 \citep{oleinik1966system}; the detailed later treatment is \citet[Lecture II, Theorems 4--5]{oleinik1968boundary}.

\begin{GSitem}{1966-113.1}{E}{Planar incompressible Navier--Stokes and Euler systems}
\GSfield{Kind; source}{definition; \GSpage{P000169}}
\GSfield{The notes}{The notes display both momentum equations for \(u,v\), the pressure gradient, \(\nu\Delta u,\nu\Delta v\), and \(u_x+v_y=0\); \(\nu=0\) is called Euler's ideal-fluid system.}
\GSfield{Statement}{For constant density \(\rho>0\) and kinematic viscosity \(\nu>0\), a smooth planar velocity \(w=(u,v)\) obeys \(\partial_tw+(w\cdot\nabla)w=-\rho^{-1}\nabla p+\nu\Delta w\), \(\nabla\cdot w=0\). The formal zero-viscosity equations are the incompressible Euler system. This algebraic substitution alone is not a convergence theorem.}
\GSfield{Reconstruction}{The vector notation condenses the two correctly written component equations. Boundary and initial conditions are a separate part of either problem.}
\GSfield{Fate}{These are the classical governing equations used in \citet[\S1]{oleinik1968boundary}.}
\GSfield{Present status}{As of 9 November 2026: standard equations; the boundary-condition issue is the next item.}
\GSfield{Assessment}{literal statement: correct; publication match: exact; verification: read the result; development: annotated; confidence: high}
\end{GSitem}
\begin{GSitem}{1966-113.2}{E}{Why vanishing viscosity is singular at a wall}
\GSfield{Kind; source}{remark; \GSpage{P000169}}
\GSfield{The notes}{The incompressible planar Navier--Stokes equations are written. Setting viscosity to zero gives Euler, but the viscous terms cannot simply be discarded because no-slip and ideal-fluid boundary conditions differ.}
\GSfield{Statement}{At a stationary impermeable wall, viscous no-slip imposes both velocity components zero, whereas impermeable Euler flow imposes only the normal component zero. If the limiting outer tangential velocity is nonzero, convergence to that Euler field cannot be uniform up to the wall. A thin transition region is therefore required in the classical matched-asymptotic picture.}
\GSfield{Reconstruction}{This observation does not prove a Navier--Stokes-to-Prandtl limit, a boundary-layer thickness theorem, or formation of a layer for all initial data.}
\GSfield{Fate}{This is the classical Prandtl motivation reviewed by \citet[\S1]{oleinik1968boundary}.}
\GSfield{Present status}{As of 9 November 2026: the boundary-condition mismatch is established; a rigorous inviscid limit requires additional hypotheses.}
\GSfield{Assessment}{literal statement: correct; publication match: exact; verification: read the result; development: annotated; confidence: high}
\end{GSitem}
\begin{GSitem}{1966-113.3}{E}{The Prandtl system and outer pressure}
\GSfield{Kind; source}{construction; \GSpage{P000170}}
\GSfield{The notes}{The tangential equation retains \(\nu u_{yy}\), continuity remains, and the pressure is a prescribed function \(p(t,x)\) constant in the normal direction. Although \(p_y=0\) is crossed out in the display, the prose explicitly retains its meaning.}
\GSfield{Statement}{In boundary-layer coordinates at a flat wall, the Prandtl model is \(u_t+uu_x+vu_y=\nu u_{yy}-\rho^{-1}p_x(t,x)\), \(u_x+v_y=0\). Matching to the outer tangential Euler velocity \(U(t,x)\) imposes \(U_t+UU_x=-\rho^{-1}p_x\). The formal scaling \(y=\sqrt{\nu}\,Y\), with the normal velocity scaled accordingly, explains the retained normal diffusion.}
\GSfield{Reconstruction}{The outer Bernoulli equation and the distinction between physical and rescaled variables are supplied. The formal derivation is not a proved approximation theorem.}
\GSfield{Fate}{The displayed model is the one studied in \citet[p.~751, equations (1)--(3)]{oleinik1966system} and \citet[\S1]{oleinik1968boundary}.}
\GSfield{Present status}{As of 9 November 2026: standard boundary-layer equations.}
\GSfield{Assessment}{literal statement: correct with missing hypotheses; publication match: exact; verification: read the result; development: annotated; confidence: high}
\end{GSitem}
\begin{GSitem}{1966-113.4}{E}{Correctly posed initial and inflow data}
\GSfield{Kind; source}{problem; \GSpage{P000170}}
\GSfield{The notes}{The wall conditions are \(u=0\), \(v=v_0(x,t)\); the notebook writes inconsistent arguments in \(u|_{x=0}=u_0(x)\) and \(u|_{t=0}=u_0(x,t)\), and requires \(u\to U(t,x)\) as \(y\to\infty\).}
\GSfield{Statement}{On \(0<x<X\), \(y>0\), \(t>0\), prescribe \(u(t,x,0)=0\), \(v(t,x,0)=v_0(t,x)\), initial profile \(u(0,x,y)=u_0(x,y)\), inflow profile \(u(t,0,y)=u_1(t,y)\), and \(\lim_{y\to\infty}u(t,x,y)=U(t,x)\). These data must agree at their intersections, satisfy the wall compatibility equations, and match \(U\) at infinity.}
\GSfield{Reconstruction}{Distinct \(u_0,u_1\) and their variables are reconstructed from the PDE and checked in Oleinik's paper. The notebook's variable mistakes are faithful errors, not proof of a different problem.}
\GSfield{Fate}{The correctly indexed data are explicit in \citet[p.~751]{oleinik1966system}; the note was received 1 March 1966 and published before this October report.}
\GSfield{Present status}{As of 9 November 2026: the standard initial/inflow Prandtl problem under a specified solution class.}
\GSfield{Assessment}{literal statement: false as written; publication match: exact; verification: read the result; development: annotated; confidence: high}
\end{GSitem}
\begin{GSitem}{1966-113.5}{E}{Local existence and uniqueness in the monotone class}
\GSfield{Kind; source}{theorem; \GSpage{P000170}}
\GSfield{The notes}{The solution is said to be unique if it exists, and to exist on \(0\le t\le T\) over the whole body; separation may occur later but not instantaneously.}
\GSfield{Statement}{A sufficient classical setting is \(U>c>0\), smooth compatible data with \(u_0,u_1>0\) for \(y>0\), \(\partial_yu_0,\partial_yu_1>0\), convergence to \(U\) at infinity, and two-sided tail bounds \(c_1(U-u_j)\le\partial_yu_j\le c_2(U-u_j)\), together with the bounded weighted derivatives and corner compatibility in Oleinik's Theorem 4. There is then a unique classical monotone solution for sufficiently short time on a finite streamwise interval.}
\GSfield{Reconstruction}{The notebook omits this crucial monotone solution class. The complete derivative and compatibility package is specified by the cited theorem rather than an assertion about all smooth data. The Crocco transformation uses monotonicity to replace the normal coordinate by the velocity.}
\GSfield{Fate}{\citet[pp.~751--754]{oleinik1966system} already announced this theory; \citet[Theorem 4, pp.~28--29]{oleinik1968boundary} supplies its detailed sufficient assumptions.}
\GSfield{Present status}{As of 9 November 2026: established in the monotone class; no arbitrary-data well-posedness statement is valid.}
\GSfield{Assessment}{literal statement: correct with missing hypotheses; publication match: partial; verification: read the result; development: annotated; confidence: high}
\end{GSitem}
\begin{GSitem}{1966-113.6}{E}{All-time existence on a short inflow interval}
\GSfield{Kind; source}{theorem; \GSpage{P000170}}
\GSfield{The notes}{Near the beginning of the flow around the body, the solution exists for all time. The inequality \(p_x\le0\) appears separately beneath this statement.}
\GSfield{Statement}{Under the uniform-in-time data bounds and compatibility assumptions of Oleinik's Theorem 5, one can choose a sufficiently small streamwise length \(X_1>0\) so that a monotone classical solution exists for all \(t\ge0\), \(0\le x\le X_1\). This differs from global existence on the entire body. A favourable pressure gradient is an additional recorded regime, not a license to omit all other assumptions.}
\GSfield{Reconstruction}{The relation of the isolated inequality to the previous sentence is not grammatically explicit. The detailed cited short-inflow theorem is the verified comparison; no universal all-time theorem for every vessel is claimed.}
\GSfield{Fate}{The result announced in the 1966 theory is proved in \citet[Theorem 5, pp.~29--32]{oleinik1968boundary}.}
\GSfield{Present status}{As of 9 November 2026: established with the stated short-inflow and uniform-data restrictions.}
\GSfield{Assessment}{literal statement: correct with missing hypotheses; publication match: partial; verification: read the result; development: annotated; confidence: high}
\end{GSitem}
\begin{GSitem}{1966-113.7}{F}{The scope of stability of a boundary layer}
\GSfield{Kind; source}{claim; \GSpage{P000170}}
\GSfield{The notes}{The final sentence says a solution is always stable if it exists and is a “nast.” boundary layer; the abbreviation is not safely expanded.}
\GSfield{Statement}{Oleinik-type monotone solutions have stability/continuous-dependence statements in the corresponding controlled class and norms. Existence of a smooth Prandtl solution alone does not imply Sobolev stability: linearization around suitable nonmonotone shear profiles with a nondegenerate critical point has high-frequency instability.}
\GSfield{Reconstruction}{The note gives neither a norm nor the meaning of the qualifying abbreviation. Thus the exact historical stability theorem is not reconstructed. The later counter-scope is included to prevent reading “always” as an unrestricted theorem.}
\GSfield{Fate}{The monotone framework is in \citet{oleinik1968boundary}; the nonmonotone obstruction is the result of \citet[Theorem 1, author version pp.~4--5]{gerardvaretdormy2010illposed}, first circulated in 2009 and published in 2010.}
\GSfield{Present status}{As of 9 November 2026: the literal sentence is underdetermined; general Sobolev stability is false, while qualified monotone theory is established.}
\GSfield{Assessment}{literal statement: not determinable; publication match: partial; verification: read the result; development: annotated; confidence: medium}
\end{GSitem}

\GStalk{114}{1966-11-14}{M. S. Agranovich}{P000170-P000174}{Pseudodifferential heading and nine geometry problems}
P000171 begins “Agranovich, Lecture 3 on pseudodifferential operators”, but no mathematical exposition of that lecture follows. Instead P000171--P000172 contain a numbered differential-geometry problem list, without an explicit speaker attribution for the list. It includes Helgason's spherical support question, entire Monge--Ampere and minimal-graph problems, metric reconstruction, symmetric-space eigenfunctions, infinite volume, two geodesic questions and contact structures. These are associated with this record by manuscript position, not individually attributed to Agranovich. P000173 begins Arnold's following record and is excluded.

\begin{GSitem}{1966-114.1}{F}{Third pseudodifferential lecture: heading only}
\GSfield{Kind; source}{remark; \GSpage{P000171}}
\GSfield{The notes}{The heading names Agranovich's third lecture on pseudodifferential operators; the next written heading is “Problems of differential geometry”.}
\GSfield{Statement}{There is no surviving mathematical statement of the pseudodifferential lecture on these pages.}
\GSfield{Reconstruction}{Neither the title nor the catalogue overlap licenses importing the content of Agranovich's publications or the subsequent Arnold pages.}
\GSfield{Fate}{No particular theorem or new result can be matched to the heading.}
\GSfield{Present status}{As of 9 November 2026: a documented lecture heading without reconstructible mathematical content.}
\GSfield{Assessment}{literal statement: not determinable; publication match: none found; verification: unresolved lead; development: none; confidence: high}
\end{GSitem}
\begin{GSitem}{1966-114.2}{C}{Helgason's support problem for geodesic spheres}
\GSfield{Kind; source}{problem; \GSpage{P000171}; Unattributed differential-geometry discussion}
\GSfield{The notes}{On a complete negatively curved manifold, \(f\in C^\infty_0\); if every sphere enclosing a given ball \(B\) has zero surface integral of \(f\), prove \(f=0\) outside \(B\). The name Helgason is written twice.}
\GSfield{Statement}{Let \(M\) be a complete simply connected smooth Riemannian manifold with negative sectional curvature, \(B\) a closed geodesic ball, and \(f\in C_c^\infty(M)\). If \(\int_{S(y,r)}f\,d\sigma=0\) for every geodesic sphere whose closed ball \(D(y,r)\) contains \(B\), must \(\operatorname{supp}f\subset B\)?}
\GSfield{Reconstruction}{Simple connectivity is supplied from the exact published formulation, making geodesic spheres globally regular by the Hadamard setting. “Sphere containing \(B\)” means its enclosed ball contains \(B\), not that the hypersurface contains the ball as a subset. The transform is spherical integration, not the geodesic ray transform.}
\GSfield{Fate}{This is exactly \citet[Problem 4.1]{quinto2008helgason}. That work proves the real-analytic metric case through its support theorems, including Theorem 5.1; it does not remove analyticity for an arbitrary smooth negatively curved metric.}
\GSfield{Present status}{As of 9 November 2026: the real-analytic case is verified; no resolution of the full smooth-metric statement was located in the targeted searches. The classification concerns this precise remaining regularity gap, not generic integral geometry.}
\GSfield{Assessment}{literal statement: correct with missing hypotheses; publication match: exact; verification: read the result; development: annotated; confidence: medium; problem \hyperref[prob:GP-1966-11]{GP-1966-11}}
\GSfield{Worth developing}{Worth revisiting through smooth local support estimates, but the available evidence does not justify a bounded new theorem in this run.}
\end{GSitem}
\begin{GSitem}{1966-114.3}{E}{Constant-curvature support theorem}
\GSfield{Kind; source}{theorem; \GSpage{P000171}; Unattributed differential-geometry discussion}
\GSfield{The notes}{The notes say Helgason has done the preceding problem for constant curvature.}
\GSfield{Statement}{For real hyperbolic space, a compactly supported smooth function whose surface integrals vanish on all geodesic spheres enclosing a fixed geodesic ball is supported in that ball.}
\GSfield{Reconstruction}{This is the already known special case of the preceding question. The broader real-analytic generalization is later and is kept distinct from what the note says was already proved.}
\GSfield{Fate}{\citet[\S4, discussion after Problem 4.1]{quinto2008helgason} identifies Helgason's earlier constant-curvature proof and its reduction in the disk model.}
\GSfield{Present status}{As of 9 November 2026: established; it does not settle arbitrary variable smooth curvature.}
\GSfield{Assessment}{literal statement: correct; publication match: exact; verification: read the result; development: annotated; confidence: high}
\end{GSitem}
\begin{GSitem}{1966-114.4}{B}{Entire convex solutions of the constant Monge--Ampere equation}
\GSfield{Kind; source}{problem; \GSpage{P000171}; Unattributed differential-geometry discussion}
\GSfield{The notes}{Calabi's problem asks whether a convex function on \(\mathbb R^n\) satisfying \(\det D^2u=1\) must be a polynomial of degree two; the notes mention \(n=2\) and Calabi's \(n\le5\) result.}
\GSfield{Statement}{Every convex entire \(C^2\) solution of \(\det D^2u=1\) on \(\mathbb R^n\) is \(u(x)=\tfrac12x^\top Ax+b\cdot x+c\), where \(A\) is positive definite and \(\det A=1\). The modern theorem extends to convex viscosity solutions.}
\GSfield{Reconstruction}{“Entire” and convexity are essential. The conclusion permits affine terms and arbitrary positive definite Hessian of determinant one; it does not say the Hessian is the identity.}
\GSfield{Fate}{The all-dimensional result is the J\"orgens--Calabi--Pogorelov theorem. \citet[Introduction and Theorem 1.1, pp.~549--550]{caffarellili2003extension} gives the exact modern statement and records the sequence from J\"orgens's \(n=2\) result through Calabi to Pogorelov's all-dimensional resolution. The all-dimensional paper \citep{pogorelov1972improper} was received in June 1971 and published in 1972; its metadata, rather than its full proof, was checked.}
\GSfield{Present status}{As of 9 November 2026: settled affirmatively in every dimension.}
\GSfield{Assessment}{literal statement: correct; publication match: exact; verification: read the result; development: annotated; confidence: high; problem \hyperref[prob:GP-1966-12]{GP-1966-12}}
\end{GSitem}
\begin{GSitem}{1966-114.5}{E}{The known low-dimensional Monge--Ampere cases}
\GSfield{Kind; source}{theorem; \GSpage{P000171}; Unattributed differential-geometry discussion}
\GSfield{The notes}{The Calabi question is followed by “proved for \(n=2\) and for \(n\le5\), Calabi”.}
\GSfield{Statement}{Before this 1966 list, quadratic rigidity for convex entire \(\det D^2u=1\) solutions was known in dimension two by J\"orgens and through dimension five by Calabi. These are special cases of the later all-dimensional theorem, but are separate historical results in the notes.}
\GSfield{Reconstruction}{The notebook's punctuation is not read as attributing J\"orgens's two-dimensional theorem to Calabi.}
\GSfield{Fate}{\citet[Introduction, p.~549]{caffarellili2003extension} identifies the two precedents; \citet[bibliography]{pogorelov1972improper} dates them to 1954 and 1958.}
\GSfield{Present status}{As of 9 November 2026: established before the seminar.}
\GSfield{Assessment}{literal statement: correct; publication match: exact; verification: read the result; development: annotated; confidence: high}
\end{GSitem}
\begin{GSitem}{1966-114.6}{B}{The dimension-dependent Bernstein problem}
\GSfield{Kind; source}{problem; \GSpage{P000171}; Unattributed differential-geometry discussion}
\GSfield{The notes}{For a smooth entire solution of a displayed equation described as the minimal-surface equation, is \(u\) linear? The final parenthesis says “Done” with a reference to De Giorgi. The uncancelled formula visibly has an extra square on the second-derivative sum.}
\GSfield{Statement}{For \(u\in C^\infty(\mathbb R^n)\) satisfying \(\operatorname{div}(\nabla u/\sqrt{1+|\nabla u|^2})=0\), equivalently \((1+|\nabla u|^2)\Delta u-\sum_{i,j}u_i u_j u_{ij}=0\), must \(u\) be affine? The answer is yes for \(n\le7\) and no for every \(n\ge8\).}
\GSfield{Reconstruction}{The notebook's raised square must be deleted to obtain the PDE identified by its prose. The graph lies in ambient \(\mathbb R^{n+1}\), so ambient and domain dimensions must not be confused. “Linear” is interpreted as affine.}
\GSfield{Fate}{\citet[Introduction, Theorem B and \S~IV]{bombieridegiorgigiusti1969minimal} construct the high-dimensional counterexamples and explain the preceding positive results. The original appeared in 1969; the verified Russian translation is 1971, pp.~84--103.}
\GSfield{Present status}{As of 9 November 2026: settled with the sharp dimension split. The unqualified “Done” does not establish the all-dimensional affirmative assertion.}
\GSfield{Assessment}{literal statement: false as written; publication match: exact; verification: read the result; development: annotated; confidence: high; problem \hyperref[prob:GP-1966-13]{GP-1966-13}}
\end{GSitem}
\begin{GSitem}{1966-114.7}{B}{Recovering the metric from a Sobolev pullback isometry}
\GSfield{Kind; source}{problem; \GSpage{P000171-P000172}; Unattributed differential-geometry discussion}
\GSfield{The notes}{The notes give the \(H^1\) inner product \(\int\varphi\psi+\int\langle d\varphi,d\psi\rangle\) and the Dirichlet energy. An \(H^1\) isometry inducing a diffeomorphism between \(M\) and \(M'\) should imply a Riemannian isometry.}
\GSfield{Statement}{Let \(F:M\to M'\) be a smooth diffeomorphism of closed smooth Riemannian manifolds. If the pullback \(u\mapsto u\circ F\) preserves the full \(H^1\) inner product for every smooth \(u,v\), then \(F^*g'=g\). An arbitrary abstract Hilbert-space isometry does not suffice.}
\GSfield{Reconstruction}{“Inducing a diffeomorphism” is interpreted as composition pullback. Elementary verification \ours: testing the inner product against \(1\) shows \(F\) preserves volume. Subtracting the resulting \(L^2\) identity leaves equality of Dirichlet forms. In local coordinates their difference is \(\int q^{ij}\partial_i u\,\partial_jv=0\) for all compactly supported \(u,v\). Integration by parts gives \(\partial_j(q^{ij}\partial_i u)=0\); choose \(du=0\) and an arbitrary Hessian at a point to obtain \(q^{ij}=0\), hence equality of inverse metrics. A second check via equality of the principal symbols of the associated Laplacians gives the same conclusion.}
\GSfield{Fate}{This precise pullback version follows by the displayed elementary argument. The later stronger reconstruction from an order isomorphism intertwining heat semigroups is \citet[Theorem 1.1]{arendt2017diffusion}, first circulated in 2011 and published in 2017; it is a partial, not exact historical match.}
\GSfield{Present status}{As of 9 November 2026: settled under the composition interpretation. No solution of an unspecified abstract Hilbert-space reconstruction problem is claimed.}
\GSfield{Assessment}{literal statement: correct with missing hypotheses; publication match: partial; verification: read the result; development: annotated; confidence: medium; problem \hyperref[prob:GP-1966-14]{GP-1966-14}}
\end{GSitem}
\begin{GSitem}{1966-114.8}{E}{Laplace eigenvalues are invariants, not complete metric data}
\GSfield{Kind; source}{remark; \GSpage{P000172}; Unattributed differential-geometry discussion}
\GSfield{The notes}{The line following the Sobolev question says “Invariants: eigenvalues of the Laplace operator”.}
\GSfield{Statement}{A Riemannian isometry intertwines the Laplace--Beltrami operators and preserves their spectra with multiplicities. Conversely, equality of spectra alone does not imply isometry: an abstract unitary spectral identification need not preserve multiplication, order or the composition action of a map.}
\GSfield{Reconstruction}{The source calls these invariants, not complete invariants. The converse warning prevents the preceding question from being read as an affirmative inverse-spectral theorem.}
\GSfield{Fate}{\citet[Introduction]{arendt2017diffusion} recalls Milnor's 1964 isospectral nonisometric examples and explains the additional order structure used in reconstruction.}
\GSfield{Present status}{As of 9 November 2026: spectral invariance is established; isospectral nonisometric manifolds are known.}
\GSfield{Assessment}{literal statement: correct; publication match: partial; verification: read the result; development: annotated; confidence: high}
\end{GSitem}
\begin{GSitem}{1966-114.9}{E}{Evaluation map from the first Laplace eigenspace}
\GSfield{Kind; source}{construction; \GSpage{P000172}; Unattributed differential-geometry discussion}
\GSfield{The notes}{For compact irreducible symmetric \(M\), \(V\) is the eigenspace of the smallest nonzero Laplace eigenvalue and \(\varphi:M\to V^*\) is written.}
\GSfield{Statement}{For a connected closed Riemannian manifold, the first positive Laplace eigenspace \(V\) is finite dimensional and smooth. Evaluation defines \(\varphi(x)(f)=f(x)\) in \(V^*\), after choosing the \(L^2\) inner product if one wants \(V^*\cong V\). The map is equivariant under isometries. If \(M\) is an irreducible compact symmetric space, its induced symmetric quadratic tensor is an invariant multiple of the metric.}
\GSfield{Reconstruction}{Evaluation is the natural construction inferred from the displayed domain and codomain; it is not written in words. No assertion that the map is injective is recorded, and none is attributed to the source.}
\GSfield{Fate}{This is the standard spectral construction; the link between functions, the Laplacian and the metric is reviewed in \citet[\S1]{arendt2017diffusion}. That reference is contextual rather than an exact source for this one-line symmetric-space item.}
\GSfield{Present status}{As of 9 November 2026: the construction is elementary; the problem or conclusion intended after it is missing.}
\GSfield{Assessment}{literal statement: correct with missing hypotheses; publication match: thematic only; verification: read the result; development: annotated; confidence: medium}
\end{GSitem}
\begin{GSitem}{1966-114.10}{F}{The unrecorded noncompact eigenfunction question}
\GSfield{Kind; source}{question; \GSpage{P000172}; Unattributed differential-geometry discussion}
\GSfield{The notes}{After the evaluation map comes the doubtful line “If not compact”.}
\GSfield{Statement}{No specific noncompact manifold, spectral realization, eigenvalue, function space or desired conclusion is supplied. In general a complete noncompact manifold need not have a first positive \(L^2\) eigenvalue or a finite-dimensional space analogous to the compact \(V\).}
\GSfield{Reconstruction}{The sentence is retained as an unresolved fragment; continuous-spectrum harmonic analysis is only one possible intended continuation.}
\GSfield{Fate}{No exact question could be matched; the compact spectral reconstruction reference is not a solution of this fragment.}
\GSfield{Present status}{As of 9 November 2026: not identifiable as a definite mathematical problem.}
\GSfield{Assessment}{literal statement: not determinable; publication match: none found; verification: unresolved lead; development: annotated; confidence: low}
\end{GSitem}
\begin{GSitem}{1966-114.11}{B}{Infinite volume under positive Ricci curvature}
\GSfield{Kind; source}{question; \GSpage{P000172}; Unattributed differential-geometry discussion}
\GSfield{The notes}{For complete noncompact \(M^3\) with positive Ricci tensor, does \(V\) tend to infinity? A two-dimensional result of Cohn--Vossen is mentioned.}
\GSfield{Statement}{Every complete connected noncompact Riemannian manifold with \(\operatorname{Ric}\ge0\) has infinite Riemannian volume. In particular \(\operatorname{Vol}(B(p,R))\to\infty\) as \(R\to\infty\) in the three-dimensional positive-Ricci case.}
\GSfield{Reconstruction}{The letter \(V\) is read as volume, supported by the curvature context and the Cohn--Vossen reference. Completeness and noncompactness are retained; positive scalar curvature alone would be a different assertion.}
\GSfield{Fate}{The stronger all-dimensional nonnegative-Ricci assertion is stated explicitly in \citet[p.~659, opening summary]{yau1976function}. That first page was opened and visually inspected; no unverified theorem number is assigned.}
\GSfield{Present status}{As of 9 November 2026: settled affirmatively, with weaker curvature than the question requests.}
\GSfield{Assessment}{literal statement: correct with missing hypotheses; publication match: exact; verification: read the result; development: annotated; confidence: high; problem \hyperref[prob:GP-1966-15]{GP-1966-15}}
\end{GSitem}
\begin{GSitem}{1966-114.12}{F}{An ambiguous infinite-geodesic question}
\GSfield{Kind; source}{question; \GSpage{P000172}; Unattributed differential-geometry discussion}
\GSfield{The notes}{For compact \(M^n\), the notes ask for infinitely many \ru{пересек.} geodesics, followed by “actually distinct”; a preceding short word is crossed out.}
\GSfield{Statement}{No unique global proposition can be extracted: closedness is not legibly stated, and “intersecting” does not specify mutual intersections, intersections through a fixed point, endpoints or periodicity. If only distinct local geodesic germs through one point are meant and \(n\ge2\), infinitely many exist immediately by their tangent directions.}
\GSfield{Reconstruction}{The scan was inspected. It does not justify replacing the written abbreviation by “closed” and thereby importing the famous infinitely-many-closed-geodesics conjecture. No problem-index entry is manufactured from that possible interpretation.}
\GSfield{Fate}{Searches for closed-geodesic problems provide only thematic alternatives; none resolves the actual uncertain wording.}
\GSfield{Present status}{As of 9 November 2026: unidentified question because a defining qualifier is missing or doubtful.}
\GSfield{Assessment}{literal statement: not determinable; publication match: none found; verification: unresolved lead; development: annotated; confidence: low}
\end{GSitem}
\begin{GSitem}{1966-114.13}{C}{Smooth type of a manifold with equally long closed geodesics}
\GSfield{Kind; source}{conjecture; \GSpage{P000172}; Unattributed differential-geometry discussion}
\GSfield{The notes}{The next numbered problem asks whether compact simply connected \(M\), all of whose geodesics are closed and have the same length, is diffeomorphic to a symmetric space “of the first type”. A rank-like expression is crossed out.}
\GSfield{Statement}{The standard intended topology question is: if every unit-speed prime geodesic of a closed simply connected Riemannian manifold has the same period, must the manifold be diffeomorphic to a compact rank-one symmetric space, namely a sphere, complex or quaternionic projective space, or the Cayley plane?}
\GSfield{Reconstruction}{Rank one is a contextual identification, not a perfectly legible manuscript word. “Same length” means prime periods, since arbitrary iterates have different lengths. This topology problem is distinct from the Blaschke metric-rigidity conjecture and from the Berger conjecture about equality of periods.}
\GSfield{Fate}{The cohomological restrictions below are established. \citet[definition on p.~1 and Main Theorem on p.~3]{lisong2026zoll} use the stronger Zoll hypothesis that every unit-speed prime geodesic is simple closed and all have the same period. Under that hypothesis and the additional entire-Grauert-tube assumption, they prove metric rigidity of the complex-projective cohomology type, up to scale. Simplicity is not explicit in the notebook, nor is an entire Grauert tube assumed; the result supplies no general diffeomorphism classification under its recorded hypotheses.}
\GSfield{Present status}{As of 9 November 2026: no full smooth classification was verified; the remaining issue is the diffeomorphism type beyond the known cohomology restrictions. The recent additional-hypothesis theorem is not counted as a solution of that issue.}
\GSfield{Assessment}{literal statement: correct with missing hypotheses; publication match: partial; verification: read the result; development: annotated; confidence: medium; problem \hyperref[prob:GP-1966-16]{GP-1966-16}}
\GSfield{Worth developing}{Worth pursuing through precise topological obstructions or restricted geometric classes, but no new bounded contribution is established here.}
\end{GSitem}
\begin{GSitem}{1966-114.14}{E}{Nonround metrics with all geodesics closed}
\GSfield{Kind; source}{claim; \GSpage{P000172}; Unattributed differential-geometry discussion}
\GSfield{The notes}{There is an analytic manifold diffeomorphic to a sphere but not isometric to it, in the preceding closed-geodesic context.}
\GSfield{Statement}{There exist nonround real-analytic metrics on \(S^2\) for which every prime geodesic is simple closed and has the same length. Consequently the conclusion sought in the preceding problem cannot be equality with the round metric merely from the closed-geodesic hypothesis.}
\GSfield{Reconstruction}{“Not isometric” is interpreted relative to a standard round sphere after scale normalization. The source does not supply the metric, and no formula for a new example is invented.}
\GSfield{Fate}{The historical nonround Zoll examples, beginning in 1903, are recalled by \citet[Introduction, p.~2]{lisong2026zoll}; they predate the seminar.}
\GSfield{Present status}{As of 9 November 2026: established and compatible with the proposed diffeomorphism conclusion.}
\GSfield{Assessment}{literal statement: correct; publication match: partial; verification: read the result; development: annotated; confidence: high}
\end{GSitem}
\begin{GSitem}{1966-114.15}{E}{Bott--Samelson cohomology restriction}
\GSfield{Kind; source}{theorem; \GSpage{P000172}; Unattributed differential-geometry discussion}
\GSfield{The notes}{The notes attribute to Bott the result that the preceding manifold has the same cohomological structure as the symmetric model.}
\GSfield{Statement}{For a closed simply connected manifold in the all-geodesics-closed setting, Bott--Samelson theory constrains its integral cohomology ring to that of a compact rank-one symmetric space. In particular it is a truncated polynomial ring of the corresponding sphere or projective-space type; this does not determine the smooth structure.}
\GSfield{Reconstruction}{The coefficient ring and distinction from diffeomorphism classification are made explicit. The attribution is broadened to the standard Bott--Samelson formulation rather than attributing every refinement to one manuscript parenthesis.}
\GSfield{Fate}{The modern formulation and reference to its classical proof are given by \citet[Introduction, pp.~1--2]{lisong2026zoll}; it is an already known partial answer, not the full smooth classification.}
\GSfield{Present status}{As of 9 November 2026: established topological restriction.}
\GSfield{Assessment}{literal statement: correct with missing hypotheses; publication match: partial; verification: read the result; development: annotated; confidence: high}
\end{GSitem}
\begin{GSitem}{1966-114.16}{B}{Contact structures on compact three-manifolds}
\GSfield{Kind; source}{question; \GSpage{P000172}; Unattributed differential-geometry discussion}
\GSfield{The notes}{The last question is whether every compact \(M^3\) has a contact structure.}
\GSfield{Statement}{Every closed orientable smooth three-manifold admits a contact structure. The statement with no orientability assumption is false: a contact plane field locally given by \(\ker\alpha\) orients a three-manifold by \(\alpha\wedge d\alpha\), since replacement by \(f\alpha\) multiplies this volume form by \(f^2>0\). Thus a closed nonorientable example such as \(\mathbb{RP}^2\times S^1\) admits none, even if coorientation is not required.}
\GSfield{Reconstruction}{The orientability correction is necessary for the literal question, not an optional modern strengthening. Compact is read as closed; boundaries require an additional convention. The local-form orientation check was independently verified by expanding \(f\alpha\wedge d(f\alpha)\). The local identity \(\alpha'=f\alpha\Rightarrow\alpha'\wedge d\alpha'=f^2\alpha\wedge d\alpha\) is checked here \ours.}
\GSfield{Fate}{Martinet's affirmative orientable theorem appeared in 1971 \citep{martinet1971contact}; \citet[pp.~165--166]{eliashberg1992contact} explicitly states both the existence theorem and the orientation consequence.}
\GSfield{Present status}{As of 9 November 2026: settled: yes for closed orientable three-manifolds, no for nonorientable ones.}
\GSfield{Assessment}{literal statement: false as written; publication match: exact; verification: read the result; development: annotated; confidence: high; problem \hyperref[prob:GP-1966-17]{GP-1966-17}}
\end{GSitem}

\GStalk{115}{1966-11-21}{V. I. Arnold}{P000173-P000178}{Singularities of smooth mappings I: stability, jets, and Boardman strata}
The first lecture moves from Morse and Whitney models through the perturbation \(z^2+\varepsilon\bar z\), smooth moduli, and Thom\textquotesingle s topological announcement to jets, transversality, and Boardman\textquotesingle s universal strata. It resumes the earlier 1965 record 72, including GP-1965-03. Arnold\textquotesingle s 1968 survey explicitly derives from his 1966 lectures and supplies unusually close publication matches. The MCCME catalogue associates P000172--P000180 with the record; the actual dated lecture starts on P000173 and ends on P000178.

\begin{GSitem}{1966-115.1}{E}{Smooth and topological left--right stability}
\GSfield{Kind; source}{definition; \GSpage{P000173}}
\GSfield{The notes}{The square with maps \(f,g\), source map \(h\), and target map \(k\) defines equivalence; homeomorphisms give topological stability, diffeomorphisms smooth stability. Local stability is phrased on neighborhoods and germs.}
\GSfield{Statement}{For smooth manifolds without boundary, two maps are left--right equivalent if \(k\circ f=g\circ h\) for diffeomorphisms \(h,k\). A map is smoothly stable if its orbit is a neighborhood of it in the Whitney \(C^\infty\) topology. For a germ, representatives may be restricted and the nearby singular point may move. Replacing coordinate changes by homeomorphisms defines topological equivalence and the corresponding stability notion.}
\GSfield{Reconstruction}{The topology and the moving-point convention are supplied; stability of an origin-fixed germ under origin-fixed perturbations would be a different assertion. This repeats 1965-72.1.}
\GSfield{Fate}{Classical framework, reproduced in \citet[\S1, Definition 1 and the definition of local stability]{arnold1968singularities}. The 1968 paper explicitly derives from 1966 lectures.}
\GSfield{Present status}{As of 9 November 2026: Established, with the hypotheses stated here.}
\GSfield{Assessment}{literal statement: correct with missing hypotheses; publication match: exact; verification: read the result; development: annotated; confidence: high}
\end{GSitem}
\begin{GSitem}{1966-115.2}{E}{A stable square and an unstable cube}
\GSfield{Kind; source}{construction; \GSpage{P000173}}
\GSfield{The notes}{The graphs \(y=x^2\) and \(y=x^3\) are labelled stable and unstable respectively.}
\GSfield{Statement}{As real smooth map germs \((\mathbb R,0)\to(\mathbb R,0)\), \(x^2\) is a fold and is stable with a moving critical point. The germ \(x^3\) is unstable: \(x^3+\varepsilon x\) has no critical point if \(\varepsilon>0\), while for \(\varepsilon<0\) its derivative has two simple zeros. These germs therefore cannot be related by source and target diffeomorphisms near the original critical point.}
\GSfield{Reconstruction}{The two perturbations spell out the diagram; they do not constitute a newly posed problem.}
\GSfield{Fate}{These elementary models are \citet[\S1, Examples 1--2]{arnold1968singularities}; the underlying fold is already Morse theory.}
\GSfield{Present status}{As of 9 November 2026: Established, with the hypotheses stated here.}
\GSfield{Assessment}{literal statement: correct; publication match: exact; verification: read the result; development: annotated; confidence: high}
\end{GSitem}
\begin{GSitem}{1966-115.3}{E}{Maximal rank and the Morse local normal forms}
\GSfield{Kind; source}{theorem; \GSpage{P000173-P000174}}
\GSfield{The notes}{The scan reads \ru{Ранг max}; maximal rank is stable. The Morse lemma supplies a nondegenerate quadratic form, followed by the regular form \(y=x_1\).}
\GSfield{Statement}{A maximal-rank germ is a submersion or an immersion and has the corresponding projection or inclusion normal form. For a scalar-valued smooth germ with \(df_0=0\) and nonsingular Hessian there are coordinates with \(f=f(0)+x_1^2+\cdots+x_k^2-x_{k+1}^2-\cdots-x_m^2\). Regular scalar germs have the form \(f=f(0)+x_1\). These are exactly the smoothly stable scalar germ types.}
\GSfield{Reconstruction}{The additive constant, smooth category, and distinction between immersions and submersions complete the shorthand. Earlier occurrence: 1965-72.2.}
\GSfield{Fate}{The standard Morse statement is recalled in \citet[\S1, Examples 3--4 and Theorem 1, p.~3]{arnold1968singularities}; it is exposition of an older theorem.}
\GSfield{Present status}{As of 9 November 2026: Established, with the hypotheses stated here.}
\GSfield{Assessment}{literal statement: correct with missing hypotheses; publication match: exact; verification: read the result; development: annotated; confidence: high}
\end{GSitem}
\begin{GSitem}{1966-115.4}{E}{Global stability and density for functions on a compact manifold}
\GSfield{Kind; source}{theorem; \GSpage{P000173-P000174}}
\GSfield{The notes}{For compact \(M\), stable maps \(M\to\mathbb R\) form an open dense set. Conditions M1 and M2 are nondegenerate critical points and distinct critical values.}
\GSfield{Statement}{Let \(M\) be a compact smooth manifold without boundary. A smooth function \(f:M\to\mathbb R\) is smoothly stable if and only if all critical points are nondegenerate and no two critical points have the same value. Such functions form an open dense subset of \(C^\infty(M,\mathbb R)\).}
\GSfield{Reconstruction}{Boundarylessness and the Whitney topology are explicit here; the notebook gives compactness. Distinct critical values are a global condition, not part of the local Morse lemma.}
\GSfield{Fate}{The compact global Morse criterion is recalled in \citet[\S1, Theorem 1, p.~3]{arnold1968singularities}.}
\GSfield{Present status}{As of 9 November 2026: Established, with the hypotheses stated here.}
\GSfield{Assessment}{literal statement: correct with missing hypotheses; publication match: exact; verification: read the result; development: annotated; confidence: high}
\end{GSitem}
\begin{GSitem}{1966-115.5}{E}{Whitney regular, fold, and cusp local forms}
\GSfield{Kind; source}{theorem; \GSpage{P000174}}
\GSfield{The notes}{The three displayed forms are \((x_1,x_2)\), \((x_1,x_2^2)\), and \((x_1,x_1x_2+x_2^3)\).}
\GSfield{Statement}{For generic smooth maps of surfaces without boundary, every point has one of the following local coordinate forms: a regular point \((u,v)\mapsto(u,v)\); a fold \((u,v)\mapsto(u,v^2)\); or a cusp \((u,v)\mapsto(u,uv+v^3)\). Each is a smoothly stable germ; these exhaust the stable surface-to-surface monogerms.}
\GSfield{Reconstruction}{Genericity concerns derivatives as well as the map. The phrase nondegenerate critical point in the notebook is read in this map-theoretic sense, not as a nonsingular scalar Hessian. Earlier occurrence: 1965-72.3.}
\GSfield{Fate}{Whitney published the plane classification in 1955; \citet[\S1, Theorem 2, pp.~3--4]{arnold1968singularities} reproduces the relevant forms. The original metadata are checked in \citet{whitney1955plane}.}
\GSfield{Present status}{As of 9 November 2026: Established, with the hypotheses stated here.}
\GSfield{Assessment}{literal statement: correct with missing hypotheses; publication match: exact; verification: read the result; development: annotated; confidence: high}
\end{GSitem}
\begin{GSitem}{1966-115.6}{E}{Global stability of maps between surfaces}
\GSfield{Kind; source}{theorem; \GSpage{P000174}}
\GSfield{The notes}{W1: critical points are nondegenerate; W2: fold images meet transversely; W3: cusp images are distinct. Stable mappings are said to be dense.}
\GSfield{Statement}{For a proper smooth map of surfaces, stability requires only regular, fold, and cusp monogerms and stable multigerms over each target point. Concretely the singular image has only ordinary transverse double crossings of fold arcs, each cusp image comes from one singular point, and no other singular branch passes through a cusp image; there are no triple fold intersections. For compact source, such maps form an open dense set.}
\GSfield{Reconstruction}{The excluded triple intersections and cusp--fold coincidences are missing from W2--W3. The literal list alone is not sufficient. Compact source (or properness with the appropriate topology) supplies the global hypothesis.}
\GSfield{Fate}{Whitney-type general position underlies the statement; see \citet[\S1, Theorem 2]{arnold1968singularities}. The precise multigerm formulation is part of \citet[\S1, especially the regular-intersection condition]{mather1969algebras}.}
\GSfield{Present status}{As of 9 November 2026: Established, with the hypotheses stated here.}
\GSfield{Assessment}{literal statement: correct with missing hypotheses; publication match: partial; verification: read the result; development: annotated; confidence: high}
\end{GSitem}
\begin{GSitem}{1966-115.7}{E}{The critical parabola and cusp discriminant}
\GSfield{Kind; source}{construction; \GSpage{P000174-P000175}}
\GSfield{The notes}{The kernel directions are vertical. The cusp is where such a direction is tangent to the critical parabola. For \(x_1>0\) the vertical restriction is monotone; for \(x_1<0\) it has two turning points. The continuation says the merger has odd order.}
\GSfield{Statement}{For \(f(u,v)=(u,uv+v^3)\), \(\det Df=u+3v^2\), so the critical curve is \(u=-3v^2\), its image is \((-3v^2,-2v^3)\), and the kernel along it is generated by \(\partial_v\). Kernel tangency occurs only at the origin. The cubic \(v^3+uv=w\) has three distinct real roots inside the cusp, one outside, and coalescing roots on its discriminant.}
\GSfield{Reconstruction}{\oursComputed{} Coefficients and the discriminant are computed directly from the written normal form. Odd order describes this cubic transition, not a sufficient criterion for stability of arbitrary odd-degree map germs. Earlier occurrence: 1965-72.4.}
\GSfield{Fate}{The Whitney cusp geometry is shown in \citet[\S1, Example 5, p.~4]{arnold1968singularities}.}
\GSfield{Present status}{As of 9 November 2026: Established, with the hypotheses stated here.}
\GSfield{Assessment}{literal statement: correct; publication match: exact; verification: read the result; development: annotated; confidence: high}
\end{GSitem}
\begin{GSitem}{1966-115.8}{A}{Resolving the complex square into folds and three cusps}
\GSfield{Kind; source}{construction; \GSpage{P000175}}
\GSfield{The notes}{The real smooth map \(z\mapsto z^2\) is unstable; the proposed perturbation is \(F_\varepsilon(z)=z^2+\varepsilon\bar z\). The drawings show three cusp points.}
\GSfield{Statement}{For real \(\varepsilon>0\), viewed as a map \(\mathbb R^2\to\mathbb R^2\), the Jacobian of \(F_\varepsilon\) is \(4|z|^2-\varepsilon^2\). Its critical set is \(|z|=\varepsilon/2\); its critical-value curve is \(\varepsilon^2(e^{2it}/4+e^{-it}/2)\). The curve has ordinary cusps at \(e^{3it}=1\); the other critical points are folds. Regular points inside the three-cusped curve have four preimages and exterior regular points have two.}
\GSfield{Reconstruction}{\oursComputed{} The real parameter normalization and critical-set calculation resolve the diagrams. They repeat the already indexed question GP-1965-03, occurrence 1965-72.5; they are not a fresh open problem.}
\GSfield{Fate}{Published explicitly by \citet[\S1, Example 6, pp.~4--5]{arnold1968singularities}; the paper records its 1966 lecture origin, while the earlier notebook occurrence predates this repetition.}
\GSfield{Present status}{As of 9 November 2026: The stated real smooth perturbation and its fold--cusp picture are settled.}
\GSfield{Assessment}{literal statement: correct; publication match: exact; verification: read the result; development: annotated; confidence: high; problem \hyperref[prob:GP-1965-03]{GP-1965-03}}
\end{GSitem}
\begin{GSitem}{1966-115.9}{F}{Large-circle winding and the trefoil caption}
\GSfield{Kind; source}{remark; \GSpage{P000175}}
\GSfield{The notes}{The caption says that large circles make two turns and adds \ru{трилистн. узел} (trefoil knot), next to several immersed planar curves.}
\GSfield{Statement}{For \(R>\varepsilon\), \(F_\varepsilon(Re^{it})\) is nonzero and has winding number two about zero: homotope the conjugate-linear term to zero without crossing zero. A knot type is not defined for this immersed planar image unless an embedding or lift to a three-dimensional space is specified; none is specified in the note.}
\GSfield{Reconstruction}{\ours{} The winding statement is unambiguous; the knot interpretation is not reconstructed from the drawing.}
\GSfield{Fate}{\citet[\S1, Example 6, pp.~4--5]{arnold1968singularities} confirms the double winding, but does not identify a three-dimensional knot in this caption.}
\GSfield{Present status}{As of 9 November 2026: The winding is elementary; the intended trefoil assertion remains unidentified from this record.}
\GSfield{Assessment}{literal statement: not determinable; publication match: partial; verification: read the result; development: annotated; confidence: low}
\end{GSitem}
\begin{GSitem}{1966-115.10}{B}{Failure of density of smoothly stable maps in dimension nine}
\GSfield{Kind; source}{theorem; \GSpage{P000175}}
\GSfield{The notes}{Thom is credited with failure of density of stable germs for \(\mathbb R^9\to\mathbb R^9\), and in higher equidimensions; continuous differential invariants are mentioned.}
\GSfield{Statement}{For every \(n\geq9\), smooth stability is not dense among smooth maps of equidimensional \(n\)-manifolds in general. There exist local configurations persisting under small perturbations whose singular germs have smooth moduli. The mechanism starts with a transverse corank-three point and a nontrivial modulus of a triple of quadratic forms in three variables.}
\GSfield{Reconstruction}{The precise mechanism is developed in the next lecture, P000184. This is a failure of density, not a claim that no stable germ exists in these dimensions. Earlier occurrence: 1965-72.10.}
\GSfield{Fate}{\citet[\S3, Theorem 1]{arnold1968singularities} gives the quadratic-modulus argument. The complete density range was published by \citet[Introduction and \S\S3--5]{mather1971nice} in 1971 (Russian translation 1974).}
\GSfield{Present status}{As of 9 November 2026: Established, with the hypotheses stated here.}
\GSfield{Assessment}{literal statement: correct with missing hypotheses; publication match: exact; verification: read the result; development: annotated; confidence: high}
\end{GSitem}
\begin{GSitem}{1966-115.11}{F}{The unspecified infinity of invariants from dimension twelve}
\GSfield{Kind; source}{claim; \GSpage{P000175}}
\GSfield{The notes}{A separate emphasized line says \ru{нач. с 12, инвар. беск.-много}: beginning with 12 there are infinitely many invariants.}
\GSfield{Statement}{The note does not specify whether this means infinitely many orbit types, arbitrarily many scalar moduli, or functional moduli, nor the class of germs to which the assertion applies. No exact theorem with these quantifiers can therefore be extracted.}
\GSfield{Reconstruction}{These three meanings are not interchangeable. The well-defined dimension-nine result is indexed separately.}
\GSfield{Fate}{The dimension-nine argument and Mather density classification were checked in \citet[\S3]{arnold1968singularities} and \citet{mather1971nice}; neither was used as an exact match for this undefined extra threshold.}
\GSfield{Present status}{As of 9 November 2026: Unresolved identification of the notebook wording; no new open problem is asserted.}
\GSfield{Assessment}{literal statement: not determinable; publication match: thematic only; verification: unresolved lead; development: annotated; confidence: low}
\end{GSitem}
\begin{GSitem}{1966-115.12}{B}{Thom topological discreteness outside infinite codimension}
\GSfield{Kind; source}{claim; \GSpage{P000175}}
\GSfield{The notes}{The space of germs is written \(A=\bigcup_{\gamma\in\Gamma}A_\gamma+A_\infty\), with countable \(\Gamma\), finite-codimension pieces of constant topological type and an infinite-codimension remainder. A constant map of a line is placed in the exceptional set.}
\GSfield{Statement}{The supported assertion is a finite-jet-controlled, countable topological classification of smooth map germs after an exceptional set of infinite jet codimension is removed. Here infinite codimension means containment, for each integer \(q\), in a jet-defined exceptional set of codimension at least \(q\). The asserted pieces are topological equivalence classes or constant-type subsets; no explicit list is provided.}
\GSfield{Reconstruction}{The note does not define its partition rigorously. We preserve the announced theorem and do not infer that every germ is topologically finitely determined, or that the listed strata are smooth equivalence orbits. Earlier occurrence: 1965-72.11.}
\GSfield{Fate}{\citet[\S1, Theorem 4, pp.~5--6]{arnold1968singularities} repeats the announcement and says its proof was unpublished. A precise later result is \citet[Introduction, p.~88]{mather1972stratifications}: for each source/target dimension pair there is a closed semialgebraic exceptional set in a finite jet space, of codimension greater than the source dimension, outside which only finitely many topological germ types occur. That paper gives a proof sketch; its Harvard preprint dates from 1971 and the opened translation appeared in 1972. The broader isotopy framework is in \citet[\S11]{mather1970topological}. These statements are a verified development, not an exact verification of the stronger countable partition outside infinite codimension.}
\GSfield{Present status}{As of 9 November 2026: the later generic finite-classification theorem is established; the exact infinite-codimension partition wording of this announcement was not independently verified in the opened sources.}
\GSfield{Assessment}{literal statement: correct with missing hypotheses; publication match: partial; verification: indirect; development: annotated; confidence: medium}
\end{GSitem}
\begin{GSitem}{1966-115.13}{E}{Finite jets and their bundle}
\GSfield{Kind; source}{definition; \GSpage{P000175-P000176}}
\GSfield{The notes}{Two germs have the same \(k\)-jet when they have contact of order \(k\); the note writes \(f(x)-g(x)=o(x^k)\). Order zero remembers the image point, order one also the differential. \(J^k(M,N)\) is a bundle over \(M\); for \(N=\mathbb R^n\) it is a vector bundle.}
\GSfield{Statement}{For smooth germs with the same source and target base points, equality of all partial derivatives of orders at most \(k\) defines equality of \(k\)-jets; equivalently in coordinates, \(f(x)-g(x)=o(\|x\|^k)\). The jet bundle has projections to \(M\times N\) and to \(M\), and \(j^kf\) is a section over \(M\). For vector-valued maps it has a natural vector-bundle structure over \(M\). In arbitrary smooth coordinates its fibres are finite-dimensional spaces of Taylor coefficients, not a canonical global algebraic variety.}
\GSfield{Reconstruction}{The norm in the remainder and the limitation of algebraic to coordinate fibres are explicit editorial qualifications.}
\GSfield{Fate}{Jet spaces and the jet-section construction are set out in \citet[\S1]{boardman1967singularities}; the definitions predate the lecture.}
\GSfield{Present status}{As of 9 November 2026: Established, with the hypotheses stated here.}
\GSfield{Assessment}{literal statement: correct with missing hypotheses; publication match: exact; verification: read the result; development: annotated; confidence: high}
\end{GSitem}
\begin{GSitem}{1966-115.14}{E}{The truncation exponent in the jet algebra}
\GSfield{Kind; source}{claim; \GSpage{P000176}}
\GSfield{The notes}{The displayed formula is \(j_0^k(n,1)=\mathbb R[x_1,\ldots,x_n]/M^k\), with \(M=(x_1,\ldots,x_n)\); the scan has no \(+1\).}
\GSfield{Statement}{The algebra of scalar \(k\)-jets at the origin is \(\mathbb R[x_1,\ldots,x_n]/(x_1,\ldots,x_n)^{k+1}\). Its basis consists of monomials of total degree at most \(k\), and its dimension is \(\binom{n+k}{k}\).}
\GSfield{Reconstruction}{The exponent is a mathematical emendation of a faithful transcription. The notebook itself says that first jets remember the differential; quotienting by \(M\) would erase it.}
\GSfield{Fate}{The standard Taylor-jet definition in \citet[\S1]{boardman1967singularities} fixes the indexing. The discrepancy is not a new mathematical claim.}
\GSfield{Present status}{As of 9 November 2026: Established, with the hypotheses stated here.}
\GSfield{Assessment}{literal statement: false as written; publication match: partial; verification: read the result; development: annotated; confidence: high}
\end{GSitem}
\begin{GSitem}{1966-115.15}{E}{Admissible finite-codimension conditions on germs}
\GSfield{Kind; source}{definition; \GSpage{P000176}}
\GSfield{The notes}{An admissible set is the inverse image of a finite-codimension subset of a finite jet space. A codimension-\(e\) variety is said to be given by \(e\) equations.}
\GSfield{Statement}{If \(S\subset J^k_{0,0}(\mathbb R^m,\mathbb R^n)\) is a smooth locally closed submanifold of codimension \(e\), the inverse image of \(S\) under the jet projection is a finite-jet condition of codimension \(e\). Locally, \(e\) equations specify this codimension only when their differentials are independent.}
\GSfield{Reconstruction}{The regularity condition on the equations is essential; the bare number of equations is not a dimension theorem.}
\GSfield{Fate}{This finite-jet interpretation of codimension accompanies \citet[\S1, pp.~5--6]{arnold1968singularities}.}
\GSfield{Present status}{As of 9 November 2026: Established, with the hypotheses stated here.}
\GSfield{Assessment}{literal statement: correct with missing hypotheses; publication match: exact; verification: read the result; development: annotated; confidence: high}
\end{GSitem}
\begin{GSitem}{1966-115.16}{E}{Rank strata and their generic codimension}
\GSfield{Kind; source}{theorem; \GSpage{P000176-P000177}}
\GSfield{The notes}{\(\Sigma^i(f)=\{x:\dim\ker df_x=i\}\). For generic maps the sets are manifolds, with codimension the product of source and target coranks, \((m-r)(n-r)\).}
\GSfield{Statement}{The matrices of rank exactly \(r\) in \(\operatorname{Hom}(\mathbb R^m,\mathbb R^n)\) form a smooth locally closed submanifold of codimension \((m-r)(n-r)\). If \(j^1f\) is transverse to this stratum, its inverse image is empty or a smooth submanifold of \(M\) of that codimension. The exact-rank stratum is locally closed; the rank-at-most locus is determinantal and can be singular.}
\GSfield{Reconstruction}{Transversality and the distinction between an exact-rank stratum and its algebraic closure are made explicit. Earlier occurrence: 1965-72.6.}
\GSfield{Fate}{Thom\textquotesingle s older generic-rank theorem appears in \citet[Part I]{thom1956singularites}; compare \citet[\S2 and Theorem 6.2]{boardman1967singularities}.}
\GSfield{Present status}{As of 9 November 2026: Established, with the hypotheses stated here.}
\GSfield{Assessment}{literal statement: correct with missing hypotheses; publication match: exact; verification: read the result; development: annotated; confidence: high}
\end{GSitem}
\begin{GSitem}{1966-115.17}{E}{Transversality of two maps}
\GSfield{Kind; source}{definition; \GSpage{P000177}}
\GSfield{The notes}{For \(f:A\to B\) and \(g:C\to B\), transversality is \(f_*T_aA+g_*T_cC=T_bB\) whenever \(f(a)=g(c)=b\). Intersecting lines in \(\mathbb R^3\) illustrate failure.}
\GSfield{Statement}{Smooth maps \(f:A\to B\), \(g:C\to B\) are transverse if the sum of their differential images is \(T_bB\) at every matched pair. Then their fibre product is a smooth submanifold of \(A\times C\) of dimension \(\dim A+\dim C-\dim B\), when nonempty.}
\GSfield{Reconstruction}{The fibre-product formulation specifies the meaning for arbitrary maps; the note supplies the tangent-space definition.}
\GSfield{Fate}{The classical tangent-space definition and its inverse-image consequence are recalled in \citet[\S2, footnote on p.~8]{arnold1968singularities}.}
\GSfield{Present status}{As of 9 November 2026: Established, with the hypotheses stated here.}
\GSfield{Assessment}{literal statement: correct; publication match: exact; verification: read the result; development: annotated; confidence: high}
\end{GSitem}
\begin{GSitem}{1966-115.18}{E}{Transverse inverse images and perturbation}
\GSfield{Kind; source}{theorem; \GSpage{P000177}}
\GSfield{The notes}{If \(f\) is transverse to an embedded submanifold \(C\subset B\), then \(f^{-1}C\) is a manifold of the same codimension. Every map can be perturbed to be transverse.}
\GSfield{Statement}{For a smooth embedded submanifold \(C\subset B\) and \(f\pitchfork C\), the preimage is empty or a smooth submanifold with \(\operatorname{codim}_A f^{-1}C=\operatorname{codim}_B C\). Thom transversality makes transverse maps residual and dense in the smooth topology; a corresponding jet-transversality theorem applies to submanifolds of \(J^k(A,B)\).}
\GSfield{Reconstruction}{The statement involving \(\Sigma^i\) is about the jet extension \(j^1f\), not about \(f:M\to N\) being transverse to a subset of a different ambient space. Openness is not asserted for arbitrary nonclosed strata.}
\GSfield{Fate}{Thom\textquotesingle s method is in \citet{thom1956singularites}; the jet-space application is used explicitly by \citet[Theorem 6.3]{boardman1967singularities}.}
\GSfield{Present status}{As of 9 November 2026: Established, with the hypotheses stated here.}
\GSfield{Assessment}{literal statement: correct with missing hypotheses; publication match: exact; verification: read the result; development: annotated; confidence: high}
\end{GSitem}
\begin{GSitem}{1966-115.19}{E}{Iterating the corank of restrictions}
\GSfield{Kind; source}{construction; \GSpage{P000177-P000178}}
\GSfield{The notes}{The restriction \(f|_{\Sigma^i(f)}\) is formed and its corank loci define \(\Sigma^{ij}(f)\); the example \(\Sigma^{110}\) is mentioned.}
\GSfield{Statement}{Provided the successive loci are smooth, define \(\Sigma^{i_1\cdots i_k}(f)\) by imposing kernel dimension \(i_k\) on the restriction of \(f\) to the previous locus. The successive kernels are nested and \(i_1\geq\cdots\geq i_k\). In the surface cusp, the critical curve has symbol beginning \((1)\), ordinary folds \((1,0)\), and the cusp \((1,1,0)\).}
\GSfield{Reconstruction}{The construction is conditional on smoothness; Boardman\textquotesingle s jet definition avoids this restriction. Earlier occurrence: 1965-72.7.}
\GSfield{Fate}{The naive iterative construction and its limitations are explicitly described by \citet[Introduction, p.~21]{boardman1967singularities}.}
\GSfield{Present status}{As of 9 November 2026: Established, with the hypotheses stated here.}
\GSfield{Assessment}{literal statement: correct with missing hypotheses; publication match: exact; verification: read the result; development: annotated; confidence: high}
\end{GSitem}
\begin{GSitem}{1966-115.20}{B}{Boardman universal strata and simultaneous genericity}
\GSfield{Kind; source}{theorem; \GSpage{P000178}}
\GSfield{The notes}{The notes credit Boardman (1965) with universal \(\Sigma^I\subset J^k(M,N)\), computed codimensions, and approximation by maps transverse to all these strata.}
\GSfield{Statement}{For a finite nonincreasing symbol \(I=(i_1,\ldots,i_k)\), the Boardman set is a smooth locally closed submanifold of the finite jet space when the symbol satisfies the dimension restrictions. Its codimension is the Boardman combinatorial number. Smooth maps whose jet extensions are transverse to all Boardman strata are dense; for them the intrinsic jet definition agrees with the iterative corank loci.}
\GSfield{Reconstruction}{The notes\textquotesingle{} algebraic subvarieties is understood in local jet coordinates as locally closed sets given by polynomial equalities and inequalities, not as closed varieties. The explicit formula is treated in record 116. Earlier occurrence: 1965-72.8.}
\GSfield{Fate}{The 1965 circulation date is the notebook\textquotesingle s testimony. The published theorem is \citet[Theorems 6.1--6.3, p.~46]{boardman1967singularities}, 1967.}
\GSfield{Present status}{As of 9 November 2026: Established, with the hypotheses stated here.}
\GSfield{Assessment}{literal statement: correct with missing hypotheses; publication match: exact; verification: read the result; development: annotated; confidence: high}
\end{GSitem}
\begin{GSitem}{1966-115.21}{B}{Boardman symbols are not complete smooth invariants}
\GSfield{Kind; source}{claim; \GSpage{P000178}}
\GSfield{The notes}{The last line says the symbols are not a complete set of invariants, first in \(\mathbb R^4\); rank drops cannot distinguish all cases.}
\GSfield{Statement}{There are smoothly inequivalent stable germs \((\mathbb R^4,0)\to(\mathbb R^4,0)\) with the same symbol \((2,0)\). Their intrinsic quadratic pencils have different real types. Therefore a Boardman symbol records successive degeneracy dimensions but not the full left--right equivalence class.}
\GSfield{Reconstruction}{The next lecture supplies explicit real quadratic-pencil representatives. Earlier occurrence: 1965-72.10.}
\GSfield{Fate}{The two real corank-two examples are given by \citet[\S2, Example 4; \S3]{arnold1968singularities}, while the systematic local-algebra classification is \citet{mather1969algebras}.}
\GSfield{Present status}{As of 9 November 2026: Established, with the hypotheses stated here.}
\GSfield{Assessment}{literal statement: correct; publication match: exact; verification: read the result; development: annotated; confidence: high}
\end{GSitem}

\GStalk{116}{1966-11-28}{V. I. Arnold}{P000179-P000186}{Singularities of smooth mappings II: intrinsic derivatives and local algebras}
The second lecture constructs Boardman ideals, states their codimension formula, and studies Morin germs and intrinsic quadratic pencils. It then explains the first smooth moduli, Mather\textquotesingle s local-algebra and infinitesimal-stability criteria, smooth preparation, and five real corank-two families. The scans reveal substantive mathematical slips in the codimension conventions and one mixed monomial; the corrected statements are separated from the notebook readings. The MCCME catalogue span P000178--P000188 overlaps both adjacent records; the dated lecture itself occupies P000179--P000186.

\begin{GSitem}{1966-116.1}{E}{Nested corank loci and the limitation of the naive definition}
\GSfield{Kind; source}{construction; \GSpage{P000179}}
\GSfield{The notes}{The lecture restarts with \(\Sigma^i(f)=\{\dim\ker df=i\}\), then \(\Sigma^{ij}(f)=\Sigma^j(f|_{\Sigma^i(f)})\), labelled Thom\textquotesingle s naive definition. Nested kernels imply \(i_1\geq i_2\geq\cdots\).}
\GSfield{Statement}{If every successive corank locus is a smooth manifold, the restriction construction defines iterated corank strata. The tangent kernels at a point satisfy \(K_{j+1}\subset K_j\); their dimensions form a nonincreasing sequence. For arbitrary smooth maps some intermediate locus need not be smooth, so this is not a universal definition.}
\GSfield{Reconstruction}{The note itself says to abandon the naive definition in favor of the jet construction; it is retained here as a separate occurrence of 1966-115.19.}
\GSfield{Fate}{\citet[Introduction and Theorem 6.2]{boardman1967singularities} gives the conditional geometric interpretation and its relation to the intrinsic construction.}
\GSfield{Present status}{As of 9 November 2026: Established, with the hypotheses stated here.}
\GSfield{Assessment}{literal statement: correct; publication match: exact; verification: read the result; development: annotated; confidence: high}
\end{GSitem}
\begin{GSitem}{1966-116.2}{B}{Jacobian extension of an ideal and source-coordinate invariance}
\GSfield{Kind; source}{definition; \GSpage{P000179}}
\GSfield{The notes}{For an ideal \(B\subset A\), adjoin the \(k\)-rowed Jacobian minors of functions in \(B\) to obtain \(\Delta_k(B)\). A chain-rule display explains coordinate invariance.}
\GSfield{Statement}{Let \(A=C^\infty_{\mathbb R^m,0}\), and let \(B\) be a finitely generated ideal. Define \(\Delta_k(B)\) to be \(B\) plus the ideal generated by all \(k\times k\) minors of Jacobian matrices of finite lists of elements of \(B\). This depends only on \(B\), not on a generating list or smooth local coordinates. One may use a generating list, because derivatives of coefficients only contribute terms already in \(B\). Coordinate invariance follows from the chain rule and Cauchy--Binet.}
\GSfield{Reconstruction}{The proof for all minors extends the notebook\textquotesingle s first-derivative calculation. The construction is stated for local germs, avoiding an undefined global ring of functions.}
\GSfield{Fate}{\citet[\S2, Definition 2.1 and Lemmas 2.2--2.3]{boardman1967singularities} develops Jacobian extensions invariantly, publishing in 1967 work already circulating in 1965.}
\GSfield{Present status}{As of 9 November 2026: Established, with the hypotheses stated here.}
\GSfield{Assessment}{literal statement: correct with missing hypotheses; publication match: exact; verification: read the result; development: annotated; confidence: high}
\end{GSitem}
\begin{GSitem}{1966-116.3}{B}{Critical Jacobian extension: the first proper ideal}
\GSfield{Kind; source}{claim; \GSpage{P000180}}
\GSfield{The notes}{As \(k\) increases, the ideals decrease. The text calls \(k\) max when the ideal is not the whole \(B\), but the displayed condition is \(\Delta_{k-1}B=A\), \(\Delta_kB\subset A\); also \(\Delta_0B=A\).}
\GSfield{Statement}{For a proper local ideal \(B\), the critical minor size is the least \(k\geq1\) for which \(\Delta_k(B)\ne A\), equivalently \(\Delta_{k-1}(B)=A\) and \(\Delta_k(B)\subsetneq A\). With source dimension \(m\), write \(\Delta^i=\Delta_{m-i+1}\); the critical superscript then encodes the corank.}
\GSfield{Reconstruction}{The words max and whole B conflict with both the monotonicity and the displayed transition. They are emended to min and whole A, without altering the diplomatic source.}
\GSfield{Fate}{The intrinsic corank/critical-extension definition is established in \citet[\S2]{boardman1967singularities}.}
\GSfield{Present status}{As of 9 November 2026: Established, with the hypotheses stated here.}
\GSfield{Assessment}{literal statement: false as written; publication match: partial; verification: read the result; development: annotated; confidence: high}
\end{GSitem}
\begin{GSitem}{1966-116.4}{E}{Successive extensions of the cubic ideal}
\GSfield{Kind; source}{construction; \GSpage{P000180}}
\GSfield{The notes}{The example begins with the ideal generated by \(x^3\), gives \(\Delta_1=(x^2)\), and has a crowded iterated expression ending in \((x)\).}
\GSfield{Statement}{In \(C^\infty_{\mathbb R,0}\), the successive first Jacobian extensions are \((x^3)\mapsto(x^2)\mapsto(x)\mapsto A\). Hence the cubic map germ has initial Boardman symbol \((1,1,0)\), although as a one-variable germ it is not a transverse, stable representative of that symbol.}
\GSfield{Reconstruction}{\oursComputed{} The iteration sequence, rather than a literal interpretation of the overloaded superscript, is fixed by differentiation. The final zero is the critical corank when the current ideal is \((x)\).}
\GSfield{Fate}{This is a direct example of the definition in \citet[\S2]{boardman1967singularities}. It illustrates algebraic symbol computation rather than a stability theorem.}
\GSfield{Present status}{As of 9 November 2026: Established, with the hypotheses stated here.}
\GSfield{Assessment}{literal statement: not determinable; publication match: partial; verification: read the result; development: annotated; confidence: high}
\end{GSitem}
\begin{GSitem}{1966-116.5}{B}{Target-coordinate independence of the Boardman ideal construction}
\GSfield{Kind; source}{question; \GSpage{P000180}}
\GSfield{The notes}{For \(B=(f_1,\ldots,f_n)\), all iterated extensions \(\Delta^{i_j}\cdots\Delta^{i_1}B\) are to be critical. Source independence is stated; target independence is marked \ru{нужно пров.} (needs checking).}
\GSfield{Statement}{For a germ \(f:(\mathbb R^m,0)\to(\mathbb R^n,0)\), the ideal is intrinsically \(B=f^*\mathfrak m_n\,A\). If \(h\) is a target diffeomorphism germ fixing zero, then \(h^*\mathfrak m_n=\mathfrak m_n\), so \((h\circ f)^*\mathfrak m_n A=B\). Thus all subsequent Jacobian extensions and the Boardman symbol are unchanged. Membership for a symbol of length \(k\) depends only on the \(k\)-jet.}
\GSfield{Reconstruction}{\ours{} The maximal-ideal argument is the missing check. It resolves a pedagogical proof obligation, not an unanswered research problem.}
\GSfield{Fate}{\citet[\S2, (2.17)--(2.18)]{boardman1967singularities} supplies the universal coordinate-independent construction.}
\GSfield{Present status}{As of 9 November 2026: Resolved by the intrinsic ideal definition; no separate open problem is introduced.}
\GSfield{Assessment}{literal statement: correct; publication match: exact; verification: read the result; development: annotated; confidence: high}
\end{GSitem}
\begin{GSitem}{1966-116.6}{B}{Boardman codimension formula and two missing conventions}
\GSfield{Kind; source}{claim; \GSpage{P000180-P000181}}
\GSfield{The notes}{The coefficient is written \(m-n+i_1\). The count \(\mu(I)\) includes nonincreasing sequences satisfying \(0\leq j_s\leq i_s\), without excluding the all-zero sequence.}
\GSfield{Statement}{For source dimension \(m\), target dimension \(n\), and an admissible nonincreasing symbol \(I=(i_1,\ldots,i_k)\), put \(\mu(I)=\#\{j_1\geq\cdots\geq j_k\geq0: j_s\leq i_s,\ j_1>0\}\). Then \[c(I)=(n-m+i_1)\mu(I)-\sum_{r=1}^{k-1}(i_r-i_{r+1})\mu(i_{r+1},\ldots,i_k).\] For length one this gives \(i(n-m+i)\), and for \(I=(1,\ldots,1)\) of length \(k\) it gives \(k(n-m+1)\).}
\GSfield{Reconstruction}{The dimension difference must be target minus source; the all-zero sequence must be excluded. Both repairs are forced already by the rank-stratum case. The supplied check script independently counts the sequences by enumeration and recursion; these finite computations are \oursComputed{} and do not replace Boardman\textquotesingle s theorem.}
\GSfield{Fate}{The formula and both conventions are \citet[Definitions 6.4--6.5, pp.~46--47]{boardman1967singularities}. The scan confirms that these are mathematical emendations of the notebook, not transcription errors.}
\GSfield{Present status}{As of 9 November 2026: Established, with the hypotheses stated here.}
\GSfield{Assessment}{literal statement: false as written; publication match: partial; verification: read the result; development: annotated; confidence: high}
\end{GSitem}
\begin{GSitem}{1966-116.7}{B}{Generic agreement of intrinsic and geometric Boardman loci}
\GSfield{Kind; source}{theorem; \GSpage{P000181}}
\GSfield{The notes}{Boardman\textquotesingle s three assertions are smooth strata, the codimension formula, and transversality for almost all maps, giving agreement with Thom\textquotesingle s naive definition.}
\GSfield{Statement}{If a smooth map has jet extension transverse to each relevant Boardman stratum, the inverse-image strata are smooth of the stated codimensions, and iterated corank loci agree pointwise with the intrinsic jet loci. Such maps are dense (indeed obtained by jet transversality). Almost all denotes genericity in function space, not a probability measure.}
\GSfield{Reconstruction}{The note writes \(f\in\Sigma^I\); the precise statement is \(x\in\Sigma^I(f)\) if and only if \(j_x^kf\in\Sigma^I\).}
\GSfield{Fate}{This is \citet[Theorems 6.1--6.3]{boardman1967singularities}, published 1967.}
\GSfield{Present status}{As of 9 November 2026: Established, with the hypotheses stated here.}
\GSfield{Assessment}{literal statement: correct with missing hypotheses; publication match: exact; verification: read the result; development: annotated; confidence: high}
\end{GSitem}
\begin{GSitem}{1966-116.8}{B}{First dimensions of transverse Boardman symbols}
\GSfield{Kind; source}{construction; \GSpage{P000181}}
\GSfield{The notes}{For equidimensional maps the table gives \(1_k\) in dimension \(k\), \((2)\) at four, \((2,1)\) at seven, \((3)\) at nine, and \((2,2)\), \((2,1,1)\) at ten. A later line mentions 16 and 4. The note asks to write the table precisely.}
\GSfield{Statement}{A Boardman stratum of codimension \(c\) can meet a transverse jet section only if the source dimension is at least \(c\). For equidimensional maps the displayed values are \(c(1_k)=k\), \(c(2)=4\), \(c(2,1)=7\), \(c(3)=9\), \(c(2,2)=c(2,1,1)=10\), and \(c(4)=16\). At \(n=9\), the fold, cusp, and corank-two loci have dimensions 8, 7, and 5 respectively.}
\GSfield{Reconstruction}{First appearance means the codimension threshold for a generic map, not existence of such a germ in arbitrary small dimension without transversality. The isolated 4 at dimension 16 is interpreted as corank four, not four equivalence classes. Values were checked by two independent counts in the supplied script (\oursComputed).}
\GSfield{Fate}{These values follow from \citet[(6.5), p.~47]{boardman1967singularities}; they are not an exhaustive classification by smooth type.}
\GSfield{Present status}{As of 9 November 2026: Established, with the hypotheses stated here.}
\GSfield{Assessment}{literal statement: correct with missing hypotheses; publication match: partial; verification: read the result; development: annotated; confidence: high}
\end{GSitem}
\begin{GSitem}{1966-116.9}{E}{Morin normal forms for corank-one singularities}
\GSfield{Kind; source}{theorem; \GSpage{P000182}}
\GSfield{The notes}{The theorem name is written Moren/Morin. A good \(\Sigma^{1_k0}\) germ is given a generalized cusp normal form; the display specializes the number of ones to the ambient dimension.}
\GSfield{Statement}{Let \(1\leq k\leq n\). A transverse corank-one germ \((\mathbb R^n,0)\to(\mathbb R^n,0)\) of symbol \((\underbrace{1,\ldots,1}_{k},0)\) is smoothly equivalent to \[(u_1,\ldots,u_{n-1},t)\longmapsto\left(u_1,\ldots,u_{n-1},\ t^{k+1}+\sum_{j=1}^{k-1}u_jt^j\right).\] It is a stable Morin germ. The notebook\textquotesingle s formula is the case \(k=n\).}
\GSfield{Reconstruction}{The transversality hypothesis cannot be replaced simply by algebraic membership of an arbitrary jet in the same Boardman locus; for example \(x^3\) alone is not a stable one-variable germ. Parameters and the index \(k\) are distinguished here.}
\GSfield{Fate}{Morin\textquotesingle s result appeared in two 1965 Comptes rendus notes; \citet[\S2, Example 6, pp.~9--10]{arnold1968singularities} states the form and cites them. This item uses that checked exposition, not an uninspected original proof.}
\GSfield{Present status}{As of 9 November 2026: Established, with the hypotheses stated here.}
\GSfield{Assessment}{literal statement: correct with missing hypotheses; publication match: exact; verification: indirect; development: annotated; confidence: high}
\end{GSitem}
\begin{GSitem}{1966-116.10}{E}{The intrinsic quadratic differential}
\GSfield{Kind; source}{construction; \GSpage{P000182-P000183}}
\GSfield{The notes}{The map \(f_{xx}:\ker df_x\to\operatorname{coker}df_x\) is defined by dividing the projected curve \(f(X(t\xi))\) by \(t^2\). The note says coordinate independence is easy and adds \ru{Посмотреть!}.}
\GSfield{Statement}{For a smooth map germ, put \(K=\ker df_x\), \(C=T_{f(x)}N/\operatorname{im}df_x\). The symmetric bilinear form \(H_f:K\times K\to C\) defined in coordinates by the class of \(D^2f_x(v,w)\) is intrinsic. The note\textquotesingle s diagonal limit is \(\tfrac12H_f(\xi,\xi)\). Under source-coordinate change the additional term lies in \(\operatorname{im}df_x\); under target-coordinate change the extra term contains \(df_x(v),df_x(w)\), which vanish.}
\GSfield{Reconstruction}{The factor \(1/2\) follows from Taylor\textquotesingle s formula. The one-line coordinate-independence proof fills the explicit checking instruction; it is not attributed to the notebook.}
\GSfield{Fate}{Intrinsic derivatives predate the lecture; \citet[\S7]{boardman1967singularities} develops them, and \citet[\S3, Definition 1 and Lemma 1]{arnold1968singularities} gives this exact quadratic construction.}
\GSfield{Present status}{As of 9 November 2026: Established, with the hypotheses stated here.}
\GSfield{Assessment}{literal statement: correct; publication match: exact; verification: read the result; development: annotated; confidence: high}
\end{GSitem}
\begin{GSitem}{1966-116.11}{E}{Why a cubic differential is not defined by the same recipe}
\GSfield{Kind; source}{remark; \GSpage{P000183}}
\GSfield{The notes}{The parenthesis says that a cubic differential cannot be defined this way.}
\GSfield{Statement}{For a general germ, the raw third Taylor coefficient restricted to \(\ker df\) and projected to \(\operatorname{coker}df\) is not coordinate invariant. Already \(f(x)=x^2\) has zero cubic coefficient in coordinate \(x\), but under \(x=u+au^2\) it becomes \(u^2+2au^3+O(u^4)\). Higher intrinsic derivatives require additional kernel/quotient data or vanishing of the relevant lower derivatives.}
\GSfield{Reconstruction}{\oursComputed{} The explicit one-variable check clarifies the intended limitation; it does not deny Boardman\textquotesingle s higher intrinsic derivatives.}
\GSfield{Fate}{\citet[\S7]{boardman1967singularities} gives the corrected higher-derivative framework. The notebook\textquotesingle s warning is mathematically appropriate for the naive recipe.}
\GSfield{Present status}{As of 9 November 2026: Established, with the hypotheses stated here.}
\GSfield{Assessment}{literal statement: correct; publication match: exact; verification: read the result; development: annotated; confidence: high}
\end{GSitem}
\begin{GSitem}{1966-116.12}{E}{Binary quadratic pencils and seven linear orbit types}
\GSfield{Kind; source}{construction; \GSpage{P000183}}
\GSfield{The notes}{The dual of the cokernel maps to quadratic forms on the kernel. For kernel and cokernel of dimension two, binary quadratic forms occupy \(\mathbb R^3\), with degenerate cone \(b^2=ac\). Seven positions, including degeneracies, give seven orbits.}
\GSfield{Statement}{Let \(q:\mathbb R^2\to\mathbb R^2\) be homogeneous quadratic. Its dual pencil is a linear map \((\mathbb R^2)^*\to\operatorname{Sym}^2((\mathbb R^2)^*)\). Up to independent invertible source and target linear changes there are seven types: three rank-two pencils (plane cutting, tangent to, or exterior to the discriminant cone), three rank-one pencils (definite, indefinite, or rank-one form), and the zero pencil. Only the pencil is classified by this seven-orbit statement, not every smooth germ with that quadratic part.}
\GSfield{Reconstruction}{Allowing the pencil image to have dimensions zero or one is essential to the count seven; a literal plane alone gives only three cases. The seven quadratic orbits must not be confused with seven smooth \(\Sigma^{20}\) germs.}
\GSfield{Fate}{The elementary real linear classification appears in \citet[\S3, Example 2, pp.~13--14]{arnold1968singularities}.}
\GSfield{Present status}{As of 9 November 2026: Established, with the hypotheses stated here.}
\GSfield{Assessment}{literal statement: correct with missing hypotheses; publication match: exact; verification: read the result; development: annotated; confidence: high}
\end{GSitem}
\begin{GSitem}{1966-116.13}{A}{Two transverse corank-two germs in four dimensions}
\GSfield{Kind; source}{construction; \GSpage{P000184}}
\GSfield{The notes}{The maps are \(y_1=x_1x_2\), \(y_2=x_1^2\pm x_2^2+x_1x_3+x_2x_4\), \(y_3=x_3\), \(y_4=x_4\); the mixed terms are explicitly added for transversality.}
\GSfield{Statement}{These two germs \(F_\pm:(\mathbb R^4,0)\to(\mathbb R^4,0)\) have rank two at zero. Their normal Jacobian block is \(\left(\begin{smallmatrix}x_2&x_1\\2x_1+x_3&\pm2x_2+x_4\end{smallmatrix}\right)\), and its linearization is an isomorphism \(\mathbb R^4\to M_2(\mathbb R)\), so their first jet sections are transverse to the corank-two stratum. Their intrinsic pencils \((xy,x^2\pm y^2)\) are inequivalent over \(\mathbb R\).}
\GSfield{Reconstruction}{Transversality was recomputed symbolically; the determinant is nonzero for either sign (\oursComputed). The signs encode a real invariant, even though the Boardman symbol is the same.}
\GSfield{Fate}{\citet[\S2, Example 4 and \S3]{arnold1968singularities} publishes equivalent corank-two examples based on the 1966 lectures.}
\GSfield{Present status}{As of 9 November 2026: Established, with the hypotheses stated here.}
\GSfield{Assessment}{literal statement: correct; publication match: exact; verification: read the result; development: annotated; confidence: high}
\end{GSitem}
\begin{GSitem}{1966-116.14}{E}{Real signs, complexification, and variation along a corank-two curve}
\GSfield{Kind; source}{remark; \GSpage{P000184}}
\GSfield{The notes}{In \(\mathbb R^5\) the corank-two set can be a curve with varying singularities. The real plus and minus cases cease to differ over \(\mathbb C\), while a parabolic case remains.}
\GSfield{Statement}{For a transverse map between five-dimensional manifolds, the corank-two stratum has dimension one; its quadratic-pencil type can vary along it. Over \(\mathbb C\), replacing one source variable by its multiple by \(i\) identifies the two nonsingular real sign types of binary quadratic pencils, whereas a tangency degeneration is still distinguished by the discriminant.}
\GSfield{Reconstruction}{The assertion is about quadratic-pencil types. No claim is made that every higher-order real or complex germ is determined by this second-order invariant.}
\GSfield{Fate}{\citet[\S3, Examples 2--3]{arnold1968singularities} explains the real pencil geometry and degeneration types.}
\GSfield{Present status}{As of 9 November 2026: Established, with the hypotheses stated here.}
\GSfield{Assessment}{literal statement: correct with missing hypotheses; publication match: exact; verification: read the result; development: annotated; confidence: high}
\end{GSitem}
\begin{GSitem}{1966-116.15}{B}{The overstatement that every corank-two germ is stable}
\GSfield{Kind; source}{claim; \GSpage{P000184}}
\GSfield{The notes}{The line says \(\Sigma^{20}\): all singularities are stable and there is a differential classification; for \(\Sigma^{30}\) this is no longer so.}
\GSfield{Statement}{A Boardman symbol \((2,0)\) by itself does not imply smooth stability. Stable corank-two germs can be classified by their local algebras, subject to the source/target dimension and unfolding conditions. Nonstable germs with the same symbol exist. For example \((x,y)\mapsto(x^2+y^2,xy)\) has algebraic symbol \((2,0)\), but corank two cannot occur in a generic surface map, so it is unstable.}
\GSfield{Reconstruction}{\ours{} The strong literal sentence is corrected to a statement about stable or sufficiently generic germs. The following table supplies an algebraic classification but does not remove the stability hypothesis.}
\GSfield{Fate}{\citet[Theorems A and B]{mather1969algebras} distinguishes algebraic classification from stable realizability; \citet[\S7]{mather1971nice} gives the stable corank-two list.}
\GSfield{Present status}{As of 9 November 2026: Established, with the hypotheses stated here.}
\GSfield{Assessment}{literal statement: false as written; publication match: partial; verification: read the result; development: annotated; confidence: high}
\end{GSitem}
\begin{GSitem}{1966-116.16}{B}{Dimension count forcing moduli of quadratic maps}
\GSfield{Kind; source}{theorem; \GSpage{P000184}}
\GSfield{The notes}{For \(n\geq3\), the group \(GL(n)\times GL(n)\) has dimension \(2n^2\), the space of \(n\) quadratic forms has dimension \(n^2(n+1)/2\), and scalar transformations provide a one-dimensional stabilizer. Hence every orbit has positive codimension.}
\GSfield{Statement}{Let \(H_n=\operatorname{Hom}(\operatorname{Sym}^2\mathbb R^n,\mathbb R^n)\), with the action \((A,B)q=B\circ q\circ A^{-1}\). The one-dimensional subgroup \((tI,t^2I)\), \(t\ne0\), acts trivially. Thus every orbit has dimension at most \(2n^2-1\), and codimension at least \(n^2(n-3)/2+1\), which is positive for \(n\geq3\). For a transverse corank-three point of a map \(\mathbb R^9\to\mathbb R^9\), these quadratic moduli obstruct smooth stability.}
\GSfield{Reconstruction}{The precise scalar action and the lower bound expand the proof already present in the notebook. Here \(n\) is the kernel/cokernel dimension, not the ambient dimension nine.}
\GSfield{Fate}{This is the argument of \citet[\S3, Lemma 3 and Theorem 1]{arnold1968singularities}, attributed there to Thom; the later complete setting is \citet{mather1971nice}.}
\GSfield{Present status}{As of 9 November 2026: Established, with the hypotheses stated here.}
\GSfield{Assessment}{literal statement: correct; publication match: exact; verification: read the result; development: annotated; confidence: high}
\end{GSitem}
\begin{GSitem}{1966-116.17}{B}{The announced dimension-dependent density classification}
\GSfield{Kind; source}{claim; \GSpage{P000184}}
\GSfield{The notes}{A dimension diagram labels the embedding region, Whitney, Morse, and cells 9,9 and 6,9; Mather is credited with announcing where moduli occur and classifying stable singularities.}
\GSfield{Statement}{For fixed source and target dimensions, there is a numerical criterion for density of proper smoothly stable maps among proper smooth maps. In Mather\textquotesingle s notation it is \(\dim M<\sigma(\dim M,\dim N)\), with an explicitly computed function \(\sigma\). The notebook diagram alone does not specify the ordering of the coordinates of every cell or the exact boundary inequalities.}
\GSfield{Reconstruction}{The numerical criterion is identified through the later primary paper; no missing line in the diagram is reconstructed. The two ordered pairs are not silently interchanged.}
\GSfield{Fate}{The announcement matured into \citet[Introduction and \S\S1--5]{mather1971nice}, published 1971; the opened Russian translation dates from 1974. Earlier occurrence: 1965-72.14.}
\GSfield{Present status}{As of 9 November 2026: Established, with the hypotheses stated here.}
\GSfield{Assessment}{literal statement: not determinable; publication match: partial; verification: read the result; development: annotated; confidence: medium}
\end{GSitem}
\begin{GSitem}{1966-116.18}{E}{The local algebra of a map germ}
\GSfield{Kind; source}{definition; \GSpage{P000185}}
\GSfield{The notes}{The local ring is written as a quotient of \(E_x\) by the ideal generated by the pullback of the maximal ideal of \(E_y\); it is called the ring of functions on the coalesced inverse image.}
\GSfield{Statement}{For a smooth map germ \(f:(\mathbb R^m,0)\to(\mathbb R^n,0)\), let \(E_m\) and \(E_n\) be the real algebras of smooth germs and \(\mathfrak m_n\) the target maximal ideal. Define \(Q(f)=E_m/(f^*\mathfrak m_n)E_m\). Its isomorphism type as a real algebra is invariant under source and target diffeomorphisms. The ideal product is necessary: the pullback subalgebra alone is not the quotient ideal.}
\GSfield{Reconstruction}{The notation \(f_*\) in the notebook denotes pullback of functions; modern notation \(f^*\) avoids conflict with differential pushforward.}
\GSfield{Fate}{The quotient is the starting definition in \citet[Introduction, pp.~223--224]{mather1969algebras}; \citet[\S4]{arnold1968singularities} gives the same exposition.}
\GSfield{Present status}{As of 9 November 2026: Established, with the hypotheses stated here.}
\GSfield{Assessment}{literal statement: correct; publication match: exact; verification: read the result; development: annotated; confidence: high}
\end{GSitem}
\begin{GSitem}{1966-116.19}{E}{Dual numbers as two coalescing points}
\GSfield{Kind; source}{construction; \GSpage{P000185}}
\GSfield{The notes}{For \(f(x)=x^2\), multiplication by \(E\) adds all higher terms, leaving classes \(a+bx\). This is the ring of functions on an infinitesimal pair of points.}
\GSfield{Statement}{For \(f(x)=x^2\), Hadamard\textquotesingle s lemma gives \(Q(f)\simeq\mathbb R[x]/(x^2)\). The family \(\mathbb R[x]/(x^2-t)\) has dimension two for every \(t\); for \(t>0\) it is \(\mathbb R\oplus\mathbb R\), and at \(t=0\) it is the nonreduced dual-number algebra. This makes coalescence precise.}
\GSfield{Reconstruction}{The parameter \(t\) and flat algebraic family formalize the diagram. The real fibre for \(t<0\) has no real points, showing why vector-space length must not be equated with the number of real preimages.}
\GSfield{Fate}{The model and its limiting multiplication are worked out in \citet[\S4, Example 1 and Propositions 1--2]{arnold1968singularities}.}
\GSfield{Present status}{As of 9 November 2026: Established, with the hypotheses stated here.}
\GSfield{Assessment}{literal statement: correct; publication match: exact; verification: read the result; development: annotated; confidence: high}
\end{GSitem}
\begin{GSitem}{1966-116.20}{E}{Local algebra length and multiplicity}
\GSfield{Kind; source}{claim; \GSpage{P000185-P000186}}
\GSfield{The notes}{The note conjecturally identifies \(\dim Q\) with multiplicity for an analytic \(f:\mathbb C^n\to\mathbb C^m\), adding \ru{По-видимому, так} (apparently so); the preparation discussion calls it the correct definition.}
\GSfield{Statement}{For a holomorphic germ \(f:(\mathbb C^n,0)\to(\mathbb C^n,0)\) with isolated fibre \(f^{-1}(0)=\{0\}\) near zero, the finite length \(\dim_\mathbb C\mathcal O_n/(f_1,\ldots,f_n)\) is the local mapping multiplicity; a nearby regular value in a sufficiently small representative has that number of preimages counted with multiplicity. For unequal dimensions the length still makes sense if finite, but it is not automatically the number of inverse images of an arbitrary ambient regular value.}
\GSfield{Reconstruction}{The equal-dimensional finite-map hypotheses are essential for the literal preimage-count interpretation. This resolves a tentative expository question rather than proposing a new conjecture.}
\GSfield{Fate}{The analytic finite-module statement is part of \citet[Chapter III]{malgrange1966ideals}; its application to the local ring is explained in \citet[\S4, Theorem 2 and examples]{arnold1968singularities}. Malgrange\textquotesingle s lectures were delivered in 1964, with a July 1965 preface and 1966 publication.}
\GSfield{Present status}{As of 9 November 2026: Established, with the hypotheses stated here.}
\GSfield{Assessment}{literal statement: correct with missing hypotheses; publication match: exact; verification: read the result; development: annotated; confidence: high}
\end{GSitem}
\begin{GSitem}{1966-116.21}{B}{Stable germs are classified by their local algebras}
\GSfield{Kind; source}{theorem; \GSpage{P000185}}
\GSfield{The notes}{Mather\textquotesingle s theorem says stable singularities are equivalent if and only if their local rings coincide.}
\GSfield{Statement}{For stable smooth monogerms \(f,g:(\mathbb R^m,0)\to(\mathbb R^n,0)\) of the same source and target dimensions, an isomorphism of the appropriate local real algebras \(Q(f)\simeq Q(g)\) determines smooth left--right equivalence. In the finite-algebra case the ordinary local algebras suffice; in general one may use their formal completions or the finite truncations in Mather\textquotesingle s Theorem A. Equality of dimensions and stability cannot be dropped.}
\GSfield{Reconstruction}{Coincide means algebra isomorphism, not literal equality of ideals or equality of vector-space dimensions. This is the map-germ classification theorem, not the later Mather--Yau theorem for hypersurface singularities.}
\GSfield{Fate}{Announced by 1966 according to this notebook and \citet[\S4]{arnold1968singularities}; the checked primary statement is \citet[Theorem A, pp.~223--225]{mather1969algebras}. The manuscript was received on 1 June 1968; the journal volume is dated 1969, while the scanned printing colophon says 1970. The bibliography follows the journal metadata and keeps that distinction explicit.}
\GSfield{Present status}{As of 9 November 2026: Established, with the hypotheses stated here.}
\GSfield{Assessment}{literal statement: correct with missing hypotheses; publication match: exact; verification: read the result; development: annotated; confidence: high}
\end{GSitem}
\begin{GSitem}{1966-116.22}{B}{Infinitesimal stability and the source--target linearization}
\GSfield{Kind; source}{theorem; \GSpage{P000185-P000186}}
\GSfield{The notes}{Every vertical field \(\varphi\) along \(f\) must be \(f_*\xi+f^*\eta\). The accompanying theorem asserts stability if and only if infinitesimal stability.}
\GSfield{Statement}{For a smooth monogerm, let \(\theta_m,\theta_n\) be the modules of source and target vector-field germs and \(\theta(f)\) the vector fields along \(f\). The linearization is \(t_f(\xi)=df\circ\xi\), \(\omega_f(\eta)=\eta\circ f\). Infinitesimal stability is the equality \(t_f\theta_m+\omega_f\theta_n=\theta(f)\); it is equivalent to smooth stability of the germ. For global maps a corresponding theorem requires properness (compact source suffices).}
\GSfield{Reconstruction}{Vector fields in this extended tangent-space criterion need not vanish at the base point. Restricting to vanishing fields would test an origin-fixed equivalence problem. Earlier occurrence: 1965-72.12.}
\GSfield{Fate}{Published in \citet{mather1969stability}; its metadata were checked. The germ linearization and stability convention were read directly in \citet[Introduction and \S1]{mather1969algebras}; \citet[\S5--\S7]{arnold1968singularities} explains the proof method.}
\GSfield{Present status}{As of 9 November 2026: Established, with the hypotheses stated here.}
\GSfield{Assessment}{literal statement: correct with missing hypotheses; publication match: exact; verification: read the result; development: annotated; confidence: high}
\end{GSitem}
\begin{GSitem}{1966-116.23}{E}{Smooth preparation as finite generation over the target}
\GSfield{Kind; source}{theorem; \GSpage{P000186}}
\GSfield{The notes}{If \(Q\) has finite dimension \(r\) and representatives \(e_1,\ldots,e_r\), every smooth germ is written \(\varphi(x)=\sum_i\psi_i(y)e_i(x)\). The heading says Weierstrass preparation.}
\GSfield{Statement}{Let \(f:(\mathbb R^m,0)\to(\mathbb R^n,0)\) be smooth and \(Q(f)\) finite-dimensional. If the classes of \(e_1,\ldots,e_r\) span \(Q(f)\) as a real vector space, then every \(\varphi\in E_m\) has \(\varphi(x)=\sum_{i=1}^r\psi_i(f(x))e_i(x)\) with \(\psi_i\in E_n\). Equivalently \(E_m\) is a finitely generated \(E_n\)-module via \(f^*\). Neither uniqueness of the coefficients nor freeness is asserted.}
\GSfield{Reconstruction}{For \(C^\infty\) germs this is the Malgrange preparation theorem and its finite-module form, not merely analytic Weierstrass preparation. Generators must mean linear spanning representatives in \(Q\), not merely generators of \(Q\) as an algebra.}
\GSfield{Fate}{\citet[Chapter V]{malgrange1966ideals} publishes the smooth preparation theorem, from 1964 lectures; \citet[\S4, Theorem 1]{arnold1968singularities} gives exactly this finite-module formulation and records Mather\textquotesingle s 1966 simplification.}
\GSfield{Present status}{As of 9 November 2026: Established, with the hypotheses stated here.}
\GSfield{Assessment}{literal statement: correct with missing hypotheses; publication match: exact; verification: read the result; development: annotated; confidence: high}
\end{GSitem}
\begin{GSitem}{1966-116.24}{B}{Five families of stable real corank-two local algebras}
\GSfield{Kind; source}{theorem; \GSpage{P000186}}
\GSfield{The notes}{The list is \(I_{a,b}\), \(II_{a,b}\), \(III_{a,b}\), \(IV_a\), \(V_a\). The first ideal has a doubtful comma between \(x,y\). The notebook\textquotesingle s IV has \((x^a,xy^{a-1},x^2+y^2)\), and V has \((x^a,x^2+y^2)\). The left margin mixes dimension conditions and row numbers.}
\GSfield{Statement}{For stable real smooth germs of type \((2,0)\) with source dimension at most target dimension, Mather\textquotesingle s final list is, using \(R=\mathbb R[[x,y]]\): \[\begin{aligned}I_{a,b}&:R/(xy,x^a+y^b),&b\geq a\geq2;\\II_{a,b}&:R/(xy,x^a-y^b),&b\geq a\geq2\ \text{both even};\\III_{a,b}&:R/(xy,x^a,y^b),&b\geq a\geq2;\\IV_a&:R/(x^2+y^2,x^a),&a\geq3;\\V_a&:R/(x^2+y^2,x^a,x^{a-1}y),&a\geq3.\end{aligned}\] They are pairwise nonisomorphic. Stable realization for given dimensions is a further condition, not guaranteed merely by writing a family name; the three-generator cases require target dimension greater than source dimension.}
\GSfield{Reconstruction}{The published labels IV and V are interchanged relative to the notebook. The first ideal must contain \(xy\), not separate \(x,y\); the three-generator elliptic family needs \(x^{a-1}y\), not \(xy^{a-1}\). At odd \(a\), the latter is redundant modulo \((x^a,x^2+y^2)\); the supplied exact algebra check detects the resulting length error (\oursComputed).}
\GSfield{Fate}{\citet[\S7, pp.~150--151 of the Russian translation]{mather1971nice} proves the list and reports that Mather\textquotesingle s autumn 1965 letter to Morin contained an error, corrected by Morin; Arnold published the corrected account in 1968. That history makes the 1966 notebook especially valuable without licensing silent correction.}
\GSfield{Present status}{As of 9 November 2026: Established, with the hypotheses stated here.}
\GSfield{Assessment}{literal statement: false as written; publication match: partial; verification: read the result; development: annotated; confidence: high}
\end{GSitem}
\begin{GSitem}{1966-116.25}{B}{Finite-codimension algebra types and the exceptional remainder}
\GSfield{Kind; source}{claim; \GSpage{P000186}}
\GSfield{The notes}{The five families are followed by remainder: codimension infinity, the codimensions have been calculated, and a statement that \(\Sigma^{20}\) is subdivided into submanifolds.}
\GSfield{Statement}{In the source-dimension-at-most-target-dimension range, the stable finite local-algebra types of corank two are the five families just listed. The formal classification also has nonfinite residual cases; such exceptional germs are excluded from the finite stable classification. The primary paper computes orbit codimensions and stable-realizability conditions. The notebook gives no numerical codimension formula here, so none is attributed to this talk.}
\GSfield{Reconstruction}{The finite-dimensional local-algebra condition, stability, and a universal finite-jet stratum are distinct. Remainder of infinite codimension is retained as the announced generic-classification qualification, not a proof that every other individual \((2,0)\) germ is absent.}
\GSfield{Fate}{\citet[\S7]{mather1971nice} supplies the classification and numerical invariants; \citet[Theorem B]{mather1969algebras} states the realization test. The precise scope is narrower than an unqualified all singularities are stable.}
\GSfield{Present status}{As of 9 November 2026: Established, with the hypotheses stated here.}
\GSfield{Assessment}{literal statement: correct with missing hypotheses; publication match: partial; verification: read the result; development: annotated; confidence: medium}
\end{GSitem}

\GStalk{117}{}{}{P000186-P000188}{Year 1967}
This is the MCCME catalogue\textquotesingle s year-heading record, not a separate mathematical lecture. Its broad page interval overlaps the end of Arnold\textquotesingle s lecture and the February plan. P000186 ends with a next-meeting notice for 13 February 1967. No distinct theorem, problem, or construction is assigned to the heading.


\GStalk{118}{}{}{P000187}{Plan for 13 February 1967}
The catalogue labels the record 13-Feb-67 but supplies no machine-readable date; the label is preserved as a title. The plan lists Kirillov on Harish-Chandra (about ten lectures), Anosov on Maslov and semiclassics (two), Kazhdan on A.~Weil (six or seven), Dynin on the index (three), Anosov on first-order Pfaff equations (one), and Bernstein on the \(KG\)-functor (two). The next meeting is assigned to Kirillov and Anosov. The Russian surname \ru{Дынин} is Dynin, not the English translation\textquotesingle s Dynkin. These are topic announcements without mathematical assertions, so the plan has no occurrence-index items. The catalogue span P000186--P000188 is wider than the plan at the top of P000187.


\GStalk{119}{1967-02-13}{A. A. Kirillov}{P000187-P000191}{Invariant eigendistributions and the first descent construction}
The first lecture on Harish-Chandra begins after the programme on P000187. It states local integrability and regular-set analyticity for conjugation-invariant, centrally finite distributions, then sketches ellipticity, descent through centralizers, and reduction near the identity to the Lie algebra. Gelfand supplies explicit central operators for the general linear group. The lecture ends on P000191 before Dynin's heading. Several compressed proof steps need correction; the unexplained sentence about characters at infinity is retained as unidentified.

\begin{GSitem}{1966-119.1}{E}{Polynomial algebra of bi-invariant differential operators}
\GSfield{Kind; source}{theorem; \GSpage{P000187-P000188}}
\GSfield{The notes}{For a real semisimple group the notes identify the bi-invariant differential operators with a polynomial ring in \(\ell\) generators, where \(\ell\) is the rank; for \(SL(n)\), there is one generator in each degree \(2,\ldots,n\). Gelfand calls all these generators ``Laplace operators.}
\GSfield{Statement}{For a connected real semisimple Lie group \(G\), the algebra of complex-valued differential operators invariant under left and right translations is \(Z(U(\mathfrak g_{\mathbb C}))\). The Harish-Chandra isomorphism identifies it with \(S(\mathfrak h_{\mathbb C})^W\), hence with a polynomial algebra on \(\operatorname{rank}_{\mathbb C}\mathfrak g_{\mathbb C}\) generators. For \(\mathfrak{sl}_n\), their orders are \(2,\ldots,n\).}
\GSfield{Reconstruction}{Connectedness and complex coefficients are supplied. ``Rank'' is complex absolute rank, not real split rank. The historical name ``Laplace operator'' here includes higher-order central operators.}
\GSfield{Fate}{The central-operator construction predates the lecture; see \citet[pp.~103--104]{gelfand1950center}.}
\GSfield{Present status}{As of 9 November 2026, this remains an established construction under the stated hypotheses.}
\GSfield{Assessment}{literal statement: correct; publication match: partial; verification: read the result; development: none; confidence: high}
\end{GSitem}
\begin{GSitem}{1966-119.2}{E}{Gelfand’s cyclic trace generators}
\GSfield{Kind; source}{construction; \GSpage{P000187}; I. M. Gelfand}
\GSfield{The notes}{Gelfand writes \(E_{ij}f(x)=\frac d{dt}|_0f(x(1+te_{ij}))\), then \(\sum E_{ii}\), \(\sum E_{ik}E_{ki}\), and \(\sum E_{ij}E_{jk}E_{ki}\). He says subsequent expressions depend on the first \(n\) and refers to ``Gelfand, 1947.}
\GSfield{Statement}{In \(U(\mathfrak{gl}_n(\mathbb C))\), set \(c_r=\sum_{i_1,\ldots,i_r}E_{i_1i_2}\cdots E_{i_ri_1}\). These central elements, for \(1\le r\le n\), generate the center. Passing to \(\mathfrak{sl}_n\) removes \(c_1\), leaving generators of degrees \(2,\ldots,n\).}
\GSfield{Reconstruction}{The path \(1+te_{ii}\) is taken in \(GL_n\); the trace-free restriction gives the \(SL_n\) statement. Noncommuting factors are retained in their written cyclic order. The bibliographic date is not silently corrected to a first-circulation date.}
\GSfield{Fate}{The matching paper appeared in 1950, received 3 July 1948: \citet[\S2, example~1, pp.~105--106; \S3, Theorem~2, p.~108]{gelfand1950center}.}
\GSfield{Present status}{As of 9 November 2026, the construction is classical. The notebook's 1947 reference remains a provenance discrepancy; earlier circulation has not been established.}
\GSfield{Assessment}{literal statement: correct; publication match: exact; verification: read the result; development: none; confidence: high}
\end{GSitem}
\begin{GSitem}{1966-119.3}{E}{Centrally finite invariant distributions}
\GSfield{Kind; source}{definition; \GSpage{P000187-P000188}}
\GSfield{The notes}{The notes impose \(T^x=T\) and \(JT=0\), where \(R=\mathbb C[D_1,\ldots,D_\ell]\) and \(\dim R/J<\infty\); maximal ideals give ordinary joint eigenvalues, whereas \((D-\lambda)^2T=0\) illustrates a generalized eigencondition.}
\GSfield{Statement}{A distribution \(T\in\mathcal D'(G)\) is invariant and centrally finite if it is conjugation invariant and some finite-codimensional ideal \(J\subset Z(U(\mathfrak g_{\mathbb C}))\) annihilates it. For a maximal ideal \((D_i-\lambda_i)_i\), this is the joint eigendistribution condition. Finite codimension permits finite generalized joint eigenspaces; it does not require a maximal ideal.}
\GSfield{Fate}{This is the finite-ideal framework of \citet[\S2, Lemma~1]{harishchandra1965algebra}; see also \citet[Definition~1.6.1]{galinalaurent2003characters}.}
\GSfield{Present status}{As of 9 November 2026, this remains an established construction under the stated hypotheses.}
\GSfield{Assessment}{literal statement: correct; publication match: exact; verification: read the result; development: none; confidence: high}
\end{GSitem}
\begin{GSitem}{1966-119.4}{E}{Local integrability and regular-set analyticity}
\GSfield{Kind; source}{theorem; \GSpage{P000188-P000189}}
\GSfield{The notes}{Kirillov's main theorem says that a conjugation-invariant distribution annihilated by a finite-codimensional central ideal is locally summable; it is analytic at regular points but need not be continuous everywhere.}
\GSfield{Statement}{For a connected real reductive group in Harish-Chandra's setting, an invariant, centrally finite distribution on an invariant open set is represented by a locally integrable function. Its restriction to the regular semisimple set is real analytic. In particular, this applies to invariant eigendistributions and to distribution characters in the admissible irreducible setting. Analyticity is not asserted across the singular set.}
\GSfield{Reconstruction}{The open set, the reductive setting, and ``regular semisimple'' are made explicit. The theorem is local; no assertion that all such distributions are characters, nor any classification of all unitary representations, follows.}
\GSfield{Fate}{Already announced in 1963 and proved by 1965; \citet[Introduction; \S1.6]{galinalaurent2003characters} reviews the group theorem.}
\GSfield{Present status}{As of 9 November 2026, this is Harish-Chandra regularity, not an open problem. A later algebraic proof is given in \citet[Theorem~1.6.2, Corollary~1.6.3]{galinalaurent2003characters}.}
\GSfield{Assessment}{literal statement: correct with missing hypotheses; publication match: exact; verification: read the result; development: none; confidence: high}
\end{GSitem}
\begin{GSitem}{1966-119.5}{F}{Unidentified normalization of characters at infinity}
\GSfield{Kind; source}{remark; \GSpage{P000188}}
\GSfield{The notes}{The sentence reads \ru{Характеры обращ. в 1 в $\infty$ --- м.б. разрывы}: characters become \(1\) at infinity and may have discontinuities. The scan supports the numeral \(1\) and the infinity sign.}
\GSfield{Statement}{No precise assertion can be recovered: neither a character class, a normalization, nor the meaning of the limiting regime is specified.}
\GSfield{Reconstruction}{This is not replaced by a decay theorem or by the normalization of a spherical function at the identity.}
\GSfield{Fate}{Neither the group regularity theorem nor its modern restatement identifies this sentence; compare \citet[\S1.6]{galinalaurent2003characters}.}
\GSfield{Present status}{As of 9 November 2026, the intended normalization remains unidentified. There is no well-posed problem to develop.}
\GSfield{Assessment}{literal statement: not determinable; publication match: none found; verification: unresolved lead; development: none; confidence: low}
\end{GSitem}
\begin{GSitem}{1966-119.6}{E}{Regular transverse sections and elliptic regularity}
\GSfield{Kind; source}{method; \GSpage{P000188-P000189}}
\GSfield{The notes}{The notes stratify conjugacy classes, choose a transverse \(\ell\)-dimensional section, assert that the radial operators have finite codimension in all differential operators, and write \(\Delta=\sum D_i^2\). After obtaining a polynomial in \(\Delta\) in \(J\), they conclude \(\Delta T=0\).}
\GSfield{Statement}{Near a regular semisimple element, conjugation invariance reduces the distribution to a Cartan slice. Its radial joint equations form an analytic elliptic system. Finite central codimension supplies polynomial annihilators for the generators. Combining equal-order powers of these annihilators with their formal adjoints gives an elliptic scalar consequence. The radial invariant polynomial algebra is a finite invariant subalgebra in the appropriate Cartan polynomial ring; it is not of finite vector-space codimension in the full algebra of differential operators.}
\GSfield{Reconstruction}{The construction of an elliptic consequence repairs the sketch. With generators of unequal orders, the displayed sum of squares need not be elliptic in the ordinary single-order sense. The immediate conclusion from \(P(\Delta)\in J\) is \(P(\Delta)T=0\), not \(\Delta T=0\), unless the latter symbol has been renamed.}
\GSfield{Fate}{The analytic regular-set step is in \citet[\S2, Lemma~1]{harishchandra1965algebra}.}
\GSfield{Present status}{As of 9 November 2026, the intended elliptic-system argument is established; the literal finite-codimension and scalar-operator shortcuts require emendation.}
\GSfield{Assessment}{literal statement: false as written; publication match: partial; verification: read the result; development: annotated; confidence: high}
\end{GSitem}
\begin{GSitem}{1966-119.7}{E}{Induction through semisimple centralizers and the Lie algebra}
\GSfield{Kind; source}{method; \GSpage{P000189}}
\GSfield{The notes}{The proof is divided into semisimple, ``nilpotent,'' and identity cases. A noncentral semisimple point descends to its smaller reductive centralizer; logarithm and exponential are used near the identity; the nilpotent case requires another argument.}
\GSfield{Statement}{Harish-Chandra descent reduces the local problem at a noncentral semisimple element to its reductive centralizer, of smaller semisimple dimension. Central elements are translated to the identity. A local analytic exponential comparison reduces the remaining group problem to the Lie algebra, where nilpotent orbits require a separate argument. In a linear algebraic group the group counterpart of a nilpotent Lie-algebra element is unipotent.}
\GSfield{Reconstruction}{The comparison near the identity includes analytic density and conjugation factors and is not the bare substitution \(X\mapsto\exp X\) for every central differential operator. The three cases are a proof plan, not a disjoint exhaustive Jordan decomposition as literally listed.}
\GSfield{Fate}{The 1965 proof explicitly uses centralizers and nilpotent analysis: \citet[Introduction]{harishchandra1965algebra}.}
\GSfield{Present status}{As of 9 November 2026, this remains an established construction under the stated hypotheses.}
\GSfield{Assessment}{literal statement: correct with missing hypotheses; publication match: partial; verification: read the result; development: none; confidence: high}
\end{GSitem}
\begin{GSitem}{1966-119.8}{E}{The conjugation map is a submersion on the descent domain}
\GSfield{Kind; source}{theorem; \GSpage{P000189-P000190}}
\GSfield{The notes}{For semisimple \(a\), the notes put \(\mathfrak z=\mathfrak g^a\), let \(\Xi\) be its connected subgroup, and define \(\varphi(x,y)=xayx^{-1}\), \(V_a(y)=\det(\operatorname{Ad}((ay)^{-1})-1)_{\mathfrak g/\mathfrak z}\), and \(\Xi'=\{V_a\ne0\}\). ``Theorem 1 (Lemma 14)'' says \(\operatorname{rank}d\varphi=\dim G\).}
\GSfield{Statement}{Let \(a\) be semisimple, \(\Xi=Z_G(a)^0\), and \(\Xi'\) the indicated open subset. Then \(\varphi:G\times\Xi'\to G\) is a smooth analytic submersion. In tangent coordinates at \((1,y)\), obtained by right multiplication \(y\exp(tZ)\) and left trivialization at \(ay\), \[d\varphi(X,Z)=(\operatorname{Ad}((ay)^{-1})-1)X+Z.\] Invertibility on \(\mathfrak g/\mathfrak z\) and the \(Z\)-term prove surjectivity.}
\GSfield{Reconstruction}{The determinant is a function, not an assertion of nonvanishing for every \(y\in\Xi\); it is nonzero at \(y=1\), so the useful neighborhood is nonempty. This is a submersion, not a globally unique representative of each conjugacy class.}
\GSfield{Fate}{This belongs to the 1965 group descent proof, \citet[pp.~457--508]{harishchandra1965group}; the notebook's Lemma~14 number was not independently verified.}
\GSfield{Present status}{As of 9 November 2026, this remains an established construction under the stated hypotheses.}
\GSfield{Assessment}{literal statement: correct; publication match: partial; verification: publisher metadata checked; development: none; confidence: high}
\end{GSitem}
\begin{GSitem}{1966-119.9}{E}{Fiber integration and distributional pullback under a submersion}
\GSfield{Kind; source}{construction; \GSpage{P000190-P000191}}
\GSfield{The notes}{For a submersion \(\pi:M\to N\) and smooth volume measures, \(\alpha\in C_c^\infty(M)\) is sent to \(f_\alpha\in C_c^\infty(N)\) by \(\int_M\alpha(F\circ\pi)d\mu=\int_NFf_\alpha d\nu\). The verbal phrase says generalized functions move ``from \(M\) to \(N\).'' The next lecture corrects the direction in a diagram.}
\GSfield{Statement}{With positive smooth densities fixed, fiber integration defines a continuous map \(\pi_!:C_c^\infty(M)\to C_c^\infty(N)\). Its transpose defines the pullback \(\pi^*:\mathcal D'(N)\to\mathcal D'(M)\), with \(\langle\pi^*T,\alpha\rangle=\langle T,\pi_!\alpha\rangle\). No properness is needed for compactly supported test functions; pushing arbitrary distributions forward would require an additional support/properness condition.}
\GSfield{Reconstruction}{The map is phrased in densities to avoid orientation assumptions. The verbal direction of transfer is separated from the displayed, correct test-function formula.}
\GSfield{Fate}{A standard distributional ingredient of the descent reviewed in \citet[Introduction; \S1.6]{galinalaurent2003characters}; the notebook itself gives its local product proof.}
\GSfield{Present status}{As of 9 November 2026, this remains an established construction under the stated hypotheses.}
\GSfield{Assessment}{literal statement: correct with missing hypotheses; publication match: partial; verification: read the result; development: none; confidence: high}
\end{GSitem}
\begin{GSitem}{1966-119.10}{E}{Mixed differential notation for conjugation}
\GSfield{Kind; source}{definition; \GSpage{P000191}}
\GSfield{The notes}{Kirillov writes \(f(x;D_1:y;D_2)\) for differentiation in both variables and uses \(f(x:y)=f(xyx^{-1})\), omitting identity operators.}
\GSfield{Statement}{For a smooth function \(f\) on \(G\), put \(F(x,y)=f(xyx^{-1})\). Then \(f(x;D_1:y;D_2)=(D_{1,x}D_{2,y}F)(x,y)\), with the order and left/right invariant conventions fixed consistently. This notation introduces the bilinear operator \(\Gamma_y\) in record 121.}
\GSfield{Fate}{Expository notation for the group proof cited as \citet{harishchandra1965group}; no independent theorem is asserted.}
\GSfield{Present status}{As of 9 November 2026, this remains an established construction under the stated hypotheses.}
\GSfield{Assessment}{literal statement: correct; publication match: partial; verification: publisher metadata checked; development: none; confidence: high}
\end{GSitem}

\GStalk{120}{1967-02-13}{A. S. Dynin}{P000191-P000192}{Pseudodifferential operators on bundles and the de Rham example}
Dynin begins after Kirillov on P000191, globalizing pseudodifferential operators by coordinate cutoffs and the off-diagonal smoothing condition. The only worked example is the operator \(d+d^*\), its Laplacian, and the Hodge identification with de Rham cohomology. The final timetable is administrative, not another mathematical result.

\begin{GSitem}{1966-120.1}{E}{Pseudodifferential operators between vector bundles}
\GSfield{Kind; source}{definition; \GSpage{P000191}}
\GSfield{The notes}{For compact smooth \(X\) and bundles \(E_1,E_2\), locally cut off operators \(\psi_iA\varphi_i\) must be ordinary matrix pseudodifferential operators; if \(\varphi\psi=0\), \(\psi A\varphi\) is smoothing.}
\GSfield{Statement}{For smooth finite-rank complex bundles on a closed smooth manifold, a continuous linear map on smooth sections is a pseudodifferential operator of order \(m\) if its localized expressions lie in the prescribed coordinate symbol class and its Schwartz kernel is smooth off the diagonal. Cutoffs with disjoint supports give smoothing compositions. The stated pointwise condition \(\varphi\psi=0\) also suffices for finite-order pseudodifferential kernels: the product \(\psi(x)\varphi(y)\) is flat along the diagonal where the supports touch.}
\GSfield{Reconstruction}{The symbol class, order, and continuity are supplied. Pointwise zero product is not erroneously replaced by a supposedly necessary positive separation of supports.}
\GSfield{Fate}{Dynin's prelecture exposition gives this localization: \citet[\S1.2]{dynin1966index}, submitted March 1965, published 1966.}
\GSfield{Present status}{As of 9 November 2026, this remains an established construction under the stated hypotheses.}
\GSfield{Assessment}{literal statement: correct with missing hypotheses; publication match: exact; verification: read the result; development: none; confidence: high}
\end{GSitem}
\begin{GSitem}{1966-120.2}{E}{The Hodge–de Rham operator and its square}
\GSfield{Kind; source}{construction; \GSpage{P000192}}
\GSfield{The notes}{The notes use all inhomogeneous differential forms, set \(\delta=d^*\), call \(d+\delta\) a Cauchy--Riemann operator, and identify \((d+\delta)^2\) with the Laplacian.}
\GSfield{Statement}{On a closed Riemannian manifold \(X\), \(D=d+d^*\) acts on \(\Omega^*(X;\mathbb C)\), exchanges even and odd forms, and satisfies \(D^2=dd^*+d^*d=\Delta\), the Hodge Laplacian, which preserves form degree. The bundles are \(\Lambda^*T^*X\otimes\mathbb C\); a metric can identify tangent and cotangent bundles.}
\GSfield{Reconstruction}{The modern name is the Hodge--de Rham Dirac operator. The notebook's use of \(\Lambda TX\) is interpreted via the metric rather than as the natural bundle of differential forms.}
\GSfield{Fate}{Already in \citet[\S1.3, example~7]{dynin1966index}.}
\GSfield{Present status}{As of 9 November 2026, this remains an established construction under the stated hypotheses.}
\GSfield{Assessment}{literal statement: correct with missing hypotheses; publication match: exact; verification: read the result; development: none; confidence: high}
\end{GSitem}
\begin{GSitem}{1966-120.3}{E}{Harmonic representatives for de Rham cohomology}
\GSfield{Kind; source}{theorem; \GSpage{P000192}}
\GSfield{The notes}{Hodge's theorem is written \(\ker\Delta|_{\Lambda^kT(X)}=H^k(X,\mathbb C)\); a heat-equation proof is attributed informally to ``Moren.}
\GSfield{Statement}{For a closed smooth Riemannian manifold, the map from harmonic complex \(k\)-forms to their de Rham classes is an isomorphism \(\mathcal H^k(X)\simeq H^k_{\mathrm{dR}}(X;\mathbb C)\). Every cohomology class has a unique harmonic representative. A manifold with boundary requires appropriate absolute or relative boundary conditions and a corresponding cohomology theory.}
\GSfield{Reconstruction}{The equality is a canonical isomorphism, not equality of sets. No specific Moren/Maurin book or heat proof is identified from the surname alone.}
\GSfield{Fate}{The Hodge example is used in \citet[\S1.3, examples~7--8]{dynin1966index}.}
\GSfield{Present status}{As of 9 November 2026, this remains an established construction under the stated hypotheses.}
\GSfield{Assessment}{literal statement: correct with missing hypotheses; publication match: exact; verification: read the result; development: none; confidence: high}
\end{GSitem}

\GStalk{121}{1967-02-20}{A. A. Kirillov}{P000193-P000195}{Enveloping-algebra construction of radial differential operators}
The second Harish-Chandra lecture constructs \(\Gamma_y\), transferring differentiation in the conjugating variable to the section. Its associated-graded calculation yields a PBW decomposition and the constant transverse term defining the radial operator. The first part of P000195 still belongs to this lecture; the dated 27 February heading starts record 122.

\begin{GSitem}{1966-121.1}{E}{Local centralizer slice and pointwise enveloping-algebra description}
\GSfield{Kind; source}{construction; \GSpage{P000193}}
\GSfield{The notes}{A semisimple \(a\) is fixed; the connected centralizer is called a section meeting every conjugacy class once. The notes write \(\mathfrak z=\{X:\operatorname{ad}_aX=X\}\) and \(D_y\in\mathfrak g\) for the left-invariant operator agreeing with a differential operator \(D\) at \(y\).}
\GSfield{Statement}{Use \(\mathfrak z=\ker(\operatorname{Ad}(a)-1)\). Left trivialization identifies the value of an order-\(m\) differential operator at \(y\) with an element \(D_y\in U_{\le m}(\mathfrak g_{\mathbb C})\); smoothly varying \(D_y\) describes \(D\). The centralizer provides a redundant local descent submanifold on the regularity domain, not a global one-point-per-orbit cross-section.}
\GSfield{Reconstruction}{Here \(\operatorname{Ad}\), rather than the Lie-algebra map \(\operatorname{ad}\), acts because \(a\in G\). The notation \(\mathfrak g\) for the enveloping algebra is explicitly clarified on P000194. Weyl-group ambiguity survives on a Cartan section.}
\GSfield{Fate}{This is the classical descent framework of \citet{harishchandra1965group}; the exact source lemma was not read.}
\GSfield{Present status}{As of 9 November 2026, this remains an established construction under the stated hypotheses.}
\GSfield{Assessment}{literal statement: correct with missing hypotheses; publication match: partial; verification: publisher metadata checked; development: none; confidence: high}
\end{GSitem}
\begin{GSitem}{1966-121.2}{E}{The representation sigma and the transfer operator Gamma}
\GSfield{Kind; source}{construction; \GSpage{P000193-P000194}}
\GSfield{The notes}{The formula is \(\Gamma_y(u\otimes v)=\sigma_y(u)v\), with \(\sigma_y(X)v=\operatorname{Ad}(y^{-1})X\,v-vX\). It is used to replace operators in the two variables of \(f(xyx^{-1})\) by one operator in \(y\).}
\GSfield{Statement}{For fixed \(y\in G\), the map \(X\mapsto L_{\operatorname{Ad}(y^{-1})X}-R_X\) is a Lie-algebra representation on \(U(\mathfrak g_{\mathbb C})\), hence extends to its enveloping algebra. Define the bilinear transfer \(\Gamma_y(u\otimes v)=\sigma_y(u)v\). With the notebook's invariant differentiation convention it encodes \(f(x;u:y;v)=f(x:y;\Gamma_y(u\otimes v))\).}
\GSfield{Reconstruction}{The left and right multiplications act on the enveloping algebra. Their commutators have opposite signs, which is essential to the representation identity. For descent at \(ay\), use \(\Gamma_{ay}\).}
\GSfield{Fate}{The construction is expository Harish-Chandra descent, \citet{harishchandra1965group}; publication matching of this formula is partial.}
\GSfield{Present status}{As of 9 November 2026, this remains an established construction under the stated hypotheses.}
\GSfield{Assessment}{literal statement: correct; publication match: partial; verification: publisher metadata checked; development: none; confidence: high}
\end{GSitem}
\begin{GSitem}{1966-121.3}{E}{PBW leading term and the transverse decomposition}
\GSfield{Kind; source}{theorem; \GSpage{P000194}}
\GSfield{The notes}{The leading term of \(\Gamma_y(X_1\cdots X_r\otimes Y_1\cdots Y_s)\) is the product of \(X'_i=(\operatorname{Ad}(y^{-1})-1)X_i\) and the \(Y_j\). A complement \(\mathfrak g=\mathfrak z+\mathfrak q\), stable under \(\mathfrak z\), is introduced, and \(\Gamma_y\) is declared an isomorphism on the corresponding generated spaces.}
\GSfield{Statement}{Choose an \(\operatorname{Ad}(a)\)-stable complement \(\mathfrak q\) to \(\mathfrak z\), also stable under the connected centralizer. Identify \(S(\mathfrak q)\) with its symmetrized PBW subspace in \(U(\mathfrak g_{\mathbb C})\). For \(y\in\Xi'\), \[\Gamma_{ay}:S(\mathfrak q_{\mathbb C})\otimes U(\mathfrak z_{\mathbb C})\longrightarrow U(\mathfrak g_{\mathbb C})\] is a filtered vector-space isomorphism: its associated graded is multiplication after the invertible transformation \((\operatorname{Ad}((ay)^{-1})-1)|_{\mathfrak q}\).}
\GSfield{Reconstruction}{The space generated by \(\mathfrak q\) must mean the PBW symmetric subspace, not the associative subalgebra: \(\mathfrak q\) need not be a Lie subalgebra. The displayed leading-term equality holds modulo lower filtration. If \(X\) commutes with \(y\), this leading term vanishes, whereas \(\sigma_y(X)=[X,\cdot]\) need not vanish.}
\GSfield{Fate}{This is the filtered-algebra mechanism of the cited 1965 descent construction, \citet{harishchandra1965group}.}
\GSfield{Present status}{As of 9 November 2026, this remains an established construction under the stated hypotheses.}
\GSfield{Assessment}{literal statement: correct with missing hypotheses; publication match: partial; verification: publisher metadata checked; development: annotated; confidence: high}
\end{GSitem}
\begin{GSitem}{1966-121.4}{E}{Radial part as the constant transverse coefficient}
\GSfield{Kind; source}{construction; \GSpage{P000194-P000195}}
\GSfield{The notes}{On \(a\Xi\), \(\Delta=\delta_{a,G/\Xi}(D)\) is defined by \(\Delta_{ay}=\alpha(D_{ay})\), where \(D-\alpha(D)\) belongs to \(\Gamma_y(\mathfrak q^+\otimes\mathfrak z)\). It is said to act identically on orbit-constant functions.}
\GSfield{Statement}{On \(a\Xi'\), expand \(D_{ay}\) through the PBW isomorphism of the preceding item and project to the degree-zero \(S(\mathfrak q)\) component. The resulting smooth \(U(\mathfrak z)\)-valued coefficients define the radial operator \(\delta_a(D)\). For a locally conjugation-invariant smooth \(f\), \[(Df)|_{a\Xi'}=\delta_a(D)(f|_{a\Xi'}).\]}
\GSfield{Reconstruction}{The omitted enveloping/symmetric-algebra notation and \(ay\)-dependence are restored. The result concerns invariant functions on the descent domain; it is not restriction of an arbitrary differential operator on arbitrary functions to an arbitrary submanifold.}
\GSfield{Fate}{A classical radial-part construction, \citet{harishchandra1965group}; exact group-paper formula verification remains metadata-level.}
\GSfield{Present status}{As of 9 November 2026, this remains an established construction under the stated hypotheses.}
\GSfield{Assessment}{literal statement: correct with missing hypotheses; publication match: partial; verification: publisher metadata checked; development: none; confidence: high}
\end{GSitem}

\GStalk{122}{1967-02-27}{A. A. Kirillov}{P000195-P000197}{Distributional descent and the induction step for regularity}
The third lecture returns to continuous fiber integration, pulls an invariant distribution back to \(G\times\Xi'\), and factors it as Haar measure times a distribution on the centralizer. It states intertwining with radial differential operators and finite central codimension, thereby reducing local integrability to a smaller group. The notebook itself rejects an attempted explicit inverse-substitution formula.

\begin{GSitem}{1966-122.1}{E}{Continuity of fiber integration}
\GSfield{Kind; source}{theorem; \GSpage{P000195-P000196}}
\GSfield{The notes}{The existence and continuity of \(\alpha\mapsto f_\alpha\) are stated in the usual compactly supported smooth topologies. The first line mistakenly repeats \(C_c^\infty(M)\) as the target; the product case integrates over the extra variable.}
\GSfield{Statement}{For a smooth submersion \(\pi:M\to N\) and smooth positive densities, fiber integration is continuous \(\pi_!:C_c^\infty(M)\to C_c^\infty(N)\) in the standard LF test-function topology. On a product with product density it is \((\pi_!\alpha)(y)=\int_P\alpha(y,z)\,dz\); coordinate products and a partition of unity give the general map.}
\GSfield{Reconstruction}{The target is \(N\), as the displayed identity and variable \(y\) already require. This repeats and proves the construction at 1966-119.9 rather than introducing a new research question.}
\GSfield{Fate}{The notebook's local product argument establishes the standard construction used in \citet{harishchandra1965group}.}
\GSfield{Present status}{As of 9 November 2026, this remains an established construction under the stated hypotheses.}
\GSfield{Assessment}{literal statement: false as written; publication match: partial; verification: publisher metadata checked; development: none; confidence: high}
\end{GSitem}
\begin{GSitem}{1966-122.2}{E}{Translation-invariant distributions are Haar measures}
\GSfield{Kind; source}{theorem; \GSpage{P000196}}
\GSfield{The notes}{A lemma says that a generalized function invariant under translations is integration against Haar measure, up to a constant.}
\GSfield{Statement}{On a connected Lie group \(G\), a left-translation-invariant distribution is a constant multiple of left Haar measure. More generally, a distribution on \(G\times V\) invariant under translations of the first factor is uniquely \(dg\otimes S\), with \(S\in\mathcal D'(V)\).}
\GSfield{Reconstruction}{The product form is the parameter version needed below. For reductive \(G\), unimodularity removes a left/right Haar distinction; the formulation still specifies left invariance.}
\GSfield{Fate}{An elementary, preexisting ingredient of the descent exposition \citep{harishchandra1965group}.}
\GSfield{Present status}{As of 9 November 2026, this remains an established construction under the stated hypotheses.}
\GSfield{Assessment}{literal statement: correct; publication match: partial; verification: publisher metadata checked; development: none; confidence: high}
\end{GSitem}
\begin{GSitem}{1966-122.3}{E}{The descended distribution sigma\_T}
\GSfield{Kind; source}{construction; \GSpage{P000195-P000196}}
\GSfield{The notes}{Pulling \(T\) back through \(\varphi(x,y)=xayx^{-1}\) gives \(T'\), translation invariant in the first variable; hence \(T'=1\otimes\sigma_T\).}
\GSfield{Statement}{For invariant \(T\) on the open image of the descent submersion, \(\varphi^*T\) is invariant under left translations of its first factor. Fix Haar density \(dg\); then a unique \(\sigma_T\in\mathcal D'(\Xi')\) satisfies \(\varphi^*T=dg\otimes\sigma_T\), under the chosen product-density convention. For \(\gamma\in C_c^\infty(G)\) with \(\int\gamma dg=1\), \[\langle\sigma_T,\beta\rangle=\langle T,\varphi_!(\gamma\otimes\beta)\rangle.\]}
\GSfield{Reconstruction}{The final formula is the valid fiber-integration interpretation, not a literal substitution of \(a^{-1}z^{-1}xz\) into a function defined only on the centralizer.}
\GSfield{Fate}{Classical Harish-Chandra descent, with the group-paper bibliographic identification \citep{harishchandra1965group}.}
\GSfield{Present status}{As of 9 November 2026, this remains an established construction under the stated hypotheses.}
\GSfield{Assessment}{literal statement: correct; publication match: partial; verification: publisher metadata checked; development: none; confidence: high}
\end{GSitem}
\begin{GSitem}{1966-122.4}{E}{The rejected explicit formula for descent}
\GSfield{Kind; source}{claim; \GSpage{P000196}}
\GSfield{The notes}{The notebook proposes \(\langle\sigma_T,\beta\rangle=\langle T,\int_G\gamma(z)\beta(a^{-1}z^{-1}xz)dz\rangle\), then adds \ru{Она не имеет смысла} (``It has no meaning''). A diagonal dashed line crosses the formula region.}
\GSfield{Statement}{As written, the formula is undefined: \(\beta\) is a test function on \(\Xi\), but \(a^{-1}z^{-1}xz\) is generally not in \(\Xi\). Arbitrary extensions of \(\beta\) do not repair it canonically. The correct test function on \(G\) is \(\varphi_!(\gamma\otimes\beta)\), defined by integration along fibers and its Jacobian, as in the preceding item.}
\GSfield{Reconstruction}{The notebook's own objection is preserved. The diagonal line's function as cancellation or connector remains uncertain; the mathematical type mismatch does not depend on that palaeographic choice.}
\GSfield{Fate}{An abandoned formula inside an established descent construction \citep{harishchandra1965group}; it is not an unsolved theorem.}
\GSfield{Present status}{As of 9 November 2026, the literal formula is invalid. The fiber-integral definition supplies the established replacement.}
\GSfield{Assessment}{literal statement: false as written; publication match: partial; verification: publisher metadata checked; development: annotated; confidence: high}
\end{GSitem}
\begin{GSitem}{1966-122.5}{E}{Intertwining radial operators with descended distributions}
\GSfield{Kind; source}{theorem; \GSpage{P000197}}
\GSfield{The notes}{The lemma is \(\sigma_{DT}=\Delta\sigma_T\), with \(\Delta=\delta_{a,G/\Xi}(D)\); the notes emphasize the use of \(\Gamma_{ay}\).}
\GSfield{Statement}{For an invariant differential operator \(D\), Harish-Chandra descent intertwines its action on invariant distributions with the corresponding radial operator on the centralizer. If distributions are defined by duality using densities, the radial operator on distributions must be defined with the compatible formal transpose and density conjugation. With those conventions, \(\sigma_{DT}=\delta_a(D)\sigma_T\).}
\GSfield{Reconstruction}{The smooth-function radial identity alone does not justify using the same coefficient expression on distributions with arbitrary densities. The compatible transpose convention is an essential supplied hypothesis.}
\GSfield{Fate}{The identity belongs to the cited 1965 group descent argument \citep{harishchandra1965group}; its exact printed lemma was not verified.}
\GSfield{Present status}{As of 9 November 2026, this remains an established construction under the stated hypotheses.}
\GSfield{Assessment}{literal statement: correct with missing hypotheses; publication match: partial; verification: publisher metadata checked; development: none; confidence: high}
\end{GSitem}
\begin{GSitem}{1966-122.6}{E}{Finite central codimension survives descent}
\GSfield{Kind; source}{theorem; \GSpage{P000197}}
\GSfield{The notes}{The notes say the central operators' images generate an ideal of finite codimension in \(Z(U(\mathfrak z))\), and display an ideal mapping from \(Z(U(\mathfrak g))\) to \(Z(U(\mathfrak z))\).}
\GSfield{Statement}{After the appropriate Harish-Chandra radial conjugation, the centralizer homomorphism identifies the relevant centers with \(S(\mathfrak h)^W\subset S(\mathfrak h)^{W_{\mathfrak z}}\). This is a finite extension. Consequently, if \(J\subset Z(U(\mathfrak g))\) has finite codimension, its extended ideal in \(Z(U(\mathfrak z))\) has finite codimension and annihilates the descended distribution.}
\GSfield{Reconstruction}{It is the image of the annihilating ideal, not the entire center, whose extension matters; the entire unital center already contains \(1\). The displayed raw radial operator includes conjugations not recorded in the notebook's schematic center diagram.}
\GSfield{Fate}{This is the finite-invariant-algebra step in classical descent \citep{harishchandra1965group}.}
\GSfield{Present status}{As of 9 November 2026, this remains an established construction under the stated hypotheses.}
\GSfield{Assessment}{literal statement: correct with missing hypotheses; publication match: partial; verification: publisher metadata checked; development: annotated; confidence: high}
\end{GSitem}
\begin{GSitem}{1966-122.7}{E}{Recovering local integrability after descent}
\GSfield{Kind; source}{method; \GSpage{P000197}}
\GSfield{The notes}{The notes conclude \(\sigma_T\in L^1_{\mathrm{loc}}\Rightarrow T\in L^1_{\mathrm{loc}}\), then translate a central \(a\) to the identity and refer to special techniques in three Harish-Chandra papers.}
\GSfield{Statement}{On the image of a surjective submersion, local product coordinates show that a distribution is locally integrable if its pullback is locally integrable. Since \(\varphi^*T=dg\otimes\sigma_T\), local integrability of \(\sigma_T\) gives that of \(T\) near the chosen semisimple orbit. Noncentral descent reduces dimension; central translation, exponential comparison, and the nilpotent analysis complete the regularity theorem.}
\GSfield{Reconstruction}{The statement is local on the covered image. This induction step does not supply the omitted nilpotent proof or identify exactly which three papers the shorthand counts.}
\GSfield{Fate}{The completed regularity proof predates the lecture: \citet[Introduction; Theorem~1]{harishchandra1965algebra}.}
\GSfield{Present status}{As of 9 November 2026, this remains an established construction under the stated hypotheses.}
\GSfield{Assessment}{literal statement: correct with missing hypotheses; publication match: partial; verification: read the result; development: none; confidence: high}
\end{GSitem}

\GStalk{123}{1967-03-06}{A. A. Kirillov}{P000198}{The Lorentzian rank-one model for invariant regularity}
A short final Kirillov note isolates the form \(x^2+y^2-z^2\), its wave operator, and invariant distributions annihilated by a polynomial in that operator. The question of local summability is the rank-one Lie-algebra case of the theorem developed in the preceding lectures. Arnold's report begins later on the same page.

\begin{GSitem}{1966-123.1}{E}{Local integrability of Lorentz-invariant polynomial wave solutions}
\GSfield{Kind; source}{question; \GSpage{P000198}}
\GSfield{The notes}{For \(q=x^2+y^2-z^2\), \(T\) is invariant under transformations preserving \(q\), and \(P(\partial_x^2+\partial_y^2-\partial_z^2)T=0\). The notes ask: is \(T\) locally summable?}
\GSfield{Statement}{Let \(T\in\mathcal D'(\mathbb R^3)\) be invariant under \(SO_0(2,1)\), and let \(0\ne P\in\mathbb C[t]\) satisfy \(P(\Box)T=0\), where \(\Box=\partial_x^2+\partial_y^2-\partial_z^2\). Then \(T\in L^1_{\mathrm{loc}}(\mathbb R^3)\). Full \(O(2,1)\)-invariance is more than needed.}
\GSfield{Reconstruction}{A linear identification with \(\mathfrak{sl}_2(\mathbb R)\) identifies \(q\), up to a nonzero scalar, with its invariant quadratic form. The adjoint group is \(SO_0(2,1)\), and the invariant constant-coefficient algebra is generated by \(\Box\); \((P)\) has finite codimension. The nonzero-polynomial hypothesis is essential.}
\GSfield{Fate}{Already covered by \citet[Theorem~1; \S2]{harishchandra1965algebra}; the notebook question is an expository specialization.}
\GSfield{Present status}{As of 9 November 2026, the answer is affirmative. It is not a surviving open regularity problem.}
\GSfield{Assessment}{literal statement: correct with missing hypotheses; publication match: exact; verification: read the result; development: none; confidence: high; problem \hyperref[prob:GP-1966-18]{GP-1966-18}}
\end{GSitem}
\begin{GSitem}{1966-123.2}{E}{Why the invariant system is elliptic away from the null cone}
\GSfield{Kind; source}{method; \GSpage{P000198}}
\GSfield{The notes}{The notes say that away from \(q=0\) the system is elliptic, referring to invariant equations plus the polynomial in the wave operator.}
\GSfield{Statement}{On each regular orbit component with \(q\ne0\), an invariant local function or distribution depends on the transverse variable \(s=q\). The radial expression is \(\Box f(q)=4q f''(q)+6f'(q)\). Thus a nonzero polynomial in \(\Box\) yields an ordinary differential equation whose leading coefficient is nonzero for \(q\ne0\). This proves analytic regularity there; the null cone, including its vertex, requires the invariant-distribution theorem.}
\GSfield{Reconstruction}{This is ellipticity after orbit reduction, or of the full invariant system; the scalar wave operator is not elliptic on \(\mathbb R^3\). \oursComputed\ The radial formula is an elementary reconstruction checked by applying the chain rule separately to all three coordinates.}
\GSfield{Fate}{The general regular-set argument appears in \citet[\S2, Lemma~1]{harishchandra1965algebra}.}
\GSfield{Present status}{As of 9 November 2026, this remains an established construction under the stated hypotheses.}
\GSfield{Assessment}{literal statement: correct with missing hypotheses; publication match: partial; verification: read the result; development: annotated; confidence: high}
\end{GSitem}

\GStalk{124}{1967-03-06}{V. I. Arnold}{P000198-P000199}{Surface diffeomorphism groups; adjacent Anosov note on operator Fourier transformation}
Arnold reports Earle--Eells on the topology of the identity component of the diffeomorphism group of an oriented surface, with genus-by-genus answers and the quasiconformal proof strategy. P000199 begins a distinct, explicit \ru{Аносов} heading and develops an operator-valued Fourier transform attributed to Maslov. The catalogue supplies no separate record for this subsection; its three items are retained here with Anosov as the item-level speaker. This grouping preserves source coverage without attributing Anosov's mathematics to Arnold.

\begin{GSitem}{1966-124.1}{E}{Diffeomorphisms of the sphere fixing three points}
\GSfield{Kind; source}{theorem; \GSpage{P000198}}
\GSfield{The notes}{For \(X=S^2\), the notes give \(\operatorname{Diff}_0X=G\times\operatorname{Diff}(S^2-a,b,c)\), where \(G\) is the conformal group and the second factor is contractible; the three points are to stay fixed.}
\GSfield{Statement}{Give smooth diffeomorphism groups the \(C^\infty\) topology. For three ordered distinct points \(a,b,c\in S^2\), the subgroup of orientation-preserving diffeomorphisms fixing them pointwise is contractible, and composition gives a homeomorphism \[\operatorname{Diff}_0(S^2)\simeq PSL_2(\mathbb C)\times\operatorname{Diff}_0(S^2;a,b,c).\] Consequently \(\operatorname{Diff}_0(S^2)\) has the homotopy type of \(SO(3)\).}
\GSfield{Reconstruction}{The puncture notation means fixed marked points, not arbitrary diffeomorphisms of an open punctured surface. The product is a homeomorphism of spaces, not generally a direct-product decomposition of groups.}
\GSfield{Fate}{\citet[Theorem~1(a)]{earleeells1967diffeomorphism}, communicated 6 February 1967, printed July 1967. The notebook explicitly reports Earle--Eells; separate public circulation before March was not established.}
\GSfield{Present status}{As of 9 November 2026, this remains an established construction under the stated hypotheses.}
\GSfield{Assessment}{literal statement: correct with missing hypotheses; publication match: exact; verification: read the result; development: none; confidence: high}
\end{GSitem}
\begin{GSitem}{1966-124.2}{E}{Diffeomorphisms of the torus fixing one point}
\GSfield{Kind; source}{theorem; \GSpage{P000198}}
\GSfield{The notes}{For \(X=T^2\), the notes give \(\operatorname{Diff}_0(T^2)=T^2\times\operatorname{Diff}(T^2-a)\), with the fixed-point factor contractible.}
\GSfield{Statement}{The evaluation at a point and the translation action identify \(\operatorname{Diff}_0(T^2)\) homeomorphically with \(T^2\times\operatorname{Diff}_0(T^2;a)\). The subgroup fixing \(a\) and homotopic to the identity is contractible. Thus \(\operatorname{Diff}_0(T^2)\simeq T^2\).}
\GSfield{Reconstruction}{Homotopy to the identity is part of the fixed-point subgroup; arbitrary mapping classes fixing the point are not included.}
\GSfield{Fate}{\citet[Theorem~1(b)]{earleeells1967diffeomorphism}; the manuscript was communicated in February and printed in July 1967; the precise first public circulation is not established.}
\GSfield{Present status}{As of 9 November 2026, this remains an established construction under the stated hypotheses.}
\GSfield{Assessment}{literal statement: correct with missing hypotheses; publication match: exact; verification: read the result; development: none; confidence: high}
\end{GSitem}
\begin{GSitem}{1966-124.3}{E}{Contractibility in genus at least two}
\GSfield{Kind; source}{theorem; \GSpage{P000198}}
\GSfield{The notes}{For genus at least \(2\), \(\operatorname{Diff}_0X\) is declared contractible.}
\GSfield{Statement}{If \(X\) is a closed, connected, oriented smooth surface of genus \(g\ge2\), its smooth diffeomorphisms homotopic to the identity form the identity component, and this topological group is contractible in the \(C^\infty\) topology.}
\GSfield{Reconstruction}{Closedness, connectedness, and the topology on the group are supplied; boundary and puncture variants are separate statements.}
\GSfield{Fate}{\citet[Theorem~1(c), Corollary]{earleeells1967diffeomorphism}; communicated to the journal in February 1967 and expressly reported as their result in these March notes.}
\GSfield{Present status}{As of 9 November 2026, this remains an established construction under the stated hypotheses.}
\GSfield{Assessment}{literal statement: correct with missing hypotheses; publication match: exact; verification: read the result; development: none; confidence: high}
\end{GSitem}
\begin{GSitem}{1966-124.4}{E}{Quasiconformal uniformization and the Teichmüller fibration}
\GSfield{Kind; source}{method; \GSpage{P000198}}
\GSfield{The notes}{The proof is said to use Ahlfors and Bers on quasiconformal maps, ``onto which any diffeomorphism contracts.}
\GSfield{Statement}{The rigorous strategy uses the contractible space of smooth conformal structures, the contractible Teichm\"uller space, and a principal action of the appropriate identity or marked-point diffeomorphism group. The resulting fibration and the manifold properties of the group imply contractibility of the relevant fiber. Quasiconformal Beltrami equations provide the required parametrization and local sections.}
\GSfield{Reconstruction}{The notes do not describe a literal deformation of every diffeomorphism ``onto quasiconformal maps''; smooth surface diffeomorphisms already are quasiconformal for suitable bounds. The fibration, rather than that phrase, is the substantive argument.}
\GSfield{Fate}{The exact proof outline is \citet[Theorem~2; \S\S3--4]{earleeells1967diffeomorphism}.}
\GSfield{Present status}{As of 9 November 2026, this remains an established construction under the stated hypotheses.}
\GSfield{Assessment}{literal statement: correct with missing hypotheses; publication match: exact; verification: read the result; development: none; confidence: high}
\end{GSitem}
\begin{GSitem}{1966-124.5}{E}{The diffeomorphism group as an infinite-dimensional Lie group}
\GSfield{Kind; source}{remark; \GSpage{P000198}}
\GSfield{The notes}{The final sentence calls \(\operatorname{Diff}\) a Lie group in a new sense of smooth structure.}
\GSfield{Statement}{For a compact smooth manifold, \(\operatorname{Diff}(X)\) admits an infinite-dimensional smooth Lie-group structure modeled locally on the Fr\'echet space of smooth vector fields. This requires a suitable infinite-dimensional differentiability calculus; it is not a finite-dimensional Lie group.}
\GSfield{Reconstruction}{The notebook specifies no calculus or regularity theorem. The contemporary reference is Leslie, as cited by Earle--Eells; Leslie's full article was not accessible in this run.}
\GSfield{Fate}{\citet[Remark~2 and reference~8]{earleeells1967diffeomorphism} cites Leslie's then-forthcoming 1967 construction.}
\GSfield{Present status}{As of 9 November 2026, this remains an established construction under the stated hypotheses.}
\GSfield{Assessment}{literal statement: correct with missing hypotheses; publication match: partial; verification: indirect; development: none; confidence: high}
\end{GSitem}
\begin{GSitem}{1966-124.6}{E}{Semiclassical Fourier transformation of the Schrödinger operator}
\GSfield{Kind; source}{construction; \GSpage{P000199}; D. V. Anosov}
\GSfield{The notes}{Under the Anosov heading, \(-\tfrac{h^2}{2}\partial_x^2+U(x)\) transforms to \(\tfrac{p^2}{2}+U(\tfrac h i\partial_p)\); the Fourier kernel is \(e^{\pm i\langle x,p\rangle/h}\).}
\GSfield{Statement}{For \(h>0\), use the unitary transform \((\mathcal F_h^+f)(p)=(2\pi h)^{-1/2}\int e^{ixp/h}f(x)\,dx\). On Schwartz functions and for polynomial \(U\), \[\mathcal F_h^+\left(-\frac{h^2}{2}\partial_x^2+U(x)\right)(\mathcal F_h^+)^{-1}=\frac{p^2}{2}+U(-ih\partial_p).\] For a more general potential the notation on the right denotes the corresponding Fourier multiplier/convolution or functional-calculus operator on its appropriate domain.}
\GSfield{Reconstruction}{The positive Fourier sign matches the notebook's \(h/i\). The polynomial/core case supplies a precise statement without assuming an arbitrary potential is a differential operator.}
\GSfield{Fate}{An elementary preexisting Fourier identity, used in Maslov's 1965 programme; only the book's metadata were checked \citep{maslov1965perturbation}.}
\GSfield{Present status}{As of 9 November 2026, this remains an established construction under the stated hypotheses.}
\GSfield{Assessment}{literal statement: correct with missing hypotheses; publication match: thematic only; verification: publisher metadata checked; development: none; confidence: high}
\end{GSitem}
\begin{GSitem}{1966-124.7}{E}{A Fourier transform with a self-adjoint operator as reciprocal Planck constant}
\GSfield{Kind; source}{construction; \GSpage{P000199}; D. V. Anosov}
\GSfield{The notes}{For Hilbert space \(H\), self-adjoint \(A\) with \(0\) not an eigenvalue, the notes replace \(1/h\) by \(A\) and insert the normalization \(|A|^{n/2}/(\sqrt{2\pi})^n\) in \(\phi^A f(p)=\int e^{iA\langle x,p\rangle}f(x)\,dx\).}
\GSfield{Statement}{A precise interpretation uses the spectral decomposition \(A=\int\lambda\,dE(\lambda)\), with \(E(\{0\})=0\). On each spectral fiber define \[\mathcal F_\lambda^+f(p)=\left(\frac{|\lambda|}{2\pi}\right)^{n/2}\int_{\mathbb R^n}e^{i\lambda\langle x,p\rangle}f(x)\,dx.\] The direct integral of these unitary transforms defines a unitary \(\mathcal F_A^+\) on \(L^2(\mathbb R^n;H)\). The notebook's integral formula is literal on a suitable spectral/Schwartz core and otherwise is understood through this extension.}
\GSfield{Reconstruction}{\ours\ The direct-integral interpretation is supplied here. It allows \(0\) in continuous spectrum: ``not an eigenvalue'' does not imply \(A^{-1}\) is bounded. Fiberwise Plancherel verifies the normalization; substitution \(q=\lambda p\) independently gives the same norm identity. No exact occurrence of this operator-valued formula was located in the checked Maslov metadata.}
\GSfield{Fate}{Maslov's 1965 book is the stated historical context, but the precise transform-to-publication match is unresolved \citep{maslov1965perturbation}.}
\GSfield{Present status}{As of 9 November 2026, the spectral construction above is well defined. The attribution is only thematic; no new general quantization theorem is inferred.}
\GSfield{Assessment}{literal statement: correct with missing hypotheses; publication match: thematic only; verification: indirect; development: annotated; confidence: medium}
\end{GSitem}
\begin{GSitem}{1966-124.8}{F}{Formal operator-valued quantization of L(x,p,A inverse)}
\GSfield{Kind; source}{method; \GSpage{P000199}; D. V. Anosov}
\GSfield{The notes}{The note ends with \(\mathcal L(x,\widehat p,A^{-1})f=\overline{\phi^A}_{p\to x}L(x,p,A^{-1})\phi^A_{x\to p}f\), after posing \(Lu=f\) with \(u(x)\in H\). No symbol estimates, domain, ordering convention, or solvability theorem is supplied.}
\GSfield{Statement}{As a formal quantization rule the expression indicates applying an \(A\)-scaled Fourier transform, a symbol, and inverse transformation with \(x\) retained as an external parameter. It does not define a unique operator for arbitrary noncommuting operator-valued \(L\). A rigorous restricted construction is possible when \(L\) is decomposable with respect to \(A\), its scalar spectral fibers satisfy stated pseudodifferential symbol bounds, and common domains and an ordering convention are fixed.}
\GSfield{Reconstruction}{The decomposability and symbol assumptions describe one possible implementation, not a recovered claim by Anosov. The unresolved issue here is the missing statement, not a demonstrated open research problem.}
\GSfield{Fate}{The Maslov book provides a thematic lead only; neither its catalogue nor the notebook identifies this general theorem \citep{maslov1965perturbation}.}
\GSfield{Present status}{As of 9 November 2026, the unrestricted construction and the intended solvability result remain unidentified from these notes.}
\GSfield{Assessment}{literal statement: not determinable; publication match: thematic only; verification: unresolved lead; development: none; confidence: low}
\end{GSitem}

\GStalk{125}{1967-03-13}{A. N. Parshin}{P000200}{Algebraic group actions and geometric orbit quotients}
Parshin's report on Mumford defines a regular algebraic action and asks when its orbit set carries an algebraic quotient structure. It records stabilizer semicontinuity, necessary restrictions on a morphism whose fibers are individual orbits, and the affine reductive closed-orbit criterion in compressed form. An unlabelled hyperbola-like sketch cannot be uniquely identified. The later projective continuation, record 128, preserves only a heading.

\begin{GSitem}{1966-125.1}{E}{Regular algebraic actions and orbits}
\GSfield{Kind; source}{definition; \GSpage{P000200}}
\GSfield{The notes}{A regular map \(\sigma:G\times X\to X\) is illustrated by an associativity square, whose top arrow is labelled \(\sigma\times1\); equivalent points form an orbit.}
\GSfield{Statement}{Over an algebraically closed field \(k\), an algebraic action of an algebraic group \(G\) on a variety \(X\) is a morphism \(\sigma:G\times X\to X\) satisfying \(\sigma(e,x)=x\) and \(\sigma(g,\sigma(h,x))=\sigma(gh,x)\). In the standard ordered product \(G\times G\times X\), the top map of the associativity square is \(1_G\times\sigma\).}
\GSfield{Reconstruction}{The identity axiom is supplied. The stated top-arrow label has the wrong factor order unless a permutation of factors is additionally declared.}
\GSfield{Fate}{Standard setup in the 1965 book \citep{mumford1965invariant}; modern definitions are in \citet[\S2.1]{berczi2016quotients}.}
\GSfield{Present status}{As of 9 November 2026, this remains an established construction under the stated hypotheses.}
\GSfield{Assessment}{literal statement: correct with missing hypotheses; publication match: exact; verification: read the result; development: none; confidence: high}
\end{GSitem}
\begin{GSitem}{1966-125.2}{E}{When does an orbit set admit an algebraic quotient?}
\GSfield{Kind; source}{question; \GSpage{P000200}}
\GSfield{The notes}{The explicit question is when \(X/G\), meaning the set of orbits, is an algebraic variety and in what sense the map \(X\to X/G\) is regular.}
\GSfield{Statement}{For an affine variety \(X\) over an algebraically closed field of characteristic zero and a reductive algebraic group \(G\), the invariant ring is finitely generated and \(q:X\to\operatorname{Spec}k[X]^G\) is a categorical quotient. Its fibers identify orbit-closure equivalence classes and contain one closed orbit; it is a geometric orbit quotient of all \(X\) precisely when all orbits are closed. For an arbitrary algebraic action, a geometric quotient exists on a suitable nonempty invariant open subset; it need not exist on all of \(X\).}
\GSfield{Reconstruction}{Affine reductive hypotheses make the notes' final criterion precise. The distinction between a categorical quotient and a quotient whose fibers are individual orbits is essential. These are answers to definite versions of the stated question, not an assertion of a single universal global quotient theorem for all varieties.}
\GSfield{Fate}{Known before the report from Mumford/Rosenlicht; compare \citet[Theorems~2.1.10, 2.2.1]{berczi2016quotients}.}
\GSfield{Present status}{As of 9 November 2026, the stated affine reductive question and the generic-open existence theorem are established. The unrestricted wording needs a chosen quotient category.}
\GSfield{Assessment}{literal statement: correct with missing hypotheses; publication match: partial; verification: read the result; development: none; confidence: high; problem \hyperref[prob:GP-1966-19]{GP-1966-19}}
\end{GSitem}
\begin{GSitem}{1966-125.3}{F}{Unlabelled orbit-space counterexample sketch}
\GSfield{Kind; source}{remark; \GSpage{P000200}}
\GSfield{The notes}{A drawing of axes and curved branches accompanies the statement that no algebraic structure can be introduced until a point is removed. Neither the group, action, variety, nor removed point is labelled.}
\GSfield{Statement}{There is insufficient information to state a unique example or a counterexample. The picture is consistent with several torus actions having nonclosed orbits, but neither the orbit partition nor the intended separation condition can be recovered.}
\GSfield{Reconstruction}{The familiar action \(t\cdot(x,y)=(tx,t^{-1}y)\) is not installed as the missing example: even there, deleting the origin has subtleties involving nonseparated orbit spaces.}
\GSfield{Fate}{The general quotient literature \citep[\S2.1]{berczi2016quotients} does not identify an unlabelled diagram.}
\GSfield{Present status}{As of 9 November 2026, this particular example remains unidentified; it does not establish an orphan open problem.}
\GSfield{Assessment}{literal statement: not determinable; publication match: none found; verification: unresolved lead; development: none; confidence: low}
\end{GSitem}
\begin{GSitem}{1966-125.4}{E}{Semicontinuity of stabilizer and orbit dimensions}
\GSfield{Kind; source}{theorem; \GSpage{P000200}}
\GSfield{The notes}{The stationary subgroup may have larger dimension than the generic stationary subgroup.}
\GSfield{Statement}{For an algebraic group acting on a variety, the function \(x\mapsto\dim G_x\) is upper semicontinuous, and \(\dim(Gx)=\dim G-\dim G_x\) is lower semicontinuous. On an irreducible variety the generic stabilizer dimension is minimal and the generic orbit dimension maximal.}
\GSfield{Reconstruction}{Dimension is meant scheme-theoretically or geometrically consistently; the characteristic-zero setting in the notes avoids nonreduced stabilizer complications.}
\GSfield{Fate}{Classical action geometry; see \citet[\S2.2, discussion preceding Theorem~2.2.1]{berczi2016quotients}.}
\GSfield{Present status}{As of 9 November 2026, this remains an established construction under the stated hypotheses.}
\GSfield{Assessment}{literal statement: correct with missing hypotheses; publication match: exact; verification: read the result; development: none; confidence: high}
\end{GSitem}
\begin{GSitem}{1966-125.5}{E}{Necessary conditions on a regular morphism with orbit fibers}
\GSfield{Kind; source}{theorem; \GSpage{P000200}}
\GSfield{The notes}{If \(X/G\) has an algebraic structure and \(X\to X/G\) is regular, all orbits must have the same dimension and, in particular, be closed.}
\GSfield{Statement}{Let \(X,Y\) be irreducible varieties over an algebraically closed field, and let \(q:X\to Y\) be a surjective morphism whose set-theoretic fibers over closed points are exactly the \(G\)-orbits. Then all such orbits are closed and have dimension \(\dim X-\dim Y\). Indeed fiber dimension is at least this number, while orbit dimension is at most its generic value, which equals it.}
\GSfield{Reconstruction}{\ours\ Irreducibility, surjectivity, and the single-orbit fiber condition make the notebook claim valid. Closedness follows directly from being a fiber, not merely from equal orbit dimensions. Independently, the generic fiber-dimension theorem and stabilizer semicontinuity give the two opposite inequalities, checking the argument.}
\GSfield{Fate}{A standard consequence of algebraic orbit and fiber dimensions in the quotient framework \citep[\S2.1]{berczi2016quotients}.}
\GSfield{Present status}{As of 9 November 2026, this remains an established construction under the stated hypotheses.}
\GSfield{Assessment}{literal statement: correct with missing hypotheses; publication match: partial; verification: read the result; development: annotated; confidence: high}
\end{GSitem}
\begin{GSitem}{1966-125.6}{E}{Closed orbits and the affine reductive geometric quotient}
\GSfield{Kind; source}{theorem; \GSpage{P000200}}
\GSfield{The notes}{For a reductive group in characteristic zero, the notes end with ``a quotient exists if and only if all orbits are closed.}
\GSfield{Statement}{If \(G\) is reductive, \(\operatorname{char}k=0\), and \(X\) is affine, the invariant quotient \(q:X\to X/\!/G\) exists whether or not every orbit is closed. It is geometric, with precisely one orbit in each fiber, if and only if every orbit is closed. Closedness alone is not the stated existence criterion for an arbitrary nonaffine variety or for an unspecified quotient category.}
\GSfield{Reconstruction}{The missing words are ``affine'' and ``geometric.'' The literal statement is false when ``quotient'' means categorical quotient. This distinguishes a valid old theorem in shorthand from a new existence problem.}
\GSfield{Fate}{The precise affine theorem is \citet[Theorem~2.2.1]{berczi2016quotients}, reflecting the prelecture Mumford theory \citep{mumford1965invariant}.}
\GSfield{Present status}{As of 9 November 2026, this remains an established construction under the stated hypotheses.}
\GSfield{Assessment}{literal statement: correct with missing hypotheses; publication match: exact; verification: read the result; development: none; confidence: high}
\end{GSitem}

\GStalk{126}{1967-03-13}{A. S. Dynin}{P000201}{Stable symbols and an axiomatic index}
Dynin's second lecture defines compactly supported stable symbols and lists homotopy invariance, additivity, products, excision, vanishing for symbols from the base, and the one-dimensional winding normalization. The first isolated uniqueness claim mentions only the first two properties and is false literally; the later full family of axioms supplies the intended setting. This is related to the earlier index lectures, but no already-indexed open question is reasserted.

\begin{GSitem}{1966-126.1}{E}{Compactly supported stable symbol classes}
\GSfield{Kind; source}{definition; \GSpage{P000201}}
\GSfield{The notes}{\(\Sigma(\mathbb R^n)\) consists of maps \(\mathbb R^n\times S^{n-1}\to GL_\infty(\mathbb C)\), equal to \(1\) outside \(K\times S^{n-1}\). The general version uses \(T_0(X)\) and a compact support in the base.}
\GSfield{Statement}{A stable invertible principal symbol on \(X\), trivial outside a compact subset of the base, is represented by a map \(S^*X\to GL_N(\mathbb C)\), stabilized by identity blocks and considered through homotopy and the relation of extending from the base. Clutching over the cotangent disk bundle yields its compactly supported \(K\)-theory symbol class in \(K_c^0(T^*X)\). A metric identifies tangent and cotangent sphere bundles when necessary.}
\GSfield{Reconstruction}{The raw maps are not themselves already all of \(K_c^0(T^*X)\); the relations and, for general bundles, bundle-valued symbols are needed. \(T_0\) is interpreted as the punctured cotangent bundle with degree-zero homogeneity, or equivalently its cosphere.}
\GSfield{Fate}{Stable symbol and \(K\)-theory framework: \citet[\S2; \S3.1]{dynin1966index}.}
\GSfield{Present status}{As of 9 November 2026, this remains an established construction under the stated hypotheses.}
\GSfield{Assessment}{literal statement: correct with missing hypotheses; publication match: partial; verification: read the result; development: none; confidence: high}
\end{GSitem}
\begin{GSitem}{1966-126.2}{E}{Homotopy invariance and direct-sum additivity do not determine the index}
\GSfield{Kind; source}{claim; \GSpage{P000201}}
\GSfield{The notes}{Above a dividing line the notes claim a unique \(\operatorname{ind}:\Sigma(\mathbb R^n)\to\mathbb Z\) with only (I) homotopy invariance and (II) direct-sum additivity.}
\GSfield{Statement}{These two axioms alone cannot normalize a nonzero index: the zero map and every integer multiple of any additive homotopy invariant satisfy them. Already for \(n=1\), \[w(\sigma)=\operatorname{wind}\det\sigma(\cdot,+1)-\operatorname{wind}\det\sigma(\cdot,-1)\] is a nonzero such invariant for compactly supported stable symbols; \(0,w,2w\) are distinct.}
\GSfield{Reconstruction}{\ours\ The counterexample uses a smooth phase that rises from \(0\) to \(2\pi\) in a compact interval on the positive component and is constant on the negative component. Determinant multiplication under direct sum proves additivity; homotopy invariance follows from integer winding. The later normalization axiom fixes this defect.}
\GSfield{Fate}{Dynin's actual index properties include a normalization: \citet[\S3.1; \S3.2.10]{dynin1966index}.}
\GSfield{Present status}{As of 9 November 2026, the isolated uniqueness statement is false; the normalized index theory remains established.}
\GSfield{Assessment}{literal statement: false as written; publication match: partial; verification: read the result; development: annotated; confidence: high}
\end{GSitem}
\begin{GSitem}{1966-126.3}{E}{Multiplicativity under external products}
\GSfield{Kind; source}{theorem; \GSpage{P000201}}
\GSfield{The notes}{A product \(\#:\Sigma(X)\times\Sigma(Y)\to\Sigma(X\times Y)\) satisfies \(\operatorname{ind}(\sigma_X\#\sigma_Y)=\operatorname{ind}\sigma_X\operatorname{ind}\sigma_Y\).}
\GSfield{Statement}{For compactly supported elliptic \(K\)-theory symbols, their external product corresponds to the graded tensor product of elliptic complexes. The analytic index is multiplicative on this product, as is the topological index. The symbol product is not merely pointwise multiplication of two unrelated square matrices on different cosphere bundles.}
\GSfield{Reconstruction}{The external-product construction is supplied; the notes state only its name and index identity. Stable triviality of tangent bundles is part of the lecture's restricted family, not a necessary hypothesis of the modern product theorem.}
\GSfield{Fate}{\citet[\S3.1.6]{dynin1966index} gives the prelecture product property.}
\GSfield{Present status}{As of 9 November 2026, this remains an established construction under the stated hypotheses.}
\GSfield{Assessment}{literal statement: correct with missing hypotheses; publication match: exact; verification: read the result; development: none; confidence: high}
\end{GSitem}
\begin{GSitem}{1966-126.4}{E}{Excision for compactly supported index classes}
\GSfield{Kind; source}{theorem; \GSpage{P000201}}
\GSfield{The notes}{If \(\operatorname{supp}\sigma\subset U\subset X\), the index on \(U\) equals the index on \(X\).}
\GSfield{Statement}{For an open embedding \(j:U\hookrightarrow X\) and a compactly supported elliptic symbol class \(a\in K_c^0(T^*U)\), extension by the identity/zero symbol class defines \(j_!a\), and \(\operatorname{ind}_X(j_!a)=\operatorname{ind}_U(a)\). The relevant support is the nontrivial part of the symbol class, not the support of an everywhere invertible matrix as a function.}
\GSfield{Reconstruction}{This is an index statement for symbols trivial at infinity; an arbitrary differential operator on a noncompact open manifold need not be Fredholm without this interpretation.}
\GSfield{Fate}{A standard topological-index property in the framework of \citet[\S\S2--3]{dynin1966index}; exact six-axiom wording was not found there.}
\GSfield{Present status}{As of 9 November 2026, this remains an established construction under the stated hypotheses.}
\GSfield{Assessment}{literal statement: correct with missing hypotheses; publication match: partial; verification: read the result; development: none; confidence: high}
\end{GSitem}
\begin{GSitem}{1966-126.5}{E}{Symbols pulled back from the base have zero index}
\GSfield{Kind; source}{theorem; \GSpage{P000201}}
\GSfield{The notes}{A commuting diagram \(T_0(X)\to X\to GL\) implies \(\operatorname{ind}\sigma=0\).}
\GSfield{Statement}{An invertible symbol that extends over the cotangent disk bundle, in particular one pulled back from a bundle automorphism on \(X\), has zero relative \(K\)-theory class and zero index. Analytically, multiplication by the corresponding invertible bundle map is an invertible order-zero realization.}
\GSfield{Reconstruction}{The compact-support/trivial-at-infinity convention from the definition is retained. This explains why raw homotopy classes on the cosphere require a quotient by extendible symbols.}
\GSfield{Fate}{An elementary consequence of the symbol construction in \citet[\S2]{dynin1966index}.}
\GSfield{Present status}{As of 9 November 2026, this remains an established construction under the stated hypotheses.}
\GSfield{Assessment}{literal statement: correct; publication match: partial; verification: read the result; development: none; confidence: high}
\end{GSitem}
\begin{GSitem}{1966-126.6}{E}{One-dimensional winding normalization}
\GSfield{Kind; source}{theorem; \GSpage{P000201}}
\GSfield{The notes}{For \(X=\mathbb R\), the Noether--Muskhelishvili formula is \[\operatorname{ind}\sigma=(2\pi)^{-1}\Delta_{\mathbb R\times S^0}\arg\det\sigma.\]}
\GSfield{Statement}{For a smooth stable invertible symbol equal to the identity off a compact subset of \(\mathbb R\), each of its two cosphere components closes to a loop. With the boundary orientation and Fourier/index sign convention chosen consistently, the index is the oriented difference of their determinant winding numbers. Reversing the Fourier or boundary orientation convention reverses the displayed sign.}
\GSfield{Reconstruction}{The unoriented notation \(\mathbb R\times S^0\) is insufficient to determine the sign; no independent sign is silently imposed on the notebook.}
\GSfield{Fate}{The normalized dimension-one calculation is \citet[\S3.2.10]{dynin1966index}.}
\GSfield{Present status}{As of 9 November 2026, this remains an established construction under the stated hypotheses.}
\GSfield{Assessment}{literal statement: correct with missing hypotheses; publication match: exact; verification: read the result; development: none; confidence: high}
\end{GSitem}
\begin{GSitem}{1966-126.7}{E}{A normalized index characterization on stably parallelizable manifolds}
\GSfield{Kind; source}{claim; \GSpage{P000201}}
\GSfield{The notes}{Below the line, the notes assert existence and uniqueness of an index for manifolds with stably trivial tangent bundle satisfying (I)--(VI).}
\GSfield{Statement}{The intended object is the normalized index homomorphism on compactly supported stable elliptic symbol classes, natural under excision and external products and normalized on the one-dimensional Bott generator. In a precise \(K\)-theoretic axiomatization, Bott periodicity and compatible pushforward determine the usual topological index, equal to the analytic index. The six shorthand lines require a specified symbol category, product construction, and stabilization before a literal uniqueness theorem can be tested.}
\GSfield{Reconstruction}{The modern normalized index statement is supplied as an interpretation. No exact proof of the notebook's six-axiom formulation, restricted to stably parallelizable manifolds, was located in the opened Dynin exposition; it is therefore not reported as a verbatim published theorem.}
\GSfield{Fate}{The 1966 exposition proves the index formula and required properties \citep[\S3.2.1]{dynin1966index}; the exact axiom list matches only partially.}
\GSfield{Present status}{As of 9 November 2026, the normalized analytic/topological index is established. The bare six-line formulation is under-specified, not evidence of an open uniqueness problem.}
\GSfield{Assessment}{literal statement: correct with missing hypotheses; publication match: partial; verification: read the result; development: none; confidence: medium}
\end{GSitem}

\GStalk{127}{1967-03-27}{A. S. Dynin}{P000202-P000204}{Boundary systems with varying partial indices and global completion}
The March 27 material begins below the separate March 20 Parshin heading on P000202. The catalogue attributes this record to Dynin, but the manuscript names Eskin on P000203; this could introduce Eskin's work or a further speaker, and the catalogue attribution is retained with that qualification. The lecture treats a continuous family of normal half-line problems even when individual partial indices jump, auxiliary potentials and boundary equations, localization, and a global completion question. The latter matches the framework and criterion of Vishik--Eskin's paper received 27 April 1967, rather than merely a later paper on the same subject.

\begin{GSitem}{1966-127.1}{E}{Normal matrix factorization and partial indices}
\GSfield{Kind; source}{construction; \GSpage{P000202-P000203}}
\GSfield{The notes}{The factorization is \(A(\omega,\eta_n)=A_-M_-M_+A_+\), with \(M_+=\operatorname{diag}((\eta_n+i)^{n_j(\omega)})\). The integers \(n_j\) may depend on \(\omega\). In the full variables \(A_\pm\) have degree zero and the order factor \(\xi_-^\alpha\) is included in \(M_-\).}
\GSfield{Statement}{For an elliptic matrix symbol, homogeneous of order \(\alpha+i\beta\) and with the regularity and normal-limit assumptions of the convolution calculus, fix a nonzero tangential direction \(\omega\). Remove the scalar order and normal-jump factors; the normalized matrix admits Wiener--Hopf factorization with integer partial indices \(n_j(\omega)\). The unreduced partial indices are \(\kappa_j=n_j+\gamma_j\), where \(\gamma_j\) arise from a chosen logarithm of the normal jump. Such pointwise factorizations need not vary continuously when the integer indices change.}
\GSfield{Reconstruction}{The normalization, half-plane analyticity of the factors, and distinction between integer and full partial indices are supplied. No continuously chosen factorization is inferred from pointwise factorization.}
\GSfield{Fate}{The pointwise factorization was already available: \citet[\S4]{vishikeskin1967elliptic}, received December 1965; the later notation is exact in \citet[\S1.3]{vishikeskin1967normally}.}
\GSfield{Present status}{As of 9 November 2026, this remains an established construction under the stated hypotheses.}
\GSfield{Assessment}{literal statement: correct with missing hypotheses; publication match: exact; verification: read the result; development: none; confidence: high}
\end{GSitem}
\begin{GSitem}{1966-127.2}{B}{Fredholm normal families despite jumping kernel dimensions}
\GSfield{Kind; source}{theorem; \GSpage{P000202}}
\GSfield{The notes}{The notes emphasize \(n_j(\omega)\) need not be constant and nevertheless assert a regularizer construction; the dimensions are denoted \(m_+(\omega)\) and \(m_-(\omega)\).}
\GSfield{Statement}{Let \(A_\omega\) be the half-line convolution family obtained by freezing the boundary point and unit tangential covector. Choose Sobolev orders avoiding the critical values \(\operatorname{Re}\kappa_j+\tfrac12+\mathbb Z\), continuously over the parameter set, and identify the resulting spaces with fixed Hilbert spaces. Under the stated symbol regularity, \(A_\omega\) is a norm-continuous Fredholm family. Its kernel and cokernel dimensions \(m_+\) and \(m_-\) can jump, but their difference is constant on each connected component and both dimensions are bounded on a compact parameter space.}
\GSfield{Reconstruction}{Uniform admissibility of the Sobolev order is essential; arbitrary elliptic symbols and arbitrary orders are not covered. Continuity concerns the Fredholm operators, not a chosen kernel basis or factorization.}
\GSfield{Fate}{The corresponding result is \citet[\S1, Lemmas~1.1--3.1 and Remark~1.1]{vishikeskin1967normally}, received 27 April 1967 and published that year.}
\GSfield{Present status}{As of 9 November 2026, the qualified family theorem is established. Its verified 1967 proof expressly removes the earlier constant-integer-index restriction.}
\GSfield{Assessment}{literal statement: correct with missing hypotheses; publication match: exact; verification: read the result; development: none; confidence: high}
\end{GSitem}
\begin{GSitem}{1966-127.3}{B}{Auxiliary potentials and boundary equations}
\GSfield{Kind; source}{construction; \GSpage{P000202}}
\GSfield{The notes}{The one-dimensional system is \(Au_++G\rho=\ell f+u_-\), \((B,u_+)_{\eta_n}+E\rho=g\), with \(\rho\in\mathbb C^M\), \(\ell-M=m_+-m_-\), unique solvability, and a continuously dependent inverse \(R_\omega\). The same glyph \(\ell\) is used for extension of \(f\) and for the number of boundary equations.}
\GSfield{Statement}{For a continuous Fredholm family \(A_\omega:H_1\to H_2\), seek a block completion \[\mathcal A_\omega=\begin{pmatrix}A_\omega&G_\omega\\B_\omega&E_\omega\end{pmatrix}:H_1\oplus\mathbb C^M\longrightarrow H_2\oplus\mathbb C^\ell\] that is invertible. A necessary index balance is \(\ell-M=\operatorname{ind}A_\omega\). Such completions exist locally in the parameter; once the block is norm-continuous and invertible, \(R_\omega=\mathcal A_\omega^{-1}\) is norm-continuous. In the convolution problem, \(G\) is a potential matrix and \(B\) a boundary functional with prescribed homogeneity and Sobolev mapping properties. The index balance alone does not ensure a single global completion.}
\GSfield{Reconstruction}{The undefined boundary-restriction subscript on \(g\) is not used to recover a new space. Abstract local completion and continuity of inversion explain the claim; global existence is separated as the later explicit question.}
\GSfield{Fate}{The same augmented system and algebraic solvability condition appear in \citet[\S1.5, equations (44.1)--(45.1)]{vishikeskin1967normally}.}
\GSfield{Present status}{As of 9 November 2026, this completion construction and its parameter continuity are standard under the stated assumptions; the global obstruction is recorded separately.}
\GSfield{Assessment}{literal statement: correct with missing hypotheses; publication match: exact; verification: read the result; development: none; confidence: high}
\end{GSitem}
\begin{GSitem}{1966-127.4}{B}{Homogeneous lifting and localization of normal inverses}
\GSfield{Kind; source}{method; \GSpage{P000203}}
\GSfield{The notes}{The normal construction is transferred to \(n\) dimensions by factoring out \(|\xi'|\); left and right regularizers are obtained by localization. The margin adds \ru{Якобы} beside the reference to Simonenko's abstract treatment of inverse operators depending on parameters.}
\GSfield{Statement}{Suppose an elliptic convolution system on a smooth bounded domain has an admissible continuous normal family and globally specified potential/boundary symbols of the required orders, whose augmented normal problem is invertible at every boundary covector. Rescaling the normal inverse by tangential homogeneity and patching local parametrices yields left and right regularizers modulo compact operators on the designated Sobolev spaces. Hence the global boundary system is Fredholm. A formal collection of pointwise inverses without symbol estimates, admissible orders, and compatible global completion is insufficient.}
\GSfield{Reconstruction}{The hypotheses are those of the matching theorem, not an assertion that arbitrary parameter dependence can be localized. The marginal reservation is retained as evidence, not interpreted as a refutation of Simonenko.}
\GSfield{Fate}{The precise half-space and bounded-domain results are in \citet[pp.~316, 322--325]{vishikeskin1967normally}. Simonenko's 1965 paper is a metadata-verified lead only \citep{simonenko1965method}.}
\GSfield{Present status}{As of 9 November 2026, the qualified Fredholm construction is established. The exact Simonenko passage meant by the speaker remains unverified.}
\GSfield{Assessment}{literal statement: correct with missing hypotheses; publication match: exact; verification: read the result; development: none; confidence: high}
\end{GSitem}
\begin{GSitem}{1966-127.5}{B}{Normal jump logarithms and admissible Sobolev orders}
\GSfield{Kind; source}{definition; \GSpage{P000203}}
\GSfield{The notes}{Under the heading \ru{Эскин}, the notes set \(a_\pm(x')=A(x',0,0,\pm1)\), \(b=a_+^{-1}a_-\), and \(\gamma=(2\pi i)^{-1}\log b\). They demand globally continuous eigenvalue branches with \(|\operatorname{Re}\gamma_i-\operatorname{Re}\gamma_j|<1\), calling this \(\widehat{\mathcal E}_\alpha\). The jump is called a measure of nonellipticity.}
\GSfield{Statement}{For a zero-order normalized elliptic symbol at a smooth boundary, the two invertible normal limits define a normal jump \(b=a_+^{-1}a_-\). Locally one can choose a matrix logarithm and spectral branches in a strip of width less than one. Global branch choices require compatibility. For nonzero symbol order, first normalize the order; in Vishik--Eskin's convention the negative limit includes the phase \(e^{i\pi(\alpha+i\beta)}\). A continuous Sobolev-order function avoiding all critical fractional indices is the functional-analytic admissibility condition.}
\GSfield{Reconstruction}{The jump measures normal factorization data, not a failure of interior ellipticity. The notebook's symbol for the class is uncertain, and its global enumeration of eigenvalue branches is stronger than a merely pointwise logarithm. The later paper uses a precisely defined class and a noncritical order condition rather than relying on this compressed name.}
\GSfield{Fate}{The normalized limits and admissibility condition are explicit in \citet[\S2, equations (1.2)--(6.2)]{vishikeskin1967normally}; local logarithmic factorization predates the lecture \citep[\S4]{vishikeskin1967elliptic}.}
\GSfield{Present status}{As of 9 November 2026, the normalized construction is established; the notebook's class label has only a partial exact identification.}
\GSfield{Assessment}{literal statement: correct with missing hypotheses; publication match: partial; verification: read the result; development: none; confidence: medium}
\end{GSitem}
\begin{GSitem}{1966-127.6}{E}{Diagonal systems and the limitation of a common ordinary space}
\GSfield{Kind; source}{remark; \GSpage{P000203-P000204}}
\GSfield{The notes}{For diagonal \(A=\operatorname{diag}(a_j)\), the notes write scalar factorization indices \(\kappa_j(x')\), \(b=e^{2\pi i\kappa_j}\), \(\gamma_j=\kappa_j\), and say the class need not contain the system. The qualification is ``operators in ordinary spaces.}
\GSfield{Statement}{For a diagonal system the full scalar factorization indices need not lie in a common real strip of width one when one insists on a fixed compatible choice and common Sobolev scale. The eigenvalues of the normal jump determine \(\kappa_j\) only modulo integers; in general \(\gamma_j\equiv\kappa_j\pmod{\mathbb Z}\), not equality for every logarithm branch. Different component orders or variable admissible orders can alter the functional setting, so this observation is a limitation of a specified common-space class, not of elliptic diagonal systems as such.}
\GSfield{Reconstruction}{The scalar index glyph is read provisionally as \(\kappa\). The equation \(\gamma_j=\kappa_j\) is meaningful only after choosing compatible branches. Pointwise branch adjustment can always narrow a finite set of fractional parts; the possible difficulty is global compatibility with the intended spaces.}
\GSfield{Fate}{The distinction between \(n_j\), \(\gamma_j\), and \(\kappa_j\) is explicit in \citet[\S1.3, p.~306; \S2.1]{vishikeskin1967normally}; the notebook's example lacks full symbols.}
\GSfield{Present status}{As of 9 November 2026, this is an understood qualification of the formulation. No particular diagonal counterexample can be uniquely recovered.}
\GSfield{Assessment}{literal statement: correct with missing hypotheses; publication match: partial; verification: read the result; development: none; confidence: medium}
\end{GSitem}
\begin{GSitem}{1966-127.7}{B}{Global completion of the boundary Fredholm family}
\GSfield{Kind; source}{question; \GSpage{P000204}}
\GSfield{The notes}{The final line asks when there globally exist \(G\) and an overwritten expression involving \(E\), ``(and \(B\)).'' The connective between \(G\) and \(E\) is not legible enough to decide among the transcription's alternatives. The preceding normal system fixes the mathematical context.}
\GSfield{Statement}{In the admissible Vishik--Eskin convolution calculus on a smooth bounded domain with compact boundary \(\Gamma\), let \(A_\gamma\), \(\gamma\in S^*\Gamma\), be the normalized continuous Fredholm family, of constant numerical index \(N\). When do globally defined finite matrices of potential and boundary symbols give an invertible normal completion using trivial auxiliary bundles? The precise criterion is \[\operatorname{Ind}(A) = N[1]\quad\hbox{in }K^0(S^*\Gamma).\] Under the paper's regularity and noncritical-order hypotheses this is necessary and sufficient, and the completion can be chosen with \(E=0\). For a connected parameter space, this says that the reduced family-index class vanishes.}
\GSfield{Reconstruction}{This identifies the definite completion question from the system on P000202; it does not resolve the overwritten connective. Trivial auxiliary bundles, admissible symbol orders, and constant numerical index are essential. Allowing arbitrary auxiliary bundles changes the obstruction. Necessity for a block with \(E\ne0\) also follows from its family index; sufficiency follows from the stronger \(E=0\) construction.}
\GSfield{Fate}{A direct contextual match is \citet[\S3, Theorems~1.3 and 2.3, pp.~328--331]{vishikeskin1967normally}, received 27 April 1967, after this lecture.}
\GSfield{Present status}{As of 9 November 2026, this precise global completion question was solved in 1967. The printed theorem matches the augmented system, jumping indices, global symbols, and obstruction; the source's ambiguous connective prevents a claim of verbatim equivalence.}
\GSfield{Assessment}{literal statement: correct with missing hypotheses; publication match: partial; verification: read the result; development: none; confidence: high; problem \hyperref[prob:GP-1966-20]{GP-1966-20}}
\end{GSitem}

\GStalk{128}{1967-03-20}{A. N. Parshin}{P000202}{Projective varieties: heading of the continued Mumford report}
Only the handwritten date 20 March and the headings ``Continuation of Parshin's report'' and ``The case of projective varieties'' survive. The next visible mathematics is below a new date, 27 March, and concerns matrix factorization; it belongs to record 127. No theorem from Mumford's projective invariant theory is supplied by this heading. Catalogue source order is preserved even though record 128 has the earlier date.


\GStalk{129}{1967-03-27}{S. I. Gelfand}{P000204-P000207}{Petrovsky lacunas and Atiyah’s algebraic-geometric formulation}
S. Gelfand explicitly reports Atiyah's Bourbaki expos\'e 319, delivered in February 1967 and printed in 1968. The notes define strict hyperbolicity, the retarded fundamental solution and stable lacunas, construct the two Petrovsky homology classes, state the even-order/even-dimension criteria, and begin the residue-integral proof. The shift from \(\mathbb R^n\) to coordinates \(x_0,\ldots,x_n\) is resolved by writing \(N=n+1\). Empty real and complex sections, projective orientation choices, and twisted canonical forms must be distinguished. The notes stop at the Herglotz--Petrovsky formula; later algebraic de Rham proof details are not attributed to the seminar.

\begin{GSitem}{1966-129.1}{E}{Strict hyperbolicity of the homogeneous symbol}
\GSfield{Kind; source}{definition; \GSpage{P000204}}
\GSfield{The notes}{For a real homogeneous polynomial \(a(\xi_0,\ldots,\xi_n)\) of degree \(m\), every nonzero spatial tuple is required to give exactly \(m\) distinct real roots in \(\xi_0\). The ambient space is first called \(\mathbb R^n\), although \(n+1\) coordinates are then used.}
\GSfield{Statement}{Write \(N=n+1\). A real homogeneous polynomial \(a\) of degree \(m\) on \(\mathbb R^N\) is strictly hyperbolic with respect to \(e_0\) if \(a(e_0)\ne0\) and, for every \(\xi'\ne0\), the polynomial \(t\mapsto a(t,\xi')\) has \(m\) distinct real roots. The corresponding constant-coefficient differential operator is \(P=a(\partial_x)\), with any Fourier convention declared consistently.}
\GSfield{Reconstruction}{Noncharacteristic time direction and the distinction between strict hyperbolicity and the weaker real-root condition are supplied. The dimension is corrected uniformly, not by changing subsequent homology degrees.}
\GSfield{Fate}{Already in \citet[\S1, p.~88]{atiyah1967hyperbolic}, lecture February 1967, print 1968.}
\GSfield{Present status}{As of 9 November 2026, this remains an established construction under the stated hypotheses.}
\GSfield{Assessment}{literal statement: correct with missing hypotheses; publication match: exact; verification: read the result; development: none; confidence: high}
\end{GSitem}
\begin{GSitem}{1966-129.2}{E}{The retarded fundamental solution}
\GSfield{Kind; source}{theorem; \GSpage{P000204}}
\GSfield{The notes}{The system \(PE=\delta\), \(E=0\) for \(x_0<0\), has a unique solution by uniqueness for the Cauchy problem.}
\GSfield{Statement}{For a real homogeneous strictly hyperbolic operator with noncharacteristic time direction, there is a unique distribution \(E\in\mathcal D'(\mathbb R^N)\) satisfying \(PE=\delta_0\) and \(\operatorname{supp}E\subset\{x_0\ge0\}\). More precisely its support lies in the closed forward propagation cone dual to the chosen hyperbolicity cone. Uniqueness uses the vanishing of a homogeneous solution with past-zero support.}
\GSfield{Reconstruction}{Existence, implicit in ``we seek'' and then ``it is denoted \(E\),'' is distinguished from the stated uniqueness argument. The propagation cone is the convex support cone, a different notion from the projective dual wavefront hypersurface below.}
\GSfield{Fate}{Classical before the report; \citet[\S1, p.~88]{atiyah1967hyperbolic}. A Fourier--Laplace construction is in \citet[\S3]{atiyahbottgarding1970lacunas}.}
\GSfield{Present status}{As of 9 November 2026, this remains an established construction under the stated hypotheses.}
\GSfield{Assessment}{literal statement: correct; publication match: exact; verification: read the result; development: none; confidence: high}
\end{GSitem}
\begin{GSitem}{1966-129.3}{E}{Analyticity away from the dual characteristic hypersurface}
\GSfield{Kind; source}{theorem; \GSpage{P000204-P000205}}
\GSfield{The notes}{The characteristic cone \(A=\{a(\xi)=0\}\) has a dual cone \(B=\{b(x)=0\}\), described by tangent hyperplanes. The solution is analytic in every connected component of its complement.}
\GSfield{Statement}{For the generic smooth projective characteristic hypersurface \(A\subset\mathbb{CP}^{N-1}\), let \(B\) be the real cone over its projective dual discriminant in the dual projective space. For real \(x\notin B\), the section \(A_x=A\cap\{x\cdot\xi=0\}\) is transverse. The retarded fundamental solution restricts to a real-analytic function on each connected component of \(\{x_0>0\}\setminus B\). The characteristic cone in covector space and its dual in physical space need not have the same defining polynomial.}
\GSfield{Reconstruction}{``Dual'' here concerns tangent hyperplanes, not the convex polar cone. The generic complex smoothness is supplied as in the next page; this avoids treating every real tangent-locus description as an exact global algebraic equation without qualification.}
\GSfield{Fate}{The precise geometric setup and analytic statement are in \citet[\S1, pp.~89--91]{atiyah1967hyperbolic}.}
\GSfield{Present status}{As of 9 November 2026, this remains an established construction under the stated hypotheses.}
\GSfield{Assessment}{literal statement: correct with missing hypotheses; publication match: exact; verification: read the result; development: none; confidence: high}
\end{GSitem}
\begin{GSitem}{1966-129.4}{E}{Lacunas and stability under perturbation}
\GSfield{Kind; source}{definition; \GSpage{P000205}}
\GSfield{The notes}{A connected component \(C\) is a lacuna if \(E\equiv0\) there; it is stable if this persists for nearby equations.}
\GSfield{Statement}{A component \(C\) of the complement of the wavefront/dual discriminant in the future half-space is a lacuna in the notebook's sense if \(E_a|_C=0\). It is stable if, for a fixed \(x\in C\), all sufficiently small real homogeneous degree-\(m\) perturbations of \(a\) have fundamental solutions vanishing in a neighborhood of \(x\); the definition is independent of the chosen point in \(C\). The comparison is local or on compact subsets of the moving component, not a requirement that the original entire component remain fixed.}
\GSfield{Reconstruction}{The perturbation class and meaning of persistence are supplied. Atiyah--Bott--G\aa rding later use a broader notion of lacuna allowing a regular continuation and call the zero-solution case a strong lacuna; their terminology must not silently replace the notebook's.}
\GSfield{Fate}{Defined by Petrovsky before 1945 \citep[pp.~290--291]{petrovsky1945lacunas}; the formulation follows \citet[pp.~89--90]{atiyah1967hyperbolic}.}
\GSfield{Present status}{As of 9 November 2026, the definition and its homological criteria are classical. Later terminology is distinguished in \citet[\S9--10]{atiyahbottgarding1970lacunas}.}
\GSfield{Assessment}{literal statement: correct with missing hypotheses; publication match: exact; verification: read the result; development: none; confidence: high}
\end{GSitem}
\begin{GSitem}{1966-129.5}{E}{Reduction to a nonsingular complex characteristic hypersurface}
\GSfield{Kind; source}{method; \GSpage{P000205}}
\GSfield{The notes}{The real polynomial defines \(A\subset\mathbb{CP}^n\), assumed to have no complex singularities. The notes say a small real change of coefficients can achieve this.}
\GSfield{Statement}{Strict hyperbolicity is open in the space of real homogeneous degree-\(m\) polynomials. The complex singular-hypersurface discriminant is a proper algebraic subset of that coefficient space. Consequently a strictly hyperbolic polynomial can be approximated by real strictly hyperbolic polynomials with nonsingular complex projective zero locus. This generic reduction is legitimate for a question defined through stability under all sufficiently small admissible perturbations.}
\GSfield{Reconstruction}{Strict hyperbolicity alone already gives smoothness of the real projective zero locus, but not complex smoothness. The reduction does not prove that the original complex variety was nonsingular.}
\GSfield{Fate}{The prelecture reduction is \citet[\S1, p.~90]{atiyah1967hyperbolic}.}
\GSfield{Present status}{As of 9 November 2026, this remains an established construction under the stated hypotheses.}
\GSfield{Assessment}{literal statement: correct with missing hypotheses; publication match: exact; verification: read the result; development: none; confidence: high}
\end{GSitem}
\begin{GSitem}{1966-129.6}{E}{The real hyperplane-section cycle beta}
\GSfield{Kind; source}{construction; \GSpage{P000205}}
\GSfield{The notes}{The notes form \(A_x=A\cap\{x\cdot\xi=0\}\), call \((A_x)_{\mathbb R}\) a cycle of dimension \(n-2\), and write \(\beta_x\in H_{n-2}(A_x)\). The phrase ``a sum of spheres'' is attached to this section.}
\GSfield{Statement}{Assume \(N>2\) and \(m,N\) even, with \(A\subset\mathbb{CP}^{N-1}\) smooth and \(x\notin B\). Hyperbolicity gives \(m/2\) spherical components of \(A_{\mathbb R}\), each of real dimension \(N-2\), with compatible orientations from the hyperbolicity direction. A trivialization of the real hyperplane line bundle over \(A_{\mathbb R}\) makes \(x\cdot\xi\) a real function. Its transverse zero locus is oriented and defines \[\beta_x=[(A_x)_{\mathbb R}]\in H_{N-3}(A_x;\mathbb C).\]}
\GSfield{Reconstruction}{The coefficient field and orientation data are restored from the identified source. Atiyah's sphere decomposition concerns \(A_{\mathbb R}\), not a stated decomposition of every \((A_x)_{\mathbb R}\); the notebook's placement of that phrase is not accepted as a general theorem about hyperplane sections.}
\GSfield{Fate}{Exact cycle construction: \citet[\S1, p.~91]{atiyah1967hyperbolic}; the sphere phrase has only a partial match.}
\GSfield{Present status}{As of 9 November 2026, this remains an established construction under the stated hypotheses.}
\GSfield{Assessment}{literal statement: correct with missing hypotheses; publication match: partial; verification: read the result; development: none; confidence: high}
\end{GSitem}
\begin{GSitem}{1966-129.7}{E}{The wave operator: an empty real section, not an empty complex quadric}
\GSfield{Kind; source}{construction; \GSpage{P000205}}
\GSfield{The notes}{For \(a=\xi_0^2-\sum_{j=1}^n\xi_j^2\) and \(x=(1,0,\ldots,0)\), the notes obtain \(\xi_0=0\), \(\sum\xi_j^2=0\), then say that in projective space this has no points.}
\GSfield{Statement}{The real projective section is empty because a sum of real squares vanishes only at the zero vector; hence \(\beta_x=0\). The complex projective section is nonempty for \(n\ge2\): for example \([\xi_0:\xi_1:\xi_2:\cdots]=[0:1:i:0:\cdots]\) lies on it. If \(N=n+1>2\) is even, the order \(m=2<N\) and the Petrovsky criterion yields the familiar stable interior lacuna of the retarded wave solution.}
\GSfield{Reconstruction}{\ours\ The explicit complex point checks the literal error independently of the real positive-definiteness argument. The parity and dimension restriction is essential: the empty real section by itself does not assert Huygens' principle in every dimension.}
\GSfield{Fate}{The wave example is classical \citep[p.~289]{petrovsky1945lacunas}; see \citet[pp.~88--89, Theorem~1]{atiyah1967hyperbolic}.}
\GSfield{Present status}{As of 9 November 2026, the corrected example is standard; the literal claim of no complex points is false for the dimensions used in the theorem.}
\GSfield{Assessment}{literal statement: false as written; publication match: partial; verification: read the result; development: annotated; confidence: high}
\end{GSitem}
\begin{GSitem}{1966-129.8}{E}{The signed relative cycle alpha and its boundary}
\GSfield{Kind; source}{construction; \GSpage{P000206}}
\GSfield{The notes}{The signed chain \(\{x\cdot\xi>0\}-\{x\cdot\xi<0\}\) on \(A_{\mathbb R}\) defines \(\alpha_x\in H_{n-1}(A,A_x)\), with \(\partial\alpha_x=2\beta_x\).}
\GSfield{Statement}{In the same oriented even-parity setting, let \(f_x\) denote the real function induced by the linear form after trivializing the real hyperplane line bundle over \(A_{\mathbb R}\). The signed fundamental chain \(\operatorname{sgn}(f_x)[A_{\mathbb R}]\) is a relative cycle of \((A,A_x)\), defining \(\alpha_x\in H_{N-2}(A,A_x;\mathbb C)\). With consistent boundary orientation, \(\partial\alpha_x=2\beta_x\), because the two sides contribute with the same boundary sign after the negative side is subtracted.}
\GSfield{Reconstruction}{The sign of a homogeneous linear form is not intrinsically defined on a real projective point without a lift/trivialization. The source supplies exactly this missing choice. \(\alpha\) is relative, not generally an absolute cycle.}
\GSfield{Fate}{The construction and factor two appear in \citet[\S1, pp.~91--92]{atiyah1967hyperbolic}.}
\GSfield{Present status}{As of 9 November 2026, this remains an established construction under the stated hypotheses.}
\GSfield{Assessment}{literal statement: correct with missing hypotheses; publication match: exact; verification: read the result; development: none; confidence: high}
\end{GSitem}
\begin{GSitem}{1966-129.9}{E}{Constancy of Petrovsky classes within a chamber}
\GSfield{Kind; source}{theorem; \GSpage{P000206}}
\GSfield{The notes}{When \(x\) moves in one connected component of the cone complement, \(\alpha_x\) and \(\beta_x\) do not change; the notation becomes \(\alpha_C,\beta_C\).}
\GSfield{Statement}{Over a component \(C\) of the complement of the dual discriminant, the smooth hyperplane sections and pairs \((A,A_x)\) form locally trivial families. Their homology groups are identified by continuation along paths, under which the geometrically defined classes \(\alpha_x\) and \(\beta_x\) are locally constant. In particular their vanishing is independent of the chosen \(x\in C\). This is a statement in the corresponding homology local systems, not literal equality of elements of different groups without identification.}
\GSfield{Reconstruction}{The local-system interpretation supplies the missing comparison maps; no global choice of an entire homology basis or absence of monodromy is asserted.}
\GSfield{Fate}{Stated before the two criteria in \citet[\S1, p.~92]{atiyah1967hyperbolic}.}
\GSfield{Present status}{As of 9 November 2026, this remains an established construction under the stated hypotheses.}
\GSfield{Assessment}{literal statement: correct with missing hypotheses; publication match: exact; verification: read the result; development: none; confidence: high}
\end{GSitem}
\begin{GSitem}{1966-129.10}{E}{Stable-lacuna criterion below the dimension}
\GSfield{Kind; source}{theorem; \GSpage{P000206}}
\GSfield{The notes}{If the order \(m\) and the dimension \(n+1\) are even and \(m<n+1\), \(\beta_C=0\) is necessary and sufficient for a stable lacuna.}
\GSfield{Statement}{Let \(N>2\) and \(m\) be even, \(m<N\), and let \(a\) be real homogeneous and strictly hyperbolic with smooth complex characteristic hypersurface. A component \(C\subset\{x_0>0\}\setminus B\) is a stable lacuna in the zero-solution sense if and only if \(\beta_C=0\) in \(H_{N-3}(A_x;\mathbb C)\).}
\GSfield{Reconstruction}{Complex homology coefficients, complex smoothness, strict hyperbolicity, and \(N>2\) are supplied from the report's precise source. An integral-homology torsion obstruction is not part of this theorem.}
\GSfield{Fate}{Petrovsky's theorem, restated in \citet[Theorem~1, p.~92]{atiyah1967hyperbolic}, circulated in the February 1967 lecture before this report.}
\GSfield{Present status}{As of 9 November 2026, the criterion is established. Atiyah's June 1967 addendum additionally says that in this low-order regime every lacuna is stable; the addendum postdates the seminar. Later broader formulations appear in \citet[\S10]{atiyahbottgarding1970lacunas}.}
\GSfield{Assessment}{literal statement: correct with missing hypotheses; publication match: exact; verification: read the result; development: none; confidence: high}
\end{GSitem}
\begin{GSitem}{1966-129.11}{E}{Stable-lacuna criterion at or above the dimension}
\GSfield{Kind; source}{theorem; \GSpage{P000206}}
\GSfield{The notes}{If \(m\ge n+1\), with \(m,n+1\) even, \(\alpha_C=0\) is necessary and sufficient for a stable lacuna.}
\GSfield{Statement}{Under the same strict-hyperbolicity and smoothness assumptions, let \(N>2\) and \(m\) be even with \(m\ge N\). Then \(C\) is a stable lacuna if and only if \(\alpha_C=0\) in \(H_{N-2}(A,A_x;\mathbb C)\). The necessary condition \(\beta_C=0\), obtained from \(\partial\alpha_C=2\beta_C\), is weaker and cannot replace the relative-class condition in this regime.}
\GSfield{Reconstruction}{The parity, homology field, and regularity are retained from the exact source. This statement addresses stability, not the classification of all accidental lacunas or all singular wavefronts.}
\GSfield{Fate}{The exact statement is \citet[Theorem~2, p.~92]{atiyah1967hyperbolic}, after Petrovsky and before the March report.}
\GSfield{Present status}{As of 9 November 2026, the criterion is classical. Atiyah's June 1967 addendum excludes nontrivial stable lacunas in this even-dimensional high-order regime; later terminology calls the notebook's zero regions strong lacunas \citep[pp.~185--188]{atiyahbottgarding1970lacunas}.}
\GSfield{Assessment}{literal statement: correct with missing hypotheses; publication match: exact; verification: read the result; development: none; confidence: high}
\end{GSitem}
\begin{GSitem}{1966-129.12}{E}{The projective canonical bundle and Euler form}
\GSfield{Kind; source}{construction; \GSpage{P000206}}
\GSfield{The notes}{The notes write \(K_{\mathbb{CP}^n}\simeq\mathcal O(-n-1)\), assert a unique holomorphic \(n\)-form with one pole of order \(n+1\), and give \(\omega=\sum(-1)^j x_j\,dx_0\wedge\cdots\widehat{dx_j}\cdots\wedge dx_n\).}
\GSfield{Statement}{With \(N=n+1\), \(K_{\mathbb{CP}^{N-1}}\simeq\mathcal O(-N)\) and \(H^0(K(N))\simeq\mathbb C\), whereas \(H^0(K)=0\). Thus a nonzero generator is unique up to scalar as a twisted canonical section, or as a meromorphic top form with poles of order at most \(N\) along a fixed hyperplane. On the affine cone its homogeneous representative is \[\omega=\sum_{j=0}^{N-1}(-1)^j\xi_j\,d\xi_0\wedge\cdots\wedge\widehat{d\xi_j}\wedge\cdots\wedge d\xi_{N-1}.\]}
\GSfield{Reconstruction}{The coordinates are renamed \(\xi\) to distinguish symbol space from the physical variable. The cone expression has homogeneity \(N\); it is not an ordinary global holomorphic top form on projective space. A fixed polar hyperplane and uniqueness up to scalar replace the literal ``unique holomorphic form with a pole.}
\GSfield{Fate}{The exact twisted-bundle formulation is \citet[\S2, p.~93]{atiyah1967hyperbolic}.}
\GSfield{Present status}{As of 9 November 2026, this remains an established construction under the stated hypotheses.}
\GSfield{Assessment}{literal statement: correct with missing hypotheses; publication match: exact; verification: read the result; development: none; confidence: high}
\end{GSitem}
\begin{GSitem}{1966-129.13}{E}{Poincare residue with the required homogeneity}
\GSfield{Kind; source}{construction; \GSpage{P000207}}
\GSfield{The notes}{The residue form is introduced by \(\omega/a=(da/a)\wedge\omega_a+\Theta\). An overwritten parenthesis appears to say that \(\Theta\) is zero on \(a=0\).}
\GSfield{Statement}{For a smooth hypersurface \(A=\{a=0\}\), the local decomposition of a meromorphic top form with a simple pole is \(\omega/a=(da/a)\wedge\eta+\Theta\), where \(\Theta\) is holomorphic and \(\omega_a=\eta|_A\) is the residue. With homogeneous projective representatives, \(\omega_a\) is a section of \(K_A\otimes\mathcal O(N-m)\); since \(K_A\simeq\mathcal O(m-N)\), the twist is essential unless \(m=N\). Multiplying by a homogeneous polynomial of degree \(m-N\), when \(m\ge N\), produces an ordinary holomorphic top form on \(A\).}
\GSfield{Reconstruction}{The pullback of the holomorphic remainder to the hypersurface vanishes by degree; its coefficient functions need not vanish there. The exact overwritten parenthesis is not resolved by this interpretation. Local residues glue with the displayed twist, rather than as an untwisted form in every degree.}
\GSfield{Fate}{The residue and line-bundle twists are explicit in \citet[\S2, p.~93]{atiyah1967hyperbolic}.}
\GSfield{Present status}{As of 9 November 2026, this remains an established construction under the stated hypotheses.}
\GSfield{Assessment}{literal statement: correct with missing hypotheses; publication match: exact; verification: read the result; development: none; confidence: high}
\end{GSitem}
\begin{GSitem}{1966-129.14}{E}{The Herglotz–Petrovsky residue integral}
\GSfield{Kind; source}{theorem; \GSpage{P000207}}
\GSfield{The notes}{For \(m\ge n+1\) the final formula is \(E(x)=c\int_{A_{\mathbb R}}\operatorname{sgn}(x\cdot\xi)(x\cdot\xi)^{m-n-1}\omega_a\), followed by a reference to the first volume of generalized functions.}
\GSfield{Statement}{Under the smooth, strictly hyperbolic setting above, with \(m,N\) even and \(m\ge N>2\), the analytic value of the retarded fundamental solution for \(x_0>0\), away from the dual discriminant, is \[E(x)=c_{m,N}\int_{A_{\mathbb R}}\operatorname{sgn}(f_x)\,(x\cdot\xi)^{m-N}\omega_a.\] Orientations and the real line-bundle trivialization are those used for \(\alpha_x\). The polynomial factor cancels the projective twist of the residue; the normalization constant depends on the Fourier convention, the chosen generator, and the normalization of \(P\). The fundamental solution is understood distributionally on the wavefront.}
\GSfield{Reconstruction}{The parity and dimension restrictions are inherited from the theorem, not repeated in the final notebook line. The shorthand constant is not assigned an unsupported numerical value. The formula supplies the announced proof route, but the omitted converse uses substantial algebraic geometry and is not proved in these pages.}
\GSfield{Fate}{The exact formula is \citet[\S2, equation (2.1), p.~93]{atiyah1967hyperbolic}; its classical origin is described in \citet[Chapter~III]{petrovsky1945lacunas}.}
\GSfield{Present status}{As of 9 November 2026, this residue representation is established. Generalized formulas and their conventions are developed in \citet[\S7]{atiyahbottgarding1970lacunas}.}
\GSfield{Assessment}{literal statement: correct with missing hypotheses; publication match: exact; verification: read the result; development: none; confidence: high}
\end{GSitem}

\section{Developments}\label{dev:1966}
\subsection{GP-1966-09: what correlation bounds actually prove}
\label{dev:canonical}

\paragraph{Historical question and the bounded scope.}
This section is a complete critical annotation of item 1966-100.11, whose recorded variance argument was published by Minlos (class A). It makes no claim to solve a previously open C/D problem. Item 1966-100.2 retains class F for the broader programme’s incomplete quantifiers. The notes posit bulk limits, control by distance to the boundary, and clustering; for
\[
 F_\Omega=\sum_{i\ne j}f(X_i-X_j),\qquad f\in L^1(\mathbb R^d),
\]
they infer a variance proportional to particle number. The strongest defensible historical formulation is conditional: find quantitative hypotheses on the canonical correlations which make this implication valid. The larger existence-and-clustering programme leaves its exceptional thermodynamic parameters unspecified, and therefore is not one sharply formulated all-density conjecture.

There is a useful published anchor beyond the shorthand. \citet[Lecture III, \S4, pp.~162--164, especially Problem~3]{minlos1968lectures} asks for precise inequalities in the representative-observable argument. On p.~164 the published account explicitly restores an integrable mixing majorant and a boundary-strip condition; these must not be portrayed as omissions persisting in that paper. The separate short-distance issue remains: condition (3.24) is an \(L^1\) assumption, while the variance calculation contains \(f^2\). The following Remark~1 discusses an exponential large-deviation estimate, which is a stronger question than the concentration proved below. The lectures originated in September 1965 and were amplified in the autumn 1966 seminar; their 1968 publication date is not a discovery date for the elementary moment argument.

Our contribution is a checked clarification: an explicit sufficient fourth-order correlation bound, a variance estimate with constants, a truncation argument covering \(L^1\), and counterexamples to stronger intermediate claims. It is not a claim of a new Gibbs-state theorem or a new general result in probability. All supplied proofs and computations carry the required editorial marks.

\paragraph{Exact literature scope.}
For stable finite-range interactions in a sufficiently dilute canonical regime, \citet[Theorem~2.5 and Corollary~2.6]{pulvirentitsagkarogiannis2015finite} control the truncated \emph{two-point} correlation and retain a finite-size contribution alongside spatial decay. The preprint appeared in September 2014 and the article in 2015. \citet[Theorems~2.1 and 2.7]{kunatsagkarogiannis2018density}, first circulated in 2016 and published in 2018, give convergent density expansions and correlation estimates within an explicit small-density domain. These are substantial realizations of the gas-phase programme. Neither citation, by itself, proves the four-point assumption below for arbitrary density, nor does a two-point estimate automatically control the covariance of two pair observables. We do not infer a theorem outside the hypotheses actually checked.

\paragraph{Framework and quantifiers.}
Let \(\Omega\subset\mathbb R^d\) be a measurable bounded region of positive volume \(V=|\Omega|\). Let \(\mathcal X\) be a simple point process with finitely many points in \(\Omega\) almost surely. Its symmetric factorial moment densities \(r_k\), \(1\leq k\leq4\), are defined by
\[
 \mathbb E\!\sum_{i_1,\ldots,i_k\;\mathrm{distinct}}
 h(X_{i_1},\ldots,X_{i_k})
 =\int_{\Omega^k}h(x_1,\ldots,x_k)r_k(x_1,\ldots,x_k)\,dx_1\cdots dx_k
\]
for nonnegative measurable \(h\). Thus the sums are ordered, exactly as on P000157. The particle number may be fixed, as in the canonical ensemble, or random. Define \(f_{ij}=f(x_i-x_j)\) for a real even measurable kernel \(f\). Write
\[
 B_{12;34}=r_4(x_1,x_2,x_3,x_4)-r_2(x_1,x_2)r_2(x_3,x_4).
\]
This is the covariance density for two disjoint ordered pairs. It is not identical to the fourth factorial cumulant: lower connected terms also contribute.

For a sequence of regions, all constants below and the same majorant \(g\) must be independent of the region. Assume, almost everywhere,
\begin{align}
 0\leq r_2&\leq C_2,&0\leq r_3&\leq C_3,\label{eq:bounds23}\\
 |B_{12;34}|&\leq C_4\sum_{a=1}^2\sum_{b=3}^4g(x_a-x_b)+\frac{C_0}{V},
 &g&\geq0,\quad g\in L^1(\mathbb R^d).\label{eq:pairmix}
\end{align}
The nonnegative constants may depend on the fixed physical parameters. No translation invariance, regular boundary, or infinite-volume state is needed for the variance implication itself. Such assumptions become relevant when deriving these bounds from a Gibbs model or when identifying the limiting mean.

\begin{lemma}[Exact collision expansion]\ours\ 
For an even kernel on a finite configuration,
\[
 F_\Omega^2=2\sum_{i\ne j}f_{ij}^2
 +4\sum_{i,j,k\;\mathrm{distinct}}f_{ij}f_{jk}
 +\sum_{i,j,k,\ell\;\mathrm{distinct}}f_{ij}f_{k\ell}.
\]
If the corresponding absolute integrals are finite, then
\begin{equation}
 \operatorname{Var}F_\Omega
 =2\int_{\Omega^2}f_{12}^2r_2
 +4\int_{\Omega^3}f_{12}f_{23}r_3
 +\int_{\Omega^4}f_{12}f_{34}B_{12;34}.
 \label{eq:exactvariance}
\end{equation}
\end{lemma}
\begin{proof}
Expand the product of two sums over ordered unequal pairs. Their union has two, three, or four labels. For two labels, the second pair equals the first or its reversal, giving coefficient two. With three labels, the shared label occupies one of two places in each pair, giving four possibilities. Symmetry of \(f\) and relabelling put them in the displayed common form. Four distinct labels require no multiplicity factor. The factorial moment identities integrate the three terms. Subtract
\((\mathbb EF_\Omega)^2=\int_{\Omega^4}f_{12}f_{34}r_2(x_1,x_2)r_2(x_3,x_4)\)
to obtain \eqref{eq:exactvariance}. Absolute integrability, verified in the following proof under its assumptions, legitimizes all rearrangements. This is a re-derivation of the notebook identity, not an attribution of a new combinatorial formula.
\end{proof}

\begin{theorem}[An explicit finite-volume variance bound]\ours\ 
Assume \eqref{eq:bounds23}--\eqref{eq:pairmix} and \(f\in L^1(\mathbb R^d)\cap L^2(\mathbb R^d)\), real and even. Then
\begin{align}
 \operatorname{Var}F_\Omega&\leq V K(f),\label{eq:variancebound}\\
 K(f)&=2C_2\|f\|_2^2+
 \bigl(4C_3+4C_4\|g\|_1+C_0\bigr)\|f\|_1^2.\label{eq:K}
\end{align}
Consequently, for \(\epsilon>0\),
\[
 \mathbb P\!\left(\left|\frac{F_\Omega-\mathbb EF_\Omega}{V}\right|>\epsilon\right)
 \leq\frac{K(f)}{V\epsilon^2}.
\]
\end{theorem}
\begin{proof}
Integrating first in a difference variable gives
\[
 \int_{\Omega^2}|f_{12}|^2\leq V\|f\|_2^2,
 \qquad
 \int_{\Omega^3}|f_{12}f_{23}|\leq V\|f\|_1^2.
\]
In any summand of the spatial part of \eqref{eq:pairmix}, the edges \(12\), \(34\), and \(ab\), with \(a\in\{1,2\}\), \(b\in\{3,4\}\), form a tree on four vertices. Successively integrate a leaf over \(\Omega\), enlarging that integral to \(\mathbb R^d\). Each removed edge contributes its \(L^1\) norm; integration of the last vertex contributes \(V\). Hence, for each of the four cross edges,
\[
 \int_{\Omega^4}|f_{12}f_{34}|g(x_a-x_b)
 \leq V\|f\|_1^2\|g\|_1.
\]
The canonical correction is bounded separately:
\[
 \frac{C_0}{V}\int_{\Omega^4}|f_{12}f_{34}|
 \leq\frac{C_0}{V}(V\|f\|_1)^2
 =C_0V\|f\|_1^2.
\]
These estimates also establish the required absolute moments: the unsubtracted four-point integral has the additional integrable product term bounded by \(C_2^2V^2\|f\|_1^2\). Apply \eqref{eq:exactvariance} and the triangle inequality. Chebyshev’s inequality proves the probability bound.
\end{proof}

The conclusion is \(O(V)\), not an asymptotic equivalent with a positive coefficient. Neither convergence of the normalized covariance integrals nor positivity of a limit was assumed; \(f=0\) alone rules out an unconditional positive coefficient. If \(N_\Omega/V\to\rho\in(0,\infty)\) with deterministic canonical particle numbers, the same result gives concentration after normalization by \(N_\Omega\).

\begin{corollary}[The integrable-kernel conclusion survives truncation]\ours\ 
Under the same uniform correlation hypotheses, let \(V\to\infty\) and let \(f\in L^1(\mathbb R^d)\) be real and even. Then
\[
 \frac{F_\Omega-\mathbb EF_\Omega}{V}\longrightarrow0
 \quad\hbox{in probability},
\]
even when \(\operatorname{Var}F_\Omega\) is infinite; the centered convergence also holds in \(L^1\). More explicitly, put
\(f_M=f\mathbf1_{\{|f|\leq M\}}\). For every \(M>0\) and \(\epsilon>0\),
\begin{equation}
 \mathbb P\!\left(\left|\frac{F_\Omega-\mathbb EF_\Omega}{V}\right|>\epsilon\right)
 \leq\frac{4K(f_M)}{V\epsilon^2}
 +\frac{4C_2}{\epsilon}\|f-f_M\|_1.
 \label{eq:truncation}
\end{equation}
\end{corollary}
\begin{proof}
The truncated kernel belongs to both spaces because
\(\|f_M\|_2^2\leq M\|f\|_1\). If \(R_\Omega=F_\Omega(f)-F_\Omega(f_M)\), then the factorial two-point formula implies
\[
 \mathbb E|R_\Omega|\leq C_2V\|f-f_M\|_1,
 \qquad
 \mathbb E|R_\Omega-\mathbb ER_\Omega|
 \leq2C_2V\|f-f_M\|_1.
\]
Split the centered observable into its truncated part and the centered remainder, and apply the union bound with threshold \(\epsilon/2\) to each. Chebyshev controls the first part and Markov controls the second, giving exactly \eqref{eq:truncation}. Cauchy--Schwarz also gives the stronger bound
\[
 \frac{\mathbb E|F_\Omega-\mathbb EF_\Omega|}{V}
 \leq\sqrt{\frac{K(f_M)}V}+2C_2\|f-f_M\|_1.
\]
First let \(V\to\infty\) at fixed \(M\), then let \(M\to\infty\). The order of limits is part of both conclusions. Convergence to a fixed thermodynamic value additionally requires convergence of \(\mathbb EF_\Omega/V\); concentration about a moving expectation does not establish that limit.
\end{proof}

\begin{proposition}[A sufficient condition identifying the mean]\ours\ 
Let \(V_n=|\Omega_n|\to\infty\), \(r_{2,n}\leq C_2\), and define
\[
 A_n(z)=\frac1{V_n}\int_{\mathbb R^d}
 \mathbf1_{\Omega_n}(x)\mathbf1_{\Omega_n}(x-z)
 r_{2,n}(x,x-z)\,dx.
\]
If \(A_n(z)\to R(z)\) for almost every \(z\), then for \(f\in L^1\),
\[
 \frac{\mathbb EF_{\Omega_n}}{V_n}\longrightarrow
 m=\int_{\mathbb R^d}f(z)R(z)\,dz.
\]
Combined with the preceding corollary, this yields \(F_{\Omega_n}/V_n\to m\) in probability. For the ideal canonical gas, a van Hove exhaustion and \(N_n/V_n\to\rho\) give \(R(z)=\rho^2\).
\end{proposition}
\begin{proof}
Change variables \(z=x-y\) in the expectation. Its normalized value is \(\int fA_n\). Since \(0\leq A_n\leq C_2\), dominated convergence applies. In the ideal gas,
\[
 A_n(z)=\frac{(N_n)_2}{V_n^2}
 \frac{|\Omega_n\cap(\Omega_n+z)|}{V_n}.
\]
For every fixed \(z\), the van Hove condition makes the last fraction tend to one. The first factor tends to \(\rho^2\). This also exhibits why pointwise correlation convergence near one fixed origin is not by itself the required moving-volume statement.
\end{proof}

\begin{example}[Finite expectation and infinite variance in an ideal gas]\ours\ 
Take \(\Omega=[0,L]^d\), \(L>2\), and \(N\geq2\) independent uniformly distributed positions, viewed as an unordered simple configuration. Coincidences have probability zero. The potential is zero, hence stable and of finite range. For
\[
 f(z)=|z|^{-\alpha}\mathbf1_{\{0<|z|<1\}},\qquad
 \frac d2\leq\alpha<d,\qquad f(0)=0,
\]
the kernel is real, even, nonnegative and integrable, because the radial integral near zero is proportional to \(\int_0^1r^{d-1-\alpha}\,dr\). Its square is not integrable. The expectation of the ordered-pair sum is finite, but its second moment and variance are infinite.
\end{example}
\begin{proof}
The factorial density is \(r_2=N(N-1)/V^2>0\). The expectation is bounded by \(r_2V\|f\|_1\). Restrict \(x_1\) to the interior box \([1,L-1]^d\); a unit ball of differences around zero remains inside the vessel. Therefore \(\int_{\Omega^2}f_{12}^2r_2=\infty\). All summands of \(F_\Omega\) are nonnegative, and the two-label terms alone imply \(\mathbb EF_\Omega^2\geq2\int f_{12}^2r_2=\infty\). A finite first moment and infinite second moment give infinite variance. Every hypothesis used here has been specified; no interacting-model approximation is involved.
\end{proof}

\begin{remark}[A hard core does not by itself repair arbitrary \(L^1\) kernels]\ours\ 
If a positive hard-core radius \(a\) is required, use two particles with potential \(+\infty\) for distances below \(a\) and zero otherwise, in a cube of side \(L>6a\). Its canonical partition function is finite and positive, and its two-point density is a positive constant on allowed pairs. The even kernel
\[
 f(z)=\bigl||z|-2a\bigr|^{-3/4}
 \mathbf1_{\{0<||z|-2a|<a/4\}}
\]
is in \(L^1\) but not \(L^2\): in spherical coordinates the only singularity is a codimension-one shell, whose radial exponents are \(-3/4\) and \(-3/2\). The shell lies strictly outside the core. A positive-measure interior set of first-particle positions accommodates the whole shell, so the previous second-moment divergence proof applies. Values exactly on the shell can be set to zero without changing any integral. Thus a hard core removes only a singularity confined to forbidden separations, not the general defect in an unrestricted \(L^1\) assumption.
\end{remark}

\begin{proposition}[The canonical correction cannot be discarded]\ours\ 
For the ideal canonical gas,
\[
 r_k=\frac{(N)_k}{V^k},\qquad
 r_4-r_2^2=-\frac{N(N-1)(4N-6)}{V^4}.
\]
For \(N\geq2\), with \(\rho_N=N/V\),
\[
 |r_4-r_2^2|\leq\frac{4\rho_N^3}{V},\qquad
 r_2\leq\rho_N^2,\quad r_3\leq\rho_N^3.
\]
Consequently \eqref{eq:pairmix} holds with \(C_4=0\), \(C_0=4\bar\rho^3\) whenever \(\rho_N\leq\bar\rho\). For a fixed \((N,V)\), the variance estimate becomes
\[
 \operatorname{Var}F_\Omega\leq
 2\rho_N^2V\|f\|_2^2+8\rho_N^3V\|f\|_1^2.
\]
\end{proposition}
\begin{proof}
The factorial density follows from the uniform labelled density. Subtract
\((N)_2^2\) from \((N)_4\), factor out \(N(N-1)\), and use
\((N-2)(N-3)-N(N-1)=6-4N\).
The inequality follows from \((N-1)(4N-6)\leq4N^2\). Substitute these constants into \eqref{eq:K}. The \(V^{-1}\) covariance density is spatially constant: it is the effect of fixing the total number, not a failure caused by an interaction tail. A strictly exponential spatial majorant without a finite-size term cannot represent this constraint uniformly over growing vessels.
\end{proof}

\begin{remark}[The effective-volume quantifier in the scan is correct]\ours\ 
For each fixed \(k\leq N\), put
\[
 \rho_{N,k}=\frac{(N)_k^{1/k}}{V},\qquad
 v_{N,k}=\frac{V}{(N)_k^{1/k}}.
\]
Then \(r_k=\rho_{N,k}^k\) exactly. If \(V/N\to v>0\), then \(v_{N,k}\to v\) for every fixed \(k\). The notebook explicitly allows the adjustment to depend on \(k\); this elementary test supports that reading. One common \(\rho_N\) cannot exactly match both \(r_1\) and \(r_2\), since it would require simultaneously \(\rho_N=N/V\) and \(\rho_N^2=N(N-1)/V^2\). This does not disprove the actual order-dependent claim and proves no interacting boundary estimate.
\end{remark}

\paragraph{Independent checks and what they found.}
\oursComputed\ The supplied standard-library program \texttt{check\_pair\_partitions.py} checks the ordered-pair expansion on configurations with \(N=0,\ldots,8\), using an even kernel with both signs and exact rational arithmetic. It independently constructs the factorial marginals of a nonuniform simple process on seven sites by enumerating all \(128\) configurations, then compares the direct variance to the correlation-density formula. It also verifies the canonical polynomial identity symbolically by coefficient arithmetic and at \(205\) rational parameter pairs. All comparisons are exact and pass. Testing \(N=0,1,2,3\) ensures that falling factorials and the two-label term are not being handled by an unjustified large-\(N\) shortcut.

The analytic argument was separately re-derived: each of the four cross terms is a three-edge tree, so the coefficient of \(C_4\|g\|_1\) is four; centering the truncation remainder costs a factor two, producing the coefficient \(4C_2/\epsilon\) in \eqref{eq:truncation}. The independent review confirmed these constants and emphasized that they must be uniform in volume. It also confirmed the order of limits and the need for a separate mean-convergence hypothesis. These checks support the complete proofs above; finite computations are not used as substitutes for the general proofs. The review file is included with the code and outputs.

\paragraph{Remaining obstruction and a feasible further run.}
The missing model-specific step is to derive \eqref{eq:pairmix}, or a comparably integrable bound on its left-hand side, for a chosen interacting canonical gas with constants uniform in the vessel. The four-point object must be obtained with its canonical normalization intact. A two-point exponential-decay estimate alone is insufficient. The second task is to prove the averaged bulk convergence used for the mean; the third, if desired, is to control enough higher cumulants to obtain the exponential deviation estimate discussed in the Remark~1 following the published Problem~3. The variance bound gives only the displayed Chebyshev rate.

A bounded follow-up should fix a particular stable finite-range potential and a certified dilute expansion domain, derive the required two-pair covariance from that expansion, and calculate explicit constants and boundary remainders. It could then establish the complete quantitative statement for that model and regime. It should not promise an all-density continuation across phase coexistence or a universal classification of exceptional parameters. The evidence does not show that the original programme was abandoned: later rigorous canonical expansions develop parts of it, while the elementary probability implication and the model-dependent correlation analysis remain logically distinct.

\section{Problem index}\label{problems:1966}
Earlier-year occurrence identifiers are retained as text; current-year identifiers link to their entries.

\subsection{GP-1965-03: The folds and cusps of the perturbed complex square}\label{prob:GP-1965-03}
For the real smooth map \(F_\varepsilon:\mathbb C\to\mathbb C\), \(F_\varepsilon(z)=z^2+\varepsilon\bar z\), \(\varepsilon>0\), determine the critical set, critical values, local singularity types, and the regular-fibre cardinalities.
\par\noindent\textit{Occurrences:} 1965-72.5, \hyperref[item:1966-115.8]{1966-115.8}.
\par\noindent\textit{Status (2026-11-09):} As of 9 November 2026: settled. The critical circle is \(|z|=\varepsilon/2\); its image is the three-cusped curve \(\varepsilon^2(e^{2it}/4+e^{-it}/2)\), with cusps where \(e^{3it}=1\) and folds elsewhere. There are four preimages in its interior and two in its exterior. The explicit treatment is \citet[\S1, Example 6, pp.~4--5]{arnold1968singularities}, published in 1968 from 1966 lectures. This repeats the earlier indexed question; it does not create a new problem. See \citep{arnold1968singularities}.

\subsection{GP-1966-01: Algebraic classification of Anosov diffeomorphisms}\label{prob:GP-1966-01}
Is every Anosov diffeomorphism of a closed manifold topologically conjugate to an affine Anosov map of an infranilmanifold?
\par\noindent\textit{Occurrences:} \hyperref[item:1966-94.4]{1966-94.4}.
\par\noindent\textit{Status (2026-11):} As of November 2026, unresolved in general in the checked literature; see Spatzier’s Conjecture 2.3. The weaker manifold-realization question also remains open. The smooth-category replacement is too strong. See \citep{smale1967dynamical,spatzier2016rigidity,deremainkar2023anosov,farrellgogolev2012exotic}.

\subsection{GP-1966-02: Structural stability of Morse--Smale systems}\label{prob:GP-1966-02}
Are Morse--Smale maps and flows \(C^1\) structurally stable?
\par\noindent\textit{Occurrences:} \hyperref[item:1966-94.8]{1966-94.8}.
\par\noindent\textit{Status (2026-11):} Affirmatively solved; 1968 circulation, 1969 Russian and 1970 English publication. See \citep{palissmale1970stability}.

\subsection{GP-1966-03: Necessity of Morse--Smale conditions}\label{prob:GP-1966-03}
Must structural stability imply the Morse--Smale conditions?
\par\noindent\textit{Occurrences:} \hyperref[item:1966-94.9]{1966-94.9}.
\par\noindent\textit{Status (2026-11):} No; hyperbolic toral automorphisms and their suspension flows provide counterexamples. See \citep{smale1967dynamical}.

\subsection{GP-1966-04: Density of Morse--Smale systems}\label{prob:GP-1966-04}
Are Morse--Smale systems dense on every closed manifold?
\par\noindent\textit{Occurrences:} \hyperref[item:1966-94.10]{1966-94.10}.
\par\noindent\textit{Status (2026-11):} No in general; an open family of Anosov systems supplies an obstruction. See \citep{smale1967dynamical}.

\subsection{GP-1966-05: Density of structurally stable fields}\label{prob:GP-1966-05}
Are structurally stable vector fields dense on every compact manifold?
\par\noindent\textit{Occurrences:} \hyperref[item:1966-94.11]{1966-94.11}.
\par\noindent\textit{Status (2026-11):} No in general, by Smale’s 1966 counterexample, already published at the time of the question. See \citep{smale1966nondensity}.

\subsection{GP-1966-06: Density of Axiom A}\label{prob:GP-1966-06}
Are Axiom A diffeomorphisms \(C^1\)-dense on every closed manifold?
\par\noindent\textit{Occurrences:} \hyperref[item:1966-94.16]{1966-94.16}.
\par\noindent\textit{Status (2026-11):} No in general; Abraham--Smale, 1970. See \citep{abrahamsmale1970nongenericity}.

\subsection{GP-1966-07: Periodic density from hyperbolic nonwandering set}\label{prob:GP-1966-07}
Does hyperbolicity of \(\Omega(f)\) force \(\Omega(f)=\overline{\operatorname{Per}(f)}\)?
\par\noindent\textit{Occurrences:} \hyperref[item:1966-94.17]{1966-94.17}.
\par\noindent\textit{Status (2026-11):} No in general; Kurata, 1979. See \citep{kurata1979hyperbolic}.

\subsection{GP-1966-08: Uniqueness of the vanishing identity-viscosity limit}\label{prob:GP-1966-08}
Under what hypotheses do the solutions of \(u_t+A(u)u_x=\epsilon u_{xx}\) have a unique limit as \(\epsilon\downarrow0\)?
\par\noindent\textit{Occurrences:} \hyperref[item:1966-96.6]{1966-96.6}.
\par\noindent\textit{Status (2026-11):} As of November 2026: Bianchini--Bressan solve the smooth strictly hyperbolic small-total-variation case. The notes do not determine a universal large-data conjecture. See \citep{bianchinibressan2005viscosity}.

\subsection{GP-1966-09: Representative pair observables: moment hypotheses and concentration}\label{prob:GP-1966-09}
A precise question extracted from the representativeness argument is which moment and mixing hypotheses justify concentration of \(N^{-1}\sum_{i\ne j}f(x_i-x_j)\), and whether an arbitrary real even \(L^1\) kernel permits the claimed finite linear variance. Minlos’s later Problem 3 asks for quantitative inequalities in this argument. This entry does not index an unspecified all-density existence conjecture.
\par\noindent\textit{Occurrences:} \hyperref[item:1966-100.11]{1966-100.11}.
\par\noindent\textit{Status (2026-11):} As of November 2026: the finite-variance claim under bare \(L^1\) is false even for an ideal gas. The chapter supplies a complete conditional correction with an explicit four-point bound, an \(L^1\) truncation conclusion, and counterexamples. These are checked critical annotations, not a claim that a previously open C/D problem has been solved. Establishing the assumed correlation bound in an interacting model remains model- and regime-dependent. See \citep{minlos1968lectures,pulvirentitsagkarogiannis2015finite,kunatsagkarogiannis2018density}.

\subsection{GP-1966-10: Realizability of a full correlation hierarchy}\label{prob:GP-1966-10}
Characterize the symmetric locally finite factorial moment measures of probability laws on locally finite configurations, distinguishing existence from uniqueness.
\par\noindent\textit{Occurrences:} \hyperref[item:1966-106.6]{1966-106.6}.
\par\noindent\textit{Status (2026-11):} As of 9 November 2026: Lenard's symmetry, positivity and normalization characterization supplies existence; determinacy requires additional conditions. See \citep{lenard1975characterization,lenard1973correlation}.

\subsection{GP-1966-11: Helgason spherical support problem on a smooth Hadamard manifold}\label{prob:GP-1966-11}
For complete simply connected negatively curved smooth \(M\), do zero integrals of \(f\in C_c^\infty(M)\) over all geodesic spheres enclosing a closed ball \(B\) imply \(\operatorname{supp}f\subset B\)?
\par\noindent\textit{Occurrences:} \hyperref[item:1966-114.2]{1966-114.2}.
\par\noindent\textit{Status (2026-11):} As of 9 November 2026: verified for real-analytic metrics; a general smooth-metric resolution was not found. The precise source is Quinto's Problem 4.1, not a ray-transform support theorem. See \citep{quinto2008helgason}.

\subsection{GP-1966-12: Entire convex Monge--Ampere rigidity}\label{prob:GP-1966-12}
Must every convex entire solution of \(\det D^2u=1\) be quadratic?
\par\noindent\textit{Occurrences:} \hyperref[item:1966-114.4]{1966-114.4}.
\par\noindent\textit{Status (2026-11):} As of 9 November 2026: yes in all dimensions, by the J\"orgens--Calabi--Pogorelov theorem; convex viscosity solutions are included. See \citep{caffarellili2003extension}.

\subsection{GP-1966-13: Entire minimal-graph Bernstein problem}\label{prob:GP-1966-13}
For which \(n\) is every smooth entire minimal graph over \(\mathbb R^n\) an affine hyperplane?
\par\noindent\textit{Occurrences:} \hyperref[item:1966-114.6]{1966-114.6}.
\par\noindent\textit{Status (2026-11):} As of 9 November 2026: exactly \(n\le7\); nonaffine entire minimal graphs exist for \(n\ge8\). See \citep{bombieridegiorgigiusti1969minimal}.

\subsection{GP-1966-14: Metric rigidity of a Sobolev composition isometry}\label{prob:GP-1966-14}
If a smooth diffeomorphism of closed Riemannian manifolds induces an isometry of their full \(H^1\) inner products by pullback, must it preserve their metrics?
\par\noindent\textit{Occurrences:} \hyperref[item:1966-114.7]{1966-114.7}.
\par\noindent\textit{Status (2026-11):} As of 9 November 2026: yes, by volume and Dirichlet-form recovery; the composition condition is essential. See \citep{arendt2017diffusion}.

\subsection{GP-1966-15: Infinite volume of a complete noncompact positive-Ricci three-manifold}\label{prob:GP-1966-15}
Does a complete connected noncompact Riemannian three-manifold with positive Ricci curvature have infinite volume?
\par\noindent\textit{Occurrences:} \hyperref[item:1966-114.11]{1966-114.11}.
\par\noindent\textit{Status (2026-11):} As of 9 November 2026: yes; Yau's theorem covers nonnegative Ricci curvature in every dimension. See \citep{yau1976function}.

\subsection{GP-1966-16: Diffeomorphism type of equal-period closed-geodesic manifolds}\label{prob:GP-1966-16}
Must a closed simply connected Riemannian manifold whose prime unit-speed geodesics all have the same period be diffeomorphic to a compact rank-one symmetric space?
\par\noindent\textit{Occurrences:} \hyperref[item:1966-114.13]{1966-114.13}.
\par\noindent\textit{Status (2026-11):} As of 9 November 2026: cohomology is constrained to the rank-one symmetric models; the unrestricted smooth classification was not found resolved. The result of \citet[definition on p.~1 and Main Theorem on p.~3]{lisong2026zoll} assumes simple closed prime geodesics of a common period, as well as an entire Grauert tube, and proves metric rigidity only for the complex-projective cohomology type. Those additional hypotheses do not remove the gap for the notebook's question. See \citep{lisong2026zoll}.

\subsection{GP-1966-17: Existence of contact structures in dimension three}\label{prob:GP-1966-17}
Which closed smooth three-manifolds admit a contact plane field?
\par\noindent\textit{Occurrences:} \hyperref[item:1966-114.16]{1966-114.16}.
\par\noindent\textit{Status (2026-11):} As of 9 November 2026: exactly the orientable ones: necessity follows from the contact volume orientation and sufficiency from Martinet's theorem. See \citep{martinet1971contact,eliashberg1992contact}.

\subsection{GP-1966-18: Invariant distributions for the Lorentz form}\label{prob:GP-1966-18}
For \(SO_0(2,1)\)-invariant \(T\in\mathcal D'(\mathbb R^3)\) and nonzero \(P\in\mathbb C[t]\), does \(P(\partial_x^2+\partial_y^2-\partial_z^2)T=0\) force \(T\in L^1_{\mathrm{loc}}\)?
\par\noindent\textit{Occurrences:} \hyperref[item:1966-123.1]{1966-123.1}.
\par\noindent\textit{Status (2026-11):} Solved before this occurrence: the affirmative answer is the rank-one case of \citet[Theorem~1]{harishchandra1965algebra}. Checked as of 9 November 2026. No identical precise question was located in the supplied earlier index. See \citep{harishchandra1965algebra}.

\subsection{GP-1966-19: Geometric orbit quotients for affine reductive actions}\label{prob:GP-1966-19}
For an affine variety over an algebraically closed field of characteristic zero with a reductive group action, when does the invariant-theoretic categorical quotient have individual orbits as its fibers?
\par\noindent\textit{Occurrences:} \hyperref[item:1966-125.2]{1966-125.2}.
\par\noindent\textit{Status (2026-11):} Solved before the March 1967 report: exactly when all orbits are closed. The modern verification is \citet[Theorem~2.2.1 and Definition~2.1.4]{berczi2016quotients}, in the framework of \citet{mumford1965invariant}. The broader notebook question is recorded with its quotient-category ambiguity. Status checked 9 November 2026.

\subsection{GP-1966-20: Global boundary completion with trivial auxiliary bundles}\label{prob:GP-1966-20}
For the admissible continuous normal Fredholm family of an elliptic convolution system over \(S^*\Gamma\), when can finitely many globally defined potential terms and boundary equations, with trivial auxiliary bundles, produce an invertible completed normal family?
\par\noindent\textit{Occurrences:} \hyperref[item:1966-127.7]{1966-127.7}.
\par\noindent\textit{Status (2026-11):} Solved in 1967 under the Vishik--Eskin regularity and noncritical-order hypotheses: the family index must be the class of a signed trivial bundle, and this condition is sufficient with \(E=0\); see \citet[\S3, Theorems~1.3 and 2.3]{vishikeskin1967normally}. Received 27 April 1967; earlier circulation beyond the present seminar evidence was not established. Status checked 9 November 2026. The overwritten connective in the notebook remains uncertain. See \citep{vishikeskin1967normally}.

\section{Corrections to the transcription and translation}\label{corr:1966}
The evidence distinguishes a transcription correction from a mathematical emendation of wording faithfully preserved in the notebook.
\subsection*{P000140 (both; faithful error): Hypersurface exponent}
\noindent\textit{Transcription reads.} \(3k-1\) (provisional)\par
\noindent\textit{Critical reading.} Retain \(3k-1\) diplomatically; annotate the intended \(6k-1\).\par
\noindent\textit{Evidence.} The circled exponent in the scan begins with 3. The Brieskorn graph criterion rules out odd \(k\) for the written family; the published family has \(6k-1\).\par
\subsection*{P000142 (both; faithful error): Cohomology reduction}
\noindent\textit{Transcription reads.} Vanishing on \(\Gamma/\Gamma_p\) finishes the proof.\par
\noindent\textit{Critical reading.} Add control of \(H^1(\Gamma_p,L/pL)^{\Gamma/\Gamma_p}\).\par
\noindent\textit{Evidence.} The displayed finite quotient does not by itself exhaust the inflation--restriction sequence; this is a missing proof step, not a changed glyph.\par
\subsection*{P000143 (both; faithful error): Equivalence in structural stability}
\noindent\textit{Transcription reads.} Diffeomorphic equivalence\par
\noindent\textit{Critical reading.} Topological conjugacy for maps; topological orbit equivalence for flows.\par
\noindent\textit{Evidence.} The notebook later explicitly specifies an \(\epsilon\)-homeomorphism on P000144. Smooth equivalence would preserve derivatives at periodic points and is not the intended stable notion.\par
\subsection*{P000145 (both; ambiguity): Name and year beside generic density}
\noindent\textit{Transcription reads.} A name interpreted as Goto or Pugh, 1964\par
\noindent\textit{Critical reading.} Keep the name/date uncertain; compare Pugh’s 1967 theorem without rewriting the scan.\par
\noindent\textit{Evidence.} The scan does not securely resolve the attribution. The later theorem is a mathematical match, not palaeographic proof.\par
\subsection*{P000146 (both; faithful error): Ambient uniform algebra}
\noindent\textit{Transcription reads.} \(C(S)\), all strongly continuous functions\par
\noindent\textit{Critical reading.} Use \(C_b(S)\) when referring to a uniform-norm Banach algebra.\par
\noindent\textit{Evidence.} Norm continuity on a closed infinite-dimensional ball does not imply boundedness; the generated regular algebra is bounded.\par
\subsection*{P000146 (both; faithful error): Dirichlet boundary described as weakly closed}
\noindent\textit{Transcription reads.} \(\Sigma\): closed surface (weak)\par
\noindent\textit{Critical reading.} Preserve the word; annotate that norm-sphere boundaries are not weakly closed.\par
\noindent\textit{Evidence.} The scan has the parenthetical \ru{слаб}. An orthonormal sequence on the unit sphere converges weakly to 0, so literal weak closedness is incompatible with the basic spherical example.\par
\subsection*{P000147 (both; faithful error): Scalar uniqueness}
\noindent\textit{Transcription reads.} Always unique if the scalar flux is convex.\par
\noindent\textit{Critical reading.} Unique in the entropy/admissible class; not among all weak solutions.\par
\noindent\textit{Evidence.} The statement is faithfully present in the scan. The Burgers expansion shock and rarefaction with the same Riemann data distinguish the two readings.\par
\subsection*{P000148 (both; faithful error): Associated-bundle arrow}
\noindent\textit{Transcription reads.} \ru{Вложение} \(R_G\to K(B_G)\)\par
\noindent\textit{Critical reading.} A homomorphism, inducing an isomorphism after completion.\par
\noindent\textit{Evidence.} The word is clear in the scan. The explicit nonzero element in \(\bigcap I^n\subset R(C_6)\) shows that injectivity is not valid for all compact groups.\par
\subsection*{P000148 (both; faithful error): Canonical bundle on projective space}
\noindent\textit{Transcription reads.} Fibre \(\mathbb{CP}^1\), total space \(\mathbb{CP}^{n+1}\)\par
\noindent\textit{Critical reading.} Tautological line fibre \(\mathbb C\); unit-circle total space \(S^{2n+1}\).\par
\noindent\textit{Evidence.} The projective-space symbols are visible in the scan. This is a mathematical slip in the source, not an OCR substitution.\par
\subsection*{P000150 (both; faithful error): Lower-right corner of the comparison square}
\noindent\textit{Transcription reads.} \(K_T(X)\)\par
\noindent\textit{Critical reading.} \(K(X_T)\) in the lower Borel row; \(K_T(X)\) identifies with the upper-right equivariant term.\par
\noindent\textit{Evidence.} The scan really labels the lower-right entry \(K_T(X)\). With that label the right vertical arrow is not the intended completion comparison. The corrected label is a mathematical emendation.\par
\subsection*{P000155 (both; ambiguity): Effective-volume quantifier}
\noindent\textit{Transcription reads.} For each \(\tau_k\), there exists \(v_N\); parenthesis doubtful.\par
\noindent\textit{Critical reading.} Retain dependence on \(k\): \(v_{N,k}\), not one common adjustment for all orders.\par
\noindent\textit{Evidence.} The scan visibly ends the parenthesis with \ru{и от k}. The ideal-gas factorial densities independently confirm why this quantifier matters.\par
\subsection*{P000156 (both; faithful error): Perturbative parameter}
\noindent\textit{Transcription reads.} Proved only for small \(v\) and \(\beta\).\par
\noindent\textit{Critical reading.} For the cited gas regime, small density/activity, hence large \(v=1/\rho\).\par
\noindent\textit{Evidence.} The source really has \ru{малых v}. Its own definition \(v=1/\rho\) on the preceding page and the cited gas-phase results show the inversion.\par
\subsection*{P000156 (both; ambiguity): Last mean formula}
\noindent\textit{Transcription reads.} \(\bar\tau_k\) with uncertain subscript\par
\noindent\textit{Critical reading.} \(\bar\tau_2\) in the pair-observable formula.\par
\noindent\textit{Evidence.} The last glyph is overwritten and resembles k; all preceding three integrals on this page have the second correlation function. This is an interpretive emendation, not a confident new glyph reading.\par
\subsection*{P000157 (both; faithful error): Variance conclusion under an L1 hypothesis}
\noindent\textit{Transcription reads.} \(DF_N\sim dN\) from the assumptions displayed.\par
\noindent\textit{Critical reading.} Add second-moment and integrable finite-volume clustering hypotheses for an \(O(N)\) bound; an asymptotic equivalent requires further convergence.\par
\noindent\textit{Evidence.} The scan’s leading variance term is \(2\iint f^2\tau_2\). An \(L^1\) pair kernel need not make it finite. The 1968 account restores integrable decay, but this separate short-distance defect remains.\par
\subsection*{P000158 (both; provenance): Lecturer and collaborator}
\noindent\textit{Transcription reads.} Prof. Lodzhs. [?], with P. Szego\par
\noindent\textit{Critical reading.} Identify the lecturer editorially as Lee Lorch and the collaborator as Peter Szego; preserve uncertainty in the diplomatic name.\par
\noindent\textit{Evidence.} The catalogue has Lee Lorch; the joint 1963 paper and Muldoon’s 1968 references identify Peter Szego. Neither supports substituting Gabor Szego.\par
\subsection*{P000158 (both; faithful error): Sturm converse}
\noindent\textit{Transcription reads.} Increasing coefficient gives closer zeros, and conversely.\par
\noindent\textit{Critical reading.} Keep the words diplomatically; delete the unrestricted converse from a theorem statement.\par
\noindent\textit{Evidence.} The explicit positive-coefficient counterexample in item 1966-102.1 has decreasing zero gaps but a nonmonotone coefficient.\par
\subsection*{P000158 (both; faithful error): Strict higher differences}
\noindent\textit{Transcription reads.} Weak derivative signs imply strict alternating differences.\par
\noindent\textit{Critical reading.} Add an actual nondegeneracy hypothesis; do not repair the formula solely by changing the index range.\par
\noindent\textit{Evidence.} The scan clearly has weak signs in the hypothesis and strict signs in the conclusion. The constant coefficient 1 meets all written assumptions but gives constant gaps.\par
\subsection*{P000158 (both; ambiguity): The endpoint lambda=-1}
\noindent\textit{Transcription reads.} Formal endpoint integral associated with reciprocal derivatives.\par
\noindent\textit{Critical reading.} Distinguish the divergent integral from a regularized endpoint limit; retain the unresolved abbreviation.\par
\noindent\textit{Evidence.} A simple zero makes the absolute-value integrand comparable to 1/|x-a|. The supplied limit has a sum of absolute reciprocal derivatives, not an unverified signed difference.\par
\subsection*{P000159 (both; faithful error): Ring over a subring}
\noindent\textit{Transcription reads.} An R(G)-module, a ring over a subring.\par
\noindent\textit{Critical reading.} An R(G)-algebra; the structure map need not be injective.\par
\noindent\textit{Evidence.} For a free action, representation bundles may become trivial or satisfy additional relations. A unital ring map alone does not identify its domain with a subring.\par
\subsection*{P000159-P000160 (both; faithful error): Kernel minus cokernel}
\noindent\textit{Transcription reads.} Finite kernel and cokernel give two bundles over the parameter space.\par
\noindent\textit{Critical reading.} Use a stabilized virtual index unless both dimensions are locally constant.\par
\noindent\textit{Evidence.} The scalar family \(A_t=t\) has jumping kernel and cokernel dimensions; the same problem occurs for general elliptic families.\par
\subsection*{P000160 (both; ambiguity): Projection argument}
\noindent\textit{Transcription reads.} The displayed homotopy already gives locally constant kernel dimension.\par
\noindent\textit{Critical reading.} Retain the ambiguous projection notation; add a global finite-dimensional complement and a surjective augmentation for a complete argument.\par
\noindent\textit{Evidence.} The scan does not state the requisite complement condition. Compactness of the base supplies a finite common complement; the reconstruction proves the resulting kernel-bundle assertion.\par
\subsection*{P000161 (both; transcription): opening analytic-index paragraph}
\noindent\textit{Transcription reads.} никаких связей с топологией / no connections with topology\par
\noindent\textit{Critical reading.} неясна связь с топологией / the connection with topology is unclear\par
\noindent\textit{Evidence.} Direct inspection shows the singular words неясна связь, not никаких связей.\par
\subsection*{P000161 (both; transcription): opening analytic-index paragraph}
\noindent\textit{Transcription reads.} главные члены / leading terms\par
\noindent\textit{Critical reading.} младшие члены / lower-order terms\par
\noindent\textit{Evidence.} The scan clearly shows the initial м and the word младшие; this is also the mathematically relevant order-perturbation issue.\par
\subsection*{P000162 (both; faithful error): sentence following (iii)}
\noindent\textit{Transcription reads.} At \(s=0\) this is not the identity operator\par
\noindent\textit{Critical reading.} \(A^0=I\); the continued diagonal density at zero is not restriction of its delta kernel\par
\noindent\textit{Evidence.} The sentence is present in the scan. The correction separates two noncommuting operations; it is a mathematical emendation.\par
\subsection*{P000162 (both; ambiguity): order-one reduction parenthesis}
\noindent\textit{Transcription reads.} Take m=1 (multiply by an invertible operator)\par
\noindent\textit{Critical reading.} An order reduction must specify what invariant it preserves; individual complex powers are not preserved by arbitrary multiplication\par
\noindent\textit{Evidence.} Faithfully written shorthand. Index preservation under suitable elliptic order reduction does not identify the original powers.\par
\subsection*{P000163 (both; ambiguity): uniform mixing gloss}
\noindent\textit{Transcription reads.} uniformly over all pairs U\par
\noindent\textit{Critical reading.} Keep the diplomatic wording; the later precise formula fixes one local event before taking the remote-event supremum\par
\noindent\textit{Evidence.} Scan confirms the strong gloss. Minlos1968 formula(2.5), pp151–152, has more careful quantifiers.\par
\subsection*{P000165 (both; ambiguity): mean-energy ratio}
\noindent\textit{Transcription reads.} Numerator has no 1/N! while the preceding Z includes it\par
\noindent\textit{Critical reading.} Use the same Gibbs reference-measure convention in numerator and denominator\par
\noindent\textit{Evidence.} The ratio is correct only with consistent normalization; an abbreviated da may conceal the quotient measure.\par
\subsection*{P000165 (both; ambiguity): pressure derivative}
\noindent\textit{Transcription reads.} p=∂\_V log Z\par
\noindent\textit{Critical reading.} Physical pressure is β\textasciicircum{}{-1}∂\_V log Z; the written quantity may be reduced pressure βp\par
\noindent\textit{Evidence.} The manuscript gives no unit convention; adjacent p=τ1 supports reduced-pressure interpretation.\par
\subsection*{P000165 (both; faithful error): entropy definition}
\noindent\textit{Transcription reads.} S=∫π logπ\par
\noindent\textit{Critical reading.} S=−∫π logπ for Gibbs entropy\par
\noindent\textit{Evidence.} The plus sign is visibly in the source. Substitution of π=Z\textasciicircum{}{-1}e\textasciicircum{}{-βH} checks the sign independently.\par
\subsection*{P000165 (both; faithful error): entropy per-particle identity}
\noindent\textit{Transcription reads.} S/N=(log Z)/N−β\textasciicircum{}{-1}∂\_β((log Z)/N)\par
\noindent\textit{Critical reading.} S/N=(log Z)/N−β∂\_β((log Z)/N)\par
\noindent\textit{Evidence.} The source visibly has 1/β. It contradicts the standard Gibbs identity even after choosing an entropy sign convention.\par
\subsection*{P000166 (both; faithful error): generating-function example}
\noindent\textit{Transcription reads.} graph p(q) Lagrangian iff p=∂S/∂q\par
\noindent\textit{Critical reading.} Locally iff; globally iff p·dq is exact\par
\noindent\textit{Evidence.} The scan has no local qualifier. A closed angular form on an annulus supplies the elementary obstruction.\par
\subsection*{P000166 (en; translation): word after map to the circle}
\noindent\textit{Transcription reads.} katsikl [?]\par
\noindent\textit{Critical reading.} cocycle / cohomology class [contextual reading]\par
\noindent\textit{Evidence.} The subsequent H\textasciicircum{}1(M,Z) and degree formula determine the mathematical meaning; the Russian letter-level uncertainty remains.\par
\subsection*{P000167 (both; provenance): top heading}
\noindent\textit{Transcription reads.} Minlos. Lecture5. Phase transitions\par
\noindent\textit{Critical reading.} Record this as an additional heading only, separate from Gelfand's duality and Nirenberg's heat-kernel report\par
\noindent\textit{Evidence.} The scan supplies no phase-transition lecture body before the next speaker heading; the MCCME list has no separately allocated record here.\par
\subsection*{P000167 (both; faithful error): finite-group multiplicity claim}
\noindent\textit{Transcription reads.} Tensor multiplicities recover a finite group\par
\noindent\textit{Critical reading.} Multiplicity-only interpretation is false; actual intertwiner and tensor data are needed\par
\noindent\textit{Evidence.} The sentence is overwritten, so its reading is qualified. D8 and Q8 have the same complex fusion rules but are nonisomorphic.\par
\subsection*{P000167 (both; faithful error): arbitrary matrices in duality equations}
\noindent\textit{Transcription reads.} Every family of arbitrary linear C\_i satisfying the equations equals T\_i(g)\par
\noindent\textit{Critical reading.} Require unit-normalized unitary tensor-natural automorphisms and compatible duals\par
\noindent\textit{Evidence.} The zero family satisfies the bare homogeneous equations; for S1, C\_k=z\textasciicircum{}k with nonunit |z| shows the missing reality condition.\par
\subsection*{P000167 (both; faithful error): heat-kernel exponent}
\noindent\textit{Transcription reads.} u=exp[−g/t+o(1/t)], g=d²\par
\noindent\textit{Critical reading.} g=d²/4 for the written generator ∑a\textasciicircum{}{ij}∂\_{ij}\par
\noindent\textit{Evidence.} The same scanned page gives the Euclidean Gaussian exp[−|x−y|²/(4t)], which fixes the discrepancy.\par
\subsection*{P000168 (both; provenance): Vishik Lecture1 and extra professor heading}
\noindent\textit{Transcription reads.} Kohn–Nirenberg definition; boundedness theorem; Prof.Lodzhs. on Sturm–Liouville\par
\noindent\textit{Critical reading.} Preserve as headings; do not reconstruct a complete theorem or assign the extra heading to Vishik\par
\noindent\textit{Evidence.} The source contains no symbol class or theorem statement; the professor's name is uncertain.\par
\subsection*{P000169 (both; transcription): last abbreviation}
\noindent\textit{Transcription reads.} отлип. / otlip.\par
\noindent\textit{Critical reading.} отлич. / differ\par
\noindent\textit{Evidence.} Direct scan inspection shows отлич.; the sentence describes strong differences between viscous and inviscid velocity in a narrow region.\par
\subsection*{P000170 (both; faithful error): initial and inflow profiles}
\noindent\textit{Transcription reads.} u|x=0=u0(x), u|t=0=u0(x,t)\par
\noindent\textit{Critical reading.} u|x=0=u1(t,y), u|t=0=u0(x,y), with corner compatibility\par
\noindent\textit{Evidence.} The variables are inconsistent in the scan itself. Oleinik1966 p751 gives the correctly indexed initial/inflow problem.\par
\subsection*{P000170 (both; ambiguity): crossed-out p\_y=0 line}
\noindent\textit{Transcription reads.} p\_y=0 is crossed out in the display\par
\noindent\textit{Critical reading.} Preserve cancellation but retain the following explicit statement p=p(t,x)\par
\noindent\textit{Evidence.} The prose immediately says pressure is constant along the normal; cancellation does not remove that model assumption.\par
\subsection*{P000170 (both; faithful error): unqualified existence and stability lines}
\noindent\textit{Transcription reads.} Unique if it exists; always stable\par
\noindent\textit{Critical reading.} These statements require the monotone solution class, compatibility, decay and specified stability norms\par
\noindent\textit{Evidence.} Oleinik1966 and Oleinik1968 Theorems4–5 supply omitted assumptions; the final abbreviation does not specify a norm.\par
\subsection*{P000171 (both; faithful error): minimal-surface equation}
\noindent\textit{Transcription reads.} (1+Σu\_i²)(Σu\_jj²)−Σu\_ij u\_i u\_j=0\par
\noindent\textit{Critical reading.} (1+Σu\_i²)Σu\_jj−Σu\_ij u\_i u\_j=0\par
\noindent\textit{Evidence.} The extra raised square is visible in the uncancelled manuscript line. Deleting it is required by the note's identification with minimal graphs.\par
\subsection*{P000171 (both; faithful error): Bernstein question parenthesis}
\noindent\textit{Transcription reads.} Done (reference in De Giorgi)\par
\noindent\textit{Critical reading.} Affirmative only in restricted dimensions; modern sharp answer is yes n≤7, no n≥8\par
\noindent\textit{Evidence.} The note supplies no dimension restriction. Bombieri–De Giorgi–Giusti1969 resolves the higher-dimensional negative cases.\par
\subsection*{P000171-P000172 (both; provenance): attribution of geometry problems}
\noindent\textit{Transcription reads.} Problems placed after Agranovich's PDO Lecture3 heading\par
\noindent\textit{Critical reading.} Associate by manuscript position but keep the problem-list speaker unattributed\par
\noindent\textit{Evidence.} No speaker name introduces the geometry list and no PDO lecture text appears; page overlap is not evidence that every problem was Agranovich's.\par
\subsection*{P000172 (both; ambiguity): problem7 geodesics}
\noindent\textit{Transcription reads.} intersect. geodesics / пересек.\par
\noindent\textit{Critical reading.} Retain uncertain intersecting-geodesic wording; do not silently emend to closed geodesics\par
\noindent\textit{Evidence.} The scan has пересек. after a cancelled word and no legible замкнутых in this problem; problem8 does explicitly write geodesics closed.\par
\subsection*{P000172 (both; ambiguity): number following7}
\noindent\textit{Transcription reads.} 3.[?]\par
\noindent\textit{Critical reading.} 8 [sequential numbering, manuscript glyph ambiguous]\par
\noindent\textit{Evidence.} Its position between7 and9 identifies list position; the written glyph itself remains 3/8-like.\par
\subsection*{P000172 (both; ambiguity): symmetric-space type in problem8}
\noindent\textit{Transcription reads.} of the1st type; cancelled rank-like expression\par
\noindent\textit{Critical reading.} Rank-one symmetric space [mathematical reconstruction, not a certain diplomatic reading]\par
\noindent\textit{Evidence.} The surrounding Bott cohomology and sphere examples support the CROSS identification, but the cancelled word cannot be recovered with certainty.\par
\subsection*{P000172 (both; faithful error): contact existence question}
\noindent\textit{Transcription reads.} every compact M3\par
\noindent\textit{Critical reading.} every closed orientable smooth M3; without orientability the answer is no\par
\noindent\textit{Evidence.} Contact local forms induce a consistent orientation in dimension3; Martinet proves sufficiency for orientable closed manifolds.\par
\subsection*{P000173 (both; ambiguity): upper-page stability examples}
\noindent\textit{Transcription reads.} Unclear Rangmax\par
\noindent\textit{Critical reading.} \ru{Ранг} max (maximal rank)\par
\noindent\textit{Evidence.} Visual inspection of the handwritten Russian resolves the word; the next Morse formula is consistent.\par
\subsection*{P000174 (both; faithful error): W1--W3 global surface stability list}
\noindent\textit{Transcription reads.} Fold images transverse; cusp images distinct\par
\noindent\textit{Critical reading.} Require only ordinary transverse double fold crossings; cusp images must meet no other critical image\par
\noindent\textit{Evidence.} The scan faithfully has the shorter list. Triple crossings and cusp--fold coincidences are not stable multigerms; additional conditions belong to the mathematical statement.\par
\subsection*{P000175 (both; ambiguity): trefoil caption}
\noindent\textit{Transcription reads.} Trefoil knot\par
\noindent\textit{Critical reading.} Retain caption, but no knot type asserted without a specified spatial lift\par
\noindent\textit{Evidence.} The scan shows planar curves and the abbreviation \ru{трилистн. узел}; it specifies no embedding into a three-dimensional space.\par
\subsection*{P000175 (both; ambiguity): threshold 12}
\noindent\textit{Transcription reads.} Beginning with 12, infinitely many invariants\par
\noindent\textit{Critical reading.} Retain wording; number/type of moduli is unspecified\par
\noindent\textit{Evidence.} The scan confirms the threshold but supplies no quantifiers distinguishing infinitely many types, scalar parameters, and functional invariants.\par
\subsection*{P000176 (both; faithful error): scalar jet algebra}
\noindent\textit{Transcription reads.} \(\mathbb R[x_1,\ldots,x_n]/M^k\)\par
\noindent\textit{Critical reading.} \(\mathbb R[x_1,\ldots,x_n]/M^{k+1}\)\par
\noindent\textit{Evidence.} The scan has \(k\), without \(+1\). A first jet must retain linear terms; Boardman 1967 §1 uses the usual \(k\)-jet indexing.\par
\subsection*{P000176 (both; faithful error): algebraic structure of jets}
\noindent\textit{Transcription reads.} \(J^k(M,N)\) is an algebraic variety\par
\noindent\textit{Critical reading.} Polynomial algebraic coordinates exist in each jet fibre; an arbitrary smooth global bundle need not be algebraic\par
\noindent\textit{Evidence.} Smooth changes of base coordinates are not assumed algebraic. This qualifies the notebook claim, not the transcription.\par
\subsection*{P000176 (both; faithful error): codimension and equations}
\noindent\textit{Transcription reads.} Codimension \(e\) means \(e\) equations\par
\noindent\textit{Critical reading.} Locally \(e\) independent equations define a smooth codimension-\(e\) locus\par
\noindent\textit{Evidence.} The number of equations alone does not imply codimension or smoothness.\par
\subsection*{P000177 (both; faithful error): transversality to a rank stratum}
\noindent\textit{Transcription reads.} \(f\pitchfork\Sigma\)\par
\noindent\textit{Critical reading.} \(j^1f\pitchfork\Sigma\)\par
\noindent\textit{Evidence.} The stratum is in the jet bundle rather than the target; the scan uses the abbreviated map symbol.\par
\subsection*{P000178 (both; faithful error): Boardman universal loci}
\noindent\textit{Transcription reads.} Algebraic subvarieties\par
\noindent\textit{Critical reading.} Smooth locally closed Boardman strata in local polynomial jet coordinates\par
\noindent\textit{Evidence.} Exact-rank loci are not closed; Boardman 1967 Theorem 6.1 uses locally closed submanifolds.\par
\subsection*{P000180 (both; faithful error): critical Jacobian extension}
\noindent\textit{Transcription reads.} \(k=\max\) when the ideal is not the whole \(B\)\par
\noindent\textit{Critical reading.} \(k=\min\) with \(\Delta_k(B)\ne A\)\par
\noindent\textit{Evidence.} The descending chain and the scan display \(\Delta_{k-1}B=A\), \(\Delta_kB\subsetneq A\). Both \(\max\) and \(B\) conflict with that display.\par
\subsection*{P000180 (both; ambiguity): cubic ideal example}
\noindent\textit{Transcription reads.} Crowded superscript on the iteration ending in \((x)\)\par
\noindent\textit{Critical reading.} Successive first extensions: \((x^3)\), \((x^2)\), \((x)\), \(A\)\par
\noindent\textit{Evidence.} The written exponents are ambiguous; differentiation fixes the mathematical sequence without claiming a definitive diplomatic superscript reading.\par
\subsection*{P000180 (both; faithful error): Boardman codimension coefficient}
\noindent\textit{Transcription reads.} \(m-n+i_1\)\par
\noindent\textit{Critical reading.} \(n-m+i_1\)\par
\noindent\textit{Evidence.} The scan confirms \(m-n\) while source is \(M^m\) and target \(N^n\). Boardman 1967 Definition 6.5 has target minus source; the rank-stratum case independently forces the sign.\par
\subsection*{P000181 (both; faithful error): definition of \(\mu\)}
\noindent\textit{Transcription reads.} All nonincreasing sequences \(0\le j_s\le i_s\)\par
\noindent\textit{Critical reading.} Exclude the all-zero sequence: \(j_1>0\)\par
\noindent\textit{Evidence.} The scan omits the exclusion. Boardman 1967 Definition 6.4 explicitly includes \(j_1>0\); otherwise even the length-one formula is wrong.\par
\subsection*{P000181 (both; ambiguity): dimension 16 in the informal table}
\noindent\textit{Transcription reads.} 16 and an isolated 4 / four types\par
\noindent\textit{Critical reading.} Read tentatively as first occurrence of corank four, \(c(4)=16\)\par
\noindent\textit{Evidence.} The scan does not support a confident count of four smooth equivalence classes. The codimension computation supports corank four but remains an editorial reading.\par
\subsection*{P000182 (both; ambiguity): Morin normal form index}
\noindent\textit{Transcription reads.} Ambient \(n\) and length \(k\) of \(1\)-string are intermingled\par
\noindent\textit{Critical reading.} For \(1\le k\le n\) use \(t^{k+1}+\sum_{j=1}^{k-1}u_jt^j\); the written full formula is \(k=n\)\par
\noindent\textit{Evidence.} The scan and following examples distinguish a Boardman symbol from ambient dimension; Arnold 1968 §2 supplies the clarified form.\par
\subsection*{P000182 (both; faithful error): composition with the diagonal}
\noindent\textit{Transcription reads.} \(\Delta\circ\bar\psi=\psi\)\par
\noindent\textit{Critical reading.} \(\bar\psi\circ\Delta=\psi\)\par
\noindent\textit{Evidence.} The written arrows give \(\Delta\) from the vector space to its symmetric square and \(\bar\psi\) onward. Standard composition order reverses the displayed expression.\par
\subsection*{P000183 (both; faithful error): seven orbits / preceding corank-two language}
\noindent\textit{Transcription reads.} Seven orbits of \(\Sigma^{(2,0)}\)\par
\noindent\textit{Critical reading.} Seven linear orbits of binary quadratic pencils, not of all smooth germs\par
\noindent\textit{Evidence.} The seven-way picture classifies a plane, line, or zero image in the three-dimensional space of quadratic forms. Higher-order germs can share a pencil.\par
\subsection*{P000184 (both; faithful error): all \(\Sigma^{(2,0)}\) singularities stable}
\noindent\textit{Transcription reads.} All singularities stable\par
\noindent\textit{Critical reading.} Stable or suitably generic corank-two germs admit the cited algebra classification; arbitrary symbol membership does not imply stability\par
\noindent\textit{Evidence.} The scan is unqualified. The germ \((x^2+y^2,xy):\mathbb R^2\to\mathbb R^2\) has symbol \((2,0)\) and is not stable, as checked by the codimension-four rank stratum.\par
\subsection*{P000184 (both; ambiguity): dimension diagram}
\noindent\textit{Transcription reads.} Cells \((9,9)\) and \((6,9)\)\par
\noindent\textit{Critical reading.} Retain diagram and do not infer a missing coordinate convention or inequality\par
\noindent\textit{Evidence.} Visual inspection does not fix every cell label and boundary; the later nice-dimensions theorem is cited separately.\par
\subsection*{P000185 (both; provenance): map on local functions}
\noindent\textit{Transcription reads.} \(f_*\)\par
\noindent\textit{Critical reading.} \(f^*\) (pullback of functions)\par
\noindent\textit{Evidence.} This is a notation clarification: composition with \(f\) goes from target functions to source functions. The scan faithfully uses the star below.\par
\subsection*{P000185 (both; faithful error): dimension of \(Q\) equals multiplicity}
\noindent\textit{Transcription reads.} Analytic maps \(\mathbb C^n\to\mathbb C^m\); dimension equals number of inverse images\par
\noindent\textit{Critical reading.} For the direct nearby-value count require a finite holomorphic germ of equal source and target dimension\par
\noindent\textit{Evidence.} The note says apparently so and does not state these restrictions; real or unequal-dimensional examples require a different interpretation.\par
\subsection*{P000186 (both; provenance): preparation theorem heading}
\noindent\textit{Transcription reads.} Weierstrass preparation for smooth germs\par
\noindent\textit{Critical reading.} Malgrange smooth preparation / its finite-module consequence\par
\noindent\textit{Evidence.} The stated category is \(C^\infty\); analytic Weierstrass alone does not give that category. Malgrange 1966, chapter V, and Arnold 1968 §4 identify the result.\par
\subsection*{P000186 (both; ambiguity): first local-algebra family}
\noindent\textit{Transcription reads.} \((x,y,x^a+y^b)\)\par
\noindent\textit{Critical reading.} \((xy,x^a+y^b)\)\par
\noindent\textit{Evidence.} The scan has a mark between \(x\) and \(y\) that can resemble a comma; taking two separate generators would collapse the algebra to \(\mathbb R\). The checked final Mather table confirms \(xy\).\par
\subsection*{P000186 (both; faithful error): elliptic three-generator family}
\noindent\textit{Transcription reads.} \((x^a,xy^{a-1},x^2+y^2)\)\par
\noindent\textit{Critical reading.} \((x^a,x^{a-1}y,x^2+y^2)\)\par
\noindent\textit{Evidence.} The scan confirms the mixed exponent on \(y\). Mather 1971 §7 gives \(x^{a-1}y\). At \(a=3\) the literal algebra has length 6, while the corrected family has length 5; two independent exact algorithms and a normal-form argument agree.\par
\subsection*{P000186 (both; provenance): names IV and V}
\noindent\textit{Transcription reads.} IV has three generators; V has two\par
\noindent\textit{Critical reading.} Record notebook labels, but use final Mather convention IV for two generators and V for three\par
\noindent\textit{Evidence.} The scan and the published final table differ in labels. This is a convention/provenance discrepancy rather than a false theorem by itself.\par
\subsection*{P000186 (both; ambiguity): left margin of algebra list}
\noindent\textit{Transcription reads.} \(n\ge m1.;\ n>m3.;\ldots\)\par
\noindent\textit{Critical reading.} Dimension conditions followed by row numbers \(1,3,\ldots\)\par
\noindent\textit{Evidence.} The visual layout separates row enumerators from the preceding inequalities. Do not read a spurious \(n\ge m+1\) into concatenated text.\par
\subsection*{P000187 (en; translation): seminar plan, index lectures}
\noindent\textit{Transcription reads.} Dynkin\par
\noindent\textit{Critical reading.} Dynin\par
\noindent\textit{Evidence.} The Russian transcription and scan read \ru{Дынин}, with no \ru{к}; this refers to the index lecturer named in the plan.\par
\subsection*{P000187 (both; provenance): Gelfand reference beneath cyclic operators}
\noindent\textit{Transcription reads.} see Gelfand, 1947\par
\noindent\textit{Critical reading.} Keep 1947 diplomatically; identify the matching published paper as Gelfand 1950, received 3 July 1948.\par
\noindent\textit{Evidence.} The scan has 1947. Math-Net and the opened original paper establish publication and receipt dates, but do not rule out earlier circulation.\par
\subsection*{P000188 (both; ambiguity): characters at infinity}
\noindent\textit{Transcription reads.} Characters become 1 at infinity\par
\noindent\textit{Critical reading.} No emendation supported; retain as unidentified.\par
\noindent\textit{Evidence.} Scan inspection confirms the compressed wording, the numeral 1, and the infinity sign. Neither the character normalization nor the limiting variable is stated.\par
\subsection*{P000189 (both; faithful error): radial operator algebra}
\noindent\textit{Transcription reads.} finite codimension in the ring of all differential operators\par
\noindent\textit{Critical reading.} Finite-module/invariant-polynomial statement on a Cartan slice; not finite vector-space codimension in all differential operators.\par
\noindent\textit{Evidence.} The same assertion is visible in the scan. Arbitrary variable coefficients and arbitrary orders preclude the literal finite-codimension assertion.\par
\subsection*{P000189 (both; faithful error): elliptic sum and annihilator}
\noindent\textit{Transcription reads.} Delta=sum D\_i\textasciicircum{}2 is elliptic; P(Delta) is in J; thus Delta T=0\par
\noindent\textit{Critical reading.} Use an elliptic consequence of the full joint system; conclude P(Delta)T=0, or explicitly rename P(Delta).\par
\noindent\textit{Evidence.} The scan repeats the same Delta. Unequal generator orders make the displayed sum-of-squares ellipticity claim insufficient.\par
\subsection*{P000189 (both; faithful error): three cases for group elements}
\noindent\textit{Transcription reads.} nilpotent group element\par
\noindent\textit{Critical reading.} Unipotent group element, or nilpotent element after passage to the Lie algebra.\par
\noindent\textit{Evidence.} The scan says nilpotent; this is a terminology emendation, not a transcription correction.\par
\subsection*{P000190 (both; faithful error): direction of transfer}
\noindent\textit{Transcription reads.} generalized functions from M to N\par
\noindent\textit{Critical reading.} Test functions are integrated from M to N; dual distributions pull back from N to M.\par
\noindent\textit{Evidence.} The defining integral and the explicit reverse distribution arrow on P000196 determine the direction.\par
\subsection*{P000192 (both; faithful error): bundle of differential forms}
\noindent\textit{Transcription reads.} Lambda T(X)\par
\noindent\textit{Critical reading.} Lambda T\textasciicircum{}*(X), or explicitly use the Riemannian metric to identify it with Lambda T(X).\par
\noindent\textit{Evidence.} The scan and transcription omit the dual; the exterior differential naturally acts on covariant differential forms.\par
\subsection*{P000192 (both; ambiguity): name of d+delta}
\noindent\textit{Transcription reads.} Cauchy–Riemann operator\par
\noindent\textit{Critical reading.} Hodge–de Rham operator (historical terminology retained in notes).\par
\noindent\textit{Evidence.} The displayed square and all-degree form domain identify d+d\textasciicircum{}*. The name in the scan is faithfully translated.\par
\subsection*{P000193 (both; faithful error): centralizer definition}
\noindent\textit{Transcription reads.} ad\_a X=X\par
\noindent\textit{Critical reading.} Ad(a)X=X for a group element a.\par
\noindent\textit{Evidence.} The scan has the lower-case notation. The earlier page explicitly defines Ad\_a(z)=aza\textasciicircum{}{-1}.\par
\subsection*{P000193 (both; faithful error): section through conjugacy classes}
\noindent\textit{Transcription reads.} section meeting each conjugacy class in one element\par
\noindent\textit{Critical reading.} A local centralizer descent submanifold, generally redundant; uniqueness requires an additional local quotient and fails globally.\par
\noindent\textit{Evidence.} The wording is in the scan. The submersion construction and ordinary Weyl conjugacy demonstrate the omitted qualification.\par
\subsection*{P000193-P000194 (both; ambiguity): pointwise operator and generated transverse space}
\noindent\textit{Transcription reads.} D\_y in g; q and z generate subalgebras of the enveloping algebra\par
\noindent\textit{Critical reading.} D\_y in U(g); use the symmetrized PBW space S(q) and U(z).\par
\noindent\textit{Evidence.} P000194 itself labels g as the enveloping algebra. An arbitrary invariant complement q is not a Lie subalgebra, so its associative generated algebra cannot serve as the stated PBW factor.\par
\subsection*{P000196 (both; faithful error): first displayed map}
\noindent\textit{Transcription reads.} C\_c\textasciicircum{}infty(M) -> C\_c\textasciicircum{}infty(M)\par
\noindent\textit{Critical reading.} C\_c\textasciicircum{}infty(M) -> C\_c\textasciicircum{}infty(N)\par
\noindent\textit{Evidence.} The scan faithfully repeats M; the immediately following variable y, integral over N, and product example fix the target.\par
\subsection*{P000196 (both; faithful error): explicit sigma\_T formula}
\noindent\textit{Transcription reads.} beta(a\textasciicircum{}{-1}z\textasciicircum{}{-1}xz) for beta defined on Xi\par
\noindent\textit{Critical reading.} Use phi\_!(gamma tensor beta); the displayed expression is not defined for general x,z.\par
\noindent\textit{Evidence.} The scan contains the explicit objection “Она не имеет смысла”; the domain mismatch verifies it independently.\par
\subsection*{P000197 (both; ambiguity): finite-codimension central ideal}
\noindent\textit{Transcription reads.} the ideal generated by the image of Z(g) has finite codimension\par
\noindent\textit{Critical reading.} The ideal generated by the image of the finite-codimensional annihilating ideal J has finite codimension in Z(z).\par
\noindent\textit{Evidence.} The diagram includes ideals under the centers. Taken literally the whole-center assertion is vacuous because 1 is in its image.\par
\subsection*{P000198-P000199 (both; provenance): speaker ownership}
\noindent\textit{Transcription reads.} Catalogue record 124 is Arnold; P000199 begins Anosov\par
\noindent\textit{Critical reading.} Keep the catalogue record, but attribute P000199 mathematical items to D. V. Anosov.\par
\noindent\textit{Evidence.} The scan has a separate explicit Anosov heading. No additional record is present in the supplied 38-record catalogue.\par
\subsection*{P000198 (both; ambiguity): sphere and torus stabilizer notation}
\noindent\textit{Transcription reads.} Diff(S\textasciicircum{}2-a,b,c), Diff(T\textasciicircum{}2-a)\par
\noindent\textit{Critical reading.} Diffeomorphisms in the specified identity class fixing the marked points, not unrestricted diffeomorphisms of punctured surfaces.\par
\noindent\textit{Evidence.} The scan says the points stay in place. Earle–Eells Theorem 1 uses pointwise marked-point stabilizers.\par
\subsection*{P000199 (both; ambiguity): operator inverse and Fourier normalization}
\noindent\textit{Transcription reads.} 0 is not an eigenvalue; use A inverse\par
\noindent\textit{Critical reading.} A inverse may be unbounded; define the transform spectrally with the inserted factor |A|\textasciicircum{}(n/2)/(2 pi)\textasciicircum{}(n/2).\par
\noindent\textit{Evidence.} The scan confirms both the weak spectral hypothesis and the separately inserted normalization. The implication to a bounded inverse would be invalid.\par
\subsection*{P000200 (both; faithful error): associativity square top arrow}
\noindent\textit{Transcription reads.} sigma times 1\par
\noindent\textit{Critical reading.} 1\_G times sigma for the displayed order G times G times X.\par
\noindent\textit{Evidence.} The scan clearly has sigma times 1. The domain of the action is G times X, so the labelled top map is ill typed in the displayed factor order.\par
\subsection*{P000200 (both; ambiguity): quotient criterion}
\noindent\textit{Transcription reads.} quotient exists iff all orbits closed\par
\noindent\textit{Critical reading.} For an affine reductive action, the invariant quotient is geometric iff all orbits are closed; the categorical quotient already exists.\par
\noindent\textit{Evidence.} The scan does not state affine or specify geometric. These hypotheses are required by the standard GIT theorem.\par
\subsection*{P000201 (both; faithful error): first uniqueness theorem}
\noindent\textit{Transcription reads.} Unique index follows from homotopy invariance and direct-sum additivity alone\par
\noindent\textit{Critical reading.} Normalization and additional functorial axioms are required; zero and integer multiples are counterexamples to uniqueness.\par
\noindent\textit{Evidence.} The scan isolates only I and II above the dividing line. The lower block adds III–VI, including winding normalization.\par
\subsection*{P000201 (both; ambiguity): one-dimensional winding formula}
\noindent\textit{Transcription reads.} Delta arg on R times S\textasciicircum{}0 without orientation\par
\noindent\textit{Critical reading.} Specify opposite component orientations and the Fourier/index sign convention.\par
\noindent\textit{Evidence.} Both components have a winding number; a numerical signed index requires their oriented difference.\par
\subsection*{P000202-P000204 (both; provenance): record date, speaker, and source order}
\noindent\textit{Transcription reads.} Catalogue record 127: Dynin, 27 March; record 128: Parshin, 20 March\par
\noindent\textit{Critical reading.} Retain catalogue record order 127 before 128, but record that the handwritten 20 March Parshin heading physically precedes the 27 March factorization lecture; Eskin is separately named on P000203.\par
\noindent\textit{Evidence.} The dates 20/III and 27/III are clearly visible on P000202. Neither the March 27 heading nor the Eskin subheading establishes that the catalogue speaker is mistaken.\par
\subsection*{P000203 (both; ambiguity): normal-limit symbol for nonzero order}
\noindent\textit{Transcription reads.} a\_-=A(x prime,0,0,-1)\par
\noindent\textit{Critical reading.} Either first normalize A to order zero, or include the order phase in the negative normal limit according to the chosen convention.\par
\noindent\textit{Evidence.} The scan omits normalization; Vishik–Eskin 1967 equation (1.2) has exp(i pi(alpha+i beta)) multiplying the negative limit.\par
\subsection*{P000203 (both; faithful error): description of b}
\noindent\textit{Transcription reads.} measure of nonellipticity\par
\noindent\textit{Critical reading.} Normal jump or factorization datum of an already elliptic interior symbol.\par
\noindent\textit{Evidence.} The scan first calls A elliptic and then uses the phrase nonellipticity. The invertible normal limits do not measure interior nonellipticity.\par
\subsection*{P000204 (both; ambiguity): gamma versus scalar partial index}
\noindent\textit{Transcription reads.} gamma\_j=kappa\_j\par
\noindent\textit{Critical reading.} Equality only with a specified compatible logarithm branch; generally gamma\_j is congruent to kappa\_j modulo integers.\par
\noindent\textit{Evidence.} The visible exponential b=exp(2 pi i kappa\_j) fixes the index only modulo integers. Vishik–Eskin distinguish kappa\_j=n\_j+gamma\_j.\par
\subsection*{P000204 (both; ambiguity): main global question}
\noindent\textit{Transcription reads.} G in E / G not in E / G and E\par
\noindent\textit{Critical reading.} Retain the connective as uncertain. Reconstruct a global potential/boundary completion question from the preceding system, without claiming a diplomatic resolution.\par
\noindent\textit{Evidence.} Full-page scan inspection does not distinguish the overwritten connective. The matching Vishik–Eskin theorem is contextual evidence, not a reading of this glyph.\par
\subsection*{P000204 (both; faithful error): ambient dimension}
\noindent\textit{Transcription reads.} Operator in R\textasciicircum{}n, then coordinates x\_0,...,x\_n\par
\noindent\textit{Critical reading.} Use N=n+1 throughout this report; preserve the subsequent n-based projective and homology indices.\par
\noindent\textit{Evidence.} The scan explicitly lists n+1 coordinates and the theorem on P000206 calls n+1 the dimension.\par
\subsection*{P000205 (both; ambiguity): sphere decomposition beside beta}
\noindent\textit{Transcription reads.} The real hyperplane section is a sum of spheres\par
\noindent\textit{Critical reading.} The matching source decomposes A\_R into m/2 spheres; it does not state this decomposition for every (A\_x)\_R. Retain the notebook wording with a scope warning.\par
\noindent\textit{Evidence.} Atiyah Bourbaki 319 p.91 states the decomposition before introducing the hyperplane section. The scan attaches the phrase to the latter.\par
\subsection*{P000205 (both; faithful error): wave equation projective section}
\noindent\textit{Transcription reads.} The section sum xi\_j squared=0 has no projective points\par
\noindent\textit{Critical reading.} Its real projective locus is empty. Its complex locus contains [1:i:0:...] when there are at least two coordinates.\par
\noindent\textit{Evidence.} The scan calls A complex earlier on the same page. The displayed complex point is a direct counterexample to the unrestricted assertion.\par
\subsection*{P000206 (both; ambiguity): signed chain on projective space}
\noindent\textit{Transcription reads.} sign(x dot xi) without a choice of representative\par
\noindent\textit{Critical reading.} Trivialize the real hyperplane line bundle on A\_R, or make an equivalent compatible lift and orientation choice.\par
\noindent\textit{Evidence.} Changing a homogeneous representative xi to -xi reverses the sign; Atiyah pp.91–92 supplies the missing trivialization.\par
\subsection*{P000206 (both; ambiguity): homology coefficients and theorem hypotheses}
\noindent\textit{Transcription reads.} alpha and beta in unspecified homology; m and n+1 even\par
\noindent\textit{Critical reading.} Use complex coefficients, N=n+1>2, strict hyperbolicity, and a smooth complex characteristic hypersurface.\par
\noindent\textit{Evidence.} These conditions are explicit in Atiyah pp.90–92; torsion-sensitive integral homology is not the stated criterion.\par
\subsection*{P000206 (both; faithful error): unique holomorphic form with a pole}
\noindent\textit{Transcription reads.} A unique holomorphic n-form with one pole of order n+1; no global sections hence uniqueness\par
\noindent\textit{Critical reading.} K\_CP\textasciicircum{}n(n+1) has a one-dimensional space of sections. Equivalently fix a polar hyperplane; uniqueness is up to a scalar.\par
\noindent\textit{Evidence.} Atiyah p.93 states the twisted section-space assertion. Vanishing of H\textasciicircum{}0(K) alone does not establish the notebook assertion.\par
\subsection*{P000207 (both; ambiguity): residue remainder and projective twisting}
\noindent\textit{Transcription reads.} Theta=0 on a=0; omega\_a used as an ordinary form\par
\noindent\textit{Critical reading.} Theta is a holomorphic remainder with zero pullback by degree, not necessarily vanishing coefficients; retain the overwritten last glyph as uncertain. The residue carries twist O(N-m).\par
\noindent\textit{Evidence.} Scan inspection leaves the overwritten parenthesis uncertain. The local residue identity and Atiyah p.93 determine the valid invariant formulation.\par
